% ============================================================
%  دادگاه در سایه — وجدان منفصل مسئول
%  سند پیشنهادی برای تأسیس نظام شبه‌قضایی اپوزیسیون ایران
%  نسخه ۰.۹ — پیش‌نویس برای بحث نخبگانی
%  ────────────────────────────────────────
%  کامپایل: XeLaTeX (الزامی)
%  حداقل دو بار کامپایل لازم است (برای فهرست مطالب و ارجاعات)
% ============================================================
\documentclass[12pt,a4paper]{book}

% ════════════════════════════════════════
%  ۱. بسته‌های پایه (قبل از xepersian)
%     ⚠ ترتیب بسته‌ها مهم است!
%     ⚠ xepersian باید آخرین بسته باشد!
% ════════════════════════════════════════

% --- رنگ و گرافیک ---
\usepackage{xcolor}
\usepackage{graphicx}
\usepackage{colortbl}          % ← اصلاح: برای \cellcolor
\usepackage{needspace}   % ← برای \needspace


% --- صفحه‌آرایی ---
\usepackage{geometry}
\usepackage{fancyhdr}
\usepackage{setspace}
\usepackage{pdflscape}         % ← اصلاح: برای محیط landscape

% --- جدول ---
\usepackage{longtable}
\usepackage{booktabs}
\usepackage{multirow}          % ← اضافه: برای \multirow در جداول
\usepackage{multicol}
\usepackage{array}             % ← اضافه: برای ستون‌های سفارشی

% --- TikZ (با تمام کتابخانه‌های لازم) ---
\usepackage{tikz}
\usetikzlibrary{
	positioning,                 % ← اصلاح: برای 'of' syntax
	arrows.meta,                 % ← اصلاح: برای {Stealth}
	calc,                        % ← اضافه: برای محاسبات مختصاتی
	shapes.geometric,            % ← اصلاح: برای ellipse و سایر اشکال
	shapes.symbols,              % ← اضافه
	decorations.pathreplacing,   % ← اضافه
	mindmap,                     % ← اصلاح: برای نمودار mindmap
	backgrounds,                 % ← اضافه
	fit                          % ← اضافه
}
\usepackage{pgf-pie}
% --- جعبه‌های رنگی ---
\usepackage[skins,breakable]{tcolorbox}

% --- لیست‌ها ---
\usepackage{enumitem}
\usepackage{amssymb}           % ← برای $\bigstar$, $\blacksquare$, $\blacktriangle$
\usepackage{bm}                % ← برای $\boldsymbol{}$
% --- متفرقه ---
\usepackage{float}
\usepackage{footnote}
\makesavenoteenv{longtable}
\usepackage{etoolbox}
\usepackage{pdfpages}

% --- عنوان‌بندی ---
\usepackage{titlesec}
% این بسته باید قبل از xepersian لود شود:
\usepackage{supertabular}

% --- لینک‌ها (قبل از xepersian) ---
\usepackage{hyperref}

% ════════════════════════════════════════
%  ۲. fontspec و تنظیمات قلم
%     ⚠ باید قبل از xepersian باشد
%     ⚠ از قلم‌هایی استفاده شده که
%       در اکثر سیستم‌ها موجودند
% ════════════════════════════════════════
\usepackage{fontspec}

% --- بررسی وجود قلم لاتین ---
% اگر TeX Gyre Termes ندارید، از Times New Roman یا Latin Modern استفاده شود
\IfFontExistsTF{TeX Gyre Termes}{%
	\setmainfont{TeX Gyre Termes}%  % اگر نصب است
}{%
	\IfFontExistsTF{Times New Roman}{%
		\setmainfont{Times New Roman}% % جایگزین اول
	}{%
		\setmainfont{Latin Modern Roman}% % جایگزین نهایی (همیشه در TeX هست)
	}%
}

% ════════════════════════════════════════
%  ۳. xepersian — آخرین بسته (الزامی!)
% ════════════════════════════════════════
\usepackage{xepersian}

% --- قلم فارسی ---
% IRLotus باید روی سیستم نصب باشد
% جایگزین‌ها: B Nazanin, Vazirmatn, XB Niloofar
\settextfont[Scale=1.05]{IRLotus}
\setdigitfont{IRLotus}

% --- قلم لاتین در متن فارسی ---
% ⚠ setlatintextfont باید بعد از xepersian باشد
\IfFontExistsTF{TeX Gyre Termes}{%
	\setlatintextfont[Scale=0.95]{TeX Gyre Termes}%
}{%
	\IfFontExistsTF{Times New Roman}{%
		\setlatintextfont[Scale=0.95]{Times New Roman}%
	}{%
		\setlatintextfont[Scale=0.95]{Latin Modern Roman}%
	}%
}

% ════════════════════════════════════════
%  ۴. تنظیمات صفحه
% ════════════════════════════════════════
\geometry{
	top=30mm,
	bottom=30mm,
	right=30mm,
	left=25mm,
	headheight=30pt             % ← اصلاح: از 15pt به 30pt
}
\setlength{\parindent}{0pt}
\setlength{\parskip}{8pt}
\onehalfspacing

% ════════════════════════════════════════
%  ۵. رنگ‌های اصلی سند (پالت ملایم)
% ════════════════════════════════════════
\definecolor{cPrimary}{HTML}{1B2A4A}    % سرمه‌ای ملایم
\definecolor{cAccent}{HTML}{C0392B}     % قرمز — هشدارها
\definecolor{cLight}{HTML}{F5F5F0}      % پس‌زمینه‌ی روشن
\definecolor{cGray}{HTML}{7F8C8D}       % خاکستری متن فرعی
\definecolor{cGold}{HTML}{B7950B}       % طلایی — تأکیدهای مثبت
\definecolor{cWarn}{HTML}{E67E22}       % نارنجی — هشدار متوسط
\definecolor{cGreen}{HTML}{27AE60}      % سبز — موفقیت
\definecolor{cConsolid}{HTML}{6C3483}   % بنفش مات — تثبیت و تداوم

% ════════════════════════════════════════
%  ۶. استایل عنوان‌ها
% ════════════════════════════════════════
% اول: فاصله‌ها
\titlespacing*{\section}{0pt}{12pt}{4pt}
\titlespacing*{\subsection}{0pt}{10pt}{3pt}
\titlespacing*{\paragraph}{0pt}{6pt}{1em}

% دوم: فرمت‌ها (باید بعد از spacing باشد)
\titleformat{\chapter}[display]
{\normalfont\Huge\bfseries\color{cPrimary}}
{\chaptertitlename\ \thechapter}{20pt}{\Huge}
\titleformat{\section}
{\normalfont\Large\bfseries\color{cPrimary}}
{\thesection}{1em}{}
\titleformat{\subsection}
{\normalfont\large\bfseries\color{cPrimary!80}}
{\thesubsection}{1em}{}
\titleformat{\paragraph}[runin]
{\normalfont\normalsize\bfseries\color{cPrimary!70}}
{}{}{}

% ════════════════════════════════════════
%  ۷. جعبه‌های اطلاعاتی
% ════════════════════════════════════════

\newtcolorbox{principleBox}[1][]{%
	enhanced, breakable,
	colback=cLight,
	colframe=cPrimary,
	coltitle=white,
	fonttitle=\bfseries,
	title=#1,
	boxrule=0.5pt,
	arc=2pt,
	left=8pt, right=8pt, top=6pt, bottom=6pt,
	shadow={0.5mm}{-0.5mm}{0mm}{black!15}
}

\newtcolorbox{warningBox}[1][]{%
	enhanced, breakable,
	colback=cAccent!5,
	colframe=cAccent!70,
	coltitle=white,
	fonttitle=\bfseries,
	title=#1,
	boxrule=0.5pt,
	arc=2pt,
	left=8pt, right=8pt, top=6pt, bottom=6pt,
}

\newtcolorbox{noteBox}[1][]{%
	enhanced, breakable,
	colback=cGold!5,
	colframe=cGold!70,
	coltitle=white,
	fonttitle=\bfseries,
	title=#1,
	boxrule=0.5pt,
	arc=2pt,
	left=8pt, right=8pt, top=6pt, bottom=6pt,
}

\newtcolorbox{infoBox}[1][]{%
	enhanced, breakable,
	colback=blue!3,
	colframe=cPrimary!50,
	coltitle=white,
	fonttitle=\bfseries,
	title=#1,
	boxrule=0.5pt,
	arc=2pt,
	left=8pt, right=8pt, top=6pt, bottom=6pt,
}

% ════════════════════════════════════════
%  ۸. سربرگ و پابرگ
% ════════════════════════════════════════
\pagestyle{fancy}
\fancyhf{}
\fancyhead[RO]{\small\color{cGray} دادگاه در سایه — وجدان منفصل مسئول}
\fancyhead[LE]{\small\color{cGray} سند پیشنهادی}
\fancyfoot[C]{\small\color{cGray}\thepage}
\renewcommand{\headrulewidth}{0.3pt}
\renewcommand{\footrulewidth}{0pt}

% ════════════════════════════════════════
%  ۹. تنظیم لینک‌ها
% ════════════════════════════════════════
\hypersetup{
	colorlinks=true,
	linkcolor=cPrimary,
	urlcolor=cPrimary!70,
	citecolor=cGold,
	pdfauthor={گروه تدوین},
	pdftitle={دادگاه در سایه — وجدان منفصل مسئول},
	pdfsubject={پیشنهاد تأسیس نظام شبه‌قضایی اپوزیسیون ایران}
}

% ════════════════════════════════════════
%  ۱۰. فرمان‌های کمکی
% ════════════════════════════════════════

% --- تأکید کلیدواژه ---
\newcommand{\keyword}[1]{\textbf{\textcolor{cPrimary}{#1}}}

% --- متن هشدار ---
\newcommand{\warn}[1]{\textcolor{cAccent}{#1}}

% --- متن انگلیسی (با پشتیبانی bidi) ---
% ⚠ \lr از xepersian می‌آید و جهت را به چپ‌به‌راست تغییر می‌دهد
\newcommand{\en}[1]{\lr{\textit{#1}}}
% ────────────────────────────────────────
% اگر می‌خواهید متن انگلیسی italic نباشد:
% \newcommand{\en}[1]{\lr{#1}}
% ────────────────────────────────────────

% ════════════════════════════════════════
%  ۱۱. تنظیمات اضافی bidi و RTL
% ════════════════════════════════════════

% longtable در محیط RTL گاهی مشکل دارد — این خط کمک می‌کند:
\AtBeginEnvironment{longtable}{\small}

% اطمینان از عملکرد صحیح شماره‌گذاری فارسی در جداول:
\SepMark{.}


% ============================================================
%  شروع سند
% ============================================================
\begin{document}
	
	% ////////////////////////////////////////////////////////////
	%  صفحه‌ی عنوان
	% ////////////////////////////////////////////////////////////
\begin{titlepage}
	\begin{tikzpicture}[remember picture, overlay]
		\fill[cPrimary] (current page.south west) rectangle (current page.north east);
		\fill[cGold, opacity=0.8]
		([yshift=6cm]current page.south west) rectangle
		([yshift=6.15cm]current page.south east);
	\end{tikzpicture}
	
	\vspace*{4cm}
	\begin{center}
		% ← هر خط جداگانه \textcolor دارد
		{\fontsize{38}{46}\selectfont\bfseries
			\textcolor{white}{\rl{ دادگاه در سایه / دادگاه در تبعید!}}}\\[12pt]
		{\fontsize{18}{24}\selectfont
			\textcolor{cGold}{وجدان منفصل مسئول}}\\[30pt]
		{\fontsize{14}{20}\selectfont
			\textcolor{white}{سند پیشنهادی برای تأسیس نظام شبه‌قضایی}}\\[4pt]
		{\fontsize{14}{20}\selectfont
			\textcolor{white}{اپوزیسیون دموکراتیک ایران}}\\[50pt]
		{\large
			\textcolor{white}{مهدی سالم؛ ریچموند هیل}}\\[22pt]
		{\large
			\textcolor{white!90}{پیش‌نویس نسخه‌ی ۰.۹ — برای بحث نخبگانی}}\\[8pt]
		{\textcolor{white!80}{\today}}
	\end{center}
	
	\vfill
	\begin{center}
		{\small
			\textcolor{white!70}{این سند محرمانه نیست. بازتولید با ذکر منبع آزاد است.}}\\[2pt]
		{\small
			\textcolor{white!70}{\en{Creative Commons BY-SA 4.0}}}
	\end{center}
\end{titlepage}
	
	% ////////////////////////////////////////////////////////////
	%  فهرست مطالب
	% ////////////////////////////////////////////////////////////
	\frontmatter
	\tableofcontents
	\clearpage
	
	% ////////////////////////////////////////////////////////////
	%  خلاصه‌ی اجرایی (۳ صفحه)
	% ////////////////////////////////////////////////////////////
	\mainmatter
	\chapter*{خلاصه‌ی اجرایی}
	\addcontentsline{toc}{chapter}{خلاصه‌ی اجرایی}
	\thispagestyle{fancy}
	
	\begin{principleBox}[گزاره‌ی محوری]
		جنبش گذارخواهی ایران در مقطع کنونی با \keyword{خلأ مشروعیت‌بخشی حقوقی}
		مواجه است: تصمیم‌هایی با ابعاد قضایی — از حمایت از تحریم‌ها تا مستندسازی
		جرایم حکومتی — بدون هیچ \keyword{فرایند دادرسی ساختارمند} اتخاذ می‌شوند.
		نبودِ چنین فرایندی هم اپوزیسیون را در برابر جنگ روانی جمهوری اسلامی آسیب‌پذیر
		می‌کند و هم ظرفیت گفت‌وگوی سازنده‌ی درون‌جنبشی را تضعیف می‌نماید.
		
		\medskip
		\textbf{پیشنهاد:} تأسیس یک \keyword{نهاد شبه‌قضایی سایه}
		(\en{Shadow Quasi-Judicial Body})
		با عنوان \textbf{«دادگاه در سایه — وجدان منفصل مسئول»}
		که وظایف زیر را بر عهده گیرد:
		
		\begin{enumerate}[nosep,label=\textcolor{cPrimary}{\arabic*.}]
			\item مستندسازی نظام‌مند جرایم رژیم بر پایه‌ی معیارهای حقوق بین‌الملل؛
			\item داوری و تنظیم‌بخشی تصمیم‌های حساس جنبش بر اساس «پاسخ مشروع و متناسب»؛
			\item فراهم‌سازی بستر گفت‌وگوی نخبگانی برای مسائل پرمناقشه
			(تحریم، مداخله‌ی بشردوستانه، اخراج یا تعقیب وابستگان حکومتی و …)؛
			\item ثبت و بایگانی آراء و دلایل برای بازنگری پس از گذار.
		\end{enumerate}
	\end{principleBox}
	
\begin{figure}[H]
	\noindent\makebox[\textwidth][c]{%
		\begin{tikzpicture}[
			% ── استایل‌ها — اندازه‌ها کاهش یافته ──
			pillar/.style={
				draw=#1, fill=#1!8, line width=1pt,
				rounded corners=5pt,
				text width=3cm, minimum height=2cm,
				align=center, font=\footnotesize,
				inner sep=6pt
			},
			icon/.style={
				circle, draw=#1, fill=#1!15, line width=1pt,
				minimum size=0.8cm, font=\normalsize,
				inner sep=0pt
			},
			connector/.style={
				-{Stealth}, thick, #1!60, line width=0.8pt
			}
			]
			
			% ═══════════════════════════════════════
			%  مرکز — شعار اصلی (کوچک‌تر)
			% ═══════════════════════════════════════
			\node[
			ellipse, draw=cPrimary, fill=cPrimary!12,
			line width=1.5pt,
			minimum width=6cm, minimum height=2.6cm,
			align=center, inner sep=8pt
			] (center) {
				{\normalsize\bfseries\textcolor{cPrimary}{\RL{دادخواهی اندیشیده، مستند، و انسانی}}}\\[4pt]
				{\footnotesize\textcolor{cGold}{\RL{صدای بی‌صدایان، وجدان بیدار ملت}}}
			};
			
			% ═══════════════════════════════════════
			%  ستون ۱: فریاد مظلوم — پایین وسط
			% ═══════════════════════════════════════
			\node[pillar=cAccent, below=2.2cm of center] (p1) {
				{\small\bfseries\textcolor{cAccent}{\RL{فریاد مظلوم}}}\\[3pt]
				{\scriptsize\RL{انعکاس جهانی ظلم سیاسی}}\\[2pt]
				{\tiny\textcolor{cGray}{\RL{صدای قربانیان را به گوش}}}\\
				{\tiny\textcolor{cGray}{\RL{جهان می‌رسانیم — با سند}}}
			};
			\node[icon=cAccent, anchor=south] (i1) at (p1.north) {};
			\node[font=\normalsize] at (i1) {\RL{✊}};
			
			% ═══════════════════════════════════════
			%  ستون ۲: چندصدایی — چپ وسط
			% ═══════════════════════════════════════
			\node[pillar=cConsolid, left=2.8cm of center] (p5) {
				{\small\bfseries\textcolor{cConsolid}{\RL{چندصدایی}}}\\[3pt]
				{\scriptsize\RL{وجدان جمعی فراگیر}}\\[2pt]
				{\tiny\textcolor{cGray}{\RL{نمایندگی طیف‌های متنوع}}}\\
				{\tiny\textcolor{cGray}{\RL{فراجناحی و فراقومی}}}
			};
			\node[icon=cConsolid, anchor=west] (i5) at (p5.east) {};
			\node[font=\normalsize] at (i5) {\RL{🤝}};
			
			% ═══════════════════════════════════════
			%  ستون ۳: عمق نظری — راست وسط
			% ═══════════════════════════════════════
			\node[pillar=cPrimary, right=2.8cm of center] (p2) {
				{\small\bfseries\textcolor{cPrimary}{\RL{عمق نظری}}}\\[3pt]
				{\scriptsize\RL{پایه‌های حقوقی استوار}}\\[2pt]
				{\tiny\textcolor{cGray}{\RL{ابتنا بر حقوق بین‌الملل}}}\\
				{\tiny\textcolor{cGray}{\RL{و عدالت انتقالی جهان}}}
			};
			\node[icon=cPrimary, anchor=east] (i2) at (p2.west) {};
			\node[font=\normalsize] at (i2) {\RL{⚖}};
			
			% ═══════════════════════════════════════
			%  ستون ۴: عمل‌گرایی — بالا چپ
			% ═══════════════════════════════════════
			\node[pillar=cWarn, above left=2cm and 1.2cm of center] (p3) {
				{\small\bfseries\textcolor{cWarn}{\RL{عمل‌گرایی}}}\\[3pt]
				{\scriptsize\RL{پا در میدان واقعیت}}\\[2pt]
				{\tiny\textcolor{cGray}{\RL{با دستانی گره‌کرده}}}\\
				{\tiny\textcolor{cGray}{\RL{نه فقط نظریه‌پردازی}}}
			};
			\node[icon=cWarn, anchor=north] (i3) at (p3.south) {};
			\node[font=\normalsize] at (i3) {\RL{⚡}};
			
			% ═══════════════════════════════════════
			%  ستون ۵: دقت وسواسی — بالا راست
			% ═══════════════════════════════════════
			\node[pillar=green!60!black, above right=2cm and 1.2cm of center] (p4) {
				{\small\bfseries\textcolor{green!60!black}{\RL{دقت وسواسی}}}\\[3pt]
				{\scriptsize\RL{مستندسازی و بازبینی}}\\[2pt]
				{\tiny\textcolor{cGray}{\RL{غیراحساسی، ساختارمند}}}\\
				{\tiny\textcolor{cGray}{\RL{و قابل بازنگری}}}
			};
			\node[icon=green!60!black, anchor=north] (i4) at (p4.south) {};
			\node[font=\normalsize] at (i4) {\RL{🔍}};
			
			% ═══════════════════════════════════════
			%  خطوط اتصال — از آیکون‌ها به مرکز
			% ═══════════════════════════════════════
			\draw[connector=cAccent] (i1) -- (center);
			\draw[connector=cConsolid] (i5) -- (center);
			\draw[connector=cPrimary] (i2) -- (center);
			\draw[connector=cWarn] (i3) -- (center);
			\draw[connector=green!60!black] (i4) -- (center);
			
			% ═══════════════════════════════════════
			%  حلقه‌ی اتصال خارجی — خط‌چین
			% ═══════════════════════════════════════
			\draw[dashed, cGray!30, line width=0.6pt]
			(p3.east) to[bend left=12] (p4.west);
			\draw[dashed, cGray!30, line width=0.6pt]
			(p4.south west) to[bend left=12] (p2.north);
			\draw[dashed, cGray!30, line width=0.6pt]
			(p2.south) to[bend left=12] (p1.east);
			\draw[dashed, cGray!30, line width=0.6pt]
			(p1.west) to[bend left=12] (p5.south);
			\draw[dashed, cGray!30, line width=0.6pt]
			(p5.north) to[bend left=12] (p3.south west);
			
			% ═══════════════════════════════════════
			%  بیانیه — زیر نمودار
			% ═══════════════════════════════════════
			\node[
			below=0.6cm of p1,
			text width=12cm, align=center,
			font=\footnotesize
			] (motto) {
				\textcolor{cPrimary}{\rule{3cm}{0.3pt}}\\[4pt]
				{\small\bfseries\textcolor{cPrimary}{%
						\RL{«ما نه از سر خشم، که از سر خِرد دادخواهی می‌کنیم.}}}\\[2pt]
				{\small\bfseries\textcolor{cPrimary}{%
						\RL{نه برای انتقام، که برای ثبت حقیقت.}}}\\[2pt]
				{\small\bfseries\textcolor{cPrimary}{%
						\RL{نه به نمایندگی یک جناح، که به وکالت از وجدان ملت.»}}}\\[4pt]
				\textcolor{cPrimary}{\rule{3cm}{0.3pt}}
			};
			
		\end{tikzpicture}%
	}%
	\caption{پنج ستون راهنمای عمل دادگاه در سایه}
\end{figure}	

	
	«ما نه از سر خشم، که از سر خِرد دادخواهی می‌کنیم.
	نه برای انتقام، که برای ثبت حقیقت.
	نه به نمایندگی یک جناح، که به وکالت از وجدان ملت.»
	
	
	
	\section*{چرا «دادگاه در سایه»؟}
	
	جنبش‌های گذارخواه در سراسر تاریخ معاصر — از
	\en{Free French} تا \en{ANC} آفریقای جنوبی —
	با ضرورت اتخاذ تصمیم‌هایی مواجه بوده‌اند که ماهیت قضایی دارند
	اما پیش از دست‌یابی به قدرت سیاسی صورت می‌گیرند.
	وقتی یک جنبش از تحریم فلان نهاد ورزشی حمایت می‌کند، عملاً
	\keyword{حکمی صادر کرده} که بر زندگی حرفه‌ای هزاران ورزشکار اثر می‌گذارد.
	وقتی خواستار اخراج فرزندان مقام‌های حکومتی از کشور ثالث می‌شود،
	\keyword{مجازاتی اعمال کرده} که بدون دادرسی است.
	
	\medskip
	نبودِ دادرسی به معنای نبودِ هر سه عنصر زیر است:
	
	\begin{itemize}[nosep]
		\item \textbf{مبنای حقوقی:} چرا این اقدام مجاز است؟
		\item \textbf{تناسب:} آیا هزینه‌ی وارد بر بی‌گناهان، از فایده‌ی مورد انتظار کمتر است؟
		\item \textbf{مستندسازی:} آیا استدلال‌ها و مخالفت‌ها ثبت شده‌اند
		تا در آینده قابل بازنگری باشند؟
	\end{itemize}
	
	\begin{warningBox}[خلأ موجود و بهره‌برداری جمهوری اسلامی]
		جمهوری اسلامی به‌صورت سازمان‌یافته از \keyword{نبودِ فرایند}
		بهره‌برداری تبلیغاتی می‌کند:
		\begin{itemize}[nosep]
			\item «اپوزیسیون بدون هیچ محاکمه‌ای، مردم را تحریم می‌کند.»
			\item «گذارخواهان از جنگ و مداخله‌ی خارجی استقبال می‌کنند.»
			\item «آن‌ها حتی به ورزشکاران و هنرمندان رحم نمی‌کنند.»
		\end{itemize}
		
		\medskip
		وجود یک نهاد معتبر شبه‌قضایی با آراء مستدل و مستند و ثبت نظرات موافق و مخالف،
		\textbf{قوی‌ترین دفاعیه} در برابر این جنگ رسانه‌ای است.
	\end{warningBox}
	
	\section*{مخاطبان اصلی سند}
	
	\begin{enumerate}[nosep,label=\textcolor{cPrimary}{\Alph*.}]
		\item \textbf{فعالان سیاسی گذارخواه و براندازانه‌ی ایرانی}
		— برای درک ضرورت ساختارمندسازی؛
		\item \textbf{حقوق‌دانان و پژوهشگران ایرانی و بین‌المللی}
		— برای مشارکت در طراحی؛
		\item \textbf{نهادهای بین‌المللی حقوق بشر}
		— برای همکاری فنی و اعتباربخشی؛
		\item \textbf{افکار عمومی ایرانی}
		— برای اطمینان از وجود نظارت اخلاقی بر تصمیم‌های سخت.
	\end{enumerate}
	
	\section*{ساختار سند}
	
	این سند شامل ۱۱ فصل اصلی و ۱۱ پیوست است:
	
	\begin{longtable}{r p{10cm}}
		\toprule
		\textbf{بخش} & \textbf{موضوع} \\
		\midrule
		\endhead
		خلاصه‌ی اجرایی & همین بخش \\
		فصل ۱ & تشریح مسئله و جدیت نیاز \\
		فصل ۲ & انتظارات مثبت و شرایط کمک‌کننده \\
		فصل ۳ & آسیب‌شناسی و احتیاط‌های لازم \\
		فصل ۴ & دغدغه‌های سازماندهی و تکثر دیدگاهی \\
		فصل ۵ & رویه‌های اجرایی و قضایی معتبر \\
		فصل ۶ & مباحث مالی و هزینه‌ای \\
		فصل ۷ & ارتباط با مراجع بین‌المللی \\
		فصل ۸ & نمونه‌های مشابه و مطالعه‌ی موردی \\
		فصل ۹ & هضم و ادغام در نظام قضایی پس از گذار \\
		فصل ۱۰ & افراد پیشنهادی \\
		فصل ۱۱ & جمع‌بندی نهایی و نتیجه‌گیری \\
		\midrule
		پیوست الف & واژه‌نامه‌ی تخصصی \\
		پیوست ب & فهرست اسناد بین‌المللی مرجع \\
		پیوست ج & نمونه‌ی فرم ارجاع پرونده \\
		پیوست د & نمونه‌ی کیفرخواست سایه \\
		پیوست ه‍ & نمونه‌ی رأی مشورتی \\
		پیوست و & استراتژی ارتباطی و رسانه‌ای \\
		پیوست ز & زیرساخت فنی و امنیت دیجیتال \\
		پیوست ح & کتاب‌نامه و منابع \\
		پیوست ط & پرسش‌های متداول \\
		پیوست ی & نقشه‌ی راه اجرایی مفصل \\
		پیوست ک & فهرست جامع موضوعات قابل طرح، بررسی و داوری \\
		\bottomrule
	\end{longtable}

%  فصل ۱: تشریح مسئله و جدیت نیاز
	\chapter{تشریح مسئله و جدیت نیاز}\label{chap:problem}
	
	\section{خلأ نهادی در جنبش گذارخواهی ایران}
	
	از جنبش سبز ۱۳۸۸ تا خیزش ژینا ۱۴۰۱ و ادامه‌ی آن، به خصوص جنبش دی ماه 1404 که با سرکوب شدید و وحشیانیه‌ی حاکمیت فراتر از هر گونه تصوری نزد جامعه صورت گرتف. جنبش مردمی ایران با رشد کمّی و کیفی چشمگیری همراه بوده، اما یک عنصر ساختاری همچنان غایب است: \keyword{سازوکار تصمیم‌گیری حقوقی‌ـ‌اخلاقی}
	برای مسائلی که ماهیت قضایی دارند و گرفتار هیجان انقلابی و ضرورت‌های تصمیم‌های سیاسی می‌شوند. این مسائل را می‌توان در سه دسته طبقه‌بندی کرد:
	
	\subsection{دسته‌ی اول: تصمیم‌های با آثار کیفری‌ـ‌حقوقی}
	
	\begin{itemize}[nosep]
		\item مستندسازی جرایم جنگی و جرایم علیه بشریت
		(کشتار آبان ۹۸، سرکوب ژینا ۱۴۰۱، قتل و عام 1404، شکنجه‌ی زندانیان سیاسی)؛
		\item اعلام جرم علیه اشخاص حقیقی (فرماندهان سپاه، قضات، بازجویان)؛
		\item اعلام جرم علیه اشخاص حقوقی (نهادهای حکومتی، شرکت‌های وابسته به سپاه)؛
		\item درخواست صدور حکم دستگیری بین‌المللی از طریق
		\en{Universal Jurisdiction} در کشورهای ثالث.
	\end{itemize}
	
	\begin{noteBox}[نمونه‌ی واقعی]
		پرونده‌ی حمید نوری در سوئد نشان داد که تعقیب قضایی مقام‌های
		جمهوری اسلامی در خارج از کشور \textbf{عملی} است، اما فقدان
		یک نهاد هماهنگ‌کننده سبب شد تا رسیدن به این نقطه بیش از سه دهه
		طول بکشد. دادگاه در سایه می‌تواند این فرایند را \keyword{نظام‌مند و تسریع} کند.
	\end{noteBox}
	
	\subsection{دسته‌ی دوم: تصمیم‌های با آثار سیاسی‌ـ‌اجتماعی}
	
	این‌ها تصمیم‌هایی هستند که مستقیماً کیفری نیستند اما بر زندگی میلیون‌ها نفر
	اثر می‌گذارند و نیازمند \keyword{توجیه حقوقی‌ـ‌اخلاقی} هستند:
	
	\begin{longtable}{r p{7cm} p{6.5cm}}
		\toprule
		\textbf{ردیف} & \textbf{نوع تصمیم} & \textbf{گروه متأثر} \\
		\midrule
		\endhead
		۱ & حمایت از تحریم‌های اقتصادی &
		مردم عادی، طبقه‌ی متوسط، بیماران نیازمند دارو \\
		۲ & حمایت از تحریم‌های ورزشی و فرهنگی &
		ورزشکاران، هنرمندان، دانشگاهیان بدون دغدغه‌ی سیاسی \\
		۳ & درخواست اخراج وابستگان مقام‌ها از کشورهای ثالث &
		خانواده‌ها شامل افراد بی‌تقصیر احتمالی \\
		۴ & حمایت از مداخله‌ی بشردوستانه &
		کل جمعیت ایران، امنیت منطقه‌ای \\
		۵ & فراخوان اعتصاب عمومی &
		کارگران، کسبه‌ی خُرد، معلمان \\
		۶ & فراخوان نافرمانی مدنی خاص &
		شهروندانی که ریسک بازداشت را متحمل می‌شوند \\
		\bottomrule
	\end{longtable}
	
	\begin{warningBox}[مسئله‌ی کلیدی]
		هر یک از این تصمیم‌ها بدون فرایند دادرسی، صرفاً یک
		\textbf{اعلام موضع سیاسی} است. با فرایند دادرسی، به یک
		\textbf{توصیه‌ی مستدل حقوقی‌ـ‌اخلاقی} تبدیل می‌شود که هم
		اعتبار بیشتری دارد و هم در برابر بهره‌برداری تبلیغاتی مقاوم‌تر است.
	\end{warningBox}

	\subsection{دسته‌ی سوم: تصمیم‌های گفتمانی و مرزگذاری}
	\vspace{-5pt}%	
	این دسته شامل مباحثی است که اگرچه مستقیماً «تصمیم» به‌معنای رایج نیستند،
	اما مرزهای گفتمانی جنبش را تعیین می‌کنند و بدون داوری ساختارمند،
	به \keyword{شکاف‌های فلج‌کننده} تبدیل می‌شوند:
	
	\begin{itemize}[nosep]
		\item آیا همکاری با دولت‌های خارجی خاص (ایالات متحده، اسرائیل، عربستان)
		مجاز، مطلوب، یا ممنوع است؟
		\item مرز «خاکستری‌ها» کجاست؟ کسانی که نه کاملاً حکومتی‌اند و نه
		کاملاً اپوزیسیون — ورزشکارانی که پرچم جمهوری اسلامی را حمل
		می‌کنند اما شعار سیاسی نمی‌دهند.
		\item آیا عدالت انتقالی باید شامل عفو مشروط هم باشد؟
		اگر آری، با چه سازوکاری؟
		\item دفاع مشروع (\en{legitimate defense}) تا کجا می‌تواند
		شامل خشونت قهری شود؟
	\end{itemize}
	\vspace{-20pt}%
	\section{چرا وضعیت «اضطراری» است؟}
	\vspace{-10pt}%
	\nopagebreak[5]%
	\begin{principleBox}[پنج فاکتور فوریت]
		\begin{enumerate}[nosep, label=\textcolor{cAccent}{\arabic*.}]
			\item \textbf{شتاب رویدادها:}
			فاصله‌ی بین خیزش‌های مردمی در حال کاهش است
			(۱۳۸۸ → ۱۳۹۶ → ۱۳۹۸ → ۱۴۰۱). پنجره‌ی فرصت برای
			ساختارسازی نهادی ممکن است کوتاه باشد.
			
			\item \textbf{بهره‌برداری حکومت از خلأ:}
			جمهوری اسلامی دستگاه عظیم رسانه‌ای و سایبری خود را
			صرف نشان دادن اپوزیسیون به‌عنوان «فاقد مشروعیت و ساختار» می‌کند.
			هر روز تأخیر، یک روز بیشتر فرصت برای این بهره‌برداری است.
			
			\item \textbf{فرسودگی مستندسازی:}
			شاهدان دارند پیر می‌شوند، مدارک از بین می‌روند، خاطره‌ها کم‌رنگ
			می‌شوند. مستندسازی ساختارمند جرایم یک مسئله‌ی زمان‌حساس است.
			
			\item \textbf{تصمیم‌های برگشت‌ناپذیر:}
			برخی تصمیم‌ها — مثل حمایت از یک مداخله‌ی نظامی — اگر اشتباه
			گرفته شوند، قابل جبران نیستند. فرایند دادرسی سدّ حفاظتی
			در برابر شتاب‌زدگی است.
			
			\item \textbf{اعتبار بین‌المللی:}
			دولت‌ها و نهادهای بین‌المللی تمایل بیشتری به تعامل با
			اپوزیسیونی دارند که \keyword{ساختار نهادی شفاف}
			داشته باشد.
		\end{enumerate}
	\end{principleBox}
	
	\section{جمهوری اسلامی و جنگ روانی ساختارمند}
	
	برای درک ضرورت این نهاد، باید \keyword{مدل جنگ رسانه‌ای} جمهوری اسلامی را شناخت.
	حکومت در مواجهه با هر اقدام اپوزیسیون از یک \keyword{ماتریس واکنش سه‌لایه‌ای} استفاده می‌کند:
		\begin{figure}[H]
		\centering
		\begin{tikzpicture}[
			node distance=1.5cm,
			box/.style={
				rectangle, rounded corners=3pt, draw=cPrimary, thick,
				fill=cLight, text width=5.5cm, minimum height=1.5cm,
				align=center, font=\small
			},
			arrow/.style={-{Stealth}, thick, cPrimary}
			]
			\node[box] (L1) {{\RL{\textbf{لایه‌ی ۱: انکار}}}\\[3pt]
				{\RL{«جرمی رخ نداده؛ اپوزیسیون دروغ می‌گوید.»}}};
			
			\node[box, below=of L1] (L2) {{\RL{\textbf{لایه‌ی ۲: تخریب فرایند}}}\\[3pt]
				{\RL{«حتی اگر رخ داده، اپوزیسیون صلاحیت قضاوت ندارد.»}}};
			
			\node[box, below=of L2] (L3) {{\RL{\textbf{لایه‌ی ۳: مظلوم‌نمایی}}}\\[3pt]
				{\RL{«اپوزیسیون به مردم/ورزشکاران/هنرمندان آسیب می‌زند.»}}};
			
			\draw[arrow] (L1) -- (L2);
			\draw[arrow] (L2) -- (L3);
			
			\node[box, fill=cGold!10, draw=cGold, text width=5cm,
			right=3cm of L2]
			(defense) {{\RL{\textbf{پاسخ دادگاه در سایه:}}}\\[3pt]
				{\RL{مستندسازی شفاف + فرایند دادرسی معتبر + ثبت آراء موافق و مخالف}}};
			
			\draw[-{Stealth}, thick, cGold] (defense) -- (L1);
			\draw[-{Stealth}, thick, cGold] (defense) -- (L2);
			\draw[-{Stealth}, thick, cGold] (defense) -- (L3);
		\end{tikzpicture}
		\caption{ماتریس واکنش سه‌لایه‌ای جمهوری اسلامی و پاسخ نهاد پیشنهادی}
	\end{figure}
	\vspace{-30pt}%
	\subsection{تحلیل هر لایه و پاسخ نهادی}
	
	\paragraph{\textbf{لایه‌ی انکار:}}
	جمهوری اسلامی به‌صورت نظام‌مند وقوع جرایم خود را انکار می‌کند —
	از تیراندازی مستقیم به معترضان آبان ۹۸ گرفته تا مسمومیت دانش‌آموزان.
	دادگاه در سایه با ایجاد \keyword{بانک شواهد ساختارمند} شامل
	شهادت‌نامه‌های تأییدشده، تصاویر ماهواره‌ای، گزارش‌های پزشکی قانونی،
	و تحلیل منابع باز (\en{OSINT})، لایه‌ی انکار را نه از طریق
	شعار بلکه از طریق \keyword{فرایند حقوقی استاندارد} خنثی می‌کند.
	
	\paragraph{\textbf{لایه‌ی تخریب فرایند:}}
	مؤثرترین ابزار جمهوری اسلامی، زیر سؤال بردن صلاحیت اپوزیسیون
	برای هرگونه قضاوت است. این لایه فقط با \keyword{ساختار نهادی معتبر}
	قابل پاسخ‌گویی است: هیئت‌داوران مستقل، آیین دادرسی شفاف، حق دفاع
	(حتی برای متهمانِ غایب)، و انتشار آراء با نظرات اقلیت.
	
	\paragraph{\textbf{لایه‌ی مظلوم‌نمایی:}}
	وقتی جنبش بدون فرایند از تحریم ورزشی حمایت می‌کند، حکومت
	ورزشکاران آسیب‌دیده را ابزار تبلیغاتی می‌کند.
	اگر همین تصمیم از مسیر دادرسی عبور کند — با بررسی هزینه‌ـ‌فایده،
	ثبت نظرات مخالف، و ارائه‌ی راهکارهای جایگزین — بهره‌برداری
	تبلیغاتی \keyword{به‌شدت دشوارتر} می‌شود.

	\section{ابعاد حقوق بشری و مبانی حقوقی}

	\subsection{اسناد بین‌المللی پشتیبان}
	
	نهاد پیشنهادی بر مبنای اسناد حقوقی بین‌المللی زیر طراحی شده
	و مشروعیت خود را از آن‌ها وام می‌گیرد:
	
	\begin{longtable}{r p{6cm} p{8cm}}
		\toprule
		\textbf{ردیف} & \textbf{سند} & \textbf{ارتباط با نهاد پیشنهادی} \\
		\midrule
		\endhead
		۱ & اعلامیه‌ی جهانی حقوق بشر (۱۹۴۸) &
		مواد ۸، ۱۰ و ۱۱: حق دادرسی عادلانه —
		حتی نهاد سایه باید این حق را محترم بشمارد. \\
		۲ & میثاق بین‌المللی حقوق مدنی و سیاسی (\en{ICCPR}) &
		مواد ۱۴ و ۱۹: دادرسی منصفانه و آزادی بیان —
		مبنای آیین دادرسی نهاد. \\
		۳ & اساسنامه‌ی رم (۱۹۹۸) &
		تعاریف جرایم علیه بشریت و جرایم جنگی —
		چارچوب طبقه‌بندی جرایم مستندشده. \\
		۴ & اصول بنیادین نورنبرگ (۱۹۵۰) &
		مسئولیت فردی فرماندهان —
		مبنای اعلام جرم علیه اشخاص. \\
		۵ & قطعنامه‌های شورای حقوق بشر سازمان ملل درباره‌ی ایران &
		به‌ویژه گزارش‌های گزارشگر ویژه —
		منابع مکمل شواهد. \\
		۶ & اصول مسئولیت حمایت (\en{R2P}، ۲۰۰۵) &
		چارچوب بررسی مجاز بودن مداخله‌ی بشردوستانه. \\
		۷ & اصول ژوئنه/اورنتلیشر (عدالت انتقالی) &
		حق قربانیان برای دانستن حقیقت، عدالت، جبران خسارت،
		و تضمین عدم تکرار. \\
		\bottomrule
	\end{longtable}
	\vspace{-30pt}%
	\subsection{مبنای نظری: «وجدان منفصل مسئول» چیست؟}
	
	عنوان فرعی این سند — \keyword{وجدان منفصل مسئول} — ناظر به
	مفهومی است که باید به‌دقت تبیین شود:
	
	\begin{principleBox}[تعریف]
		\textbf{وجدان منفصل مسئول} (\en{Responsible Detached Conscience})
		به نهادی اطلاق می‌شود که:
		
		\begin{itemize}[nosep]
			\item \textbf{منفصل} است: مستقل از فشار سیاسی جناح‌ها و احساسات
			عمومی تصمیم می‌گیرد — مانند قاضی مستقل در یک نظام قضایی سالم؛
			\item \textbf{وجدان} است: نقش داوری اخلاقی دارد، نه صرفاً حقوقیِ تکنیکال
			— دغدغه‌ی کرامت انسانی و عدالت را مرکز کار قرار می‌دهد؛
			\item \textbf{مسئول} است: آراء و توصیه‌هایش ثبت، منتشر و قابل بازنگری‌اند
			— در برابر تاریخ و مردم پاسخ‌گوست.
		\end{itemize}
	\end{principleBox}
	
	این مفهوم از سنت فلسفی هانا آرنت
	(درباره‌ی قضاوت و مسئولیت در شرایط بحران) و سنت حقوقی
	دوست دادگاه (\en{amicus curiae}) الهام گرفته
	و با شرایط خاص جنبش ایران تطبیق داده شده است.
	
	\section{جمع‌بندی فصل اول}
	
	\begin{principleBox}[خلاصه‌ی استدلال فصل ۱]
		\begin{enumerate}[nosep, label=\textcolor{cPrimary}{\arabic*.}]
			\item جنبش گذارخواهی ایران تصمیم‌هایی با ماهیت قضایی اتخاذ می‌کند
			بدون داشتن سازوکار قضایی.
			\item این خلأ هم به اعتبار جنبش آسیب می‌زند و هم ابزار تبلیغاتی
			به جمهوری اسلامی می‌دهد.
			\item مستندسازی جرایم بدون ساختار نهادی، ناقص و فرسایشی خواهد بود.
			\item فوریت مسئله ناشی از شتاب رویدادها، فرسودگی شواهد، و
			ضرورت اعتبار بین‌المللی است.
			\item مبانی حقوقی بین‌المللی کافی برای تأسیس چنین نهادی وجود دارد.
		\end{enumerate}
	\end{principleBox}
	
	\clearpage
	
	
	% ////////////////////////////////////////////////////////////
	%  فصل ۲: انتظارات مثبت و شرایط کمک‌کننده
	% ////////////////////////////////////////////////////////////
	\chapter{انتظارات مثبت و شرایط کمک‌کننده}\label{chap:expect}
	
	\section{دستاوردهای قابل انتظار}
	
	تأسیس دادگاه در سایه، در صورت اجرای صحیح، می‌تواند
	دستاوردهای زیر را به ارمغان بیاورد:
	
	\subsection{دستاوردهای کوتاه‌مدت (۰ تا ۱۲ ماه)}
	
	\begin{enumerate}[nosep, label=\textcolor{cGold}{\alph*.}]
		\item \textbf{مرجع واحد مستندسازی:}
		ایجاد یک بانک اطلاعاتی یکپارچه از جرایم حکومتی با
		استانداردهای حقوق بین‌الملل، قابل ارائه به دادگاه‌های ملی
		و بین‌المللی.
		
		\item \textbf{اعتبار رسانه‌ای:}
		وقتی یک نهاد ساختارمند با اعضای شناخته‌شده‌ی آکادمیک
		«رأی» صادر می‌کند، بازتاب رسانه‌ای آن چندین برابر
		یک بیانیه‌ی سیاسی معمول است.
		
		\item \textbf{کاهش فرسایش درون‌جنبشی:}
		بسیاری از مناقشات داخلی اپوزیسیون (مثلاً درباره‌ی تحریم‌ها)
		ناشی از نبودِ یک مرجع داوری معتبر است.
		داشتن یک نهاد مشورتی مستقل می‌تواند سطح تنش را کاهش دهد.
		
		\item \textbf{سپر دفاعی در برابر جنگ روانی:}
		هر توصیه‌ی نهاد، با مستندات و آراء مخالف منتشر می‌شود.
		این شفافیت، اصلی‌ترین سلاح در برابر تبلیغات حکومتی است.
	\end{enumerate}
	
	\subsection{دستاوردهای میان‌مدت (۱ تا ۳ سال)}
	
	\begin{enumerate}[nosep, label=\textcolor{cGold}{\alph*.}]
		\item \textbf{انباشت حقوقی:}
		مجموعه‌ای از «آراء مشورتی» شکل می‌گیرد که
		\keyword{رویه‌ی قضایی سایه} (\en{Shadow Case Law})
		را بنیان می‌گذارد — قابل استفاده در دوره‌ی گذار.
		
		\item \textbf{تقویت پرونده‌های بین‌المللی:}
		مستندات نهاد می‌تواند به‌عنوان
		\en{amicus curiae} یا مدرک تکمیلی
		در پرونده‌های صلاحیت جهانی مورد استفاده قرار گیرد.
		
		\item \textbf{جذب اعتماد بین‌المللی:}
		دولت‌ها و نهادهای بین‌المللی تمایل بیشتری به تعامل با
		اپوزیسیونی ساختارمند دارند. وجود یک نهاد شبه‌قضایی
		معتبر، \keyword{سرمایه‌ی اعتماد} (\en{trust capital})
		تولید می‌کند.
		
		\item \textbf{تربیت نسل جدید حقوقی:}
		مشارکت حقوق‌دانان جوان ایرانی در فرایندهای نهاد،
		بستر آموزشی ارزشمندی برای دوره‌ی پس از گذار فراهم می‌کند.
		
		\item \textbf{مدیریت منطقه‌ی خاکستری:}
		ایجاد چارچوبی برای برخورد با افراد و نهادهای «خاکستری»
		— کسانی که نه کاملاً حکومتی‌اند و نه کاملاً اپوزیسیون —
		از طریق معیارهای شفاف و قابل دفاع.
	\end{enumerate}
	
	\subsection{دستاوردهای بلندمدت (پس از گذار)}
	
	\begin{enumerate}[nosep, label=\textcolor{cGold}{\alph*.}]
		\item \textbf{تسهیل عدالت انتقالی:}
		مستندات و آراء نهاد، پایه‌ی ارزشمندی برای
		کمیسیون حقیقت‌یاب و دادگاه‌های انتقالی خواهد بود.
		
		\item \textbf{الگوی نهادسازی:}
		تجربه‌ی این نهاد می‌تواند الگویی برای ساختن
		نظام قضایی مستقل در ایران آینده باشد.
		
		\item \textbf{حافظه‌ی تاریخی:}
		آرشیو کامل آراء، مباحثات، و شواهد،
		یک \keyword{حافظه‌ی نهادی} (\en{institutional memory})
		بی‌بدیل برای نسل‌های آینده خواهد بود.
		
		\item \textbf{پیشگیری از عدالت انتقامی:}
		یکی از بزرگ‌ترین خطرات پس از سقوط هر نظام استبدادی،
		«عدالت خیابانی» و انتقام‌جویی بدون فرایند است.
		وجود سوابق دادرسی مستند، احتمال این خطر را کاهش می‌دهد.
	\end{enumerate}
	
	\section{شرایط کمک‌کننده برای موفقیت}
	
	\begin{principleBox}[پیش‌شرط‌های ضروری]
		\begin{enumerate}[nosep, label=\textcolor{cPrimary}{\Roman*.}]
			\item \textbf{استقلال واقعی:}
			نهاد نباید وابسته به هیچ حزب، جناح سیاسی، یا دولت خارجی باشد.
			بودجه باید از منابع متنوع و شفاف تأمین شود.
			
			\item \textbf{تخصص بالا:}
			اعضای هیئت‌داوران باید دارای سوابق آکادمیک و حرفه‌ای
			قابل تأیید در حوزه‌ی حقوق بین‌الملل، حقوق بشر، یا
			حقوق کیفری باشند.
			
			\item \textbf{شفافیت بی‌قیدوشرط:}
			تمام آراء، مذاکرات، منابع مالی، و تعارض منافع
			باید به‌صورت عمومی قابل دسترسی باشند.
			
			\item \textbf{تکثرگرایی:}
			ترکیب نهاد باید بازتاب‌دهنده‌ی تنوع قومی، جنسیتی،
			نسلی، و فکری ایران باشد.
			
			\item \textbf{فروتنی نهادی:}
			نهاد باید همواره تأکید کند که مرجع مشورتی است و
			جایگزین نظام قضایی رسمی آینده نیست.
			
			\item \textbf{همکاری بین‌المللی:}
			ارتباط فعال با نهادهای بین‌المللی حقوقی و حقوق بشری.
		\end{enumerate}
	\end{principleBox}
	
	\section{سناریوهای مثبت}
	
	برای ملموس‌تر کردن دستاوردها، سه سناریوی فرضی ارائه می‌شود:
	
	\begin{noteBox}[سناریوی ۱: تحریم ورزشی]
		فرض کنید جریانی در اپوزیسیون خواستار تعلیق عضویت جمهوری اسلامی
		در فیفا شود.
		
		\medskip
		\textbf{بدون دادگاه در سایه:}
		بیانیه‌ای سیاسی منتشر می‌شود. حکومت فوراً ورزشکاران آسیب‌دیده را
		جلوی دوربین می‌آورد. افکار عمومی دوپاره می‌شود.
		برخی از چهره‌های اپوزیسیون مجبور به عقب‌نشینی می‌شوند.
		
		\medskip
		\textbf{با دادگاه در سایه:}
		درخواست به نهاد ارجاع می‌شود. هیئت داوران بررسی می‌کند:
		\begin{itemize}[nosep]
			\item آیا تعلیق ورزشی بر اساس سابقه‌ی فیفا
			(تعلیق آفریقای جنوبی، روسیه) توجیه‌پذیر است؟
			\item هزینه‌ی ورزشکاران بی‌تقصیر چقدر است؟
			\item آیا جایگزین‌های کم‌هزینه‌تر (مثلاً تحریم مقام‌های
			ورزشی حکومتی بدون تعلیق کامل) بررسی شده‌اند؟
			\item نظرات موافق و مخالف اعضا چیست؟
		\end{itemize}
		
		\medskip
		نتیجه: رأی مستدل صادر می‌شود — مثلاً «تعلیق کامل در حال حاضر
		نامتناسب است، اما تحریم هدفمند مقام‌های ورزشی توصیه می‌شود.»
		این رأی هم قابل دفاع است و هم فضای مانور تبلیغاتی حکومت را
		به حداقل می‌رساند.
	\end{noteBox}
	
	\begin{noteBox}[سناریوی ۲: اخراج وابستگان مقام‌ها از کشور ثالث]
		فرض کنید کمپینی خواستار اخراج فرزندان یک مقام ارشد سپاه
		از کانادا شود.
		
		\medskip
		\textbf{بدون دادگاه در سایه:}
		کمپین رسانه‌ای پُرسروصدا اما بدون مبنای حقوقی روشن.
		حکومت فوراً روایت «ظلم به خانواده‌های بی‌گناه» را تولید می‌کند.
		دولت کانادا از نبودِ فرایند حقوقی شاکی می‌شود.
		
		\medskip
		\textbf{با دادگاه در سایه:}
		\begin{itemize}[nosep]
			\item هیئت داوران ابتدا \keyword{نقش فردی} هر شخص را بررسی می‌کند:
			آیا خودِ فرزندان در فعالیت‌های حکومتی نقش دارند؟
			آیا از ثروت‌های حاصل از فساد بهره‌مند شده‌اند؟
			\item بر اساس اصل \keyword{مسئولیت فردی}
			تفکیک می‌کند بین کسانی که ذی‌نفع فعال هستند و
			کسانی که صرفاً رابطه‌ی خویشاوندی دارند.
			\item رأی صادره ممکن است شامل توصیه‌های متفاوت برای
			اعضای مختلف یک خانواده باشد.
			\item مبنای حقوقی روشن (مثلاً قانون مگنیتسکی)
			در رأی ذکر می‌شود.
		\end{itemize}
		
		\medskip
		نتیجه: دولت کانادا یک سند حقوقی معتبر دریافت می‌کند،
		نه صرفاً یک درخواست سیاسی.
	\end{noteBox}
	
	\begin{noteBox}[سناریوی ۳: مداخله‌ی بشردوستانه]
		فرض کنید در پی یک سرکوب خونین گسترده، بخشی از اپوزیسیون
		خواستار مداخله‌ی نظامی بین‌المللی شود.
		
		\medskip
		\textbf{بدون دادگاه در سایه:}
		بحث شدید درون‌جنبشی. اتهام «وطن‌فروشی» از یک سو و
		اتهام «همدستی با حکومت» از سوی دیگر. جمهوری اسلامی
		حداکثر بهره‌برداری تبلیغاتی را می‌کند.
		
		\medskip
		\textbf{با دادگاه در سایه:}
		\begin{itemize}[nosep]
			\item هیئت ویژه‌ای بر اساس چارچوب \en{R2P}
			(سند ۲۰۰۵ اجلاس جهانی سازمان ملل) بررسی می‌کند:
			\item آیا چهار شرط مداخله محقق شده‌اند؟
			\begin{enumerate}[nosep, label=\textcolor{cAccent}{\roman*.}]
				\item نسل‌کشی یا پاکسازی قومی در حال وقوع یا قریب‌الوقوع است؟
				\item آیا دولت از محافظت از شهروندان ناتوان یا ناخواسته است؟
				\item آیا تمام ابزارهای مسالمت‌آمیز امتحان شده‌اند؟
				\item آیا مداخله‌ی نظامی بدتر از وضع موجود نخواهد بود؟
			\end{enumerate}
			\item رأی مستدل با ذکر دقیق شواهد و تحلیل هر شرط صادر می‌شود.
			\item نظرات اقلیت به‌صورت کامل ثبت و منتشر می‌شوند.
			\item حتی اگر رأی مثبت باشد، \keyword{قیود و شروط}
			(نوع مداخله، محدوده، مدت، تضمین‌های پس از مداخله)
			مشخص می‌شوند.
		\end{itemize}
		
		\medskip
		نتیجه: جنبش به‌جای یک موضع‌گیری احساسی، یک سند تحلیلی
		چندلایه‌ای دارد که هم برای نخبگان ایرانی و هم برای
		مخاطبان بین‌المللی قابل فهم و دفاع است.
	\end{noteBox}
	
	\section{جمع‌بندی فصل دوم}
	
	\begin{principleBox}[خلاصه‌ی فصل ۲]
		دادگاه در سایه صرفاً یک نهاد حقوقی نیست؛ یک
		\keyword{زیرساخت استراتژیک} برای جنبش گذارخواهی است که:
		\begin{itemize}[nosep]
			\item اعتبار تولید می‌کند؛
			\item تنش درونی را مدیریت می‌کند؛
			\item سپر دفاعی در برابر جنگ رسانه‌ای فراهم می‌کند؛
			\item حافظه‌ی نهادی برای عدالت انتقالی آینده می‌سازد.
		\end{itemize}
		موفقیت آن مشروط به استقلال، تخصص، شفافیت، و فروتنی نهادی است.
	\end{principleBox}
	
	\clearpage
	
	
	% ////////////////////////////////////////////////////////////
	%  فصل ۳: آسیب‌شناسی و احتیاط‌های لازم
	% ////////////////////////////////////////////////////////////
	\chapter{آسیب‌شناسی و احتیاط‌های لازم}\label{chap:risk}
	
	\begin{warningBox}[تذکر مقدماتی]
		هیچ نهادی از خطر مصون نیست. صداقت اقتضا می‌کند که پیش از
		ترویج ایده، \keyword{بدترین سناریوها} را نیز بررسی کنیم.
		تنها با شناخت دقیق آسیب‌ها می‌توان سازوکارهای پیش‌گیری طراحی کرد.
	\end{warningBox}
	
	\section{آسیب‌شناسی ساختاری}
	
	\subsection{خطر ۱: مصادره‌ی سیاسی}
	
	\textbf{توصیف:} یک جناح سیاسی تلاش کند نهاد را ابزار مشروعیت‌بخشی
	به مواضع خود کند — مثلاً جناح سلطنت‌طلب یا جناح جمهوری‌خواه
	اکثریت هیئت داوران را به خود اختصاص دهد.
	
	\begin{principleBox}[سازوکار پیش‌گیری]
		\begin{itemize}[nosep]
			\item سیستم نامزدی و تأیید دومرحله‌ای:
			نامزدها توسط \keyword{کمیته‌ی گزینش مستقل} پیشنهاد و توسط
			\keyword{شورای ناظران بیرونی} (شامل حقوق‌دانان بین‌المللی
			غیرایرانی) تأیید شوند.
			\item ممنوعیت عضویت هم‌زمان در هر حزب یا سازمان سیاسی.
			\item دوره‌ی خدمت محدود (مثلاً ۳ سال بدون تمدید فوری).
			\item افشای کامل تعارض منافع قبل از هر پرونده.
		\end{itemize}
	\end{principleBox}
	
	\subsection{خطر ۲: بی‌اعتباری ناشی از فقدان ضمانت اجرا}
	
	\textbf{توصیف:} چون نهاد «سایه» است و قدرت اجرایی ندارد،
	ممکن است به‌سرعت به نهادی کاغذی تبدیل شود که کسی به
	آراء آن وقعی نمی‌گذارد.
	
	\begin{principleBox}[سازوکار پیش‌گیری]
		\begin{itemize}[nosep]
			\item از ابتدا صادقانه اعلام شود که نهاد \keyword{مشورتی} است.
			ارزش آن نه در ضمانت اجرا بلکه در
			\keyword{اعتبار اخلاقی و فنی} آراء آن است.
			\item تمرکز بر کیفیت آراء: هر رأی باید آن‌قدر مستدل و مستند
			باشد که رسانه‌ها، نهادهای بین‌المللی، و فعالان مدنی
			خودبه‌خود به آن ارجاع دهند.
			\item ایجاد مکانیزم «پذیرش داوطلبانه»
			(\en{voluntary compliance}):
			جریان‌های سیاسی از پیش متعهد شوند که آراء نهاد را
			به‌عنوان مرجع مشورتی بپذیرند.
			\item همکاری با رسانه‌های معتبر برای پوشش گسترده‌ی آراء.
		\end{itemize}
	\end{principleBox}
	
	\subsection{خطر ۳: نفوذ جمهوری اسلامی}
	
	\textbf{توصیف:} حکومت تلاش خواهد کرد عوامل خود را در ساختار
	نهاد نفوذ دهد — چه به‌صورت مستقیم و چه از طریق افراد واسط.
	
	\begin{principleBox}[سازوکار پیش‌گیری]
		\begin{itemize}[nosep]
			\item \keyword{بررسی امنیتی چندلایه}
			(\en{multi-layer vetting})
			برای تمام اعضا و کارکنان.
			\item استفاده از \keyword{امنیت عملیاتی}
			(\en{OPSEC}) در ارتباطات.
			\item ممنوعیت سفر به ایران یا تماس با نهادهای حکومتی
			برای اعضای فعال (مگر با مجوز ویژه).
			\item سیستم \keyword{افشاگری محافظت‌شده}
			(\en{protected whistleblowing})
			برای گزارش موارد مشکوک.
			\item بررسی دوره‌ای مستقل توسط ناظران بیرونی.
		\end{itemize}
	\end{principleBox}
	
	\subsection{خطر ۴: تبدیل شدن به ابزار فشار بر بی‌گناهان}
	
	\textbf{توصیف:} ممکن است نهاد به‌جای هدف قرار دادن ناقضان واقعی
	حقوق بشر، ناخواسته به ابزاری برای تسویه‌حساب‌های شخصی یا
	فشار بر افراد بی‌تقصیر تبدیل شود.
	
	\begin{principleBox}[سازوکار پیش‌گیری]
		\begin{itemize}[nosep]
			\item \keyword{آزمون تناسب اجباری}
			(\en{mandatory proportionality test})
			در هر پرونده: آیا هزینه‌ی وارد بر افراد ثالث
			بی‌تقصیر از فایده‌ی عمومی کمتر است؟
			\item حق دفاع برای هر شخص نام‌بُرده — حتی غیابی:
			وکیل تسخیری (\en{amicus curiae})
			موظف است بهترین استدلال ممکن را به نفع طرف مقابل ارائه دهد.
			\item مرحله‌ی تجدیدنظر اجباری برای تمام آرائی که
			مستقیماً بر شخص حقیقی اثر می‌گذارند.
			\item \keyword{حق اعتراض عمومی:}
			هر شهروند بتواند نسبت به رأی صادره اعتراض مستند ارائه دهد.
		\end{itemize}
	\end{principleBox}
	
	\subsection{خطر ۵: انتظارات نامعقول}
	
	\textbf{توصیف:} افکار عمومی ممکن است از نهاد انتظار عملکردی
	مانند یک دادگاه واقعی داشته باشد و وقتی آراء اجرا نمی‌شوند
	سرخورده شود.
	
	\begin{principleBox}[سازوکار پیش‌گیری]
		\begin{itemize}[nosep]
			\item \keyword{کمپین ارتباطی شفاف} از ابتدا: توضیح دقیق
			ماهیت مشورتی نهاد و تفاوت آن با دادگاه‌های رسمی.
			\item استفاده از عنوان‌هایی مانند «رأی مشورتی» و «توصیه‌نامه»
			به‌جای «حکم» و «رأی قطعی».
			\item گزارش‌دهی دوره‌ای درباره‌ی تأثیرات واقعی آراء صادره
			— چه مثبت و چه محدود.
		\end{itemize}
	\end{principleBox}
	
	\section{آسیب‌شناسی اخلاقی}
	
	\subsection{خطر ۶: ادعای مشروعیتی که وجود ندارد}
	
	\textbf{توصیف:} نهاد ممکن است به‌مرور ادعای مشروعیتی فراتر
	از ظرفیت خود داشته باشد — مثلاً ادعای «نمایندگی مردم ایران».
	
	\begin{warningBox}[قاعده‌ی سخت]
		نهاد \textbf{هرگز} نباید ادعای نمایندگی مردم ایران کند.
		مشروعیت آن صرفاً ناشی از:
		\begin{itemize}[nosep]
			\item تخصص اعضا؛
			\item شفافیت فرایند؛
			\item کیفیت استدلال؛
			\item و پذیرش داوطلبانه‌ی طرف‌ها
		\end{itemize}
		است. نمایندگی مردم تنها از طریق انتخابات آزاد محقق می‌شود.
	\end{warningBox}
	
	\subsection{خطر ۷: خطای قضاوت}
	
	\textbf{توصیف:} حتی بهترین قضات اشتباه می‌کنند. در شرایط
	انقلابی و فقدان اطلاعات کامل، احتمال خطا بیشتر نیز هست.
	
	\begin{principleBox}[سازوکار پیش‌گیری]
		\begin{itemize}[nosep]
			\item \keyword{اصل فروتنی معرفتی}
			(\en{epistemic humility}):
			هر رأی باید به‌صراحت محدودیت‌های اطلاعاتی خود را ذکر کند.
			\item \keyword{بازنگری دوره‌ای:}
			آراء صادره هر ۱۲ ماه یکبار بازبینی شوند.
			اگر اطلاعات جدیدی در دست باشد، رأی اصلاح شود.
			\item \keyword{ثبت نظرات اقلیت:}
			نظر مخالف اعضای اقلیت به‌صورت کامل ثبت شود
			تا اگر در آینده درست از آب درآمد، قابل بازیابی باشد.
			\item \keyword{عذرخواهی نهادی:}
			اگر خطایی اثبات شد، نهاد رسماً عذرخواهی و اصلاح کند.
		\end{itemize}
	\end{principleBox}
	
	\section{ماتریس ریسک — خلاصه}
	
	% ── اصلاح‌شده: \cellcolor اکنون با colortbl کار می‌کند ──
	\begin{longtable}{r p{3.5cm} c c p{4cm}}
		\toprule
		\textbf{ردیف} & \textbf{خطر} & \textbf{احتمال} & \textbf{شدت} &
		\textbf{اولویت پیش‌گیری} \\
		\midrule
		\endhead
		۱ & مصادره‌ی سیاسی & بالا & بالا & \cellcolor{cAccent!15} بحرانی \\
		۲ & بی‌اعتباری & متوسط & بالا & \cellcolor{cAccent!15} بحرانی \\
		۳ & نفوذ حکومت & متوسط & بالا & \cellcolor{cAccent!15} بحرانی \\
		۴ & فشار بر بی‌گناهان & متوسط & بالا & \cellcolor{cAccent!15} بحرانی \\
		۵ & انتظارات نامعقول & بالا & متوسط & \cellcolor{cGold!20} مهم \\
		۶ & ادعای مشروعیت کاذب & پایین & بالا & \cellcolor{cGold!20} مهم \\
		۷ & خطای قضاوت & بالا & متوسط & \cellcolor{cGold!20} مهم \\
		\bottomrule
	\end{longtable}
	
	\section{جمع‌بندی فصل سوم}
	
	\begin{warningBox}[پیام اصلی]
		دادگاه در سایه یک ابزار قدرتمند است و مانند هر ابزار قدرتمند،
		می‌تواند هم سازنده و هم مخرب باشد. تفاوت را
		\keyword{طراحی نهادی دقیق} و \keyword{فرهنگ فروتنی}
		رقم می‌زند. این فصل نه برای منصرف‌کردن از ایده، بلکه
		برای تضمین طراحی مسئولانه نوشته شده است.
	\end{warningBox}
	
	\clearpage
	
	
	% ////////////////////////////////////////////////////////////

	\addcontentsline{toc}{chapter}{خلاصه‌ی اجرایی}
	\thispagestyle{fancy}
	
	\begin{principleBox}[گزاره‌ی محوری]
	جنبش گذارخواهی ایران در مقطع کنونی با \keyword{خلأ مشروعیت‌بخشی حقوقی}
	مواجه است: تصمیم‌هایی با ابعاد قضایی — از حمایت از تحریم‌ها تا مستندسازی
	جرایم حکومتی — بدون هیچ \keyword{فرایند دادرسی ساختارمند} اتخاذ می‌شوند.
	نبودِ چنین فرایندی هم اپوزیسیون را در برابر جنگ روانی جمهوری اسلامی آسیب‌پذیر
	می‌کند و هم ظرفیت گفت‌وگوی سازنده‌ی درون‌جنبشی را تضعیف می‌نماید.
	
	\medskip
	\textbf{پیشنهاد:} تأسیس یک \keyword{نهاد شبه‌قضایی سایه}
	(\en{Shadow Quasi-Judicial Body})
	با عنوان \textbf{«دادگاه در سایه — وجدان منفصل مسئول»}
	که وظایف زیر را بر عهده گیرد:
	
	\begin{enumerate}[nosep,label=\textcolor{cPrimary}{\arabic*.}]
	  \item مستندسازی نظام‌مند جرایم رژیم بر پایه‌ی معیارهای حقوق بین‌الملل؛
	  \item داوری و تنظیم‌بخشی تصمیم‌های حساس جنبش بر اساس «پاسخ مشروع و متناسب»؛
	  \item فراهم‌سازی بستر گفت‌وگوی نخبگانی برای مسائل پرمناقشه
	        (تحریم، مداخله‌ی بشردوستانه، اخراج یا تعقیب وابستگان حکومتی و …)؛
	  \item ثبت و بایگانی آراء و دلایل برای بازنگری پس از گذار.
	\end{enumerate}
	\end{principleBox}
	
	\section*{چرا «دادگاه در سایه»؟}
	
	جنبش‌های گذارخواه در سراسر تاریخ معاصر — از
	\en{Free French} تا \en{ANC} آفریقای جنوبی —
	با ضرورت اتخاذ تصمیم‌هایی مواجه بوده‌اند که ماهیت قضایی دارند
	اما پیش از دست‌یابی به قدرت سیاسی صورت می‌گیرند.
	وقتی یک جنبش از تحریم فلان نهاد ورزشی حمایت می‌کند، عملاً
	\keyword{حکمی صادر کرده} که بر زندگی حرفه‌ای هزاران ورزشکار اثر می‌گذارد.
	وقتی خواستار اخراج فرزندان مقام‌های حکومتی از کشور ثالث می‌شود،
	\keyword{مجازاتی اعمال کرده} که بدون دادرسی است.
	
	\medskip
	نبودِ دادرسی به معنای نبودِ هر سه عنصر زیر است:
	
	\begin{itemize}[nosep]
	  \item \textbf{مبنای حقوقی:} چرا این اقدام مجاز است؟
	  \item \textbf{تناسب:} آیا هزینه‌ی وارد بر بی‌گناهان، از فایده‌ی مورد انتظار کمتر است؟
	  \item \textbf{مستندسازی:} آیا استدلال‌ها و مخالفت‌ها ثبت شده‌اند
	        تا در آینده قابل بازنگری باشند؟
	\end{itemize}
	
	\begin{warningBox}[خلأ موجود و بهره‌برداری جمهوری اسلامی]
	جمهوری اسلامی به‌صورت سازمان‌یافته از \keyword{نبودِ فرایند}
	بهره‌برداری تبلیغاتی می‌کند:
	\begin{itemize}[nosep]
	  \item «اپوزیسیون بدون هیچ محاکمه‌ای، مردم را تحریم می‌کند.»
	  \item «گذارخواهان از جنگ و مداخله‌ی خارجی استقبال می‌کنند.»
	  \item «آن‌ها حتی به ورزشکاران و هنرمندان رحم نمی‌کنند.»
	\end{itemize}
	
	\medskip
	وجود یک نهاد معتبر شبه‌قضایی با آراء مستدل و مستند و ثبت نظرات موافق و مخالف،
	\textbf{قوی‌ترین دفاعیه} در برابر این جنگ رسانه‌ای است.
	\end{warningBox}
	
	\section*{مخاطبان اصلی سند}
	
	\begin{enumerate}[nosep,label=\textcolor{cPrimary}{\Alph*.}]
	  \item \textbf{فعالان سیاسی گذارخواه و براندازانه‌ی ایرانی}
	        — برای درک ضرورت ساختارمندسازی؛
	  \item \textbf{حقوق‌دانان و پژوهشگران ایرانی و بین‌المللی}
	        — برای مشارکت در طراحی؛
	  \item \textbf{نهادهای بین‌المللی حقوق بشر}
	        — برای همکاری فنی و اعتباربخشی؛
	  \item \textbf{افکار عمومی ایرانی}
	        — برای اطمینان از وجود نظارت اخلاقی بر تصمیم‌های سخت.
	\end{enumerate}
	
	\section*{ساختار سند}
	
	\begin{longtable}{r r p{10cm}}
	\toprule
	\textbf{فصل} & \textbf{شماره} & \textbf{موضوع} \\
	\midrule
	\endhead
	خلاصه‌ی اجرایی & --- & همین بخش \\
	فصل ۱ & \ref{chap:problem} & تشریح مسئله و جدیت نیاز \\
	فصل ۲ & \ref{chap:expect} & انتظارات مثبت و شرایط کمک‌کننده \\
	فصل ۳ & \ref{chap:risk} & آسیب‌شناسی و احتیاط‌های لازم \\
	فصل ۴ & \ref{chap:org} & دغدغه‌های سازماندهی و تکثر دیدگاهی \\
	فصل ۵ & \ref{chap:proc} & رویه‌های اجرایی و قضایی معتبر \\
	فصل ۶ & \ref{chap:finance} & مباحث مالی و هزینه‌ای \\
	فصل ۷ & \ref{chap:intl} & ارتباط با مراجع بین‌المللی \\
	فصل ۸ & \ref{chap:cases} & نمونه‌های مشابه و مطالعه‌ی موردی \\
	فصل ۹ & \ref{chap:transition} & هضم و ادغام در نظام قضایی پس از گذار \\
	فصل ۱۰ & \ref{chap:persons} & افراد پیشنهادی \\
	فصل ۱۱ & \ref{chap:conclusion} & جمع‌بندی و نتیجه‌گیری \\
	\bottomrule
	\end{longtable}
	
	\section*{اصول حاکم بر کل سند}
	
	\begin{principleBox}[هفت اصل بنیادین]
	\begin{enumerate}[nosep, label=\textcolor{cGold}{\Roman*.}]
	  \item \textbf{دعوت به اندیشیدن، نه تحمیل حکم:}
	        این دادگاه مرجع مشورتی و وجدانی است، نه مقام حاکمیتی.
	  \item \textbf{مجوزدهی دیرانه و بااحتیاط:}
	        هرگونه توصیه‌ی قهرآمیز تنها پس از بررسی همه‌جانبه صادر می‌شود.
	  \item \textbf{شفافیت مستند:}
	        آراء، نظرات اقلیت، و دلایل مخالف به‌صورت کامل ثبت و منتشر می‌شوند.
	  \item \textbf{تناسب و ضرورت:}
	        هر اقدامِ توصیه‌شده باید هر دو آزمون تناسب
	        (\en{proportionality}) و ضرورت
	        (\en{necessity}) را بگذراند.
	  \item \textbf{احترام به کرامت انسانی:}
	        حتی وابستگان حکومت از حقوق بنیادین برخوردارند؛
	        هیچ توصیه‌ای نباید ناقض اعلامیه‌ی جهانی حقوق بشر باشد.
	  \item \textbf{استقلال از جناح‌بندی‌های سیاسی:}
	        اعضا به‌صورت فردی و بر پایه‌ی تخصص خدمت می‌کنند.
	  \item \textbf{موقتی بودن:}
	        این نهاد با تأسیس نظام قضایی رسمی پس از گذار منحل یا ادغام می‌شود.
	\end{enumerate}
	\end{principleBox}
	
	% پایان خلاصه‌ی اجرایی — حدوداً ۳ صفحه
	\clearpage
	
	% ////////////////////////////////////////////////////////////
	%  فصل ۱: تشریح مسئله و جدیت نیاز
	% ////////////////////////////////////////////////////////////
	\mainmatter



% ////////////////////////////////////////////////////////////
%  فصل ۴: دغدغه‌های سازماندهی و تکثر دیدگاهی
% ////////////////////////////////////////////////////////////
	\chapter{دغدغه‌های سازماندهی و تکثر دیدگاهی}\label{chap:org}
	
	\section{ساختار پیشنهادی نهاد}
	
	ساختار نهاد باید هم‌زمان \keyword{چابک} و \keyword{موجه} باشد.
	پیشنهاد زیر مبتنی بر مطالعه‌ی تطبیقی نهادهای مشابه
	(فصل \ref{chap:cases}) طراحی شده است.
	
		% ── نمودار ساختار سازمانی — اصلاح‌شده ──

\begin{figure}[H]
	\centering
	\begin{tikzpicture}[
		% ← FIX: فاصله‌ی عمودی و افقی کاهش یافت
		node distance=1.6cm and 2cm,
		box/.style={
			rectangle, rounded corners=3pt, draw=cPrimary, thick,
			fill=cLight,
			text width=4.5cm,           % ← کمی باریک‌تر
			minimum height=1.2cm,
			align=center, font=\small
		},
		smallbox/.style={
			rectangle, rounded corners=2pt, draw=cGray,
			fill=white,
			text width=3.2cm,           % ← کمی باریک‌تر
			minimum height=0.9cm,
			align=center, font=\footnotesize
		},
		arrow/.style={-{Stealth}, thick, cPrimary}
		]
		
		% ── سطح ۱: شورای ناظران ──
		\node[box, fill=cGold!10, draw=cGold] (supervisory)
		{{\RL{\textbf{شورای ناظران بیرونی}}}\\[2pt]
			{\footnotesize\RL{(۵ تا ۷ نفر — اکثراً غیرایرانی)}}};
		
		% ── سطح ۲: هیئت عالی ──
		\node[box, below=of supervisory] (high)
		{{\RL{\textbf{هیئت عالی داوران}}}\\[2pt]
			{\footnotesize\RL{(۹ تا ۱۵ نفر — ایرانی و غیرایرانی)}}};
		
		% ── سطح ۳: شعبات تخصصی ──
		% ← FIX: فاصله‌ی افقی از 1.5cm به 0.8cm کاهش یافت
		\node[smallbox, below left=1.8cm and 0.8cm of high] (branch1)
		{{\RL{شعبه‌ی مستندسازی}}\\{\RL{جرایم}}};
		\node[smallbox, below=1.8cm of high] (branch2)
		{{\RL{شعبه‌ی داوری}}\\{\RL{تصمیم‌های سیاسی}}};
		\node[smallbox, below right=1.8cm and 0.8cm of high] (branch3)
		{{\RL{شعبه‌ی گفتمان}}\\{\RL{و مرزگذاری}}};
		
		% ── سطح ۴: دبیرخانه ──
		% ← FIX: فاصله‌ی عمودی از 5cm به 4.5cm کاهش یافت
		\node[box, below=4.5cm of high] (secretariat)
		{{\RL{\textbf{دبیرخانه‌ی اجرایی}}}\\[2pt]
			{\footnotesize\RL{(مدیر اجرایی + تیم پژوهش و ارتباطات)}}};
		
		% ── پیکان‌ها ──
		\draw[arrow, cGold] (supervisory) -- (high)
		node[midway, left, font=\footnotesize] {\RL{نظارت}};
		\draw[arrow] (high) -- (branch1);
		\draw[arrow] (high) -- (branch2);
		\draw[arrow] (high) -- (branch3);
		\draw[arrow] (branch1) -- (secretariat);
		\draw[arrow] (branch2) -- (secretariat);
		\draw[arrow] (branch3) -- (secretariat);
		
		% ── حاشیه — مشاوران ویژه ──
		% ← FIX: از left=4cm به left=2.2cm تغییر یافت
		\node[smallbox, left=2.2cm of high, draw=cGold!60, fill=cGold!5]
		(advisors)
		{{\RL{مشاوران ویژه}}\\[2pt]
			{\footnotesize\RL{(موردی — حسب نیاز)}}};
		\draw[-{Stealth}, dashed, cGold!60] (advisors) -- (high);
		
		% ── حاشیه — وکیل تسخیری ──
		% ← FIX: از right=4cm به right=2.2cm تغییر یافت
		\node[smallbox, right=2.2cm of high, draw=cAccent!60, fill=cAccent!5]
		(amicus)
		{{\RL{دفتر وکیل تسخیری}}\\[2pt]
			{\footnotesize\RL{(دوست دادگاه)}}};
		\draw[-{Stealth}, dashed, cAccent!60] (amicus) -- (high);
		
	\end{tikzpicture}
	\caption{نمودار ساختار سازمانی پیشنهادی دادگاه در سایه}
\end{figure}
	
	\subsection{شورای ناظران بیرونی}
		
			\begin{principleBox}[مشخصات]
			\begin{itemize}[nosep]
			\item \textbf{ترکیب:} ۵ تا ۷ نفر، اکثراً حقوق‌دانان و قضات
			بازنشسته‌ی بین‌المللی (ترجیحاً غیرایرانی برای تضمین استقلال).
				\item \textbf{وظایف:}
				\begin{itemize}[nosep]
				\item نظارت بر فرایند انتخاب اعضای هیئت عالی؛
				\item بررسی سالانه‌ی عملکرد نهاد؛
			\item رسیدگی به شکایات مربوط به تعارض منافع یا سوءرفتار؛
			\item تأیید نهایی اصلاحات آیین‌نامه‌ای.
			\end{itemize}
		\item \textbf{دوره‌ی خدمت:} ۴ سال، حداکثر یک بار تمدید.
	\item \textbf{حق‌الزحمه:} نمادین (برای تأکید بر استقلال).
	\end{itemize}
	\end{principleBox}
	
	\subsection{هیئت عالی داوران}
		
			\begin{principleBox}[مشخصات]
			\begin{itemize}[nosep]
				\item \textbf{ترکیب:} ۹ تا ۱۵ نفر شامل:
				\begin{itemize}[nosep]
				\item حداقل ۳ متخصص حقوق بین‌الملل کیفری؛
				\item حداقل ۲ متخصص حقوق بشر؛
				\item حداقل ۱ متخصص علوم سیاسی / روابط بین‌الملل؛
				\item حداقل ۱ جامعه‌شناس یا روان‌شناس اجتماعی؛
				\item حداقل ۱ نماینده‌ی جامعه‌ی مدنی داخل ایران
			(با رعایت امنیت)؛
			\item حداقل ۲ عضو غیرایرانی.
			\end{itemize}
				\item \textbf{شرایط عضویت:}
				\begin{itemize}[nosep]
				\item مدرک دکتری یا معادل حرفه‌ای در حوزه‌ی مرتبط؛
				\item حداقل ۱۰ سال سابقه‌ی حرفه‌ای یا آکادمیک؛
			\item عدم عضویت در هیچ حزب سیاسی فعال؛
			\item تأیید شورای ناظران بیرونی.
			\end{itemize}
		\item \textbf{دوره‌ی خدمت:} ۳ سال، بدون تمدید فوری
	(امکان بازگشت پس از ۲ سال وقفه).
	\end{itemize}
	\end{principleBox}
	\newpage
	\subsection{شعبات سه‌گانه‌ی تخصصی}
		
		\begin{longtable}{r p{4cm} p{8cm}}
		\toprule
		\textbf{شعبه} & \textbf{حوزه‌ی صلاحیت} & \textbf{نمونه موضوعات} \\
		\midrule
		\endhead
		\textbf{مستندسازی جرایم} & اعلام جرم، جمع‌آوری شواهد،
		طبقه‌بندی بر اساس اساسنامه‌ی رم &
		کشتار آبان ۹۸، شکنجه‌ی زندانیان سیاسی، اعدام‌های مواد مخدر،
		سرکوب اقلیت‌های قومی و مذهبی \\
		\midrule
		\textbf{داوری تصمیم‌های سیاسی} & بررسی تناسب و مشروعیت
		اقدامات پیشنهادی جنبش &
		تحریم‌های اقتصادی، ورزشی، فرهنگی؛ اخراج وابستگان؛
		فراخوان اعتصاب \\
		\midrule
		\textbf{گفتمان و مرزگذاری} & تحلیل مسائل اختلاف‌آفرین
		و ارائه‌ی چارچوب گفتمانی &
		مداخله‌ی بشردوستانه، رابطه با دولت‌های خارجی، جایگاه
	«خاکستری‌ها»، عفو مشروط \\
	\bottomrule
	\end{longtable}
	
	\subsection{دفتر وکیل تسخیری (دوست دادگاه)}
		
		\begin{principleBox}[یکی از مهم‌ترین نوآوری‌های طرح]
		در هر پرونده‌ای که نام شخص حقیقی یا حقوقی مطرح شود،
			\keyword{دفتر وکیل تسخیری} موظف است:
			\begin{enumerate}[nosep, label=\textcolor{cPrimary}{\arabic*.}]
			\item بهترین استدلال ممکن را \textbf{به نفع} طرف مقابل ارائه دهد؛
		\item حقوق بنیادین طرف مقابل را یادآوری کند؛
		\item هرگونه ضعف در شواهد دادستان را برجسته کند.
		\end{enumerate}
		
		\medskip
	وجود این دفتر مهم‌ترین تضمین عدالت فرایندی و مؤثرترین
	سپر در برابر اتهام «دادگاه نمایشی» است.
	\end{principleBox}
	
	\section{مدیریت تکثر دیدگاهی}
	
	اپوزیسیون ایران طیف بسیار گسترده‌ای را شامل می‌شود:
	سلطنت‌طلبان، جمهوری‌خواهان، فدرالیست‌ها، چپ‌ها، لیبرال‌ها،
	ملی‌ـ‌مذهبی‌ها، و گروه‌های قومی‌ـ‌زبانی متنوع.
	مدیریت این تکثر یکی از بزرگ‌ترین چالش‌های نهاد است.
	
	\subsection{اصل بنیادین: تکثر در ورودی، وحدت در فرایند}
		
			\begin{principleBox}[توضیح اصل]
			\begin{itemize}[nosep]
			\item \textbf{تکثر در ورودی:}
			هر جریان فکری حق دارد نماینده‌ای در فرایند ارائه‌ی
			دیدگاه داشته باشد. جلسات استماع باید باز و فراگیر باشند.
			\item \textbf{وحدت در فرایند:}
			اما فرایند دادرسی یکسان، معیارهای حقوقی یکسان، و
			استانداردهای شواهد یکسان بر همه اعمال می‌شود.
		\item \textbf{تکثر در خروجی:}
	آراء اکثریت و اقلیت هر دو ثبت و منتشر می‌شوند.
	\end{itemize}
	\end{principleBox}
	
	\subsection{سازوکارهای عملی مدیریت تکثر}
		
		\begin{enumerate}[nosep, label=\textcolor{cPrimary}{\arabic*.}]
		\item \textbf{جلسات استماع عمومی}
		(\en{Public Hearings}):
		قبل از صدور هر رأی مهم، جلسه‌ی استماع برگزار شود
		که نمایندگان دیدگاه‌های مختلف بتوانند شفاهاً و
		کتباً نظر دهند.
		
		\item \textbf{لایحه‌ی دوست دادگاه}
		(\en{Amicus Brief}):
		هر شخص حقیقی یا حقوقی بتواند لایحه‌ی مکتوب
		به نهاد تقدیم کند.
		
		\item \textbf{کمیته‌های مشورتی قومی‌ـ‌زبانی:}
		کمیته‌هایی از نمایندگان کُرد، بلوچ، ترک‌آذری،
		عرب، ترکمن، لُر و سایر اقوام تشکیل شود
		تا در موضوعات مربوط به حقوق اقلیت‌ها مشاوره دهند.
		
		\item \textbf{نشست‌های مشورتی نسلی:}
		ترکیب هیئت باید شامل نمایندگان نسل‌های مختلف باشد —
		به‌ویژه نسل «ژینا» که بار اصلی جنبش را بر دوش دارد.
		
		\item \textbf{مکانیزم اعتراض ساختارمند:}
		هر رأی صادره تا ۳۰ روز پس از انتشار قابل اعتراض
	مستند باشد. هیئت موظف است اعتراضات مستدل را بررسی
	و در صورت لزوم رأی را اصلاح کند.
	\end{enumerate}
	
	\subsection{مسئله‌ی زبان}
		
			\begin{noteBox}[سیاست زبانی پیشنهادی]
			\begin{itemize}[nosep]
			\item \textbf{زبان رسمی فرایند:} فارسی و انگلیسی.
			\item \textbf{حق ارائه‌ی لایحه به زبان مادری:}
			هر شخص بتواند لایحه‌ی خود را به هر زبانی
			(کُردی، ترکی، عربی، بلوچی و …) ارائه دهد؛
			ترجمه‌ی رسمی بر عهده‌ی دبیرخانه خواهد بود.
			\item \textbf{آراء نهایی:}
		به فارسی و انگلیسی منتشر و در صورت
	ارتباط با اقلیتی خاص، به زبان آن اقلیت نیز ترجمه شوند.
	\end{itemize}
	\end{noteBox}
	
	\section{جمع‌بندی فصل چهارم}
		
		\begin{principleBox}[خلاصه‌ی فصل ۴]
			سازماندهی نهاد باید بر سه اصل استوار باشد:
			\begin{enumerate}[nosep, label=\textcolor{cPrimary}{\arabic*.}]
			\item \textbf{لایه‌بندی:}
			جداسازی نظارت (شورای ناظران)، داوری (هیئت عالی)،
			و اجرا (دبیرخانه) برای جلوگیری از تمرکز قدرت.
			\item \textbf{تکثرگرایی ساختاری:}
			تنوع تخصصی، قومی، جنسیتی، نسلی و ملیتی در تمام سطوح.
			\item \textbf{شفافیت بی‌قیدوشرط:}
		تمام فرایندها، منابع مالی، و آراء به‌صورت عمومی
		در دسترس باشند.
		\end{enumerate}
		
	وکیل تسخیری (دوست دادگاه) مهم‌ترین نوآوری ساختاری
	طرح است و تضمین‌کننده‌ی عدالت فرایندی.
	\end{principleBox}
	
	\clearpage
	
	
	% ////////////////////////////////////////////////////////////
	%  فصل ۵: رویه‌های اجرایی و قضایی معتبر
	% ////////////////////////////////////////////////////////////
	\chapter{رویه‌های اجرایی و قضایی معتبر}\label{chap:proc}
		
		\begin{principleBox}[اصل راهنما]
			رویه‌های این نهاد باید هم‌زمان دو شرط متعارض‌نما را برآورده کنند:
			\begin{itemize}[nosep]
			\item \textbf{استحکام حقوقی:} به‌اندازه‌ای دقیق باشند که
			در نگاه حقوق‌دانان بین‌المللی معتبر شمرده شوند؛
		\item \textbf{انعطاف عملیاتی:} به‌اندازه‌ای سبک باشند که
		یک نهاد سایه‌ی بدون بودجه‌ی دولتی بتواند آن‌ها را اجرا کند.
		\end{itemize}
		الگوی اصلی: ترکیبی از آیین دادرسی دیوان کیفری بین‌المللی (برای استحکام)
	و رویه‌ی دادگاه‌های مردمی (\en{People's Tribunals})
	(برای انعطاف).
	\end{principleBox}
	
	\section{طبقه‌بندی پرونده‌ها}
	
	پرونده‌های قابل رسیدگی در نهاد به چهار دسته تقسیم می‌شوند
	که هر کدام رویه‌ی متفاوتی دارند:
		
		\begin{longtable}{r p{3cm} p{3cm} p{2.5cm} p{3cm}}
		\toprule
		\textbf{نوع} & \textbf{موضوع} & \textbf{سطح رسمیت} &
		\textbf{مدت بررسی} & \textbf{خروجی} \\
		\midrule
		\endhead
		\textbf{نوع الف} &
		مستندسازی جرایم حکومتی &
		بالا (استاندارد بین‌المللی) &
		۳ تا ۱۲ ماه &
		کیفرخواست سایه \\
		\midrule
		\textbf{نوع ب} &
		داوری تصمیم سیاسی (تحریم، اخراج و …) &
		متوسط &
		۱ تا ۳ ماه &
		رأی مشورتی \\
		\midrule
		\textbf{نوع ج} &
		تحلیل گفتمانی (مداخله‌ی نظامی، عفو و …) &
		متوسط — تحلیلی &
		۲ تا ۶ ماه &
		توصیه‌نامه‌ی تحلیلی \\
		\midrule
		\textbf{نوع د} &
		اظهارنظر فوری (بحران‌های آنی) &
		سبک — اضطراری &
		۲۴ تا ۷۲ ساعت &
	بیانیه‌ی فوری \\
	\bottomrule
	\end{longtable}
	
	\section{رویه‌ی پرونده‌های نوع الف: مستندسازی جرایم}
	
	این مهم‌ترین و سنگین‌ترین نوع پرونده است و باید نزدیک‌ترین
	رویه به استانداردهای بین‌المللی را داشته باشد.
	
	\subsection{مرحله‌ی ۱: ارجاع}
		
		\begin{enumerate}[nosep, label=\textcolor{cPrimary}{\arabic*.}]
			\item \textbf{منابع ارجاع مجاز:}
			\begin{itemize}[nosep]
			\item هر عضو هیئت عالی داوران؛
			\item هر سازمان حقوق بشری شناخته‌شده؛
		\item هر شهروند ایرانی یا خارجی با ارائه‌ی شواهد اولیه؛
		\item دبیرخانه‌ی نهاد بر اساس رصد رسانه‌ای و تحقیق مستقل.
		\end{itemize}
		
			\item \textbf{فرم ارجاع استاندارد:} شامل:
			\begin{itemize}[nosep]
			\item شرح مختصر وقایع؛
			\item شواهد اولیه (اسناد، تصاویر، شهادت‌نامه)؛
			\item طبقه‌بندی اولیه‌ی جرم بر اساس اساسنامه‌ی رم؛
		\item اشخاص حقیقی و حقوقی مظنون؛
		\item درخواست مشخص از نهاد.
		\end{itemize}
		
		\item \textbf{بررسی اولیه‌ی قابلیت پذیرش:}
		دبیرخانه ظرف ۱۴ روز بررسی می‌کند که آیا ارجاع
	حداقل‌های لازم را دارد یا خیر. این بررسی صرفاً شکلی
	است و به ماهیت نمی‌پردازد.
	\end{enumerate}
	
	\subsection{مرحله‌ی ۲: تحقیقات مقدماتی}
		
			\begin{principleBox}[گام‌های تحقیقات مقدماتی]
			\begin{enumerate}[nosep, label=\textcolor{cPrimary}{\alph*.}]
			\item \textbf{تشکیل تیم تحقیق:}
			۲ تا ۴ نفر از پژوهشگران دبیرخانه به همراه
			حداقل ۱ عضو هیئت داوران به‌عنوان ناظر.
			
				\item \textbf{جمع‌آوری شواهد:}
				\begin{itemize}[nosep]
				\item شهادت‌نامه‌های مکتوب و ویدئویی
				(با رعایت پروتکل حفاظت از شاهد)؛
				\item تحلیل منابع باز (\en{OSINT})
				شامل تصاویر ماهواره‌ای، ویدئوهای شبکه‌های اجتماعی،
				و اسناد فاش‌شده؛
				\item گزارش‌های سازمان‌های بین‌المللی
				(گزارشگر ویژه‌ی سازمان ملل و …)؛
				\item تحلیل پزشکی قانونی از راه دور
				در صورت امکان؛
			\item اسناد مالی و تجاری مرتبط
			(برای اثبات زنجیره‌ی فرماندهی و منفعت).
			\end{itemize}
			
			\item \textbf{احراز هویت و اعتبارسنجی شواهد:}
			هر مدرک باید حداقل از دو منبع مستقل تأیید شود
			(قاعده‌ی تأیید متقابل).
			شهادت‌نامه‌ها باید با سوگندنامه‌ی مکتوب همراه باشند.
			
			\item \textbf{گزارش مقدماتی:}
			تیم تحقیق گزارشی تهیه می‌کند شامل:
		خلاصه‌ی شواهد، نقاط قوت و ضعف، طبقه‌بندی حقوقی
	پیشنهادی، و توصیه درباره‌ی ادامه یا مختومه‌سازی پرونده.
	\end{enumerate}
	\end{principleBox}
	
	\subsection{مرحله‌ی ۳: کیفرخواست سایه}
		
		\begin{enumerate}[nosep, label=\textcolor{cPrimary}{\arabic*.}]
		\item \textbf{تهیه‌ی پیش‌نویس کیفرخواست:}
		بر اساس گزارش مقدماتی، دادستان نهاد
		پیش‌نویس کیفرخواست را
			با ساختار مشابه کیفرخواست‌های بین‌المللی تهیه می‌کند:
			\begin{itemize}[nosep]
			\item شرح وقایع؛
			\item شواهد؛
			\item طبقه‌بندی حقوقی (جرم علیه بشریت، جرم جنگی، و …)؛
		\item شناسایی متهمان و نقش هر یک در زنجیره‌ی فرماندهی؛
		\item درخواست مشخص (اعلام مسئولیت، توصیه‌ی تحریم و …).
		\end{itemize}
		
		\item \textbf{فعال‌سازی وکیل تسخیری:}
		دفتر وکیل تسخیری موظف است پاسخ مکتوبی
			به کیفرخواست تهیه کند شامل:
			\begin{itemize}[nosep]
			\item نقد شواهد (آیا کافی و موثق هستند؟)؛
		\item دفاعیات احتمالی متهم؛
		\item ملاحظات حقوق بشری (حتی برای متهمان).
		\end{itemize}
		
		\item \textbf{جلسه‌ی رسیدگی:}
		هیئت عالی داوران (یا هیئت منتخب ۵ نفره) هر دو
		سند را بررسی می‌کنند. جلسه می‌تواند حضوری
		یا مجازی باشد.
		
			\item \textbf{صدور رأی:}
			\begin{itemize}[nosep]
			\item رأی با اکثریت ساده صادر می‌شود؛
			\item نظر اقلیت به‌صورت کامل ضمیمه می‌شود؛
		\item رأی شامل یافته‌های واقعی، طبقه‌بندی حقوقی،
	و توصیه‌های عملی است.
	\end{itemize}
	\end{enumerate}
	
	\subsection{مرحله‌ی ۴: انتشار و پیگیری}
		
		\begin{enumerate}[nosep, label=\textcolor{cPrimary}{\arabic*.}]
		\item \textbf{انتشار عمومی:}
		رأی نهایی به فارسی و انگلیسی منتشر می‌شود.
		\item \textbf{ارسال به مراجع بین‌المللی:}
		دبیرخانه رأی را به نهادهای بین‌المللی مرتبط
		و دادگاه‌های ملی با صلاحیت جهانی ارسال می‌کند.
		\item \textbf{بایگانی:}
		تمام مستندات در آرشیو دیجیتال امن ذخیره می‌شوند.
		\item \textbf{رصد تأثیر:}
	دبیرخانه هر ۶ ماه گزارشی از وضعیت پیگیری هر
	کیفرخواست سایه منتشر می‌کند.
	\end{enumerate}
	
	\section{رویه‌ی پرونده‌های نوع ب: داوری تصمیم سیاسی}
		
		\begin{principleBox}[آزمون سه‌مرحله‌ای تناسب]
		هر تصمیم سیاسی پیشنهادی باید از سه آزمون عبور کند:
		
		\medskip
			\textbf{آزمون ۱ — مشروعیت:}
			\begin{itemize}[nosep]
			\item آیا هدف اقدام مشروع است؟
		(حفاظت از حقوق بشر، فشار بر حکومت ناقض)
		\item آیا مبنای حقوقی بین‌المللی دارد؟
		\end{itemize}
		
		\medskip
			\textbf{آزمون ۲ — ضرورت:}
			\begin{itemize}[nosep]
		\item آیا جایگزین کم‌هزینه‌تری بررسی شده؟
		\item آیا اقدام مورد نظر آخرین گزینه‌ی معقول است؟
		\end{itemize}
		
		\medskip
			\textbf{آزمون ۳ — تناسب:}
			\begin{itemize}[nosep]
			\item آیا هزینه‌ی وارد بر افراد بی‌تقصیر از فایده‌ی
			مورد انتظار کمتر است؟
			\item آیا اثرات جانبی قابل کنترل هستند؟
		\item آیا مکانیزم جبران برای آسیب‌دیدگان بی‌تقصیر
	پیش‌بینی شده؟
	\end{itemize}
	\end{principleBox}
	
	\subsection{فرایند داوری تصمیم سیاسی}
		
		\begin{enumerate}[nosep, label=\textcolor{cPrimary}{\arabic*.}]
		\item \textbf{ثبت درخواست:}
		هر جریان سیاسی، سازمان مدنی، یا شهروند می‌تواند
		درخواست بررسی یک تصمیم سیاسی را ثبت کند.
		
		\item \textbf{تحقیق سریع:}
		تیم پژوهش دبیرخانه ظرف ۲ هفته گزارش مقدماتی
		شامل بررسی سوابق تطبیقی و تحلیل هزینه‌ـ‌فایده تهیه می‌کند.
		
		\item \textbf{جلسه‌ی استماع:}
		حداقل ۲ نفر موافق و ۲ نفر مخالف اقدام پیشنهادی
		دعوت می‌شوند تا نظرات خود را ارائه دهند.
		
		\item \textbf{اجرای آزمون سه‌مرحله‌ای:}
		هیئت داوران آزمون فوق را اجرا و نتایج را مستند می‌کند.
		
		\item \textbf{صدور رأی مشورتی:}
		رأی شامل یکی از چهار نتیجه‌ی ممکن است:
		
			% ── جدول نتایج — اصلاح‌شده با نمادهای ساده ──
			\begin{longtable}{r p{3.5cm} p{7.5cm}}
			\toprule
			\textbf{نتیجه} & \textbf{عنوان} & \textbf{معنا} \\
			\midrule
			\endhead
			{\color{cGold}\textbf{★}} &
			\textbf{توصیه‌شده} &
			اقدام هر سه آزمون را با موفقیت گذرانده؛
			نهاد اجرای آن را توصیه می‌کند. \\
			\midrule
			{\color{cPrimary}\textbf{■}} &
			\textbf{مجاز با قیود} &
			اقدام در اصل مشروع است اما نیاز به اصلاحاتی دارد؛
			قیود مشخص ذکر می‌شوند. \\
			\midrule
			{\color{cGray}\textbf{▲}} &
			\textbf{تعلیق} &
			شواهد کافی برای قضاوت نیست؛
			بررسی به زمان دیگری موکول می‌شود. \\
			\midrule
			{\color{cAccent}\textbf{✕}} &
			\textbf{توصیه‌نشده} &
			اقدام یک یا چند آزمون را نگذرانده؛
		نهاد اجرای آن را توصیه نمی‌کند. \\
		\bottomrule
		\end{longtable}
		
	\item \textbf{مهلت اعتراض:}
	۲۱ روز برای ارائه‌ی اعتراض مستند.
	\end{enumerate}
	
	\section{رویه‌ی پرونده‌های نوع ج: تحلیل گفتمانی}
		
		\begin{enumerate}[nosep, label=\textcolor{cPrimary}{\arabic*.}]
		\item \textbf{شناسایی موضوع:}
		شعبه‌ی گفتمان و مرزگذاری، موضوعات اختلافی کلیدی
		را شناسایی و اولویت‌بندی می‌کند.
		
		\item \textbf{تشکیل کارگروه تخصصی:}
		برای هر موضوع، کارگروهی از ۳ تا ۵ متخصص
		(حداقل یک نفر از هر طرف بحث) تشکیل می‌شود.
		
		\item \textbf{تهیه‌ی سند مبنا:}
		سندی تحلیلی که تمام دیدگاه‌های موجود را بی‌طرفانه
		ترسیم می‌کند.
		
		\item \textbf{نشست مشورتی گسترده:}
		نشستی با حضور ۲۰ تا ۵۰ نفر از نخبگان و فعالان
		برای بحث ساختارمند.
		
		\item \textbf{صدور توصیه‌نامه:}
			سندی مفصل شامل:
			\begin{itemize}[nosep]
			\item تحلیل هر دیدگاه با نقاط قوت و ضعف؛
			\item موضع پیشنهادی نهاد (در صورت رسیدن به اجماع)؛
		\item گزارش اختلاف‌نظرها (در صورت عدم اجماع)؛
	\item چارچوب پیشنهادی برای تصمیم‌گیری آتی.
	\end{itemize}
	\end{enumerate}
	
	\section{رویه‌ی پرونده‌های نوع د: اظهارنظر فوری}
		
		\begin{warningBox}[رویه‌ی اضطراری]
		در بحران‌های آنی (مثلاً سرکوب خونین ناگهانی، اعدام فوری
		زندانی سیاسی)، نهاد باید بتواند ظرف ۲۴ تا ۷۲ ساعت واکنش نشان دهد:
			
			\begin{enumerate}[nosep, label=\textcolor{cAccent}{\arabic*.}]
			\item \textbf{فعال‌سازی پروتکل اضطراری:}
			مدیر اجرایی یا هر ۳ عضو هیئت عالی می‌توانند
			پروتکل اضطراری را فعال کنند.
			
			\item \textbf{جلسه‌ی فوری:}
			حداقل ۵ عضو هیئت عالی ظرف ۱۲ ساعت
			جلسه‌ی مجازی تشکیل می‌دهند.
			
			\item \textbf{بیانیه‌ی فوری:}
			شامل: شرح مختصر وقایع، ارزیابی حقوقی اولیه،
			و توصیه‌ی فوری (مثلاً توقف اعدام، تحریم فوری شخص خاص).
			
			\item \textbf{بازنگری:}
		بیانیه‌ی فوری ظرف ۳۰ روز مورد بازنگری کامل قرار می‌گیرد
		و در صورت لزوم اصلاح یا تکمیل می‌شود.
		\end{enumerate}
		
		\medskip
		\textbf{قاعده‌ی حفاظتی:} بیانیه‌های فوری نمی‌توانند شامل
	توصیه‌ی اقدام قهری شدید (مانند مداخله‌ی نظامی) باشند.
	چنین توصیه‌هایی فقط از مسیر رویه‌ی کامل (نوع ب یا ج) مجاز هستند.
	\end{warningBox}
	
	\section{استانداردهای شواهد}
		
			\begin{principleBox}[سه سطح استاندارد]
			\begin{longtable}{r p{3cm} p{4cm} p{4.5cm}}
			\toprule
			\textbf{سطح} & \textbf{استاندارد} & \textbf{کاربرد} & \textbf{توضیح} \\
			\midrule
			\endhead
			۱ & فراتر از شک معقول &
			پرونده‌های نوع الف (کیفرخواست سایه) &
			بالاترین سطح اثبات — مشابه دادگاه کیفری \\
			\midrule
			۲ & ترجیح ادله &
			پرونده‌های نوع ب و ج (داوری و تحلیل) &
			سطح متوسط — مشابه دادگاه مدنی \\
			\midrule
			۳ & دلایل معقول برای باور &
			پرونده‌های نوع د (اظهارنظر فوری) &
		پایین‌ترین سطح مجاز — فقط برای واکنش اضطراری \\
	\bottomrule
	\end{longtable}
	\end{principleBox}
	
	\section{حفاظت از شاهدان و قربانیان}
		
			\begin{principleBox}[پروتکل حفاظت]
			\begin{enumerate}[nosep, label=\textcolor{cPrimary}{\arabic*.}]
			\item \textbf{ناشناس‌سازی:}
			هویت شاهدان داخل ایران هرگز فاش نمی‌شود.
			در آراء منتشرشده از کُدهای ناشناس استفاده می‌شود.
			
			\item \textbf{رمزنگاری:}
			تمام ارتباطات با شاهدان از طریق کانال‌های رمزنگاری‌شده
			سرتاسری صورت می‌گیرد.
			
			\item \textbf{حداقل‌سازی داده:}
			فقط اطلاعات ضروری از شاهدان جمع‌آوری می‌شود.
			
			\item \textbf{پشتیبانی روان‌شناختی:}
			شاهدان و قربانیان به خدمات مشاوره‌ی روان‌شناختی
			(مجازی و رایگان) دسترسی داشته باشند.
			
			\item \textbf{هشدار ریسک:}
		قبل از هر مشارکت، ریسک‌های احتمالی
	(بازداشت خانواده در ایران و …) صادقانه توضیح داده شود.
	\end{enumerate}
	\end{principleBox}
	
	\section{جمع‌بندی فصل پنجم}
		
			\begin{principleBox}[خلاصه‌ی فصل ۵]
			\begin{enumerate}[nosep, label=\textcolor{cPrimary}{\arabic*.}]
			\item چهار نوع پرونده با چهار سطح رسمیت تعریف شده‌اند.
			\item آزمون سه‌مرحله‌ای تناسب، ستون فقرات داوری تصمیم‌های سیاسی است.
			\item وکیل تسخیری در تمام پرونده‌های نوع الف اجباری است.
			\item سه سطح استاندارد شواهد برای سه نوع خروجی تعریف شده‌اند.
			\item حفاظت از شاهدان اولویت مطلق دارد.
		\item پروتکل اضطراری اجازه‌ی واکنش سریع می‌دهد اما
	توصیه‌ی اقدام قهری شدید را ممنوع می‌کند.
	\end{enumerate}
	\end{principleBox}
	
	\clearpage
	
	
	% ////////////////////////////////////////////////////////////
	%  فصل ۶: مباحث مالی و هزینه‌ای
	% ////////////////////////////////////////////////////////////
	\chapter{مباحث مالی و هزینه‌ای}\label{chap:finance}
	
	\section{اصول مالی حاکم}
		
			\begin{principleBox}[چهار اصل مالی]
			\begin{enumerate}[nosep, label=\textcolor{cGold}{\Roman*.}]
			\item \textbf{شفافیت کامل:}
			تمام درآمدها و هزینه‌ها به‌صورت سالانه و عمومی منتشر شوند.
			\item \textbf{تنوع منابع:}
			هیچ منبع واحد نباید بیش از ۲۵٪ بودجه‌ی سالانه را تأمین کند
			(برای جلوگیری از وابستگی).
			\item \textbf{ممنوعیت دریافت از دولت‌ها:}
			نهاد مستقیماً از هیچ دولتی بودجه دریافت نمی‌کند.
			دریافت از بنیادهای مستقلِ کشورهای دموکراتیک مجاز است
			مشروط به عدم اعمال نفوذ بر محتوای آراء.
			\item \textbf{صرفه‌جویی ساختاری:}
		نهاد باید با حداقل بوروکراسی و حداکثر استفاده از
	فناوری‌های دیجیتال اداره شود.
	\end{enumerate}
	\end{principleBox}
	
	\section{برآورد بودجه‌ی سالانه}
		
		\begin{warningBox}[تذکر]
		ارقام زیر برآورد اولیه هستند و بر اساس تجربه‌ی نهادهای مشابه
	و نرخ‌های بازار تخمین زده شده‌اند. بودجه‌ی واقعی به
	مقیاس عملیات بستگی دارد.
	\end{warningBox}
		
		\begin{longtable}{r p{5.5cm} r r}
		\toprule
		\textbf{ردیف} & \textbf{قلم هزینه} &
		\textbf{حداقل (€)} & \textbf{مطلوب (€)} \\
		\midrule
		\endhead
		\multicolumn{4}{r}{\textbf{الف. نیروی انسانی}} \\
		\midrule
		۱ & مدیر اجرایی (تمام‌وقت) & ۶۰٬۰۰۰ & ۸۰٬۰۰۰ \\
		۲ & هماهنگ‌کننده‌ی پژوهش (تمام‌وقت) & ۵۰٬۰۰۰ & ۶۵٬۰۰۰ \\
		۳ & مسئول ارتباطات و رسانه (تمام‌وقت) & ۴۵٬۰۰۰ & ۶۰٬۰۰۰ \\
		۴ & دستیاران پژوهشی (۲ نفر پاره‌وقت) & ۳۰٬۰۰۰ & ۵۰٬۰۰۰ \\
		۵ & مترجم فارسی‌ـ‌انگلیسی (پاره‌وقت) & ۱۵٬۰۰۰ & ۲۵٬۰۰۰ \\
		۶ & مشاور حقوقی بین‌المللی (قراردادی) & ۲۰٬۰۰۰ & ۴۰٬۰۰۰ \\
		۷ & حق‌الزحمه‌ی داوران (نمادین، جلسه‌ای) & ۱۵٬۰۰۰ & ۳۰٬۰۰۰ \\
		\midrule
		\multicolumn{4}{r}{\textbf{ب. عملیاتی}} \\
		\midrule
		۸ & فناوری اطلاعات و امنیت سایبری & ۱۵٬۰۰۰ & ۳۰٬۰۰۰ \\
		۹ & میزبانی وب، بانک اطلاعاتی، رمزنگاری & ۵٬۰۰۰ & ۱۰٬۰۰۰ \\
		۱۰ & سفر و اقامت (جلسات حضوری، سالی ۲ بار) & ۲۰٬۰۰۰ & ۴۰٬۰۰۰ \\
		۱۱ & اجاره‌ی دفتر (مجازی یا مشترک) & ۵٬۰۰۰ & ۱۵٬۰۰۰ \\
		\midrule
		\multicolumn{4}{r}{\textbf{ج. انتشارات و ارتباطات}} \\
		\midrule
		۱۲ & طراحی و انتشار آراء و گزارش‌ها & ۸٬۰۰۰ & ۱۵٬۰۰۰ \\
		۱۳ & رویدادهای عمومی و جلسات استماع & ۱۰٬۰۰۰ & ۲۵٬۰۰۰ \\
		۱۴ & تولید محتوای رسانه‌ای (ویدئو، پادکست) & ۵٬۰۰۰ & ۱۵٬۰۰۰ \\
		\midrule
		\multicolumn{4}{r}{\textbf{د. حقوقی و حسابرسی}} \\
		\midrule
		۱۵ & ثبت قانونی نهاد (در کشور میزبان) & ۵٬۰۰۰ & ۵٬۰۰۰ \\
		۱۶ & حسابرسی مالی سالانه‌ی مستقل & ۵٬۰۰۰ & ۱۰٬۰۰۰ \\
		۱۷ & بیمه‌ی مسئولیت مدنی & ۳٬۰۰۰ & ۸٬۰۰۰ \\
		\midrule
		\multicolumn{4}{r}{\textbf{ه‍. ذخیره‌ی احتیاطی}} \\
		\midrule
		۱۸ & صندوق ذخیره (۱۰٪) & ۳۲٬۰۰۰ & ۵۲٬۰۰۰ \\
		\midrule
		\midrule
		& \textbf{جمع کل سالانه} &
	\textbf{۳۴۸٬۰۰۰} & \textbf{۵۷۵٬۰۰۰} \\
	\bottomrule
	\end{longtable}
	
	\section{منابع تأمین مالی}
		
		\begin{enumerate}[nosep, label=\textcolor{cPrimary}{\arabic*.}]
			\item \textbf{بنیادهای بین‌المللی حقوق بشر و دموکراسی:}
			\begin{itemize}[nosep]
			\item بنیاد ملی دموکراسی آمریکا (\en{NED})
			\item بنیادهای جامعه‌ی باز (\en{Open Society Foundations})
			\item صندوق سیگرید راوزینگ (\en{Sigrid Rausing Trust})
			\item بنیاد اوک (\en{Oak Foundation})
			\item بنیاد فورد — بخش حقوق بشر
		\item بنیاد هاینریش بل (\en{Heinrich Böll Stiftung})
		\item بنیاد فریدریش ناومان (\en{Friedrich Naumann Foundation})
		\end{itemize}
		
			\item \textbf{کمک‌های مردمی:}
			\begin{itemize}[nosep]
			\item کمپین‌های تأمین مالی جمعی؛
			\item عضویت حمایتی ماهانه
			(مثلاً ۵ تا ۵۰ یورو در ماه)؛
			\item \textbf{سقف پذیرش فردی:}
		حداکثر ۵٬۰۰۰ یورو در سال از هر شخص حقیقی
		(برای جلوگیری از نفوذ).
		\end{itemize}
		
			\item \textbf{درآمد خدماتی:}
			\begin{itemize}[nosep]
			\item انتشار گزارش‌های تحلیلی اشتراکی
			(برای پژوهشگران و مؤسسات)؛
			\item برگزاری کارگاه‌های آموزشی
		(مستندسازی حقوق بشر، حقوق بین‌الملل کیفری)؛
	\item مشاوره‌ی حقوقی به نهادهای مدنی ایرانی.
	\end{itemize}
	\end{enumerate}
	
	\section{سازوکارهای نظارت مالی}
		
			\begin{principleBox}[نظارت چهارلایه]
			\begin{enumerate}[nosep, label=\textcolor{cPrimary}{\arabic*.}]
			\item \textbf{حسابرسی داخلی:}
			خزانه‌دار نهاد (عضو دبیرخانه) ماهانه گزارش مالی
			به هیئت عالی ارائه می‌دهد.
			\item \textbf{حسابرسی خارجی:}
			شرکت حسابرسی مستقل سالانه حساب‌ها را بررسی می‌کند.
			\item \textbf{نظارت شورای ناظران:}
			شورای ناظران بیرونی حق دسترسی کامل به
			اسناد مالی دارد.
			\item \textbf{انتشار عمومی:}
		گزارش مالی سالانه شامل تمام درآمدها (با ذکر منبع)
	و هزینه‌ها به‌صورت عمومی منتشر می‌شود.
	\end{enumerate}
	\end{principleBox}
		
		\begin{noteBox}[نکته‌ی کلیدی]
		\textbf{فقر مالی بهتر از وابستگی مالی است.}
		اگر نهاد نتواند بودجه‌ی مطلوب را تأمین کند، باید با بودجه‌ی حداقلی
	و اتکای بیشتر بر کار داوطلبانه فعالیت کند — نه این‌که
	شرط تنوع منابع را زیر پا بگذارد.
	\end{noteBox}
	
	\section{جمع‌بندی فصل ششم}
		
			\begin{principleBox}[خلاصه‌ی فصل ۶]
			\begin{itemize}[nosep]
			\item بودجه‌ی سالانه‌ی نهاد بین ۳۵۰ تا ۵۷۵ هزار یورو برآورد می‌شود.
			\item منابع مالی باید متنوع، شفاف، و مستقل از دولت‌ها باشند.
			\item قاعده‌ی ۲۵٪ (سقف سهم هر منبع) از وابستگی جلوگیری می‌کند.
			\item نظارت مالی چهارلایه تضمین‌کننده‌ی سلامت مالی است.
		\item در صورت کمبود بودجه، مقیاس کوچک‌تر ترجیح دارد بر
	وابستگی مالی.
	\end{itemize}
	\end{principleBox}
	
	\clearpage
	
	
	% ////////////////////////////////////////////////////////////
	%  فصل ۷: ارتباط با مراجع بین‌المللی
	% ////////////////////////////////////////////////////////////
	\chapter{ارتباط با مراجع قضایی بین‌المللی و کشورهای حامی}\label{chap:intl}
	
	\section{نقشه‌ی نهادهای بین‌المللی مرتبط}
	
		% ── نمودار اصلاح‌شده: ارتباطات نهادی ──
		\begin{figure}[H]
			\centering
			\begin{tikzpicture}[
				node distance=2.5cm and 3cm,
				center/.style={
				rectangle, rounded corners=8pt, draw=cPrimary, very thick,
			fill=cGold!10, text width=3.5cm, align=center,
			font=\small\bfseries, minimum height=1.5cm
				},
				sat/.style={
				rectangle, rounded corners=3pt, draw=cPrimary!60,
			fill=cLight, text width=3.5cm, align=center,
			font=\footnotesize, minimum height=1cm
			},
% ── ادامه‌ی نمودار فصل ۷ ──
		arrow/.style={<->, thick, cPrimary!50}
		]
		\node[center] (shadow)
		{{\RL{دادگاه در سایه}}};
		
		% نهادهای قضایی بین‌المللی — بالا
		\node[sat, above left=of shadow, xshift=-1cm] (icc)
		{{\RL{دیوان کیفری بین‌المللی}}\\(\en{ICC})};
		\node[sat, above=of shadow] (icj)
		{{\RL{دیوان بین‌المللی دادگستری}}\\(\en{ICJ})};
		\node[sat, above right=of shadow, xshift=1cm] (ohchr)
		{{\RL{دفتر کمیساریای}}\\{\RL{حقوق بشر سازمان ملل}}};
		
		% نهادهای حقوق بشری — چپ و راست
		\node[sat, left=of shadow, xshift=-1.5cm] (hrw)
		{{\RL{دیده‌بان حقوق بشر}}\\(\en{HRW})};
		\node[sat, right=of shadow, xshift=1.5cm] (amnesty)
		{{\RL{عفو بین‌الملل}}\\(\en{Amnesty})};
		
		% دادگاه‌های ملی و نهادهای سازمان ملل — پایین
		\node[sat, below left=of shadow, xshift=-1cm] (national)
		{{\RL{دادگاه‌های ملی}}\\{\RL{صلاحیت جهانی}}\\{\RL{(سوئد، آلمان، کانادا)}}};
		\node[sat, below=of shadow] (special)
		{{\RL{گزارشگر ویژه‌ی}}\\{\RL{حقوق بشر ایران}}\\{\RL{سازمان ملل}}};
		\node[sat, below right=of shadow, xshift=1cm] (ffm)
		{{\RL{هیئت حقیقت‌یاب}}\\{\RL{بین‌المللی ایران}}\\(\en{FFMI})};
		
		% پیکان‌ها
		\draw[arrow] (shadow) -- (icc);
		\draw[arrow] (shadow) -- (icj);
		\draw[arrow] (shadow) -- (ohchr);
		\draw[arrow] (shadow) -- (hrw);
		\draw[arrow] (shadow) -- (amnesty);
		\draw[arrow] (shadow) -- (national);
		\draw[arrow] (shadow) -- (special);
		\draw[arrow] (shadow) -- (ffm);
		
		\end{tikzpicture}
		\caption{نقشه‌ی ارتباطات نهادی دادگاه در سایه با مراجع بین‌المللی}
		\end{figure}
		
		\section{دیوان کیفری بین‌المللی}
		
		\subsection{وضعیت حقوقی ایران}
		
		\begin{warningBox}[محدودیت اساسی]
		ایران اساسنامه‌ی رم را امضا نکرده و عضو دیوان کیفری بین‌المللی نیست.
		بنابراین دیوان نمی‌تواند به‌صورت مستقل پرونده‌ای
		علیه مقام‌های ایرانی باز کند \textbf{مگر}:
		\begin{itemize}[nosep]
		\item شورای امنیت سازمان ملل موضوع را به دیوان
		ارجاع دهد (مشابه ارجاع سودان و لیبی) — که
		با وتوی احتمالی روسیه و چین بعید است؛
		\item جرایم در قلمرو یک کشور عضو دیوان
		ارتکاب یافته باشند (مانند حمله به هواپیمای
		اوکراینی — اوکراین صلاحیت دیوان را پذیرفته).
		\end{itemize}
		\end{warningBox}
		
		\subsection{نقش دادگاه در سایه}
		
		\begin{principleBox}[استراتژی سه‌گانه]
		\begin{enumerate}[nosep, label=\textcolor{cPrimary}{\arabic*.}]
		\item \textbf{آماده‌سازی پرونده:}
		مستندات نهاد می‌توانند پایه‌ی تحقیقات مقدماتی
		دادستان دیوان قرار گیرند — به‌ویژه در
		پرونده‌ی هواپیمای اوکراینی (پرواز ۷۵۲).
		
		\item \textbf{فشار برای ارجاع شورای امنیت:}
		آراء مستدل نهاد می‌توانند ابزار لابی‌گری حقوقی
		برای متقاعدسازی اعضای شورای امنیت باشند.
		
		\item \textbf{آماده‌سازی برای پس از گذار:}
		اگر ایران پس از گذار به دیوان بپیوندد،
		مستندات نهاد آماده‌ی ارائه خواهند بود.
		(سابقه: تأخیر ۲۰ ساله‌ی پرونده‌های کامبوج
		نشان می‌دهد آماده‌سازی زودهنگام حیاتی است.)
		\end{enumerate}
		\end{principleBox}
		
		\section{صلاحیت جهانی}
		
		\subsection{کشورهای دارای سابقه‌ی اعمال صلاحیت جهانی}
		
		\begin{longtable}{r p{2.5cm} p{4cm} p{5cm}}
		\toprule
		\textbf{ردیف} & \textbf{کشور} & \textbf{سابقه‌ی مرتبط} &
		\textbf{ظرفیت همکاری} \\
		\midrule
		\endhead
		۱ & سوئد &
		پرونده‌ی حمید نوری (محکومیت ۲۰۲۲) &
		بالاترین سابقه‌ی مستقیم در پرونده‌ی ایران.
		واحد جرایم بین‌المللی پلیس سوئد فعال. \\
		\midrule
		۲ & آلمان &
		محکومیت انور رسلان (مقام سوری، ۲۰۲۲) &
		قوی‌ترین ساختار صلاحیت جهانی در اروپا.
		دادستانی فدرال آلمان واحد ویژه‌ی جرایم
		بین‌المللی دارد. \\
		\midrule
		۳ & کانادا &
		قانون عدالت برای قربانیان مقام‌های فاسد خارجی
		(قانون مگنیتسکی) &
		جامعه‌ی بزرگ ایرانی؛ ظرفیت حقوقی مناسب؛
		سابقه‌ی تحریم مقام‌های ایرانی. \\
		\midrule
		۴ & هلند &
		میزبان دیوان کیفری و دیوان دادگستری بین‌المللی؛
		سابقه‌ی پرونده‌ی هواپیمای مالزیایی &
		محیط حقوقی مساعد؛ دسترسی به نهادهای
		بین‌المللی. \\
		\midrule
		۵ & بریتانیا &
		سابقه‌ی دادگاه پینوشه (۱۹۹۸)؛
		مقررات تحریمی حقوق بشر (۲۰۲۰) &
		ظرفیت حقوقی بالا؛ جامعه‌ی ایرانی فعال. \\
		\midrule
		۶ & آرژانتین &
		پرونده‌ی جرایم فرانکوئیسم اسپانیا؛
		بررسی جرایم میانمار &
		سابقه‌ی نوآورانه در صلاحیت جهانی؛
		کمتر در معرض فشار سیاسی. \\
		\midrule
		۷ & بلژیک &
		سابقه‌ی طولانی صلاحیت جهانی
		(هرچند محدود شده در ۲۰۰۳) &
		همچنان ظرفیت محدودی برای جرایم
		بین‌المللی دارد. \\
		\bottomrule
		\end{longtable}
		
		\subsection{نقش دادگاه در سایه در تقویت پرونده‌های صلاحیت جهانی}
		
		\begin{principleBox}[سه کارکرد]
		\begin{enumerate}[nosep, label=\textcolor{cPrimary}{\arabic*.}]
		\item \textbf{تهیه‌ی بسته‌های حقوقی آماده:}
		برای هر متهم اصلی، بسته‌ای شامل شواهد طبقه‌بندی‌شده،
		تحلیل حقوقی، و شناسایی شاهدان بالقوه تهیه شود
		که مستقیماً قابل ارائه به دادستان‌های ملی باشد.
		
		\item \textbf{شناسایی حضور متهمان:}
		رصد سفرها و اقامت مقام‌ها و وابستگان آن‌ها
		در کشورهای دارای صلاحیت جهانی.
		
		\item \textbf{هماهنگی حقوقی:}
		ایجاد شبکه‌ای از وکلای داوطلب در هر کشور هدف
		برای پیگیری شکایات.
		\end{enumerate}
		\end{principleBox}
		
		\section{هیئت حقیقت‌یاب بین‌المللی ایران}
		
		در نوامبر ۲۰۲۲، شورای حقوق بشر سازمان ملل
		\keyword{هیئت حقیقت‌یاب بین‌المللی مستقل}
		درباره‌ی ایران را تأسیس کرد.
		
		\begin{noteBox}[فرصت همکاری]
		\begin{itemize}[nosep]
		\item دادگاه در سایه می‌تواند شواهد و مستندات خود را
		به‌صورت رسمی به هیئت حقیقت‌یاب ارائه دهد.
		\item آراء مشورتی نهاد می‌توانند تحلیل حقوقی مکمل
		برای یافته‌های هیئت باشند.
		\item هیئت یک مأموریت موقت سازمان ملل است
		و ممکن است تمدید نشود. مستندات دادگاه در سایه
		تداوم کار را فارغ از تصمیمات سیاسی
		سازمان ملل تضمین می‌کند.
		\item \textbf{تفاوت کلیدی:}
		هیئت حقیقت‌یاب فقط حقیقت‌یابی می‌کند و توصیه‌ی
		سیاستی ارائه نمی‌دهد.
		دادگاه در سایه هم حقیقت‌یابی می‌کند و هم
		\keyword{توصیه‌ی عملیاتی} صادر می‌کند
		(مثلاً توصیه‌ی تحریم شخص خاص).
		این دو نهاد مکمل یکدیگرند، نه رقیب.
		\end{itemize}
		\end{noteBox}
		
		\section{گزارشگر ویژه‌ی حقوق بشر ایران}
		
		از سال ۲۰۱۱، شورای حقوق بشر سازمان ملل یک
		\keyword{گزارشگر ویژه}
		برای وضعیت حقوق بشر در ایران دارد.
		
		\begin{principleBox}[ارتباط با گزارشگر ویژه]
		\begin{enumerate}[nosep, label=\textcolor{cPrimary}{\arabic*.}]
		\item \textbf{ارائه‌ی منظم مستندات:}
		دبیرخانه‌ی نهاد هر ۶ ماه گزارشی از یافته‌های خود
		را به دفتر گزارشگر ویژه ارسال کند.
		\item \textbf{هماهنگی تحقیقات:}
		شناسایی موضوعاتی که گزارشگر ویژه در حال بررسی
		آن‌هاست و تأمین شواهد مکمل.
		\item \textbf{تقویت متقابل:}
		ارجاع به گزارش‌های گزارشگر ویژه در آراء نهاد
		و بالعکس — ایجاد زنجیره‌ی اعتبار متقابل.
		\item \textbf{حضور در جلسات شورای حقوق بشر:}
		نماینده‌ی نهاد در جلسات مربوط به ایران
		به‌عنوان ناظر سازمان غیردولتی حضور یابد
		و لوایح حقوقی ارائه دهد.
		\end{enumerate}
		\end{principleBox}
		
		\section{سازمان‌های حقوق بشری بین‌المللی}
		
		\begin{longtable}{r p{3cm} p{4cm} p{4.5cm}}
		\toprule
		\textbf{ردیف} & \textbf{سازمان} & \textbf{نوع همکاری} &
		\textbf{ملاحظات} \\
		\midrule
		\endhead
		۱ & دیده‌بان حقوق بشر &
		تبادل شواهد، بررسی مشترک،
		ارجاع متقابل &
		تیم ایران قوی‌ای دارد.
		می‌تواند اعتبار بین‌المللی نهاد را تقویت کند. \\
		\midrule
		۲ & عفو بین‌الملل &
		همکاری در مستندسازی،
		کمپین‌های مشترک &
		پوشش جهانی گسترده.
		باید استقلال نهاد حفظ شود. \\
		\midrule
		۳ & فدراسیون بین‌المللی حقوق بشر (\en{FIDH}) &
		مشاوره‌ی حقوقی، حمایت از
		پرونده‌های صلاحیت جهانی &
		سابقه‌ی طولانی همکاری با
		جنبش‌های دموکراتیک. \\
		\midrule
		۴ & گزارشگران بدون مرز (\en{RSF}) &
		مستندسازی جرایم علیه
		روزنامه‌نگاران ایرانی &
		تخصص ویژه در آزادی مطبوعات.
		مکمل کار شعبه‌ی مستندسازی. \\
		\midrule
		۵ & مرکز عدالت و حقوق بین‌الملل (\en{CEJIL}) &
		مشاوره‌ی فنی در طراحی
		آیین دادرسی &
		تجربه‌ی بالا در دادخواهی
		بین‌المللی (آمریکای لاتین). \\
		\midrule
		۶ & مرکز بین‌المللی عدالت انتقالی (\en{ICTJ}) &
		مشاوره در عدالت انتقالی
		و طراحی ساختار پس از گذار &
		تجربه‌ی جهانی در عدالت انتقالی.
		بسیار مرتبط با مأموریت بلندمدت نهاد. \\
		\midrule
		۷ & مرکز عبدالرحمن برومند (\en{ABC}) &
		تبادل مستندات،
		بانک اطلاعاتی مشترک &
		تخصص ویژه در ایران.
		بانک اطلاعاتی ارزشمند اعدام‌ها.
		نزدیک‌ترین نهاد موجود به مأموریت
		مستندسازی دادگاه در سایه. \\
		\midrule
		۸ & ایران حقوق بشر (\en{IHR}) &
		تبادل داده‌ها، همکاری در
		رصد اعدام‌ها &
		مستقر در نروژ. فعال‌ترین نهاد
		رصد اعدام در ایران. \\
		\bottomrule
		\end{longtable}
		
		\section{دولت‌های حامی: چارچوب تعامل}
		
		\begin{warningBox}[اصل بنیادین]
		دادگاه در سایه \textbf{هرگز از هیچ دولتی مستقیماً بودجه
		دریافت نمی‌کند} و \textbf{هرگز ابزار سیاست خارجی هیچ
		کشوری نمی‌شود}. با این حال، همکاری با دولت‌ها در
		چارچوب‌های مشخص ضروری و مجاز است.
		\end{warningBox}
		
		\begin{principleBox}[سطوح مجاز تعامل با دولت‌ها]
		\begin{longtable}{r p{3.5cm} p{4cm} p{4cm}}
		\toprule
		\textbf{سطح} & \textbf{نوع تعامل} & \textbf{مثال} & \textbf{قید} \\
		\midrule
		\endhead
		{\color{cGold}\textbf{۱. مجاز}} &
		ارائه‌ی مستندات و آراء
		به دولت‌ها &
		ارسال کیفرخواست سایه
		به وزارت خارجه‌ی کانادا &
		یک‌طرفه و علنی \\
		\midrule
		{\color{cGold}\textbf{۲. مجاز}} &
		حضور در جلسات پارلمانی
		به‌عنوان شاهد متخصص &
		شهادت در کمیته‌ی حقوق بشر
		پارلمان اروپا &
		فقط بر اساس دعوت رسمی \\
		\midrule
		{\color{cGold}\textbf{۳. مشروط}} &
		مشاوره‌ی فنی به دولت‌ها
		در طراحی تحریم‌ها &
		کمک به تهیه‌ی فهرست
		تحریمی مگنیتسکی &
		فقط فنی، نه سیاسی؛
		نهاد حق وتو بر لیست نهایی
		ندارد \\
		\midrule
		{\color{cAccent}\textbf{۴. ممنوع}} &
		دریافت بودجه‌ی مستقیم
		از دولت‌ها &
		--- &
		ممنوعیت مطلق \\
		\midrule
		{\color{cAccent}\textbf{۵. ممنوع}} &
		هماهنگی عملیاتی با سرویس‌های
		اطلاعاتی &
		--- &
		ممنوعیت مطلق \\
		\midrule
		{\color{cAccent}\textbf{۶. ممنوع}} &
		حمایت از اقدام نظامی
		دولت خاص بدون فرایند داخلی نهاد &
		--- &
		فقط از مسیر رأی مشورتی
		رسمی نهاد مجاز است \\
		\bottomrule
		\end{longtable}
		\end{principleBox}
		
		\section{کشور میزبان: ملاحظات حقوقی}
		
		نهاد باید در یک کشور دموکراتیک به ثبت قانونی برسد. انتخاب
		کشور میزبان اهمیت استراتژیک بالایی دارد.
		
		\begin{longtable}{r p{2.5cm} p{4.5cm} p{4.5cm}}
		\toprule
		\textbf{ردیف} & \textbf{کشور} & \textbf{مزایا} & \textbf{معایب} \\
		\midrule
		\endhead
		۱ & هلند &
		میزبان دیوان کیفری و دیوان دادگستری بین‌المللی؛
		فضای حقوقی مساعد؛
		دسترسی به نهادهای بین‌المللی؛
		سابقه‌ی میزبانی دادگاه‌های ویژه &
		هزینه‌ی بالای زندگی؛
		فاصله‌ی نسبی از دیاسپورای
		بزرگ ایرانی \\
		\midrule
		۲ & بریتانیا &
		سنت حقوقی قوی؛
		جامعه‌ی ایرانی بزرگ؛
		رسانه‌های بین‌المللی فعال &
		پس از برگزیت ارتباط کمتر
		با نهادهای اروپایی؛
		هزینه‌ی بسیار بالا (لندن) \\
		\midrule
		۳ & آلمان &
		قوی‌ترین ساختار صلاحیت جهانی؛
		بنیادهای سیاسی متعدد؛
		هزینه‌ی متوسط &
		بوروکراسی سنگین‌تر؛
		زبان آلمانی مانعی برای
		ارتباطات عمومی \\
		\midrule
		۴ & سوئد &
		سابقه‌ی پرونده‌ی نوری؛
		سنت حقوق بشری قوی؛
		جامعه‌ی ایرانی فعال &
		جمعیت کمتر؛
		نفوذ بین‌المللی محدودتر \\
		\midrule
		۵ & کانادا &
		جامعه‌ی ایرانی بسیار بزرگ؛
		قانون مگنیتسکی؛
		دوزبانه (انگلیسی/فرانسوی) &
		فاصله‌ی جغرافیایی از اروپا
		و نهادهای بین‌المللی \\
		\bottomrule
		\end{longtable}
		
		\begin{principleBox}[توصیه‌ی اولیه]
		\textbf{ثبت اصلی:} هلند (لاهه) — به دلیل نزدیکی به
		دیوان کیفری بین‌المللی، دیوان دادگستری بین‌المللی،
		و سایر نهادهای بین‌المللی.
		
		\medskip
		\textbf{دفاتر نمایندگی (مجازی یا فیزیکی):}
		لندن، برلین، تورنتو — برای دسترسی به دیاسپورای ایرانی
		و نهادهای حقوقی ملی.
		
		\medskip
		\textbf{فرم حقوقی پیشنهادی:}
		بنیاد هلندی (\en{Stichting}) — ساده‌ترین و انعطاف‌پذیرترین
		فرم حقوقی برای سازمان‌های غیرانتفاعی بین‌المللی در هلند.
		\end{principleBox}
		
		\section{جمع‌بندی فصل هفتم}
		
		\begin{principleBox}[خلاصه‌ی فصل ۷]
		\begin{enumerate}[nosep, label=\textcolor{cPrimary}{\arabic*.}]
		\item ارتباط با دیوان کیفری بین‌المللی استراتژیک اما بلندمدت است
		(به دلیل عدم عضویت ایران).
		\item صلاحیت جهانی مؤثرترین مسیر
		فعلی تعقیب قضایی است — سوئد و آلمان در اولویت.
		\item هیئت حقیقت‌یاب و گزارشگر ویژه‌ی سازمان ملل
		نزدیک‌ترین شرکای نهادی هستند.
		\item سازمان‌های حقوق بشری (به‌ویژه مرکز برومند و ایران حقوق بشر)
		مکمل طبیعی دادگاه در سایه هستند.
		\item تعامل با دولت‌ها مجاز اما مشروط به قیود سخت‌گیرانه است.
		\item هلند (لاهه) بهترین گزینه‌ی کشور میزبان است.
		\end{enumerate}
		\end{principleBox}
		
		\clearpage
		






% ////////////////////////////////////////////////////////////
%  فصل ۸: نمونه‌های مشابه و مطالعه‌ی موردی
% ////////////////////////////////////////////////////////////
	\chapter{نمونه‌های مشابه و مطالعه‌ی موردی}\label{chap:cases}
	
	\begin{principleBox}[چرا مطالعه‌ی تطبیقی؟]
	ایده‌ی «دادگاه در سایه» بی‌سابقه نیست. در تاریخ معاصر،
	نمونه‌هایی وجود دارند که هر یک درس‌هایی ارزشمند
	— هم مثبت و هم منفی — به همراه دارند.
مطالعه‌ی دقیق این نمونه‌ها ما را از اختراع دوباره‌ی چرخ
بی‌نیاز می‌کند.
\end{principleBox}

	\section{نمونه‌ی ۱: دادگاه راسل (\en{Russell Tribunal})}
	
	\subsection{معرفی}
		
			\begin{noteBox}[مشخصات]
			\begin{itemize}[nosep]
			\item \textbf{سال تأسیس:} ۱۹۶۶
			\item \textbf{بنیان‌گذاران:} برتراند راسل و ژان‌پل سارتر
			\item \textbf{موضوع:} جنایات جنگی آمریکا در ویتنام
			\item \textbf{ماهیت:} دادگاه مردمی بدون ضمانت اجرا
		\item \textbf{نتیجه:} محکومیت اخلاقی آمریکا؛
	تأثیر عمیق بر افکار عمومی جهانی
	\end{itemize}
	\end{noteBox}
	
	\subsection{درس‌های مثبت}
		
		\begin{enumerate}[nosep, label=\textcolor{cGold}{\alph*.}]
		\item \textbf{اعتبار اخلاقی:}
		حضور چهره‌های شناخته‌شده‌ی جهانی (راسل، سارتر)
		به نهاد اعتبار بالایی بخشید —
		حتی بدون ضمانت اجرا.
		\item \textbf{مستندسازی ارزشمند:}
		مجموعه شهادت‌ها و مستندات دادگاه راسل هنوز
		هم یکی از مهم‌ترین منابع تاریخی جنگ ویتنام است.
		\item \textbf{تأثیر بر افکار عمومی:}
		این دادگاه نقش مؤثری در تغییر افکار عمومی اروپا و آمریکا درباره‌ی
		جنگ ویتنام ایفا کرد و فشار عمومی برای پایان جنگ را
		تقویت نمود.
		\item \textbf{الگوسازی:}
		دادگاه راسل الگویی شد که بعدها بارها تکرار شد
	(دادگاه راسل ۲ درباره‌ی آمریکای لاتین، دادگاه راسل
	درباره‌ی فلسطین، و …).
	\end{enumerate}
	
	\subsection{درس‌های منفی و محدودیت‌ها}
		
		\begin{enumerate}[nosep, label=\textcolor{cAccent}{\alph*.}]
		\item \textbf{اتهام جانب‌داری:}
		منتقدان دادگاه راسل را «دادگاه نمایشی چپ‌گرا» خواندند.
		نبودِ وکیل مدافع برای طرف آمریکایی این اتهام را تقویت کرد.
		\textbf{درس برای ما:} وکیل تسخیری (\en{Amicus Curiae})
		ضروری است.
		
		\item \textbf{فقدان فرایند تجدیدنظر:}
		آراء دادگاه راسل قطعی و غیرقابل بازنگری بودند.
		\textbf{درس برای ما:} مکانیزم بازنگری و اعتراض باید
		از ابتدا طراحی شود.
		
		\item \textbf{وابستگی به چهره‌ها:}
		پس از درگذشت راسل (۱۹۷۰)، نهاد عملاً متوقف شد.
		\textbf{درس برای ما:} نهاد باید ساختار مستقل از اشخاص داشته باشد.
		
		\item \textbf{فقدان ارتباط با نظام حقوقی رسمی:}
		آراء دادگاه راسل هرگز به هیچ نهاد حقوقی رسمی ارجاع نشد.
	\textbf{درس برای ما:} ارتباط ساختارمند با \en{ICC}،
	دادگاه‌های ملی، و \en{UJ} از ابتدا برنامه‌ریزی شود.
	\end{enumerate}
	
	\subsection{میزان شباهت به دادگاه در سایه}
	
		% ── اصلاح: عرض ستون‌ها تنظیم و هم‌ترازی بهبود یافت ──
		\begin{longtable}{p{5cm} c c}
		\toprule
		\textbf{ویژگی} & \textbf{دادگاه راسل} & \textbf{دادگاه در سایه} \\
		\midrule
		\endhead
		ماهیت مشورتی
		& \textcolor{cGold}{\checkmark} & \textcolor{cGold}{\checkmark} \\
		وکیل تسخیری
		& \textcolor{cAccent}{$\times$} & \textcolor{cGold}{\checkmark} \\
		مکانیزم بازنگری
		& \textcolor{cAccent}{$\times$} & \textcolor{cGold}{\checkmark} \\
		ارتباط با نهادهای رسمی
		& \textcolor{cAccent}{$\times$} & \textcolor{cGold}{\checkmark} \\
		ساختار مستقل از اشخاص
		& \textcolor{cAccent}{$\times$} & \textcolor{cGold}{\checkmark} \\
		مستندسازی نظام‌مند
		& جزئی & \textcolor{cGold}{\checkmark} \\
		داوری تصمیم‌های سیاسی
	& \textcolor{cAccent}{$\times$} & \textcolor{cGold}{\checkmark} \\
	\bottomrule
	\end{longtable}
	
	
	\section{نمونه‌ی ۲: دادگاه ایران (\en{Iran Tribunal}, ۲۰۱۲)}
	
	\subsection{معرفی}
		
			\begin{noteBox}[مشخصات]
			\begin{itemize}[nosep]
			\item \textbf{سال:} ۲۰۱۲ (لاهه)
			\item \textbf{موضوع:} کشتار زندانیان سیاسی ۱۳۶۷
			\item \textbf{ماهیت:} دادگاه مردمی بین‌المللی
			\item \textbf{هیئت داوران:} شامل قاضی \en{Geoffrey Robertson QC}
			و قاضی \en{John Dugard}
			(حقوق‌دان برجسته‌ی آفریقای جنوبی)
		\item \textbf{نتیجه:} اعلام وقوع جنایت علیه بشریت؛
	توصیه‌ی ارجاع به \en{ICC}
	\end{itemize}
	\end{noteBox}
	
	\subsection{درس‌های مثبت}
		
		\begin{enumerate}[nosep, label=\textcolor{cGold}{\alph*.}]
		\item \textbf{حضور قضات بین‌المللی معتبر:}
		حضور \en{Robertson} و \en{Dugard} اعتبار بالایی
		به دادگاه بخشید.
		
		\item \textbf{مستندسازی تاریخی:}
		شهادت‌های ضبط‌شده‌ی بازماندگان، آرشیو بی‌نظیری
		از کشتار ۶۷ ایجاد کرد.
		
		\item \textbf{تأثیر بر پرونده‌ی نوری:}
		بسیاری از مستنداتی که در دادگاه ایران ۲۰۱۲ ارائه شدند،
		بعدها در پرونده‌ی حمید نوری در سوئد مورد استفاده قرار گرفتند.
		\keyword{این مهم‌ترین درس است:}
		مستندسازی امروز، پرونده‌ی قضایی فردا را ممکن می‌سازد.
		
		\item \textbf{بازتاب رسانه‌ای:}
	این دادگاه برای نخستین بار موضوع کشتار ۶۷ را
	به‌صورت ساختارمند وارد رسانه‌های بین‌المللی کرد.
	\end{enumerate}
	
	\subsection{درس‌های منفی و محدودیت‌ها}
		
		\begin{enumerate}[nosep, label=\textcolor{cAccent}{\alph*.}]
		\item \textbf{رویداد‌محور بود، نه نهادمحور:}
		دادگاه ایران ۲۰۱۲ یک رویداد یکباره بود و
		ساختار دائمی نداشت. پس از پایان جلسات، فعالیت
		عملاً متوقف شد.
		\textbf{درس برای ما:} دادگاه در سایه باید ساختار
		دائمی و فعالیت مستمر داشته باشد.
		
		\item \textbf{محدود به یک موضوع:}
		فقط کشتار ۶۷ بررسی شد و سایر جرایم پوشش نیافتند.
		\textbf{درس برای ما:} صلاحیت موضوعی نهاد باید
		فراگیر و شامل تمام جرایم حکومتی باشد.
		
		\item \textbf{فقدان مکانیزم پیگیری:}
		پس از صدور رأی، هیچ مکانیزم رسمی‌ای برای
		پیگیری اجرای توصیه‌ها وجود نداشت.
		\textbf{درس برای ما:} دبیرخانه‌ی دائمی باید
		وظیفه‌ی رصد و پیگیری آراء را بر عهده داشته باشد.
		
		\item \textbf{مشارکت محدود دیاسپورا:}
	بسیاری از ایرانیان خارج از کشور از وجود دادگاه بی‌اطلاع بودند.
	\textbf{درس برای ما:} استراتژی ارتباطی قوی از ابتدا ضروری است.
	\end{enumerate}
	
	
	\section{نمونه‌ی ۳: دادگاه اویغورها (\en{Uyghur Tribunal}, ۲۰۲۱)}
	
	\subsection{معرفی}
		
			\begin{noteBox}[مشخصات]
			\begin{itemize}[nosep]
			\item \textbf{سال:} ۲۰۲۱ (لندن)
			\item \textbf{موضوع:} نسل‌کشی و جرایم علیه بشریت
			در سین‌کیانگ چین
			\item \textbf{ریاست:} سر \en{Geoffrey Nice QC}
			(دادستان سابق پرونده‌ی میلوشویچ در \en{ICTY})
		\item \textbf{ماهیت:} دادگاه مردمی مستقل
	\item \textbf{نتیجه:} اعلام وقوع نسل‌کشی
	\end{itemize}
	\end{noteBox}
	
	\subsection{درس‌های مثبت}
		
		\begin{enumerate}[nosep, label=\textcolor{cGold}{\alph*.}]
		\item \textbf{ریاست با اعتبار جهانی:}
		\en{Geoffrey Nice} به‌واسطه‌ی نقشش در
		محاکمه‌ی میلوشویچ، اعتبار بی‌نظیری به دادگاه بخشید.
		
		\item \textbf{فرایند دقیق:}
		دادگاه اویغورها جلسات استماع عمومی، بررسی شواهد
		کارشناسی، و حق دفاع (دعوت از دولت چین برای
		ارائه‌ی دفاعیات — که البته رد شد) را رعایت کرد.
		
		\item \textbf{رأی مستدل و مفصل:}
		رأی نهایی بیش از ۶۰ صفحه بود و هر ادعا با
		شواهد مشخص پشتیبانی شده بود.
		
		\item \textbf{تأثیر سیاسی فوری:}
		رأی دادگاه به‌سرعت مورد استناد پارلمان‌های
		بریتانیا، کانادا، و هلند قرار گرفت و
		در تصمیم‌های تحریمی اتحادیه‌ی اروپا و آمریکا
		نقش ایفا کرد.
		
		\item \textbf{مستقل از دولت‌ها:}
	تأمین مالی از طریق \en{crowdfunding}
	و بدون بودجه‌ی دولتی — الگوی موفق استقلال مالی.
	\end{enumerate}
	
	\subsection{درس‌های منفی و محدودیت‌ها}
		
		\begin{enumerate}[nosep, label=\textcolor{cAccent}{\alph*.}]
		\item \textbf{واکنش تهاجمی چین:}
		چین کمپین رسانه‌ای گسترده‌ای علیه دادگاه ترتیب داد.
		\textbf{درس برای ما:} باید از ابتدا برای واکنش
		تهاجمی جمهوری اسلامی آماده بود.
		
		\item \textbf{رویدادمحور:}
		مشابه دادگاه ایران ۲۰۱۲، ساختار دائمی نداشت.
		\textbf{درس برای ما:} تداوم نهادی.
		
		\item \textbf{فقدان مکانیزم اجرا:}
		رأی دادگاه اویغورها الزام‌آور نبود و
		تأثیر واقعی‌اش بسته به اراده‌ی سیاسی دولت‌ها بود.
	\textbf{درس برای ما:} تأکید بر ماهیت مشورتی
	و تمرکز بر کیفیت استدلال.
	\end{enumerate}
	
		
	\section{نمونه‌ی ۴: کمیسیون حقیقت و آشتی آفریقای جنوبی
	(\en{TRC}, ۱۹۹۶)}
	
	\subsection{معرفی}
		
			\begin{noteBox}[مشخصات]
			\begin{itemize}[nosep]
			\item \textbf{سال:} ۱۹۹۶–۲۰۰۲
			\item \textbf{موضوع:} جرایم دوره‌ی آپارتاید
			\item \textbf{ریاست:} اسقف دزموند توتو
			\item \textbf{ماهیت:} کمیسیون رسمی دولتی پس از گذار
			\item \textbf{ویژگی کلیدی:} ترکیب حقیقت‌یابی با عفو مشروط
		\item \textbf{نتیجه:} ثبت بیش از ۲۱٬۰۰۰ شهادت؛
	عفو مشروط ۱٬۵۰۰ نفر از عاملان جرایم
	\end{itemize}
	\end{noteBox}
	
	\subsection{درس‌های مثبت}
		
		\begin{enumerate}[nosep, label=\textcolor{cGold}{\alph*.}]
		\item \textbf{مدل عفو مشروط:}
		عفو فقط در ازای اعتراف کامل و علنی به جرایم
		اعطا می‌شد. این مدل هم‌زمان حقیقت‌یابی و
		آشتی ملی را ممکن کرد.
		\textbf{درس برای ما:} دادگاه در سایه می‌تواند
		از هم‌اکنون چارچوب عفو مشروط را طراحی کند
		و به مقام‌های حکومتی پیام دهد:
		«اگر اکنون همکاری کنید، عفو مشروط ممکن خواهد بود.»
		
		\item \textbf{مشارکت قربانیان:}
		\en{TRC} فضایی برای شنیده شدن صدای قربانیان فراهم کرد.
		این فرایند خودش ارزش درمانی داشت — حتی بدون مجازات عاملان.
		\textbf{درس برای ما:} جلسات استماع عمومی نهاد باید
		فضای کرامت‌مند شنیده شدن برای قربانیان فراهم کنند.
		
		\item \textbf{مستندسازی جامع:}
		گزارش نهایی \en{TRC} (۷ جلد) یکی از جامع‌ترین
		مستندات تاریخی درباره‌ی آپارتاید است.
		\textbf{درس برای ما:} بایگانی نظام‌مند از ابتدا.
		
		\item \textbf{رهبری اخلاقی:}
		حضور دزموند توتو — چهره‌ای اخلاقی با اعتبار جهانی —
		مشروعیت فرایند را تضمین کرد.
	\textbf{درس برای ما:} اهمیت جذب چهره‌های
	اخلاقی‌ـ‌آکادمیک معتبر.
	\end{enumerate}
	
	\subsection{درس‌های منفی و محدودیت‌ها}
		
		\begin{enumerate}[nosep, label=\textcolor{cAccent}{\alph*.}]
		\item \textbf{نارضایتی بخشی از قربانیان:}
		برخی قربانیان احساس کردند عفو مشروط به معنای
		فرار عاملان از عدالت است.
		\textbf{درس برای ما:} چارچوب عفو باید شرایط
		بسیار سخت‌گیرانه‌ای داشته باشد و نظرات مخالف
		ثبت شوند.
		
		\item \textbf{تفاوت بنیادین با شرایط ایران:}
		\en{TRC} \textbf{پس از} گذار تشکیل شد و
		حمایت دولت جدید را داشت.
		دادگاه در سایه \textbf{پیش از} گذار و بدون
		حمایت دولتی تشکیل می‌شود.
		\textbf{درس برای ما:} نهاد نباید خود را \en{TRC}
		بنامد یا ادعای مشابه داشته باشد.
		باید صادقانه بگوید: «ما پیش‌درآمد \en{TRC} هستیم،
		نه خود آن.»
		
		\item \textbf{اجرای ناقص توصیه‌ها:}
		بسیاری از توصیه‌های \en{TRC} درباره‌ی
		جبران خسارت قربانیان هرگز اجرا نشدند.
	\textbf{درس برای ما:} انتظارات واقع‌بینانه
	درباره‌ی تأثیرگذاری.
	\end{enumerate}
	
	
	\section{نمونه‌ی ۵: فرانسه‌ی آزاد (\en{Free France}, ۱۹۴۰–۱۹۴۴)}
	
	\subsection{معرفی}
		
			\begin{noteBox}[مشخصات]
			\begin{itemize}[nosep]
			\item \textbf{سال:} ۱۹۴۰–۱۹۴۴
			\item \textbf{بنیان‌گذار:} ژنرال شارل دوگل
			\item \textbf{ماهیت:} حکومت در تبعید با ساختار شبه‌دولتی کامل
			شامل نظام قضایی سایه
			\item \textbf{ویژگی کلیدی:} ایجاد نهادهای حکومتی موازی
			پیش از آزادسازی فرانسه
			\item \textbf{نتیجه:} انتقال سریع قدرت پس از آزادسازی؛
		اما همراه با «پاکسازی‌های وحشیانه»
	(\en{épuration sauvage})
	\end{itemize}
	\end{noteBox}
	
	\subsection{درس‌های مثبت}
		
		\begin{enumerate}[nosep, label=\textcolor{cGold}{\alph*.}]
		\item \textbf{نهادسازی پیش از پیروزی:}
		فرانسه‌ی آزاد قوانین، دادگاه‌ها، و ساختارهای اداری
		را پیش از آزادسازی آماده کرد.
		این آمادگی نهادی انتقال قدرت را تسهیل کرد.
		\textbf{درس برای ما:} دقیقاً همین منطق —
		آمادگی نهادی پیش از گذار.
		
		\item \textbf{مشروعیت‌سازی تدریجی:}
		فرانسه‌ی آزاد با کسب شناسایی تدریجی از سوی
		متفقین، مشروعیت خود را ساخت.
	\textbf{درس برای ما:} اعتبار بین‌المللی
	گام‌به‌گام ساخته می‌شود.
	\end{enumerate}
	
	\subsection{درس‌های منفی — هشدار بسیار مهم}
		
		\begin{warningBox}[هشدار تاریخی]
		پس از آزادسازی فرانسه، حدود ۹٬۰۰۰ نفر بدون محاکمه
		و در «پاکسازی‌های وحشیانه» کشته شدند. بسیاری از این
		قربانیان بی‌گناه بودند یا جرم‌شان ناچیز بود.
		
		\medskip
		\textbf{علت:} نبودِ سازوکار قضایی آماده و کافی در لحظه‌ی
		انتقال قدرت سبب شد مردم خودسرانه «عدالت» اجرا کنند.
		
		\medskip
		\textbf{مهم‌ترین درس برای دادگاه در سایه:}
		یکی از اصلی‌ترین دلایل تأسیس این نهاد، دقیقاً
		\keyword{پیش‌گیری از عدالت وحشیانه‌ی پس از سقوط}
	جمهوری اسلامی است. اگر مردم بدانند فرایند قضایی آماده است،
	احتمال خودسری کاهش می‌یابد.
	\end{warningBox}
	
	
	\section{نمونه‌ی ۶: کنگره‌ی ملی آفریقا (\en{ANC}) — ساختارهای سایه}
	
	\subsection{معرفی}
		
			\begin{noteBox}[مشخصات]
			\begin{itemize}[nosep]
			\item \textbf{دوره:} ۱۹۶۰–۱۹۹۰ (دوره‌ی ممنوعیت)
			\item \textbf{ماهیت:} جنبش آزادی‌بخش با ساختارهای شبه‌دولتی
			در تبعید
			\item \textbf{ویژگی‌ها:} تأسیس نظام آموزشی، خدمات اجتماعی،
		و سازوکارهای قضایی داخلی
	\item \textbf{نتیجه:} گذار مسالمت‌آمیز ۱۹۹۴
	\end{itemize}
	\end{noteBox}
	
	\subsection{درس‌ها}
		
		\begin{enumerate}[nosep, label=\textcolor{cGold}{\alph*.}]
		\item \textbf{نهادسازی در تبعید ممکن است:}
		\en{ANC} نشان داد که حتی یک جنبش ممنوع
		می‌تواند ساختارهای نهادی مؤثر بسازد.
		
		\item \textbf{مشروعیت بین‌المللی:}
		\en{ANC} با جذب حمایت سازمان ملل، کشورهای
		اسکاندیناوی، و جنبش ضدآپارتاید جهانی،
		مشروعیت خود را ساخت.
		
		\item \textbf{مدیریت تحریم‌ها:}
		\en{ANC} تجربه‌ی بی‌نظیری در حمایت از تحریم‌های
		اقتصادی، ورزشی، و فرهنگی دارد —
		از جمله مدیریت هزینه‌های انسانی آن.
	\textbf{درس برای ما:} مطالعه‌ی دقیق مدل \en{ANC}
	در توجیه تحریم‌ها و مدیریت هزینه‌ها.
	\end{enumerate}
		
		\begin{warningBox}[هشدار]
		\en{ANC} در دوره‌ی تبعید مرتکب نقض حقوق بشر در
		اردوگاه‌های خود نیز شد (\en{Quatro camp}).
	\textbf{درس:} حتی جنبش‌های آزادی‌بخش می‌توانند
	عاملان نقض حقوق بشر شوند. نظارت مستقل ضروری است.
	\end{warningBox}
	
	
	\section{نمونه‌ی ۷: دادگاه خلق‌ها (\en{Permanent Peoples' Tribunal})}
	
	\subsection{معرفی}
		
			\begin{noteBox}[مشخصات]
			\begin{itemize}[nosep]
			\item \textbf{سال تأسیس:} ۱۹۷۹ (بولونیا، ایتالیا)
			\item \textbf{بنیان‌گذار:} لئولیو باسو (سناتور ایتالیایی)
			\item \textbf{ماهیت:} نهاد دائمی بین‌المللی مردمی
			\item \textbf{موضوعات:} بیش از ۴۵ جلسه در
			نقض حقوق بشر در سراسر جهان
		\item \textbf{ویژگی کلیدی:} تنها نهاد دادگاه مردمی
	با ساختار دائمی
	\end{itemize}
	\end{noteBox}
	
	\subsection{درس‌های مثبت}
		
		\begin{enumerate}[nosep, label=\textcolor{cGold}{\alph*.}]
		\item \textbf{تداوم نهادی:}
		بر خلاف دادگاه راسل و دادگاه ایران ۲۰۱۲،
		\en{PPT} ساختار دائمی دارد و بیش از ۴ دهه
		فعالیت مستمر داشته.
		\textbf{درس:} مدل سازمانی \en{PPT} الگوی
		خوبی برای تداوم است.
		
		\item \textbf{انعطاف موضوعی:}
		\en{PPT} درباره‌ی موضوعات بسیار متنوع
		(از نسل‌کشی تا حقوق محیط زیست) رسیدگی کرده.
	\textbf{درس:} دادگاه در سایه نیز باید صلاحیت
	موضوعی گسترده‌ای داشته باشد.
	\end{enumerate}
	
	
	\section{خلاصه‌ی تطبیقی نمونه‌ها}
	
		% ── اصلاح: جدول در landscape قرار گرفت تا سرریز نکند ──
			\begin{landscape}
			\begin{longtable}{r p{3cm} p{2cm} p{2cm} p{2cm} p{2cm} p{3.5cm}}
			\toprule
			\textbf{\RL{ردیف}} & \textbf{\RL{نمونه}} & \textbf{\RL{دائمی؟}} &
			\textbf{\RL{وکیل تسخیری؟}} & \textbf{\RL{فرایند بازنگری؟}} &
			\textbf{\RL{ارتباط رسمی؟}} & \textbf{\RL{تأثیر واقعی}} \\
			\midrule
			\endhead
			۱ & {\RL{دادگاه راسل}} & {\RL{خیر}} & {\RL{خیر}} & {\RL{خیر}}
			& {\RL{خیر}} & {\RL{بالا (افکار عمومی)}} \\
			۲ & {\RL{دادگاه ایران ۲۰۱۲}} & {\RL{خیر}} & {\RL{خیر}} & {\RL{خیر}}
			& {\RL{ضعیف}} & {\RL{متوسط}} \\
			۳ & {\RL{دادگاه اویغور}} & {\RL{خیر}} & {\RL{محدود}} & {\RL{خیر}}
			& {\RL{متوسط}} & {\RL{بالا}} \\
			۴ & \en{TRC} & {\RL{خیر}}\footnote{\RL{رسمی بود اما پس از پایان مأموریت منحل شد.}} &
			{\RL{بله}} & {\RL{بله}} & {\RL{بالا}} & {\RL{بسیار بالا}} \\
			۵ & {\RL{فرانسه‌ی آزاد}} & {\RL{بله}} & {\RL{محدود}} & {\RL{محدود}}
			& {\RL{بالا}} & {\RL{بسیار بالا}} \\
			۶ & {\RL{سایه}} \en{ANC} & {\RL{بله}} & {\RL{خیر}} & {\RL{محدود}}
			& {\RL{بالا}} & {\RL{بسیار بالا}} \\
			۷ & \en{PPT} & {\RL{بله}} & {\RL{محدود}} & {\RL{بله}}
			& {\RL{ضعیف}} & {\RL{متوسط}} \\
			\midrule
			& \textbf{\RL{دادگاه در سایه (هدف)}} & \textbf{\RL{بله}} &
		\textbf{\RL{بله}} & \textbf{\RL{بله}} & \textbf{\RL{بالا}} & \textbf{—} \\
	\bottomrule
	\end{longtable}
	\end{landscape}
	
	\section{جمع‌بندی فصل هشتم}
		
			\begin{principleBox}[خلاصه‌ی فصل ۸]
			\begin{enumerate}[nosep, label=\textcolor{cPrimary}{\arabic*.}]
			\item هیچ‌یک از نمونه‌های موجود دقیقاً مشابه مدل پیشنهادی نیست،
			اما هر یک درس‌های ارزشمندی دارد.
				\item \textbf{مهم‌ترین درس‌ها:}
				\begin{itemize}[nosep]
				\item وکیل تسخیری ضروری است (درس دادگاه راسل)؛
				\item تداوم نهادی حیاتی است (درس دادگاه ایران ۲۰۱۲)؛
				\item استقلال مالی ممکن است (درس دادگاه اویغور)؛
				\item نهادسازی پیش از گذار از پاکسازی وحشیانه
			جلوگیری می‌کند (درس فرانسه‌ی آزاد)؛
			\item حتی جنبش‌ها می‌توانند ناقض حقوق بشر شوند (درس \en{ANC}).
			\end{itemize}
		\item دادگاه در سایه تلاش می‌کند نقاط قوت همه‌ی این نمونه‌ها را
	ترکیب و نقاط ضعف آن‌ها را اصلاح کند.
	\end{enumerate}
	\end{principleBox}
	
	\clearpage
	
	
	% ////////////////////////////////////////////////////////////
	%  فصل ۹: هضم و ادغام در نظام قضایی پس از گذار
	% ////////////////////////////////////////////////////////////
	\chapter{هضم و ادغام در نظام قضایی پس از گذار}\label{chap:transition}
		
		\begin{principleBox}[اصل بنیادین]
		دادگاه در سایه یک \keyword{نهاد موقت} است.
		هدف نهایی آن \keyword{منحل کردن خودش}
		در لحظه‌ای است که نظام قضایی رسمی و دموکراتیک
		در ایران تأسیس شود. اما «انحلال» نباید به معنای
	«از بین رفتن» باشد — بلکه به معنای \keyword{تحول و ادغام}
	در ساختارهای رسمی است.
	\end{principleBox}
	
	\section{سه مرحله‌ی گذار}
		
		\begin{figure}[H]
			\centering
			\begin{tikzpicture}[
				node distance=0.5cm and 2cm,
				phase/.style={
				rectangle, rounded corners=5pt, draw=cPrimary, thick,
			fill=cLight, text width=4.5cm, minimum height=2.5cm,
			align=center, font=\small
			},
			arrow/.style={-{Stealth[length=5pt]}, very thick, cGold}
				]
				\node[phase] (pre) {
				\textbf{\RL{مرحله‌ی ۱}}\\
					\textbf{\RL{پیش از گذار}}\\[4pt]
					{\footnotesize \RL{وضعیت فعلی}\\
			\RL{نهاد سایه فعال}\\
			\RL{بدون قدرت اجرایی}}
				};
				\node[phase, left=of pre] (trans) {
				\textbf{\RL{مرحله‌ی ۲}}\\
					\textbf{\RL{دوره‌ی گذار}}\\[4pt]
					{\footnotesize \RL{سقوط جمهوری اسلامی}\\
			\RL{تا تأسیس نظام جدید}\\
			\RL{(تخمین: ۶ تا ۲۴ ماه)}}
				};
				\node[phase, left=of trans] (post) {
				\textbf{\RL{مرحله‌ی ۳}}\\
					\textbf{\RL{پس از تثبیت}}\\[4pt]
					{\footnotesize \RL{نظام قضایی رسمی}\\
			\RL{دموکراتیک مستقر}\\
			\RL{انحلال/ادغام نهاد}}
			};
			
		\draw[arrow] (pre) -- (trans);
		\draw[arrow] (trans) -- (post);
	\end{tikzpicture}
	\caption{\RL{سه مرحله‌ی گذار و تحول نهاد}}
	\end{figure}
	
	\section{مرحله‌ی ۱: پیش از گذار (وضعیت فعلی)}
	
	این مرحله در فصل‌های قبل به‌تفصیل شرح داده شد. خلاصه:
		
		\begin{itemize}[nosep]
		\item نهاد ماهیت مشورتی دارد.
		\item کیفرخواست‌های سایه صادر می‌کند.
		\item آراء مشورتی درباره‌ی تصمیم‌های سیاسی صادر می‌کند.
	\item مستندات را برای آینده بایگانی می‌کند.
	\item با نهادهای بین‌المللی همکاری می‌کند.
	\end{itemize}
	
	\section{مرحله‌ی ۲: دوره‌ی گذار (بحرانی‌ترین مرحله)}
		
		\begin{warningBox}[هشدار: خطر عدالت وحشیانه]
		تجربه‌ی تاریخی نشان می‌دهد که خطرناک‌ترین لحظه،
		\textbf{فاصله‌ی بین سقوط نظام قبلی و تثبیت نظام جدید} است.
			در این فاصله:
			\begin{itemize}[nosep]
			\item نظام قضایی قبلی از هم پاشیده؛
			\item نظام قضایی جدید هنوز تأسیس نشده؛
		\item خشم عمومی در اوج است؛
	\item ریسک انتقام‌جویی بدون فرایند بسیار بالاست.
	\end{itemize}
	\end{warningBox}
	
	\subsection{نقش دادگاه در سایه در دوره‌ی گذار}
		
			\begin{principleBox}[شش وظیفه‌ی کلیدی]
			\begin{enumerate}[nosep, label=\textcolor{cPrimary}{\arabic*.}]
			\item \textbf{ارائه‌ی فهرست اولویت‌دار بازداشت:}
			بر اساس کیفرخواست‌های سایه صادرشده،
			فهرستی اولویت‌بندی‌شده از افرادی که باید
			بازداشت و محاکمه شوند تهیه و به
			دولت موقت ارائه شود.
			\keyword{این فهرست مانع بازداشت‌های خودسرانه می‌شود.}
			
			\item \textbf{فراخوان آرامش و اعتماد:}
			نهاد با انتشار بیانیه‌ی عمومی اعلام می‌کند:
			«مستندات آماده است. عدالت اجرا خواهد شد.
			لطفاً از خودسری بپرهیزید.»
			
			\item \textbf{مشاوره به دولت موقت:}
				ارائه‌ی توصیه‌های فنی درباره‌ی:
				\begin{itemize}[nosep]
				\item تأسیس کمیسیون حقیقت‌یاب؛
				\item طراحی دادگاه‌های انتقالی؛
			\item سازوکار عفو مشروط؛
			\item اصلاح نظام قضایی.
			\end{itemize}
			
			\item \textbf{حفاظت از آرشیو:}
			تمام مستندات انباشته‌شده در اختیار مراجع رسمی
			جدید قرار می‌گیرد.
			
			\item \textbf{نظارت بر فرایند عدالت انتقالی:}
			نهاد در دوره‌ی گذار از ماهیت «مشورتی» به
			ماهیت «ناظر» تحول می‌یابد و بر رعایت
			حقوق بشر در فرایند انتقالی نظارت می‌کند.
			
			\item \textbf{حمایت از قربانیان:}
		ارائه‌ی مشاوره‌ی حقوقی و روان‌شناختی به
	قربانیان و خانواده‌های آن‌ها در فرایند انتقالی.
	\end{enumerate}
	\end{principleBox}
	
	\subsection{چارچوب عفو مشروط پیشنهادی}
		
		\begin{principleBox}[شرایط عفو مشروط]
		دادگاه در سایه از هم‌اکنون چارچوب زیر را پیشنهاد می‌دهد
		تا هم امکان عدالت ترمیمی فراهم شود و هم
		انگیزه‌ای برای همکاری مقام‌های حکومتی ایجاد گردد:
			
			\begin{enumerate}[nosep, label=\textcolor{cGold}{\arabic*.}]
			\item \textbf{اعتراف کامل و علنی:}
			متقاضی عفو باید تمام جرایم خود را
			به‌صورت کامل و علنی اعتراف کند.
			
			\item \textbf{انگیزه‌ی سیاسی:}
			جرم باید انگیزه‌ی سیاسی (نه شخصی یا مالی) داشته باشد.
			\warn{فساد مالی مشمول عفو نیست.}
			
			\item \textbf{تناسب:}
			میزان جرم باید در حدی باشد که عفو مشروط
				نسبت به آن «نامتناسب» به نظر نرسد.
			\warn{عاملان نسل‌کشی و شکنجه‌ی سیستماتیک
			مشمول عفو نیستند.}
			
			\item \textbf{همکاری فعال:}
			متقاضی باید اطلاعات مفید درباره‌ی ساختار
			سرکوب و سایر عاملان ارائه دهد.
			
			\item \textbf{رضایت قربانیان:}
			نظر قربانیان و خانواده‌های آن‌ها باید
			شنیده و در تصمیم‌گیری لحاظ شود
			(هرچند حق وتوی مطلق ندارند).
			
			\item \textbf{جبران خسارت:}
			متقاضی عفو باید به صندوق جبران خسارت
		قربانیان کمک مالی کند
	(متناسب با توانایی مالی خود).
	\end{enumerate}
	\end{principleBox}
	
	\section{مرحله‌ی ۳: پس از تثبیت نظام جدید}
	
	\subsection{سناریوهای ادغام و انحلال}
		
		\begin{longtable}{r p{3cm} p{8.5cm}}
		\toprule
		\textbf{سناریو} & \textbf{عنوان} & \textbf{شرح} \\
		\midrule
		\endhead
		\textbf{الف} & ادغام کامل &
		تمام مستندات، آراء، و نیروی انسانی نهاد
		در نظام قضایی رسمی جدید ادغام می‌شوند.
		اعضای هیئت عالی می‌توانند در کمیسیون حقیقت‌یاب
		یا دادگاه‌های انتقالی خدمت کنند.
		\keyword{محتمل‌ترین سناریو.} \\
		\midrule
		\textbf{ب} & تحول به نهاد ناظر &
		نهاد به‌جای انحلال، به یک نهاد نظارتی مستقل
		(مانند \en{ombudsman} یا دیده‌بان حقوق بشر ملی)
		تبدیل شود.
		\keyword{سناریوی مطلوب اگر نظام جدید استقبال کند.} \\
		\midrule
		\textbf{ج} & انحلال با بایگانی &
		نهاد منحل شود اما آرشیو آن به‌صورت عمومی
		و دائمی در دسترس باشد — مشابه آرشیو \en{TRC}
		آفریقای جنوبی یا آرشیو \en{Stasi} آلمان شرقی.
		\keyword{حداقل قابل قبول.} \\
		\midrule
		\textbf{د} & تداوم مستقل &
		اگر نظام جدید از ادغام استقبال نکند یا
		خود دچار نقض حقوق بشر شود، نهاد به‌صورت
		مستقل به فعالیت ادامه دهد.
	\keyword{سناریوی احتیاطی.} \\
	\bottomrule
	\end{longtable}
	
	\subsection{حفاظت از آرشیو}
		
			\begin{principleBox}[سیاست بایگانی بلندمدت]
			\begin{enumerate}[nosep, label=\textcolor{cPrimary}{\arabic*.}]
			\item \textbf{نسخه‌ی فیزیکی:}
			در حداقل دو کشور مختلف نگهداری شود.
			\item \textbf{نسخه‌ی دیجیتال:}
			با رمزنگاری قوی و بر روی سرورهای مستقل
			در حداقل سه موقعیت جغرافیایی متفاوت.
			\item \textbf{دسترسی عمومی:}
			آراء منتشرشده‌ی نهاد برای همیشه
			به‌صورت عمومی در دسترس باشند.
			\item \textbf{محرمانگی شاهدان:}
			اطلاعات هویتی شاهدان تا ۳۰ سال پس از انحلال
			(یا تا درخواست خود شاهد) محرمانه بمانند.
			\item \textbf{واگذاری به نهاد رسمی:}
			در سناریوهای الف و ب، آرشیو رسماً به
			نهاد جانشین تحویل داده شود.
			\item \textbf{پشتیبان بین‌المللی:}
			یک نسخه‌ی کامل از آرشیو به نهاد بین‌المللی معتبر
			(مانند \en{OHCHR} یا \en{Swisspeace})
		سپرده شود تا در صورت بروز مشکلات
	سیاسی داخلی، محفوظ بماند.
	\end{enumerate}
	\end{principleBox}
	
	\section{جمع‌بندی فصل نهم}
		
			\begin{principleBox}[خلاصه‌ی فصل ۹]
			\begin{enumerate}[nosep, label=\textcolor{cPrimary}{\arabic*.}]
			\item نهاد از ابتدا با آگاهی به موقتی بودن خود طراحی شده.
			\item بحرانی‌ترین نقش آن در \keyword{دوره‌ی گذار} است:
			جلوگیری از عدالت وحشیانه و ارائه‌ی زیرساخت
			حقوقی آماده.
			\item چارچوب عفو مشروط از هم‌اکنون طراحی شود
			تا هم ابزار عدالت ترمیمی باشد و هم
			انگیزه‌ی همکاری مقام‌ها.
			\item چهار سناریوی انحلال/ادغام با اولویت‌بندی ارائه شده.
		\item حفاظت از آرشیو اولویت مطلق دارد —
	حافظه‌ی تاریخی نباید از بین برود.
	\end{enumerate}
	\end{principleBox}
	
	\clearpage
	
	
	% ////////////////////////////////////////////////////////////
	%  فصل ۱۰: افراد پیشنهادی
	% ////////////////////////////////////////////////////////////
	\chapter{افراد پیشنهادی}\label{chap:persons}
		
		\begin{warningBox}[تذکر مهم]
		فهرست زیر صرفاً \textbf{پیشنهاد اولیه} بر اساس سوابق
		آکادمیک و حرفه‌ای شناخته‌شده است. درج نام هیچ‌کس
		به‌معنای تأیید یا مشارکت آن شخص نیست.
		هرگونه دعوت رسمی منوط به تأسیس نهاد و طی فرایند
		گزینش مستقل (فصل \ref{chap:org}) خواهد بود.
		
		\medskip
			معیارهای انتخاب:
			\begin{enumerate}[nosep, label=\textcolor{cPrimary}{\arabic*.}]
			\item سابقه‌ی آکادمیک یا حرفه‌ای شناخته‌شده در سطح بین‌المللی؛
			\item تخصص مرتبط (حقوق بین‌الملل، حقوق بشر، علوم سیاسی،
			جامعه‌شناسی، حقوق کیفری)؛
			\item استقلال از جناح‌بندی‌های سیاسی اپوزیسیون
			(تا حد ممکن)؛
		\item اعتبار اخلاقی و سابقه‌ی فعالیت حقوق بشری؛
	\item تنوع جنسیتی، نسلی، قومی و تخصصی.
	\end{enumerate}
	\end{warningBox}
	
	\section{دسته‌ی اول: حقوق‌دانان و متخصصان حقوق بین‌الملل}
	
	% ── اصلاح: \\ درون سلول‌های longtable باعث خطا می‌شود ──
	% ── از \newline یا \par به‌جای \\ استفاده شد ──
		% ── همچنین عرض ستون‌ها تنظیم شد ──
		\begin{longtable}{r p{3.2cm} p{3.8cm} p{4.5cm}}
		\toprule
		\textbf{ردیف} & \textbf{نام} & \textbf{وابستگی نهادی} &
		\textbf{تخصص و سابقه‌ی مرتبط} \\
		\midrule
		\endhead
		۱ & \textbf{پیام اخوان}
		\en{(Payam Akhavan)} &
		استاد حقوق، دانشگاه مک‌گیل
		(\en{McGill University})، کانادا &
		دادستان سابق \en{ICTY} و \en{ICTR}؛
		مشاور ارشد \en{OHCHR}؛
		نویسنده‌ی کتاب‌های متعدد درباره‌ی
		عدالت بین‌المللی؛ تجربه‌ی مستقیم
			در تعقیب جرایم علیه بشریت.
		\keyword{مناسب‌ترین فرد برای ریاست
		یا عضویت ارشد هیئت عالی.} \\
		\midrule
		۲ & \textbf{لیلا سادات}
		\en{(Leila Sadat)} &
		استاد حقوق بین‌الملل،
			دانشگاه واشنگتن
		(\en{Washington University
		in St.\ Louis}) &
		مشاور ویژه‌ی دادستان \en{ICC}؛
			مدیر مرکز مطالعات جرایم
		بین‌المللی (\en{Whitney R.\ Harris
		World Law Institute})؛
			متخصص اساسنامه‌ی رم.
		\keyword{تخصص ویژه در ارتباط
		با \en{ICC}.} \\
		\midrule
		۳ & \textbf{هادی قائمی}
		\en{(Hadi Ghaemi)} &
		مدیر اجرایی مرکز حقوق بشر
			در ایران
		(\en{Center for Human Rights
		in Iran}, نیویورک) &
		بیش از ۲۰ سال سابقه‌ی مستندسازی
		نقض حقوق بشر در ایران؛
		سابقه‌ی همکاری با \en{HRW}
		و \en{OHCHR}.
		\keyword{تخصص در مستندسازی.} \\
		\midrule
		۴ & \textbf{شیرین عبادی}
		\en{(Shirin Ebadi)} &
		برنده‌ی جایزه‌ی صلح نوبل ۲۰۰۳؛
		حقوق‌دان و قاضی سابق ایران &
		اولین قاضی زن ایران؛
		بنیان‌گذار مرکز مدافعان حقوق بشر؛
		شناخته‌شده‌ترین چهره‌ی حقوق بشری
		ایران در سطح جهانی.
			\keyword{اعتبار اخلاقی و نمادین بالا.}
			\warn{ملاحظه: مواضع سیاسی ممکن
			است بحث‌برانگیز باشد — نقش نمادین
		(شورای ناظران) مناسب‌تر از نقش
		اجرایی است.} \\
		\midrule
			۵ & \textbf{محمد‌حسین برکتی}
		\en{(Mohammad-Hossein
		Barakati)} &
		وکیل دادگستری و محقق حقوق
		بین‌الملل کیفری &
		سابقه‌ی همکاری با پرونده‌ی
		حمید نوری در سوئد؛
		تخصص در \en{Universal Jurisdiction}؛
		آشنایی مستقیم با نظام قضایی اسکاندیناوی. \\
		\midrule
		۶ & \textbf{کاوه موسوی}
		\en{(Kaveh Moussavi)} &
		وکیل دادگستری بین‌المللی،
		مستقر در لندن &
		سابقه‌ی پیگیری پرونده‌های
		کلاهبرداری و فساد بین‌المللی
		مرتبط با ایران؛
		آشنایی با نظام حقوقی بریتانیا
			و سازوکارهای تحریم.
		\keyword{مناسب برای شعبه‌ی
		داوری تصمیم‌های سیاسی.} \\
		\midrule
		۷ & \textbf{نسرین ستوده}
		\en{(Nasrin Sotoudeh)} &
		وکیل دادگستری و مدافع
		حقوق بشر (داخل ایران) &
		برنده‌ی جایزه‌ی ساخاروف (۲۰۱۲)؛
		سابقه‌ی دفاع از زندانیان سیاسی
			و حقوق زنان.
		\warn{در صورت خروج از ایران
		یا از طریق کانال امن مشارکت ممکن.}
	\keyword{نماد جامعه‌ی مدنی داخل.} \\
	\bottomrule
	\end{longtable}
	
	\section{دسته‌ی دوم: متخصصان علوم سیاسی و روابط بین‌الملل}
		
		\begin{longtable}{r p{3.2cm} p{3.8cm} p{4.5cm}}
		\toprule
		\textbf{ردیف} & \textbf{نام} & \textbf{وابستگی نهادی} &
		\textbf{تخصص و سابقه‌ی مرتبط} \\
		\midrule
		\endhead
		۸ & \textbf{عباس میلانی}
		\en{(Abbas Milani)} &
		استاد علوم سیاسی،
		دانشگاه استنفورد
		(\en{Stanford University}) &
		مدیر برنامه‌ی مطالعات ایران
		در هوور؛ مورخ سیاسی ایران؛
		نویسنده‌ی آثار متعدد درباره‌ی
		تحولات سیاسی ایران.
		\keyword{تحلیل سیاسی و تاریخی.} \\
		\midrule
		۹ & \textbf{مهرانگیز کار}
		\en{(Mehrangiz Kar)} &
		حقوق‌دان و نویسنده؛
		پژوهشگر مهمان دانشگاه هاروارد &
		متخصص حقوق زنان و آزادی‌های
		مدنی در ایران؛ سابقه‌ی زندان
		سیاسی؛ نویسنده‌ی ده‌ها کتاب
			و مقاله‌ی حقوقی.
		\keyword{تخصص در حقوق زنان
		و جامعه‌ی مدنی.} \\
		\midrule
		۱۰ & \textbf{رویا برومند}
		\en{(Roya Boroumand)} &
		مدیر اجرایی مرکز
			عبدالرحمن برومند
		(\en{Abdorrahman Boroumand
		Center}) &
		بیش از ۲۰ سال مستندسازی
		نقض حقوق بشر در ایران؛
		بانک اطلاعاتی اعدام‌ها؛
		نزدیک‌ترین تجربه به مأموریت
			مستندسازی نهاد.
		\keyword{مناسب برای شعبه‌ی
		مستندسازی جرایم.} \\
		\midrule
		۱۱ & \textbf{لادن برومند}
		\en{(Ladan Boroumand)} &
		مورخ و پژوهشگر؛
		هم‌بنیان‌گذار مرکز برومند &
		متخصص تاریخ جنبش‌های
		دموکراتیک ایران؛ دکترای تاریخ
			از \en{EHESS} پاریس.
		\keyword{بُعد تاریخی و حافظه‌ی
		جمعی.} \\
		\midrule
		۱۲ & \textbf{محمود امیری‌مقدم}
		\en{(Mahmood Amiry-Moghaddam)} &
		مدیر سازمان حقوق بشر ایران
		(\en{Iran Human Rights},
		نروژ) &
		فعال‌ترین نهاد رصد اعدام در ایران؛
		برنده‌ی جوایز بین‌المللی متعدد؛
		تجربه‌ی گسترده در مستندسازی
			و لابی‌گری بین‌المللی.
		\keyword{تخصص در رصد و
	مستندسازی اعدام.} \\
	\bottomrule
	\end{longtable}
	
	\section{دسته‌ی سوم: جامعه‌شناسان، فیلسوفان و روشنفکران}
		
		\begin{longtable}{r p{3.2cm} p{3.8cm} p{4.5cm}}
		\toprule
		\textbf{ردیف} & \textbf{نام} & \textbf{وابستگی نهادی} &
		\textbf{تخصص و سابقه‌ی مرتبط} \\
		\midrule
		\endhead
		۱۳ & \textbf{آصف بیات}
		\en{(Asef Bayat)} &
		استاد جامعه‌شناسی،
		دانشگاه ایلینوی
		(\en{University of Illinois}) &
		متخصص جنبش‌های اجتماعی
		خاورمیانه؛ نویسنده‌ی
			\en{Life as Politics} و
		\en{Revolution without
			Revolutionaries}.
		\keyword{تحلیل جامعه‌شناختی
		جنبش‌ها.} \\
		\midrule
		۱۴ & \textbf{نیره توحیدی}
		\en{(Nayereh Tohidi)} &
		استاد مطالعات زنان و
		جنسیت، دانشگاه
		ایالتی کالیفرنیا
		(\en{Cal State Northridge}) &
		متخصص مطالعات جنسیتی،
		فمینیسم ایرانی، و جنبش زنان؛
			مشاور سابق \en{UNDP}.
		\keyword{بُعد جنسیتی عدالت
		انتقالی.} \\
		\midrule
		۱۵ & \textbf{رامین جهانبگلو}
		\en{(Ramin Jahanbegloo)} &
		فیلسوف و استاد علوم سیاسی،
		دانشگاه تورنتو (سابقاً) &
		متخصص فلسفه‌ی عدم خشونت؛
		مصاحبه‌کننده با \en{Isaiah Berlin}؛
			سابقه‌ی زندان سیاسی.
		\keyword{بُعد فلسفی و اخلاقی
		دفاع مشروع و خشونت‌پرهیزی.} \\
		\midrule
		۱۶ & \textbf{فرهاد خسروخاور}
		\en{(Farhad Khosrokhavar)} &
		جامعه‌شناس، مدیر مطالعات
		\en{EHESS} پاریس &
		متخصص جامعه‌شناسی خشونت،
		رادیکالیسم، و زندان؛
		نویسنده‌ی آثار متعدد درباره‌ی
			جامعه‌ی ایران.
		\keyword{تحلیل خشونت سازمان‌یافته
		و رادیکالیسم.} \\
		\midrule
		۱۷ & \textbf{مهدی خلجی}
		\en{(Mehdi Khalaji)} &
		پژوهشگر ارشد،
		مؤسسه‌ی واشنگتن
		(\en{Washington Institute}) &
		متخصص ساختار قدرت در
		جمهوری اسلامی و
		حوزه‌های علمیه؛ سابقه‌ی
			تحصیل در حوزه‌ی قم.
			\keyword{شناخت عمیق از ساختار
		درونی حکومت — ضروری
	برای شعبه‌ی مستندسازی.} \\
	\bottomrule
	\end{longtable}
	
	\section{دسته‌ی چهارم: اعضای پیشنهادی شورای ناظران بیرونی (غیرایرانی)}
		
		\begin{longtable}{r p{3.2cm} p{3.8cm} p{4.5cm}}
		\toprule
		\textbf{ردیف} & \textbf{نام} & \textbf{وابستگی نهادی} &
		\textbf{تخصص و سابقه‌ی مرتبط} \\
		\midrule
		\endhead
		۱۸ & \textbf{سر جفری نایس}
		\en{(Sir Geoffrey Nice QC)} &
		وکیل و قاضی بین‌المللی
		(بریتانیا) &
		دادستان پرونده‌ی میلوشویچ
		در \en{ICTY}؛ رئیس دادگاه
		اویغورها (۲۰۲۱) و دادگاه
			چین (\en{China Tribunal}, ۲۰۱۹).
		\keyword{تجربه‌ی بی‌نظیر در
		رهبری دادگاه‌های مردمی.} \\
		\midrule
		۱۹ & \textbf{جان دوگارد}
		\en{(John Dugard)} &
		حقوق‌دان بین‌المللی
		(آفریقای جنوبی) &
		قاضی سابق \en{ICJ}؛
		گزارشگر ویژه‌ی سابق سازمان ملل؛
			عضو هیئت داوران دادگاه ایران ۲۰۱۲.
		\keyword{سابقه‌ی مستقیم در
		قضاوت درباره‌ی ایران.} \\
		\midrule
		۲۰ & \textbf{کریستین شینک}
		\en{(Christine Chinkin)} &
		استاد حقوق بین‌الملل،
		\en{LSE} (بریتانیا) &
		عضو هیئت حقیقت‌یاب سازمان ملل
		برای غزه (۲۰۰۹)؛ متخصص
		حقوق بشردوستانه و حقوق زنان
			در مناقشات.
		\keyword{بُعد جنسیتی و
		حقوق بشردوستانه.} \\
		\midrule
		۲۱ & \textbf{خوان مندز}
		\en{(Juan E.\ Méndez)} &
		استاد حقوق،
		\en{American University}
		(واشنگتن) &
		گزارشگر ویژه‌ی سابق سازمان ملل
		در امور شکنجه؛ مشاور ارشد
		\en{ICTJ}؛ بازمانده‌ی شکنجه
			در آرژانتین.
		\keyword{تجربه‌ی زیسته و حرفه‌ای
		در عدالت انتقالی.} \\
		\midrule
		۲۲ & \textbf{ناوی پیلای}
		\en{(Navi Pillay)} &
		قاضی سابق \en{ICC}
		و کمیسر عالی سابق
		حقوق بشر سازمان ملل
		(آفریقای جنوبی) &
		بالاترین سابقه‌ی قضایی بین‌المللی
		در فهرست؛ قاضی \en{ICTR}
			(رواندا)؛ مبارز ضدآپارتاید.
		\keyword{بالاترین اعتبار ممکن
		برای ریاست شورای ناظران.} \\
		\midrule
		۲۳ & \textbf{ویلیام شاباس}
		\en{(William Schabas)} &
		استاد حقوق بین‌الملل،
		دانشگاه میدلسکس (بریتانیا)
		و دانشگاه لیدن (هلند) &
		یکی از برجسته‌ترین متخصصان
		جهان در حقوق بین‌الملل کیفری
		و نسل‌کشی؛ نویسنده‌ی
			شرح اساسنامه‌ی رم.
		\keyword{مرجعیت آکادمیک
	بلامنازع.} \\
	\bottomrule
	\end{longtable}
	
	\section{دسته‌ی پنجم: نمایندگان اقلیت‌ها و جامعه‌ی مدنی}
		
		\begin{principleBox}[اصل تکثر]
		ترکیب نهاد باید شامل نمایندگان اقلیت‌های قومی، مذهبی،
	و جنسیتی ایران باشد. فهرست زیر صرفاً نمونه‌هایی از
	فعالان شناخته‌شده‌ی هر حوزه است.
	\end{principleBox}
		
		\begin{longtable}{r p{3.2cm} p{3.8cm} p{4.5cm}}
		\toprule
		\textbf{ردیف} & \textbf{نام / حوزه} & \textbf{نمونه‌ی پیشنهادی} &
		\textbf{ملاحظات} \\
		\midrule
		\endhead
		۲۴ & حقوق کُردها &
		فعالان حقوق بشر کُرد
		با سابقه‌ی آکادمیک &
			مشورت با سازمان‌هایی مانند
		\en{Kurdistan Human Rights
		Network (KHRN)} \\
		\midrule
		۲۵ & حقوق بلوچ‌ها &
		فعالان حقوق بلوچ
		(مثلاً فعالان مستند‌ساز
		کشتار زاهدان ۱۴۰۱) &
		اولویت ویژه با توجه به
		سرکوب شدید بلوچستان \\
		\midrule
		۲۶ & حقوق ترک‌آذری‌ها &
		فعالان حقوق زبانی و
		فرهنگی آذربایجان &
		مشورت با نهادهای فعال
		در حوزه‌ی حقوق زبانی \\
		\midrule
		۲۷ & حقوق عرب‌ها &
		فعالان حقوق عرب‌های
		خوزستان &
		مستندسازی تبعیض سیستماتیک
		و سرکوب محیط زیستی \\
		\midrule
		۲۸ & حقوق اقلیت‌های مذهبی &
		نمایندگان بهائیان، مسیحیان،
		یهودیان، زرتشتیان، دراویش &
		سابقه‌ی طولانی آزار مذهبی
		به‌ویژه علیه بهائیان \\
		\midrule
		۲۹ & حقوق \en{LGBTQ+} &
		فعالان حقوق دگرباشان
		ایرانی در خارج از کشور &
		اعدام و شکنجه‌ی سیستماتیک
		دگرباشان در ایران \\
		\midrule
		۳۰ & حقوق کارگران &
		فعالان سندیکایی ایرانی
		(داخل یا خارج) &
		مستندسازی سرکوب
	تشکل‌های کارگری \\
	\bottomrule
	\end{longtable}
	
	\section{نمودار ترکیب پیشنهادی}
	
	% ══════════════════════════════════════════════════════════════
	%  اصلاح اساسی: نمودار دایره‌ای (Pie Chart)
	%  ─────────────────────────────────────────────
	%  مشکل اصلی:
	%    بسته‌ی pgf-pie با xepersian و محیط RTL سازگاری ضعیف دارد.
	%    متن‌های فارسی درون \pie در بسیاری از نسخه‌ها بهم می‌ریزند
	%    یا جهت‌شان عوض می‌شود.
	%
	%  راه‌حل:
	%    نمودار را با TikZ خالص و بدون pgf-pie بازسازی می‌کنیم.
	%    از \RL{} برای تمام متن‌های فارسی درون TikZ استفاده می‌شود.
	%    legend جداگانه رسم می‌شود.
	% ══════════════════════════════════════════════════════════════
		
		\begin{figure}[H]
			\centering
			\begin{tikzpicture}[
			% ── رنگ‌ها ──
			slice1/.style={fill=cPrimary!80},
			slice2/.style={fill=cGold!80},
			slice3/.style={fill=cAccent!50},
			slice4/.style={fill=cPrimary!40},
			slice5/.style={fill=cGold!40},
			slice6/.style={fill=cGray!40},
			]
			% ── رسم برش‌های دایره‌ای ──
			% زاویه‌ها: 30%=108°  20%=72°  15%=54°  15%=54°  10%=36°  10%=36°
			% شروع از بالا (90 درجه) و در جهت ساعت‌گرد
			
			% برش ۱: حقوق بین‌الملل کیفری — ۳۰٪ (از 90 تا -18)
			\fill[slice1] (0,0) -- (90:3.2) arc (90:-18:3.2) -- cycle;
			\node at (36:2.2) {\footnotesize\bfseries\color{white} ۳۰\%};
			
			% برش ۲: حقوق بشر و مدنی — ۲۰٪ (از -18 تا -90)
			\fill[slice2] (0,0) -- (-18:3) arc (-18:-90:3) -- cycle;
			\node at (-54:2.2) {\footnotesize\bfseries ۲۰\%};
			
			% برش ۳: علوم سیاسی و تاریخ — ۱۵٪ (از -90 تا -144)
			\fill[slice3] (0,0) -- (-90:3) arc (-90:-144:3) -- cycle;
			\node at (-117:2.2) {\footnotesize\bfseries ۱۵\%};
			
			% برش ۴: جامعه‌شناسی و فلسفه — ۱۵٪ (از -144 تا -198)
			\fill[slice4] (0,0) -- (-144:3) arc (-144:-198:3) -- cycle;
			\node at (-171:2.2) {\footnotesize\bfseries ۱۵\%};
			
			% برش ۵: نمایندگان اقلیت‌ها — ۱۰٪ (از -198 تا -234)
			\fill[slice5] (0,0) -- (-198:3) arc (-198:-234:3) -- cycle;
			\node at (-216:2.2) {\footnotesize\bfseries ۱۰\%};
			
			% برش ۶: اعضای غیرایرانی — ۱۰٪ (از -234 تا -270 = 90)
			\fill[slice6] (0,0) -- (-234:3) arc (-234:-270:3) -- cycle;
			\node at (-252:2.2) {\footnotesize\bfseries ۱۰\%};
			
			% ── خط دور دایره ──
			\draw[thick, cPrimary!50] (0,0) circle (3);
			
				% ── راهنما (Legend) — سمت چپ نمودار ──
				\begin{scope}[shift={(-5.5,-2)}]
				\fill[slice1] (0, 2.5) rectangle ++(0.4,0.4);
				\node[anchor=east] at (-0.2, 2.7)
				{\footnotesize\RL{حقوق بین‌الملل کیفری}};
				
				\fill[slice2] (0, 1.8) rectangle ++(0.4,0.4);
				\node[anchor=east] at (-0.2, 2.0)
				{\footnotesize\RL{حقوق بشر و مدنی}};
				
				\fill[slice3] (0, 1.1) rectangle ++(0.4,0.4);
				\node[anchor=east] at (-0.2, 1.3)
				{\footnotesize\RL{علوم سیاسی و تاریخ}};
				
				\fill[slice4] (0, 0.4) rectangle ++(0.4,0.4);
				\node[anchor=east] at (-0.2, 0.6)
				{\footnotesize\RL{جامعه‌شناسی و فلسفه}};
				
				\fill[slice5] (0, -0.3) rectangle ++(0.4,0.4);
				\node[anchor=east] at (-0.2, -0.1)
				{\footnotesize\RL{نمایندگان اقلیت‌ها}};
				
				\fill[slice6] (0, -1.0) rectangle ++(0.4,0.4);
			\node[anchor=east] at (-0.2, -0.8)
			{\footnotesize\RL{اعضای غیرایرانی}};
		\end{scope}
		
	\end{tikzpicture}
	\caption{\RL{ترکیب تخصصی پیشنهادی هیئت عالی داوران و شورای ناظران}}
	\end{figure}
	
	\section{جمع‌بندی فصل دهم}
		
			\begin{principleBox}[خلاصه‌ی فصل ۱۰]
			\begin{enumerate}[nosep, label=\textcolor{cPrimary}{\arabic*.}]
			\item پنج دسته‌ی تخصصی شناسایی شده‌اند:
			حقوق بین‌الملل، علوم سیاسی، جامعه‌شناسی و فلسفه،
			اعضای بین‌المللی، و نمایندگان اقلیت‌ها.
			\item \keyword{پیام اخوان} (دادستان سابق \en{ICTY/ICTR})
			و \keyword{ناوی پیلای} (کمیسر سابق حقوق بشر سازمان ملل)
			بالاترین اعتبار حقوقی بین‌المللی را دارند.
			\item \keyword{رویا برومند} و \keyword{محمود امیری‌مقدم}
			نزدیک‌ترین تجربه به مأموریت مستندسازی را دارند.
			\item تنوع قومی، جنسیتی و نسلی باید از ابتدا تضمین شود.
		\item درج نام هیچ فرد به‌معنای تعهد یا مشارکت آن شخص نیست —
	فهرست صرفاً پیشنهاد اولیه است.
	\end{enumerate}
	\end{principleBox}
	
	\clearpage
	
	
	% ////////////////////////////////////////////////////////////
	%  فصل ۱۱: جمع‌بندی نهایی و نتیجه‌گیری
	% ////////////////////////////////////////////////////////////
	\chapter{جمع‌بندی نهایی و نتیجه‌گیری}\label{chap:conclusion}
	
	\section{مرور مسئله}
	
		جنبش گذارخواهی ایران در مقطعی تاریخی قرار دارد:
		\begin{itemize}[nosep]
		\item خیزش‌های پیاپی مردمی نشان‌دهنده‌ی اراده‌ی جمعی برای تغییر است؛
		\item اما فقدان ساختار نهادی — به‌ویژه در حوزه‌ی قضایی‌ـ‌اخلاقی —
	جنبش را آسیب‌پذیر، ناهماهنگ، و در برابر جنگ روانی
	حکومت بی‌دفاع کرده است.
	\end{itemize}
	
	این سند استدلال کرد که تأسیس یک \keyword{نهاد شبه‌قضایی سایه}
	با عنوان \textbf{«دادگاه در سایه — وجدان منفصل مسئول»}
	می‌تواند این خلأ را — نه کامل اما مؤثر — پُر کند.
	
	\section{خلاصه‌ی فصل‌ها}
		
		\begin{longtable}{r p{3cm} p{8.5cm}}
		\toprule
		\textbf{فصل} & \textbf{عنوان} & \textbf{یافته‌ی اصلی} \\
		\midrule
		\endhead
		۱ & تشریح مسئله &
		خلأ نهادی واقعی و فوری است. جمهوری اسلامی
		به‌صورت سازمان‌یافته از آن بهره‌برداری می‌کند. \\
		\midrule
		۲ & انتظارات مثبت &
		نهاد می‌تواند اعتبار، هماهنگی، سپر دفاعی، و
		حافظه‌ی نهادی تولید کند. \\
		\midrule
		۳ & آسیب‌شناسی &
		هفت خطر اصلی شناسایی و برای هر کدام
		سازوکار پیش‌گیری طراحی شد. \\
		\midrule
		۴ & سازماندهی &
		ساختار چهارلایه: شورای ناظران، هیئت عالی،
		شعبات تخصصی، دبیرخانه. وکیل تسخیری
		مهم‌ترین نوآوری. \\
		\midrule
		۵ & رویه‌ها &
		چهار نوع پرونده با چهار سطح رسمیت.
		آزمون سه‌مرحله‌ای تناسب. سه سطح استاندارد شواهد. \\
		\midrule
		۶ & مالی &
		بودجه‌ی سالانه ۳۵۰–۵۷۵ هزار یورو.
		تنوع منابع و شفافیت کامل. \\
		\midrule
		۷ & بین‌المللی &
		صلاحیت جهانی مؤثرترین مسیر فعلی.
		\en{FFMI} و گزارشگر ویژه نزدیک‌ترین شرکا.
		لاهه بهترین مکان. \\
		\midrule
		۸ & نمونه‌ها &
		هفت نمونه بررسی شد. مهم‌ترین درس:
		نهادسازی پیش از گذار از عدالت وحشیانه
		جلوگیری می‌کند. \\
		\midrule
		۹ & پس از گذار &
		سه مرحله‌ی گذار. چارچوب عفو مشروط.
		چهار سناریوی انحلال/ادغام. \\
		\midrule
		۱۰ & افراد &
	۳۰ نام پیشنهادی در پنج دسته‌ی تخصصی. \\
	\bottomrule
	\end{longtable}
	
	\section{ده اصل بنیادین — مرام‌نامه‌ی پیشنهادی}
		
			\begin{principleBox}[مرام‌نامه‌ی دادگاه در سایه — وجدان منفصل مسئول]
			\begin{enumerate}[nosep, label=\textcolor{cGold}{\Roman*.}]
			\item \textbf{دعوت به اندیشیدن:}
			ما حکم صادر نمی‌کنیم؛ دعوت به تفکر مسئولانه می‌کنیم.
			
			\item \textbf{احتیاط:}
			هر توصیه‌ی ما با بیشترین احتیاط ممکن صادر می‌شود.
			در تردید، سکوت بر شتاب‌زدگی ترجیح دارد.
			
			\item \textbf{شفافیت:}
			هر چه می‌اندیشیم، می‌نویسیم. هر چه تصمیم می‌گیریم، منتشر می‌کنیم.
			
			\item \textbf{فروتنی:}
			ما نماینده‌ی مردم ایران نیستیم.
			ما خادمان عدالت و حقیقت هستیم — با تمام محدودیت‌هایمان.
			
			\item \textbf{استقلال:}
			وابسته به هیچ حزب، دولت، یا جناحی نیستیم.
			تنها وفاداری ما به حقیقت و کرامت انسانی است.
			
			\item \textbf{تکثر:}
			همه‌ی صداها شنیده می‌شوند.
			نظر اقلیت به‌اندازه‌ی نظر اکثریت ثبت می‌شود.
			
			\item \textbf{تناسب:}
			هر اقدامی باید متناسب با هدف باشد.
			هزینه‌ی بی‌گناهان همیشه در محاسبه لحاظ می‌شود.
			
			\item \textbf{کرامت:}
			حتی بدترین متهمان از حقوق بنیادین انسانی برخوردارند.
			ما عدالت می‌خواهیم، نه انتقام.
			
			\item \textbf{مسئولیت:}
			آراء ما ثبت می‌شوند و ما در برابر تاریخ پاسخ‌گوییم.
			
			\item \textbf{موقتی بودن:}
			هدف ما منحل کردن خودمان است —
		در روزی که ایران، نظام قضایی مستقل و
	دموکراتیک خود را داشته باشد.
	\end{enumerate}
	\end{principleBox}
	
	\section{نقشه‌ی راه اجرایی}
		
		\begin{longtable}{r p{3cm} p{4cm} p{4.5cm}}
		\toprule
		\textbf{مرحله} & \textbf{بازه‌ی زمانی} & \textbf{اقدامات کلیدی} &
		\textbf{خروجی} \\
		\midrule
		\endhead
		\textbf{۱. بنیان‌گذاری} & ماه ۱–۶ &
		تشکیل کمیته‌ی مؤسس (۳–۵ نفر)؛
		تدوین اساسنامه؛ ثبت قانونی در هلند؛
		جذب بودجه‌ی اولیه &
		نهاد رسماً تأسیس شود \\
		\midrule
		\textbf{۲. ساختارسازی} & ماه ۳–۱۲ &
		انتخاب شورای ناظران؛
		تأیید اعضای هیئت عالی؛
		استخدام دبیرخانه؛
		راه‌اندازی زیرساخت فنی &
		ساختار کامل و عملیاتی \\
		\midrule
		\textbf{۳. نخستین پرونده‌ها} & ماه ۶–۱۸ &
		پذیرش ۲–۳ پرونده‌ی نوع الف
		(مستندسازی) و ۱–۲ پرونده‌ی نوع ب
		(داوری تصمیم سیاسی) &
		اولین آراء منتشر شوند \\
		\midrule
		\textbf{۴. تثبیت و گسترش} & ماه ۱۲–۳۶ &
		افزایش پرونده‌ها؛ تقویت ارتباطات
		بین‌المللی؛ انتشار گزارش سالانه؛
		بازنگری ساختاری اول &
		نهاد به‌عنوان مرجع شناخته شود \\
		\midrule
		\textbf{۵. آمادگی گذار} & مستمر &
		آماده‌سازی مستندات و ساختارها
		برای تحویل به نظام قضایی آینده &
	بسته‌ی کامل عدالت انتقالی \\
	\bottomrule
	\end{longtable}
	
	\section{فراخوان مشارکت}
		
			\begin{noteBox}[به چه کسانی نیاز داریم؟]
			\begin{enumerate}[nosep, label=\textcolor{cPrimary}{\arabic*.}]
			\item \textbf{حقوق‌دانان:} برای طراحی آیین دادرسی و تحلیل حقوقی.
			\item \textbf{پژوهشگران:} برای مستندسازی و تحقیق.
			\item \textbf{مترجمان:} فارسی‌ـ‌انگلیسی و سایر زبان‌ها.
			\item \textbf{متخصصان فناوری:} امنیت سایبری، بانک اطلاعاتی، وب.
			\item \textbf{قربانیان و شاهدان:} برای ثبت شهادت‌ها.
			\item \textbf{حامیان مالی:} بنیادها و افراد.
		\item \textbf{رسانه‌ها:} برای پوشش و انتشار.
	\item \textbf{هر شهروند ایرانی:} برای نقد، پیشنهاد، و نظارت.
	\end{enumerate}
	\end{noteBox}
	
	\section{کلام آخر}
		
		\begin{principleBox}[چرا «وجدان منفصل مسئول»؟]
		در لحظاتی از تاریخ، ملت‌ها مجبورند تصمیم‌هایی بگیرند
		که نه آسان‌اند و نه بدون هزینه. تحریم‌هایی که بر
		زندگی مردم عادی اثر می‌گذارند. اقداماتی که ممکن
		است خشونت‌آمیز باشند. مرزگذاری‌هایی که بعضی‌ها را
		بیرون می‌گذارند.
		
		\medskip
		در چنین لحظاتی، بدترین کار این است که بدون
		اندیشیدن تصمیم بگیریم — و بهترین کار این است که
		\keyword{وجدانی منفصل} (مستقل از فشار سیاسی و احساسات آنی)
		اما \keyword{مسئول} (پاسخ‌گو در برابر حقیقت و تاریخ)
		داشته باشیم.
		
		\medskip
		دادگاه در سایه قرار نیست جای هیچ‌کس تصمیم بگیرد.
		قرار است فضایی بسازد که تصمیم‌های سخت،
		\textbf{اندیشیده}، \textbf{مستند}، و \textbf{قابل دفاع} باشند.
		
		\medskip
	این حداقلِ وظیفه‌ای است که نسل ما
	در برابر نسل‌های آینده‌ی ایران دارد.
	\end{principleBox}
	
	\vfill
		
		\begin{center}
		\rule{0.5\textwidth}{0.5pt}\\[12pt]
		{\large\color{cPrimary}\textbf{پایان سند پیشنهادی}}\\[6pt]
		{\color{cGray} نسخه‌ی ۰.۹ — پیش‌نویس برای بحث نخبگانی}\\[4pt]
			{\color{cGray}\today}\\[12pt]
			{\small\color{cGold}
			«عدالت بدون فرایند، نامش انتقام است.\\
	فرایند بدون عدالت، نامش بوروکراسی است.\\
	ما هر دو را با هم می‌خواهیم.»}
	\end{center}
	
	\clearpage
	
	
	% ////////////////////////////////////////////////////////////
	%  پیوست‌ها
	% ////////////////////////////////////////////////////////////
	\appendix
% ////////////////////////////////////////////////////////////
%  پیوست الف: واژه‌نامه‌ی تخصصی
% ////////////////////////////////////////////////////////////
\chapter{واژه‌نامه‌ی تخصصی}\label{app:glossary}

\begin{longtable}{p{4cm} p{4cm} p{5cm}}
	\toprule
	\textbf{اصطلاح فارسی} & \textbf{اصطلاح انگلیسی} &
	\textbf{توضیح مختصر} \\
	\midrule
	\endhead
	دادگاه در سایه & \en{Shadow Quasi-Judicial Body} &
	نهاد شبه‌قضایی مشورتی بدون قدرت اجرایی رسمی \\
	\midrule
	وجدان منفصل مسئول & \en{Responsible Detached Conscience} &
	مفهوم بنیادین سند: داوری مستقل، اخلاقی، و پاسخ‌گو \\
	\midrule
	کیفرخواست سایه & \en{Shadow Indictment} &
	سند اعلام جرم صادره توسط نهاد بدون ضمانت اجرای رسمی \\
	\midrule
	رأی مشورتی & \en{Advisory Opinion} &
	نظر حقوقی غیرالزام‌آور درباره‌ی یک مسئله \\
	\midrule
	صلاحیت جهانی & \en{Universal Jurisdiction} &
	اصل حقوقی که به دادگاه ملی اجازه‌ی رسیدگی به
	جرایم بین‌المللی بدون ارتباط مستقیم را می‌دهد \\
	\midrule
	وکیل تسخیری & \en{Amicus Curiae} &
	وکیل مستقل مأمور ارائه‌ی بهترین استدلال به نفع طرف مقابل \\
	\midrule
	آزمون تناسب & \en{Proportionality Test} &
	بررسی تناسب بین هزینه و فایده‌ی یک اقدام \\
	\midrule
	عدالت انتقالی & \en{Transitional Justice} &
	مجموعه‌ی سازوکارهای حقوقی برای رسیدگی به
	جرایم نظام سابق پس از گذار \\
	\midrule
	مسئولیت حمایت & \en{Responsibility to Protect (R2P)} &
	اصل بین‌المللی مجوزدهنده‌ی مداخله‌ی بشردوستانه
	در صورت ناکامی دولت در حمایت از شهروندان \\
	\midrule
	عفو مشروط & \en{Conditional Amnesty} &
	بخشش مشروط به اعتراف کامل، همکاری، و جبران خسارت \\
	\midrule
	هیئت حقیقت‌یاب & \en{Fact-Finding Mission} &
	نهاد مأمور بررسی و گزارش واقعیت‌ها بدون صدور حکم \\
	\midrule
	فرسودگی شواهد & \en{Evidence Attrition} &
	از بین رفتن تدریجی مدارک و خاطرات با گذشت زمان \\
	\midrule
	امنیت عملیاتی & \en{Operational Security (OPSEC)} &
	مجموعه‌ی اقدامات حفاظتی در ارتباطات و فعالیت‌ها \\
	\midrule
	حافظه‌ی نهادی & \en{Institutional Memory} &
	دانش و مستندات انباشته‌شده در یک نهاد \\
	\bottomrule
\end{longtable}


% ////////////////////////////////////////////////////////////
%  پیوست ب: فهرست اسناد بین‌المللی مرجع
% ////////////////////////////////////////////////////////////
\chapter{فهرست اسناد بین‌المللی مرجع}\label{app:intldocs}

\begin{enumerate}[nosep, label=\textcolor{cPrimary}{\arabic*.}]
	\item اعلامیه‌ی جهانی حقوق بشر (۱۹۴۸)
	\item میثاق بین‌المللی حقوق مدنی و سیاسی (\en{ICCPR}, ۱۹۶۶)
	\item میثاق بین‌المللی حقوق اقتصادی، اجتماعی و فرهنگی (\en{ICESCR}, ۱۹۶۶)
	\item کنوانسیون منع شکنجه (\en{UNCAT}, ۱۹۸۴)
	\item اساسنامه‌ی رم دیوان کیفری بین‌المللی (۱۹۹۸)
	\item اصول نورنبرگ (۱۹۵۰)
	\item سند اجلاس جهانی ۲۰۰۵ (\en{R2P})
	\item اصول \en{Joinet/Orentlicher} (عدالت انتقالی)
	\item اصول \en{Paris Principles} (نهادهای ملی حقوق بشر)
	\item قطعنامه‌های شورای حقوق بشر سازمان ملل درباره‌ی ایران
	\item قانون \en{Global Magnitsky} (ایالات متحده، ۲۰۱۶)
	\item \en{Justice for Victims of Corrupt Foreign Officials Act} (کانادا، ۲۰۱۷)
	\item \en{Global Human Rights Sanctions Regulations} (بریتانیا، ۲۰۲۰)
	\item \en{EU Global Human Rights Sanctions Regime} (اتحادیه‌ی اروپا، ۲۰۲۰)
\end{enumerate}


% ////////////////////////////////////////////////////////////
%  پیوست ج: نمونه‌ی فرم ارجاع پرونده
% ////////////////////////////////////////////////////////////
\chapter{نمونه‌ی فرم ارجاع پرونده}\label{app:referralform}

\begin{principleBox}[فرم ارجاع — پرونده‌ی نوع الف (مستندسازی جرایم)]
	\begin{enumerate}[nosep, label=\textcolor{cPrimary}{\arabic*.}]
		\item \textbf{مشخصات ارجاع‌دهنده:}
		نام (یا کُد ناشناس) / سازمان / تاریخ
		
		\item \textbf{شرح وقایع:}
		شرح مختصر (حداکثر ۲ صفحه) از واقعه‌ی مورد نظر
		شامل: زمان، مکان، بازیگران اصلی، و قربانیان.
		
		\item \textbf{طبقه‌بندی حقوقی پیشنهادی:}
		$\square$ جرم علیه بشریت \quad
		$\square$ جرم جنگی \quad
		$\square$ نسل‌کشی \quad
		$\square$ شکنجه \quad
		$\square$ سایر (توضیح دهید)
		
		\item \textbf{اشخاص مظنون:}
		فهرست اسامی (حقیقی و حقوقی) با ذکر نقش هر شخص.
		
		\item \textbf{شواهد اولیه:}
		$\square$ شهادت‌نامه \quad
		$\square$ تصویر/ویدئو \quad
		$\square$ سند مکتوب \quad
		$\square$ گزارش سازمان بین‌المللی \quad
		$\square$ سایر
		
		\item \textbf{درخواست از نهاد:}
		$\square$ مستندسازی و بایگانی \quad
		$\square$ صدور کیفرخواست سایه \quad
		$\square$ ارجاع به مراجع بین‌المللی \quad
		$\square$ سایر
		
		\item \textbf{ملاحظات امنیتی:}
		آیا ارجاع‌دهنده یا شاهدان در معرض خطر هستند؟
		$\square$ بله \quad $\square$ خیر
		(در صورت بله، توضیح دهید.)
		
		\item \textbf{امضا و تاریخ}
	\end{enumerate}
\end{principleBox}

\clearpage

% ////////////////////////////////////////////////////////////
%  کولوفون
% ////////////////////////////////////////////////////////////
\thispagestyle{empty}
\vspace*{\fill}
\begin{center}
	{\large\color{cPrimary}\textbf{دادگاه در سایه — وجدان منفصل مسئول}}\\[8pt]
	{\color{cGray}سند پیشنهادی برای تأسیس نظام شبه‌قضایی
		اپوزیسیون دموکراتیک ایران}\\[16pt]
	
	\begin{tabular}{r l}
		نسخه: & ۰.۹ (پیش‌نویس برای بحث نخبگانی) \\
		تاریخ: & \today \\
		قالب: & \XeLaTeX \\
		مجوز: & \en{Creative Commons BY-SA 4.0} \\
		تماس: & \en{[shadow.court@protonmail.com]} \\
		& {\footnotesize (نشانی پیشنهادی — فعلاً فعال نیست)} \\
	\end{tabular}
	
	\vspace{24pt}
	{\small\color{cGray}
		این سند با هدف دعوت به اندیشیدن، نه تحمیل نتیجه، نوشته شده است.\\
		از هر بازخورد سازنده‌ای سپاسگزاریم.}
	
	\vspace{24pt}
	\begin{tikzpicture}
		\draw[cGold, thick] (0,0) -- (4,0);
	\end{tikzpicture}
	
	\vspace{12pt}
	{\footnotesize\color{cGold}
		«وقتی عدالت گنگ شود، وجدان باید سخن بگوید.»}
\end{center}
\vspace*{\fill}


% ////////////////////////////////////////////////////////////
%  پیوست د: نمونه‌ی کیفرخواست سایه
% ////////////////////////////////////////////////////////////
\chapter{نمونه‌ی کیفرخواست سایه}\label{app:indictment}

\begin{warningBox}[تذکر]
	متن زیر یک \textbf{نمونه‌ی ساختاری فرضی} است و
	صرفاً برای نشان دادن قالب پیشنهادی کیفرخواست سایه
	تهیه شده است. اسامی و جزئیات بر اساس اطلاعات
	عمومی موجود درج شده‌اند و نباید به‌عنوان
	کیفرخواست واقعی تلقی شوند.
\end{warningBox}

\begin{principleBox}[کیفرخواست سایه — نمونه‌ی ساختاری]
	
	\begin{center}
		{\large\bfseries\color{cPrimary}
			دادگاه در سایه — وجدان منفصل مسئول\\[4pt]
			شعبه‌ی مستندسازی جرایم\\[8pt]}
		{\bfseries کیفرخواست سایه شماره‌ی: \en{SCI-2024-001}}\\[4pt]
		{\color{cGray} تاریخ صدور: [تاریخ] \quad|\quad
			طبقه‌بندی: عمومی}
	\end{center}
	
	\vspace{12pt}
	
	\textbf{بخش اول: مشخصات پرونده}
	
	% ── اصلاح: از tabular به‌جای longtable برای جدول کوتاه ──
	\begin{tabular}{r p{10cm}}
		عنوان پرونده: & [عنوان توصیفی واقعه] \\[4pt]
		تاریخ وقوع: & [تاریخ دقیق یا بازه‌ی زمانی] \\[4pt]
		مکان وقوع: & [استان / شهر / محل دقیق] \\[4pt]
		طبقه‌بندی حقوقی: & $\boxtimes$ جرم علیه بشریت
		(مواد ۷(الف)(۱) و ۷(الف)(۲) اساسنامه‌ی رم) \\
		& $\square$ جرم جنگی \\
		& $\square$ نسل‌کشی \\
		& $\boxtimes$ شکنجه (ماده‌ی ۱ کنوانسیون منع شکنجه) \\[4pt]
		تعداد قربانیان شناسایی‌شده: & [عدد] \\[4pt]
		سطح استاندارد شواهد: & \en{Beyond reasonable doubt} \\
	\end{tabular}
	
	\vspace{8pt}
	\textbf{بخش دوم: شرح وقایع}
	
	[شرح مفصل واقعه بر اساس ترتیب زمانی.
	حداقل ۳ صفحه. هر ادعای واقعی با شماره‌ی مرجع
	شواهد مشخص شود.]
	
	\vspace{8pt}
	\textbf{بخش سوم: شواهد}
	
	\begin{longtable}{r p{2.5cm} p{3cm} p{2cm} p{3cm}}
		\toprule
		\textbf{شماره} & \textbf{نوع} & \textbf{منبع} &
		\textbf{تأیید مستقل} & \textbf{ارزیابی اعتبار} \\
		\midrule
		\endhead
		\en{E-001} & شهادت‌نامه & شاهد عینی [کُد] & بله (۲ منبع) & بالا \\
		\en{E-002} & ویدئو & شبکه‌ی اجتماعی & بله (\en{OSINT}) & بالا \\
		\en{E-003} & گزارش پزشکی & [کُد] & بله (۱ منبع) & متوسط \\
		\en{E-004} & تصویر ماهواره‌ای & \en{Planet Labs} & بله & بالا \\
		\en{E-005} & سند رسمی فاش‌شده & [کُد] & در حال بررسی & متوسط \\
		\bottomrule
	\end{longtable}
	
	\vspace{8pt}
	\textbf{بخش چهارم: اشخاص مظنون و زنجیره‌ی فرماندهی}
	
	\begin{longtable}{r p{3cm} p{3cm} p{2.5cm} p{3cm}}
		\toprule
		\textbf{ردیف} & \textbf{نام / سمت} & \textbf{نقش در واقعه} &
		\textbf{سطح مسئولیت} & \textbf{شواهد مرتبط} \\
		\midrule
		\endhead
		۱ & [سمت — بدون نام فعلاً] & فرماندهی مستقیم عملیات &
		مسئولیت فرماندهی
		(\en{command responsibility}) &
		\en{E-001}, \en{E-004} \\
		\midrule
		۲ & [سمت] & صدور دستور تیراندازی &
		مسئولیت مستقیم &
		\en{E-002}, \en{E-005} \\
		\midrule
		۳ & [سمت] & نظارت بر عملیات &
		مسئولیت فرماندهی &
		\en{E-001}, \en{E-003} \\
		\bottomrule
	\end{longtable}
	
	\vspace{8pt}
	\textbf{بخش پنجم: تحلیل حقوقی}
	
	[تحلیل مفصل انطباق وقایع با عناصر تشکیل‌دهنده‌ی
	جرایم مذکور بر اساس اساسنامه‌ی رم و
	رویه‌ی قضایی \en{ICC/ICTY/ICTR}.
	حداقل ۵ صفحه.]
	
	\vspace{8pt}
	\textbf{بخش ششم: پاسخ وکیل تسخیری (\en{Amicus Curiae})}
	
	[خلاصه‌ی لایحه‌ی دفاعیه‌ی تهیه‌شده توسط
	دفتر وکیل تسخیری. شامل نقد شواهد،
	دفاعیات احتمالی، و ملاحظات حقوق بشری.]
	
	\vspace{8pt}
	\textbf{بخش هفتم: رأی هیئت داوران}
	
	\begin{itemize}[nosep]
		\item \textbf{رأی اکثریت:} [خلاصه‌ی یافته‌ها و نتیجه‌گیری]
		\item \textbf{نظر اقلیت:} [خلاصه‌ی نظرات مخالف]
		\item \textbf{توصیه‌ها:}
		\begin{enumerate}[nosep]
			\item ارجاع به دادگاه ملی [کشور] بر اساس صلاحیت جهانی؛
			\item درخواست تحریم هدفمند اشخاص مذکور؛
			\item ارسال مستندات به \en{FFMI} و \en{OHCHR}.
		\end{enumerate}
	\end{itemize}
	
	\vspace{8pt}
	\textbf{بخش هشتم: امضاها}
	
	\begin{tabular}{p{5cm} p{5cm}}
		رئیس هیئت داوران: & [نام و امضا] \\[8pt]
		دادستان سایه: & [نام و امضا] \\[8pt]
		وکیل تسخیری: & [نام و امضا] \\[8pt]
		دبیر جلسه: & [نام و امضا] \\
	\end{tabular}
\end{principleBox}


% ////////////////////////////////////////////////////////////
%  پیوست ه‍: نمونه‌ی رأی مشورتی (پرونده‌ی نوع ب)
% ////////////////////////////////////////////////////////////
\chapter{نمونه‌ی رأی مشورتی}\label{app:advisory}

\begin{principleBox}[رأی مشورتی — نمونه‌ی ساختاری]
	
	\begin{center}
		{\large\bfseries\color{cPrimary}
			دادگاه در سایه — وجدان منفصل مسئول\\[4pt]
			شعبه‌ی داوری تصمیم‌های سیاسی\\[8pt]}
		{\bfseries رأی مشورتی شماره‌ی: \en{SAO-2024-001}}\\[4pt]
		{\color{cGray} تاریخ صدور: [تاریخ] \quad|\quad
			موضوع: [عنوان موضوع]}
	\end{center}
	
	\vspace{12pt}
	
	\textbf{۱. سؤال مطروحه}
	
	[بیان دقیق سؤال یا درخواستی که به نهاد ارجاع شده.
	مثال: «آیا حمایت از تعلیق عضویت تیم ملی فوتبال ایران
	از \en{FIFA} توصیه می‌شود؟»]
	
	\vspace{8pt}
	\textbf{۲. خلاصه‌ی دیدگاه‌های ارائه‌شده}
	
	\begin{longtable}{p{6cm} p{6cm}}
		\toprule
		\textbf{دیدگاه‌های موافق} & \textbf{دیدگاه‌های مخالف} \\
		\midrule
		\endhead
		[خلاصه‌ی استدلال‌های موافقان] &
		[خلاصه‌ی استدلال‌های مخالفان] \\
		\bottomrule
	\end{longtable}
	
	\vspace{8pt}
	\textbf{۳. اجرای آزمون سه‌مرحله‌ای تناسب}
	
	\textbf{آزمون مشروعیت:}
	\begin{itemize}[nosep]
		\item هدف مشروع است / نیست؟ [تحلیل]
		\item مبنای حقوقی: [ذکر اسناد مرتبط]
		\item نتیجه: $\square$ موفق \quad $\square$ ناموفق
	\end{itemize}
	
	\textbf{آزمون ضرورت:}
	\begin{itemize}[nosep]
		\item جایگزین‌های بررسی‌شده:
		\begin{enumerate}[nosep]
			\item [جایگزین ۱] — ارزیابی: [کافی / ناکافی]
			\item [جایگزین ۲] — ارزیابی: [کافی / ناکافی]
			\item [جایگزین ۳] — ارزیابی: [کافی / ناکافی]
		\end{enumerate}
		\item نتیجه: $\square$ موفق \quad $\square$ ناموفق
	\end{itemize}
	
	\textbf{آزمون تناسب:}
	\begin{itemize}[nosep]
		\item فایده‌ی مورد انتظار: [توصیف]
		\item هزینه‌ی بر افراد بی‌تقصیر: [توصیف]
		\item نسبت هزینه‌ـ‌فایده: [تحلیل]
		\item اثرات جانبی: [توصیف]
		\item مکانیزم جبران: [توصیف]
		\item نتیجه: $\square$ موفق \quad $\square$ ناموفق
	\end{itemize}
	
	\vspace{8pt}
	\textbf{۴. نتیجه‌گیری}
	
	\begin{itemize}[nosep]
		\item $\square$ \textcolor{cGold!80!black}{$\bigstar$} توصیه‌شده
		\item $\square$ \textcolor{cPrimary}{$\blacksquare$} مجاز با قیود [ذکر قیود]
		\item $\square$ \textcolor{cGray}{$\blacktriangle$} تعلیق [ذکر دلایل]
		\item $\square$ \textcolor{cAccent}{$\boldsymbol{\times}$} توصیه‌نشده [ذکر دلایل]
	\end{itemize}
	
	\vspace{8pt}
	\textbf{۵. نظر اقلیت}
	
	[ثبت کامل نظرات مخالف با ذکر دلایل]
	
	\vspace{8pt}
	\textbf{۶. قابلیت اعتراض}
	
	این رأی تا ۲۱ روز پس از انتشار قابل اعتراض مستند است.
	اعتراضات به دبیرخانه ارسال شوند.
	
\end{principleBox}


% ////////////////////////////////////////////////////////////
%  پیوست و: استراتژی ارتباطی و رسانه‌ای
% ////////////////////////////////////////////////////////////
\chapter{استراتژی ارتباطی و رسانه‌ای}\label{app:comms}

\section{اهداف ارتباطی}

\begin{principleBox}[چهار هدف]
	\begin{enumerate}[nosep, label=\textcolor{cPrimary}{\arabic*.}]
		\item \textbf{اعتبارسازی:}
		نهاد به‌عنوان مرجع معتبر حقوقی‌ـ‌اخلاقی
		شناخته شود.
		\item \textbf{دفاع در برابر جنگ رسانه‌ای:}
		پاسخ‌های پیش‌ساخته برای حملات تبلیغاتی
		جمهوری اسلامی آماده باشد.
		\item \textbf{مشارکت‌جویی:}
		جذب مشارکت حقوق‌دانان، پژوهشگران، و شهروندان.
		\item \textbf{مدیریت انتظارات:}
		توضیح شفاف ماهیت مشورتی نهاد و
		تفاوت آن با دادگاه رسمی.
	\end{enumerate}
\end{principleBox}

\section{مخاطبان و پیام‌های کلیدی}

% ── اصلاح: عرض ستون‌ها تنظیم شد ──
\begin{longtable}{r p{2.8cm} p{4.2cm} p{4cm}}
	\toprule
	\textbf{ردیف} & \textbf{مخاطب} & \textbf{پیام اصلی} &
	\textbf{کانال ارتباطی} \\
	\midrule
	\endhead
	۱ & دیاسپورای ایرانی &
	«تصمیم‌های سخت جنبش اکنون
	فرایند حقوقی دارد.
	مشارکت شما خوش‌آمد است.» &
	شبکه‌های اجتماعی فارسی‌زبان؛
	پادکست؛ رسانه‌های فارسی‌زبان
	(\en{BBC Persian}, \en{Iran International},
	\en{Manoto}) \\
	\midrule
	۲ & جامعه‌ی حقوقی بین‌المللی &
	«یک نهاد ساختارمند با فرایند
	دادرسی معتبر تأسیس شده که
	مستندات استاندارد تولید می‌کند.» &
	مجلات حقوقی؛ کنفرانس‌ها؛
	ارتباط مستقیم با دانشکده‌های حقوق \\
	\midrule
	۳ & نهادهای بین‌المللی &
	«مستندات ما آماده‌ی استفاده
	در فرایندهای رسمی است.» &
	ارسال مستقیم؛ حضور در
	جلسات \en{UNHRC}؛
	لابی‌گری حقوقی \\
	\midrule
	۴ & رسانه‌های بین‌المللی &
	«داستان‌های مستند و تحلیلی
	درباره‌ی جرایم حکومتی ایران
	با کیفیت دادگاهی.» &
	بسته‌های رسانه‌ای (\en{press kits})؛
	کنفرانس‌های مطبوعاتی؛
	مصاحبه‌ی اختصاصی \\
	\midrule
	۵ & مردم داخل ایران &
	«ما دیده‌ایم. ثبت کرده‌ایم.
	روزی عدالت اجرا خواهد شد.
	از خودسری بپرهیزید.» &
	\en{VPN}-محور؛ تلگرام؛
	پیام‌های صوتی و تصویری کوتاه \\
	\midrule
	۶ & مقام‌های جمهوری اسلامی &
	«جرایم شما مستند شده.
	همکاری و اعتراف می‌تواند
	به عفو مشروط منجر شود.» &
	انتشار عمومی آراء؛
	ارسال مستقیم (در صورت امکان) \\
	\bottomrule
\end{longtable}

\section{پاسخ به حملات تبلیغاتی پیش‌بینی‌شده}

% ══════════════════════════════════════════════════════════════
%  اصلاح اساسی: حذف \\ از داخل سلول‌های جدول
%  ─────────────────────────────────────────────
%  مشکل: استفاده از \\ درون p{} در longtable باعث خطای
%        "misplaced \noalign" یا "extra alignment tab" می‌شود.
%  راه‌حل: متن را در یک خط بنویسیم یا از / استفاده کنیم.
% ══════════════════════════════════════════════════════════════

\begin{longtable}{r p{4cm} p{7.5cm}}
	\toprule
	\textbf{ردیف} & \textbf{حمله‌ی احتمالی} &
	\textbf{پاسخ پیش‌ساخته} \\
	\midrule
	\endhead
	۱ & «دادگاه نمایشی است» &
	«آیین دادرسی ما بر اساس استانداردهای \en{ICC}
	طراحی شده. وکیل تسخیری بهترین استدلال را
	به نفع طرف مقابل ارائه می‌دهد. تمام آراء
	با نظرات اقلیت منتشر می‌شوند.
	دادگاه‌های نمایشی این‌ها را ندارند.» \\
	\midrule
	۲ & «مشروعیت ندارد» &
	«ما ادعای مشروعیت حاکمیتی نداریم.
	ما یک نهاد مشورتی متخصص هستیم —
	مشابه \en{Uyghur Tribunal} یا
	\en{Russell Tribunal}. مشروعیت ما
	از کیفیت استدلال و شفافیت فرایند
	ناشی می‌شود، نه از قدرت سیاسی.» \\
	\midrule
	۳ & «ابزار دست آمریکا/اسرائیل» &
	«بودجه‌ی ما عمومی و شفاف است.
	هیچ دولتی مستقیماً بودجه نمی‌دهد.
	قاعده‌ی ۲۵٪ از وابستگی به هر منبع
	واحد جلوگیری می‌کند. گزارش مالی
	سالانه‌ی ما را بخوانید.» \\
	\midrule
	% ── اصلاح: \\ حذف شد از عنوان حمله ──
	۴ & «به ورزشکاران/مردم آسیب می‌زنید» &
	«هر توصیه‌ی ما از آزمون سه‌مرحله‌ای
	تناسب عبور کرده. هزینه‌ی بی‌گناهان
	همیشه در محاسبه لحاظ می‌شود.
	رأی مشورتی ما درباره‌ی [موضوع خاص]
	را بخوانید — شامل تحلیل هزینه‌ـ‌فایده
	و نظرات مخالف.» \\
	\midrule
	% ── اصلاح: \\ حذف شد از عنوان حمله ──
	۵ & «خودشان هم ناقض حقوق بشرند» &
	«شورای ناظران بیرونی (اکثراً غیرایرانی)
	بر عملکرد ما نظارت دارد.
	هر شهروند می‌تواند نسبت به آراء ما
	اعتراض مستند ارائه دهد.
	اگر اشتباهی اثبات شود، رسماً عذرخواهی
	و اصلاح می‌کنیم — چیزی که جمهوری
	اسلامی هرگز نکرده.» \\
	\midrule
	% ── اصلاح: \\ حذف شد از عنوان حمله ──
	۶ & «در خدمت براندازی و خشونت هستند» &
	«اصل دوم مرام‌نامه‌ی ما:
	\textit{در تردید، سکوت بر شتاب‌زدگی ترجیح دارد.}
	ما نهاد مهار خشونت هستیم، نه ترویج آن.
	هر توصیه‌ی قهری فقط از مسیر دادرسی
	کامل و با رعایت اصل تناسب صادر می‌شود.» \\
	\bottomrule
\end{longtable}

\section{تقویم ارتباطی پیشنهادی}

\begin{longtable}{r p{4cm} p{7.5cm}}
	\toprule
	\textbf{زمان} & \textbf{رویداد} & \textbf{اقدام ارتباطی} \\
	\midrule
	\endhead
	هفته‌ی ۱ & تأسیس رسمی &
	کنفرانس مطبوعاتی (لاهه/لندن)؛
	بیانیه‌ی تأسیس به سه زبان
	(فارسی، انگلیسی، فرانسوی)؛
	مصاحبه با ۳–۵ رسانه‌ی کلیدی \\
	\midrule
	ماه ۱–۳ & دوره‌ی معرفی &
	سلسله‌پادکست «دادگاه در سایه چیست؟»
	(۶ قسمت)؛ مقاله‌ی توضیحی در
	\en{Just Security} یا \en{EJIL Talk} \\
	\midrule
	ماه ۶ & اولین رأی &
	انتشار همراه با بسته‌ی رسانه‌ای کامل؛
	ویدئوی توضیحی ۳ دقیقه‌ای \\
	\midrule
	سالگرد ژینا & سالانه &
	گزارش سالانه‌ی نهاد؛
	رویداد عمومی (حضوری/مجازی) \\
	\midrule
	روز حقوق بشر (۱۰ دسامبر) & سالانه &
	انتشار خلاصه‌ی مستندات سالانه؛
	بیانیه‌ی مشترک با نهادهای بین‌المللی \\
	\midrule
	سالگرد آبان ۹۸ & سالانه &
	انتشار به‌روزرسانی پرونده‌ی آبان؛
	یادبود قربانیان \\
	\bottomrule
\end{longtable}


% ////////////////////////////////////////////////////////////
%  پیوست ز: زیرساخت فنی و امنیت دیجیتال
% ////////////////////////////////////////////////////////////

% ////////////////////////////////////////////////////////////
%  پیوست ز: زیرساخت فنی و امنیت دیجیتال
% ////////////////////////////////////////////////////////////
\chapter{زیرساخت فنی و امنیت دیجیتال}\label{app:tech}

\section{معماری فنی پیشنهادی}

% ══════════════════════════════════════════════════════════════
%  اصلاح اساسی: نمودار TikZ با متن فارسی
%  ─────────────────────────────────────────────
%  ۱. تمام متن‌های فارسی درون \RL{} قرار گرفتند
%  ۲. جهت نودها از right=of به left=of تغییر کرد (برای RTL)
%  ۳. برچسب‌های فارسی روی پیکان‌ها در \RL{} قرار گرفتند
% ══════════════════════════════════════════════════════════════

\begin{figure}[H]
	\centering
	\begin{tikzpicture}[
		node distance=1.5cm and 2.5cm,
		box/.style={
			rectangle, rounded corners=3pt, draw=cPrimary, thick,
			fill=cLight, text width=3.8cm, minimum height=1.2cm,
			align=center, font=\footnotesize
		},
		secbox/.style={
			rectangle, rounded corners=3pt, draw=cAccent, thick,
			fill=cAccent!5, text width=3.8cm, minimum height=1.2cm,
			align=center, font=\footnotesize
		},
		arrow/.style={-{Stealth[length=4pt]}, thick, cPrimary}
		]
		% لایه‌ی عمومی
		\node[box] (web) {
			\textbf{\RL{وب‌سایت عمومی}}\\
			\RL{آراء، گزارش‌ها، اخبار}\\
			\en{Static site (Hugo/Jekyll)}
		};
		\node[box, left=of web] (social) {
			\textbf{\RL{شبکه‌های اجتماعی}}\\
			\RL{توییتر، یوتیوب، اینستاگرام}\\
			\RL{تلگرام (کانال عمومی)}
		};
		
		% لایه‌ی داخلی
		\node[secbox, below=of web] (db) {
			\textbf{\RL{بانک اطلاعاتی شواهد}}\\
			\RL{رمزنگاری} \en{AES-256}\\
			\RL{سرورهای مستقل}\\
			\RL{(سوئیس + ایسلند)}
		};
		\node[secbox, left=of db] (comm) {
			\textbf{\RL{ارتباطات داخلی}}\\
			\en{Signal / Element (Matrix)}\\
			\RL{ایمیل:} \en{ProtonMail}\\
			\RL{فایل:} \en{Tresorit}
		};
		
		% لایه‌ی بایگانی
		\node[box, below=of db] (archive) {
			\textbf{\RL{آرشیو بلندمدت}}\\
			\RL{نسخه‌ی آفلاین رمزنگاری‌شده}\\
			\RL{در ۲+ موقعیت جغرافیایی}\\
			+ \RL{نسخه‌ی} \en{LOCKSS}
		};
		\node[box, left=of archive] (backup) {
			\textbf{\RL{پشتیبان بین‌المللی}}\\
			\en{OHCHR} \RL{یا} \en{Swisspeace}\\
			\RL{نسخه‌ی اضطراری}
		};
		
		% پیکان‌ها
		\draw[arrow] (web) -- (db)
		node[midway, right, font=\tiny] {\RL{آراء منتشره}};
		\draw[arrow] (db) -- (archive)
		node[midway, right, font=\tiny] {\RL{بایگانی}};
		\draw[arrow] (comm) -- (db)
		node[midway, above, font=\tiny] {\RL{ورود شواهد}};
		\draw[arrow] (archive) -- (backup)
		node[midway, above, font=\tiny] {\RL{پشتیبان}};
		\draw[arrow, dashed, cAccent] (social) -- (comm)
		node[midway, left, font=\tiny] {\RL{مرز امنیتی}};
		
	\end{tikzpicture}
	\caption{\RL{معماری فنی سه‌لایه‌ی زیرساخت دیجیتال}}
\end{figure}

\section{الزامات امنیتی}

\begin{principleBox}[پروتکل امنیت دیجیتال]
	\begin{enumerate}[nosep, label=\textcolor{cAccent}{\arabic*.}]
		\item \textbf{رمزنگاری سرتاسری}
		(\en{End-to-End Encryption}):
		تمام ارتباطات داخلی و با شاهدان.
		
		\item \textbf{احراز هویت دوعاملی}
		(\en{2FA}):
		اجباری برای تمام اعضا و کارکنان.
		ترجیحاً کلید سخت‌افزاری (\en{YubiKey}).
		
		\item \textbf{\en{VPN} و \en{Tor}:}
		اجباری برای دسترسی به بانک اطلاعاتی.
		ارتباط با شاهدان داخل ایران فقط از طریق \en{Tor}.
		
		\item \textbf{حداقل‌سازی فراداده}
		(\en{Metadata Minimization}):
		تمام اسناد قبل از بایگانی از فراداده پاک‌سازی شوند
		(\en{ExifTool}).
		
		\item \textbf{تفکیک دسترسی}
		(\en{Principle of Least Privilege}):
		هر عضو فقط به اطلاعات مرتبط با وظیفه‌ی خود
		دسترسی داشته باشد.
		
		\item \textbf{ممیزی امنیتی سالانه:}
		حداقل یکبار در سال توسط شرکت امنیتی مستقل
		(مانند \en{Access Now} یا \en{Citizen Lab}).
		
		\item \textbf{پروتکل واکنش به رخنه}
		(\en{Breach Response Protocol}):
		در صورت نفوذ احتمالی، پروتکل از پیش تعیین‌شده
		فعال شود شامل: قطع دسترسی، اطلاع‌رسانی،
		ارزیابی خسارت، و بازیابی.
		
		\item \textbf{آموزش امنیتی:}
		تمام اعضا و کارکنان ملزم به گذراندن
		دوره‌ی آموزش امنیت دیجیتال (سالانه) هستند.
	\end{enumerate}
\end{principleBox}

\section{ابزارهای پیشنهادی}

\begin{longtable}{r p{3cm} p{4cm} p{4.5cm}}
	\toprule
	\textbf{کاربرد} & \textbf{ابزار پیشنهادی} & \textbf{جایگزین} &
	\textbf{ملاحظات} \\
	\midrule
	\endhead
	پیام‌رسانی & \en{Signal} & \en{Element (Matrix)} &
	\en{Signal} ساده‌تر؛ \en{Matrix}
	قابلیت‌های سازمانی بیشتر \\
	\midrule
	ایمیل & \en{ProtonMail} & \en{Tutanota} &
	هر دو سوئیسی/آلمانی و
	رمزنگاری‌شده \\
	\midrule
	فایل‌های مشترک & \en{Tresorit} & \en{Cryptpad} &
	\en{Tresorit} تجاری اما
	بسیار امن؛ \en{Cryptpad} متن‌باز \\
	\midrule
	بانک اطلاعاتی شواهد & \en{Uwazi (HURIDOCS)} & — &
	طراحی‌شده ویژه‌ی مستندسازی
	حقوق بشر. رایگان و متن‌باز. \\
	\midrule
	وب‌سایت & \en{Hugo + Netlify} & \en{Jekyll + GitHub Pages} &
	سایت ایستا: امن‌تر و سریع‌تر.
	بدون نیاز به سرور پویا. \\
	\midrule
	تحلیل منابع باز & \en{Hunchly} & \en{Wayback Machine} &
	\en{Hunchly} برای بایگانی
	خودکار صفحات وب.
	حیاتی برای \en{OSINT}. \\
	\midrule
	ویدئوکنفرانس & \en{Jitsi Meet} & \en{BigBlueButton} &
	هر دو متن‌باز و قابل
	میزبانی مستقل \\
	\midrule
	مدیریت پروژه & \en{Cryptpad Kanban} & \en{Taiga.io} &
	هر دو متن‌باز و
	قابل میزبانی مستقل \\
	\bottomrule
\end{longtable}


% ////////////////////////////////////////////////////////////
%  پیوست ح: کتاب‌نامه و منابع
% ////////////////////////////////////////////////////////////
\chapter{کتاب‌نامه و منابع}\label{app:biblio}

\section{کتاب‌ها}

\begin{enumerate}[nosep, label={[\arabic*]}]
	\item \en{Akhavan, Payam. \textit{Reducing Genocide to Law:
			Definition, Meaning, and the Ultimate Crime}.
		Cambridge University Press, 2012.}
	
	\item \en{Arendt, Hannah. \textit{Eichmann in Jerusalem:
			A Report on the Banality of Evil}. Viking Press, 1963.}
	
	\item \en{Bassiouni, M.\ Cherif. \textit{Introduction to
			International Criminal Law}. 2nd ed.\ Martinus Nijhoff, 2012.}
	
	\item \en{Bayat, Asef. \textit{Revolution without Revolutionaries:
			Making Sense of the Arab Spring}. Stanford University Press, 2017.}
	
	\item \en{De Greiff, Pablo (ed.). \textit{The Handbook of
			Reparations}. Oxford University Press, 2006.}
	
	\item \en{Hayner, Priscilla B. \textit{Unspeakable Truths:
			Transitional Justice and the Challenge of Truth Commissions}.
		2nd ed.\ Routledge, 2011.}
	
	\item \en{Jahanbegloo, Ramin. \textit{Introduction to Nonviolence}.
		Palgrave Macmillan, 2014.}
	
	\item \en{Khosrokhavar, Farhad. \textit{The New Arab Revolutions
			That Shook the World}. Paradigm, 2012.}
	
	\item \en{Nice, Geoffrey, and Francis Carey.
		\textit{The Uyghur Genocide: An Examination of China's
			Breaches of the 1948 Genocide Convention}.
		Uyghur Tribunal Report, 2021.}
	
	\item \en{Robertson, Geoffrey. \textit{Crimes Against Humanity:
			The Struggle for Global Justice}. 4th ed.\ Penguin, 2012.}
	
	\item \en{Schabas, William A. \textit{The International Criminal
			Court: A Commentary on the Rome Statute}. 2nd ed.\
		Oxford University Press, 2016.}
	
	\item \en{Teitel, Ruti G. \textit{Transitional Justice}.
		Oxford University Press, 2000.}
	
	\item \en{Tutu, Desmond. \textit{No Future Without Forgiveness}.
		Doubleday, 1999.}
\end{enumerate}

\section{مقالات و گزارش‌ها}

\begin{enumerate}[nosep, label={[\arabic*]}, resume]
	\item \en{Akhavan, Payam. ``Are International Criminal
		Tribunals a Disincentive to Peace?: Reconciling Judicial
		Romanticism with Political Realism.''
		\textit{Human Rights Quarterly} 31.3 (2009): 624–654.}
	
	\item \en{Amnesty International. ``Trampling Humanity:
		Mass Arrests, Disappearances and Torture Since
		Iran's November 2019 Protests.'' September 2020.}
	
	\item \en{Human Rights Watch. ``Deadly Repression:
		Iran's November 2019 Crackdown.'' January 2020.}
	
	\item \en{OHCHR. ``Report of the Special Rapporteur on the
		Situation of Human Rights in the Islamic Republic of Iran.''
		Annual Reports, 2011–2024.}
	
	\item \en{OHCHR. ``Independent International Fact-Finding
		Mission on the Islamic Republic of Iran: Interim Report.''
		2023.}
	
	\item \en{Robertson, Geoffrey. ``The Massacre of Political
		Prisoners in Iran, 1988.'' Report of an Inquiry conducted
		by Geoffrey Robertson QC. Abdorrahman Boroumand Foundation, 2010.}
	
	\item \en{Sadat, Leila N. ``Crimes Against Humanity in the
		Modern Age.'' \textit{American Journal of International
			Law} 107.2 (2013): 334–377.}
	
	\item \en{Schabas, William A. ``Genocide and the International
		Court of Justice: Finally, a Duty to Prevent the Crime of Crimes.''
		\textit{Genocide Studies and Prevention} 2.2 (2007): 101–122.}
\end{enumerate}

\section{وب‌سایت‌های مرجع}

\begin{enumerate}[nosep, label={[\arabic*]}, resume]
	\item \en{Abdorrahman Boroumand Center: \url{https://www.iranrights.org}}
	\item \en{Iran Human Rights: \url{https://iranhr.net}}
	\item \en{Center for Human Rights in Iran: \url{https://iranhumanrights.org}}
	\item \en{International Criminal Court: \url{https://www.icc-cpi.int}}
	\item \en{International Center for Transitional Justice: \url{https://www.ictj.org}}
	\item \en{Uyghur Tribunal: \url{https://uyghurtribunal.com}}
	\item \en{Permanent Peoples' Tribunal: \url{https://permanentpeoplestribunal.org}}
	\item \en{HURIDOCS (Uwazi): \url{https://huridocs.org}}
	\item \en{OHCHR Iran FFM: \url{https://www.ohchr.org/en/hr-bodies/hrc/ffm-iran}}
\end{enumerate}


% ////////////////////////////////////////////////////////////
%  پیوست ط: پرسش‌های متداول (FAQ)
% ////////////////////////////////////////////////////////////
\chapter{پرسش‌های متداول}\label{app:faq}

\begin{enumerate}[label=\textcolor{cPrimary}{\textbf{پ\arabic*.}}]
	
	\item \textbf{آیا این دادگاه واقعی است؟}
	
	خیر. دادگاه در سایه یک \keyword{نهاد مشورتی} است. آراء آن
	الزام قانونی ندارد اما ارزش حقوقی‌ـ‌اخلاقی دارد.
	مشابه دادگاه اویغورها یا دادگاه راسل.
	
	\item \textbf{چه تفاوتی با دادگاه ایران ۲۰۱۲ دارد؟}
	
	سه تفاوت اصلی:
	\begin{itemize}[nosep]
		\item ساختار دائمی (نه رویداد یکباره)؛
		\item صلاحیت موضوعی گسترده (نه فقط یک واقعه)؛
		\item وکیل تسخیری اجباری.
	\end{itemize}
	
	\item \textbf{چه کسی حق ارجاع پرونده دارد؟}
	
	هر شهروند، سازمان حقوق بشری، یا عضو هیئت عالی
	می‌تواند پرونده ارجاع دهد (فصل \ref{chap:proc}).
	
	\item \textbf{آیا مقام‌های حکومتی می‌توانند از خود دفاع کنند؟}
	
	بله. دفتر وکیل تسخیری موظف است بهترین استدلال ممکن را
	به نفع هر متهم ارائه دهد — حتی اگر متهم خودش
	مشارکت نکند.
	
	\item \textbf{آیا بودجه از دولت خارجی دریافت می‌شود؟}
	
	خیر. ممنوعیت مطلق دریافت مستقیم از دولت‌ها.
	فقط از بنیادهای مستقل و کمک‌های مردمی.
	
	\item \textbf{اگر نهاد اشتباه کند چه؟}
	
	مکانیزم بازنگری ۱۲ ماهه، حق اعتراض عمومی ۲۱ روزه،
	و قاعده‌ی عذرخواهی نهادی (فصل \ref{chap:risk}).
	
	\item \textbf{اگر جمهوری اسلامی سقوط کند، تکلیف نهاد چیست؟}
	
	چهار سناریو: ادغام کامل، تحول به نهاد ناظر، انحلال با بایگانی،
	یا تداوم مستقل (فصل \ref{chap:transition}).
	
	\item \textbf{آیا عفو مشروط به معنای فرار از عدالت است؟}
	
	خیر. عفو مشروط فقط در ازای اعتراف کامل و علنی،
	همکاری فعال، و جبران خسارت اعطا می‌شود.
	عاملان نسل‌کشی و شکنجه‌ی سیستماتیک مشمول نمی‌شوند
	(فصل \ref{chap:transition}).
	
	\item \textbf{چگونه می‌توانم مشارکت کنم؟}
	
	به‌عنوان: حقوق‌دان، پژوهشگر، مترجم، متخصص فناوری،
	شاهد/قربانی، حامی مالی، یا شهروند ناظر
	(فصل \ref{chap:conclusion}).
	
	\item \textbf{آیا اطلاعات من امن خواهد بود؟}
	
	پروتکل حفاظت شامل: رمزنگاری سرتاسری، ناشناس‌سازی،
	حداقل‌سازی داده، و ممیزی امنیتی سالانه
	(فصل \ref{chap:proc} و پیوست \ref{app:tech}).
	
\end{enumerate}


% ////////////////////////////////////////////////////////////
%  پیوست ی: نقشه‌ی راه اجرایی مفصل (گانت)
% ////////////////////////////////////////////////////////////
\chapter{نقشه‌ی راه اجرایی مفصل}\label{app:roadmap}

\begin{longtable}{r p{1cm} p{4cm} p{2.5cm} p{3cm}}
	\toprule
	\textbf{هفته} & \textbf{ماه} & \textbf{اقدام} &
	\textbf{مسئول} & \textbf{خروجی} \\
	\midrule
	\endhead
	۱–۴ & ۱ &
	تشکیل کمیته‌ی مؤسس (۳–۵ نفر)؛
	تدوین پیش‌نویس اساسنامه &
	بنیان‌گذاران &
	پیش‌نویس اساسنامه \\
	\midrule
	۵–۸ & ۲ &
	مشورت با حقوق‌دانان بین‌المللی؛
	اصلاح اساسنامه &
	کمیته‌ی مؤسس &
	اساسنامه‌ی نهایی \\
	\midrule
	۹–۱۲ & ۳ &
	ثبت قانونی در هلند (\en{Stichting})؛
	افتتاح حساب بانکی &
	مشاور حقوقی هلندی &
	ثبت رسمی \\
	\midrule
	۱۳–۱۶ & ۴ &
	دعوت از نامزدهای شورای ناظران بیرونی؛
	شروع جذب بودجه &
	کمیته‌ی مؤسس &
	تأیید ۳+ عضو شورای ناظران \\
	\midrule
	۱۷–۲۰ & ۵ &
	انتخاب اعضای هیئت عالی؛
	تأیید شورای ناظران &
	شورای ناظران &
	هیئت عالی تشکیل شود \\
	\midrule
	۲۱–۲۴ & ۶ &
	استخدام مدیر اجرایی و
	هماهنگ‌کننده‌ی پژوهش؛
	راه‌اندازی وب‌سایت &
	هیئت عالی &
	دبیرخانه عملیاتی \\
	\midrule
	۲۵–۲۸ & ۷ &
	\textbf{کنفرانس مطبوعاتی تأسیس}؛
	پذیرش اولین پرونده‌ی نوع الف &
	دبیرخانه + ارتباطات &
	اعلام عمومی تأسیس \\
	\midrule
	۲۹–۴۰ & ۸–۱۰ &
	تحقیقات مقدماتی اولین پرونده‌ها؛
	پذیرش اولین پرونده‌ی نوع ب &
	تیم پژوهش &
	گزارش‌های مقدماتی \\
	\midrule
	۴۱–۵۲ & ۱۱–۱۳ &
	صدور اولین کیفرخواست سایه؛
	صدور اولین رأی مشورتی &
	هیئت عالی &
	\textbf{اولین آراء منتشر شوند} \\
	\midrule
	۵۳–۷۸ & ۱۴–۱۸ &
	گسترش پرونده‌ها؛
	اولین جلسه‌ی استماع عمومی؛
	ایجاد شبکه‌ی وکلای داوطلب &
	تمام واحدها &
	تثبیت عملیاتی \\
	\midrule
	۷۹–۱۰۴ & ۱۹–۲۴ &
	انتشار گزارش سالانه‌ی اول؛
	حسابرسی مالی اول؛
	بازنگری ساختاری اول &
	شورای ناظران &
	\textbf{پایان سال اول} \\
	\bottomrule
\end{longtable}


% ============================================================
%  پیوست ک: فهرست جامع موضوعات قابل طرح، بررسی و داوری
%  (Comprehensive Docket of Adjudicable Matters)
% ============================================================

\chapter{فهرست جامع موضوعات قابل طرح، بررسی و داوری}
\label{app:topics}

\begin{principleBox}[درباره‌ی این فهرست]
	این پیوست \keyword{دستور کار موضوعی}
	(\en{Thematic Docket}) دادگاه در سایه را تعریف می‌کند:
	فهرستی ساختارمند از تمام موضوعاتی که نهاد صلاحیت
	بررسی آن‌ها را دارد، همراه با:
	\begin{itemize}[nosep]
		\item تشریح هر موضوع و ابعاد آن؛
		\item نوع پرونده (الف تا د)؛
		\item نمونه‌های تاریخی مشابه؛
		\item پرسش‌های کلیدی برای داوری؛
		\item ملاحظات ویژه.
	\end{itemize}
	
	\medskip
	این فهرست \textbf{باز و قابل بسط} است و با تحولات جنبش
	و شرایط بین‌المللی به‌روزرسانی خواهد شد.
\end{principleBox}

% ============================================================
%  بخش اول: تحریم‌ها — انواع، مدل‌ها، شدت
% ============================================================
\section{حوزه‌ی اول: تحریم‌ها}

\begin{warningBox}[چرا تحریم‌ها مهم‌ترین حوزه‌ی داوری هستند؟]
	تحریم‌ها ابزار اصلی فشار غیرنظامی بر جمهوری اسلامی هستند
	اما هم‌زمان \keyword{پُرهزینه‌ترین ابزار برای مردم عادی}.
	هر تصمیم درباره‌ی تحریم‌ها ذاتاً یک \keyword{معادله‌ی اخلاقی}
	است: آیا فشار بر حکومت، هزینه‌ی تحمیل‌شده بر مردم را
	توجیه می‌کند؟ دادگاه در سایه باید این معادله را
	شفاف و مستند حل کند.
\end{warningBox}

\subsection{موضوع ۱.۱: تحریم‌های اقتصادی جامع}

\begin{longtable}{r p{11cm}}
	\toprule
	\textbf{عنصر} & \textbf{شرح} \\
	\midrule
	\endhead
	\textbf{تعریف} &
	تحریم‌هایی که کلیت اقتصاد ایران را هدف می‌گیرند:
	تحریم نفتی، بانکی (\en{SWIFT})، تجاری، و مالی.
	شامل تحریم‌های یک‌جانبه‌ی آمریکا و تحریم‌های
	چندجانبه‌ی شورای امنیت. \\
	\midrule
	\textbf{نوع پرونده} &
	نوع ب (داوری تصمیم سیاسی) + نوع ج (تحلیل گفتمانی) \\
	\midrule
	\textbf{پرسش‌های کلیدی} &
	\begin{enumerate}[nosep, label=\textcolor{cPrimary}{\roman*.}]
		\item آیا جنبش باید از تشدید تحریم‌های جامع حمایت کند؟
		\item آیا استثناهای بشردوستانه (\en{humanitarian exemptions})
		کافی هستند؟
		\item تأثیر واقعی تحریم‌ها بر رفتار حکومت چقدر بوده؟
		\item آیا تحریم‌ها بیشتر به مردم آسیب می‌زنند یا به حکومت؟
		\item چه مدلی از تحریم بهینه‌ترین نسبت فشار‌ـ‌بر‌ـ‌حکومت
		به آسیب‌ـ‌بر‌ـ‌مردم را دارد؟
		\item آیا جایگزین‌های هدفمندتر وجود دارند؟
	\end{enumerate} \\
	\midrule
	\textbf{نمونه‌ی تاریخی} &
	\textbf{آفریقای جنوبی (۱۹۶۲–۱۹۹۳):}
	\en{ANC} فعالانه از تحریم‌های جامع حمایت کرد.
	نلسون ماندلا صریحاً گفت: «تحریم‌ها دردناک هستند
	اما آپارتاید دردناک‌تر است.» تحریم‌ها نقش مهمی
	در فروپاشی آپارتاید داشتند اما هزینه‌ی سنگینی
	بر سیاه‌پوستان فقیر تحمیل کردند.
	
	\medskip
	\textbf{عراق (۱۹۹۰–۲۰۰۳):}
	تحریم‌های جامع شورای امنیت بر عراق، طبق برآوردها
	به مرگ صدها هزار کودک منجر شد بدون آن‌که
	رژیم صدام را سرنگون کند. این نمونه‌ی
	\keyword{شکست فاجعه‌بار} تحریم جامع بدون
	هدف‌گذاری دقیق است.
	
	\medskip
	\textbf{درس:} تحریم جامع بدون استثناهای بشردوستانه
	مؤثر و بدون هدف‌گذاری دقیق، ممکن است
	ضدّاخلاقی و ناکارآمد باشد. \\
	\midrule
	\textbf{ملاحظات ویژه} &
	\begin{itemize}[nosep]
		\item تفکیک بین تحریم‌های هوشمند (\en{smart sanctions})
		و تحریم‌های کور (\en{blanket sanctions}).
		\item بررسی تأثیر بر دسترسی به دارو، غذا، و کالاهای اساسی.
		\item ارزیابی اثربخشی واقعی (نه فرضی) در تغییر رفتار حکومت.
		\item ارائه‌ی توصیه‌ی متمایز: مثلاً «تحریم نفتی ادامه یابد
		اما مکانیزم بشردوستانه‌ی دارویی تقویت شود.»
	\end{itemize} \\
	\bottomrule
\end{longtable}


\subsection{موضوع ۱.۲: تحریم‌های هدفمند (فردی و نهادی)}

\begin{longtable}{r p{11cm}}
	\toprule
	\textbf{عنصر} & \textbf{شرح} \\
	\midrule
	\endhead
	\textbf{تعریف} &
	تحریم‌هایی که اشخاص حقیقی (مقام‌ها، فرماندهان، قضات)
	یا اشخاص حقوقی (سپاه، بسیج، نهادهای اقتصادی وابسته)
	را هدف می‌گیرند: مسدودسازی دارایی، ممنوعیت سفر،
	محرومیت از خدمات بانکی.
	شامل: قانون \en{Magnitsky} (آمریکا)،
	\en{CAATSA}، تحریم‌های اتحادیه اروپا،
	و \en{Global Human Rights Sanctions} (بریتانیا). \\
	\midrule
	\textbf{نوع پرونده} &
	نوع الف (مستندسازی + کیفرخواست سایه) +
	نوع ب (توصیه‌ی تحریم خاص) \\
	\midrule
	\textbf{پرسش‌های کلیدی} &
	\begin{enumerate}[nosep, label=\textcolor{cPrimary}{\roman*.}]
		\item چه اشخاصی باید در فهرست تحریم قرار گیرند؟
		\item آیا شواهد کافی علیه هر شخص وجود دارد؟
		\item تحت کدام قانون تحریم قابل اعمال است؟
		\item آیا تحریم شخص خاص اثربخش خواهد بود؟
		\item آیا حقوق افراد بی‌تقصیر (خانواده) رعایت شده؟
		\item اولویت‌بندی: چه کسانی اول باید تحریم شوند؟
	\end{enumerate} \\
	\midrule
	\textbf{نمونه‌ی تاریخی} &
	\textbf{قانون مگنیتسکی (۲۰۱۲/۲۰۱۶):}
	پس از مرگ سرگئی مگنیتسکی در زندان روسیه،
	آمریکا قانونی تصویب کرد که تحریم فردی مقام‌های
	ناقض حقوق بشر را ممکن ساخت.
	این مدل بعداً توسط کانادا، بریتانیا، و اتحادیه اروپا
	تکرار شد. \keyword{مؤثرترین مدل تحریم هدفمند.}
	
	\medskip
	\textbf{میانمار (از ۲۰۲۱):}
	پس از کودتای نظامی، فهرست تحریم فردی
	فرماندهان ارتش میانمار توسط چندین کشور
	اعمال شد. دادگاه در سایه می‌تواند مشابه
	همین مدل، فهرست تحریم فردی برای مقام‌های
	ایرانی تهیه کند. \\
	\midrule
	\textbf{خروجی مورد انتظار} &
	\begin{itemize}[nosep]
		\item \keyword{فهرست تحریم سایه}
		(\en{Shadow Sanctions List}):
		فهرست اولویت‌بندی‌شده‌ی اشخاص حقیقی و حقوقی
		با مستندات حقوقی آماده‌ی ارائه به دولت‌ها.
		\item \keyword{بسته‌ی حقوقی} (\en{Sanctions Case Package})
		برای هر شخص: شامل شرح جرایم، شواهد، مبنای حقوقی،
		و درخواست مشخص از دولت هدف.
	\end{itemize} \\
	\bottomrule
\end{longtable}


\subsection{موضوع ۱.۳: تحریم‌های ورزشی}

\begin{longtable}{r p{11cm}}
	\toprule
	\textbf{عنصر} & \textbf{شرح} \\
	\midrule
	\endhead
	\textbf{تعریف} &
	تعلیق یا اخراج تیم‌ها و فدراسیون‌های ورزشی جمهوری اسلامی
	از رقابت‌های بین‌المللی (\en{FIFA}, \en{IOC},
	\en{FIVB}, \en{BWF} و …)؛
	یا اقدامات محدودتر مانند: تحریم مقام‌های ورزشی حکومتی
	بدون تعلیق ورزشکاران. \\
	\midrule
	\textbf{نوع پرونده} &
	نوع ب (داوری تصمیم سیاسی) \\
	\midrule
	\textbf{پرسش‌های کلیدی} &
	\begin{enumerate}[nosep, label=\textcolor{cPrimary}{\roman*.}]
		\item آیا تعلیق کامل (\en{full suspension}) توصیه می‌شود
		یا تحریم هدفمند مقام‌های ورزشی؟
		\item هزینه‌ی ورزشکاران حرفه‌ای بی‌تقصیر چقدر است؟
		\item آیا ورزشکارانی که پرچم جمهوری اسلامی را حمل می‌کنند
		اما شعار سیاسی نمی‌دهند «خاکستری» هستند یا «همراه»؟
		\item آیا ورزشکارانی که به دلیل فشار حکومت مجبور به
		رقابت نکردن با حریفان اسرائیلی شده‌اند،
		قربانی هستند یا مباشر؟
		\item مدل جایگزین: آیا می‌توان ورزشکاران ایرانی را
		تحت پرچم بی‌طرف (\en{Neutral Athletes})
		مشابه مدل روسیه اجازه‌ی رقابت داد؟
		\item آیا فدراسیون‌های ورزشی ایران مستقل هستند
		یا بازوی دولت؟ (شواهد: کنترل سپاه بر
		فدراسیون‌های ورزشی)
	\end{enumerate} \\
	\midrule
	\textbf{نمونه‌ی تاریخی} &
	\textbf{آفریقای جنوبی (۱۹۶۴–۱۹۹۲):}
	تعلیق کامل از \en{IOC} و اکثر فدراسیون‌ها.
	یکی از مؤثرترین ابزارهای فشار بین‌المللی بود.
	اما هزینه‌ی سنگینی بر ورزشکاران سیاه‌پوست
	تحمیل کرد که فرصت رقابت از دست دادند.
	
	\medskip
	\textbf{روسیه (از ۲۰۲۲):}
	پس از حمله به اوکراین، \en{IOC} ورزشکاران روس
	را فقط تحت پرچم بی‌طرف و بدون نماد ملی
	اجازه‌ی رقابت داد. \en{FIFA} تیم ملی روسیه را
	کاملاً تعلیق کرد. \keyword{مدل ترکیبی:}
	تفکیک بین رشته‌ها و سطوح تعلیق.
	
	\medskip
	\textbf{یوگوسلاوی (۱۹۹۲):}
	تعلیق از المپیک و جام جهانی.
	ورزشکاران مجاز به رقابت فردی
	تحت عنوان «شرکت‌کنندگان مستقل» بودند.
	
	\medskip
	\textbf{درس:} سه مدل وجود دارد:
	(الف) تعلیق کامل؛ (ب) پرچم بی‌طرف؛
	(ج) تحریم مقام‌ها بدون تعلیق ورزشکاران.
	دادگاه در سایه باید هر سه مدل را برای هر رشته
	جداگانه بررسی و توصیه‌ی متمایز ارائه دهد. \\
	\midrule
	\textbf{ملاحظات ویژه} &
	\begin{itemize}[nosep]
		\item برخی ورزشکاران ایرانی (مانند کشتی‌گیران و
		جودوکاران مجبور به \en{forfeit}) خود قربانی
		حکومت هستند — نباید دوبار مجازات شوند.
		\item ورزشکارانی که فعالانه از حکومت حمایت تبلیغاتی
		کرده‌اند باید جداگانه بررسی شوند.
		\item اصل \keyword{تفکیک نهاد از فرد}.
	\end{itemize} \\
	\bottomrule
\end{longtable}


\subsection{موضوع ۱.۴: تحریم‌های فرهنگی و هنری}

\begin{longtable}{r p{11cm}}
	\toprule
	\textbf{عنصر} & \textbf{شرح} \\
	\midrule
	\endhead
	\textbf{تعریف} &
	تحریم نهادهای فرهنگی وابسته به حکومت
	(سازمان فرهنگ و ارتباطات اسلامی،
	رایزنی‌های فرهنگی سفارتخانه‌ها، جشنواره‌ی فجر)؛
	تحریم یا محدودسازی شرکت هنرمندان و فیلم‌سازانی که
	نماینده‌ی رسمی حکومت هستند؛
	بحث درباره‌ی جایگاه هنرمندان مستقل. \\
	\midrule
	\textbf{نوع پرونده} &
	نوع ب + نوع ج \\
	\midrule
	\textbf{پرسش‌های کلیدی} &
	\begin{enumerate}[nosep, label=\textcolor{cPrimary}{\roman*.}]
		\item آیا تحریم فرهنگی جامع توصیه می‌شود یا هدفمند؟
		\item خط فاصل بین «هنرمند مستقل» و «نماینده‌ی فرهنگی حکومت» کجاست؟
		\item هنرمندانی که در جشنواره‌های حکومتی شرکت می‌کنند
		اما اعتراضات خیابانی را حمایت نکرده‌اند: وضعیت‌شان چیست؟
		\item آیا تحریم فرهنگی اصولاً مؤثر است؟
		\item آیا \keyword{بایکوت فرهنگی}
		(\en{cultural boycott}) با آزادی بیان تعارض دارد؟
	\end{enumerate} \\
	\midrule
	\textbf{نمونه‌ی تاریخی} &
	\textbf{جنبش \en{BDS} (فلسطین):}
	بایکوت فرهنگی اسرائیل مدل بحث‌برانگیزی است.
	برخی هنرمندان (راجر واترز) از آن حمایت و
	برخی (رادیوهِد) آن را رد کردند.
	\keyword{درس:} بایکوت فرهنگی همیشه بحث‌برانگیز
	خواهد بود — دادگاه در سایه باید معیارهای
	شفاف و قابل‌دفاع ارائه دهد.
	
	\medskip
	\textbf{آفریقای جنوبی:}
	بایکوت فرهنگی شامل منع اجراهای بین‌المللی
	هنرمندان در آفریقای جنوبی بود —
	نه بایکوت هنرمندان آفریقای جنوبی.
	\keyword{مدل معکوس} که شاید برای ایران هم
	مناسب‌تر باشد: به‌جای تحریم هنرمندان ایرانی،
	تحریم رویدادهای فرهنگی حکومتی. \\
	\midrule
	\textbf{ملاحظات ویژه} &
	\begin{itemize}[nosep]
		\item فرهنگ و هنر ابزار قدرت نرم حکومت هستند —
		تحریم آن‌ها مشروعیت دارد.
		\item اما هنرمندان مستقل نباید قربانی شوند.
		\item معیار پیشنهادی: \keyword{آیا شخص نماینده‌ی رسمی
			نهاد حکومتی است یا به‌عنوان فرد مستقل فعالیت می‌کند؟}
	\end{itemize} \\
	\bottomrule
\end{longtable}


\subsection{موضوع ۱.۵: تحریم‌های علمی و دانشگاهی}

\begin{longtable}{r p{11cm}}
	\toprule
	\textbf{عنصر} & \textbf{شرح} \\
	\midrule
	\endhead
	\textbf{تعریف} &
	قطع یا محدودسازی همکاری‌های دانشگاهی با نهادهای
	آموزشی وابسته به حکومت (دانشگاه‌های نظامی، دانشگاه
	امام حسین، مراکز پژوهشی وابسته به سپاه)؛
	بحث درباره‌ی همکاری با دانشگاه‌های عمومی ایران. \\
	\midrule
	\textbf{نوع پرونده} &
	نوع ب + نوع ج \\
	\midrule
	\textbf{پرسش‌های کلیدی} &
	\begin{enumerate}[nosep, label=\textcolor{cPrimary}{\roman*.}]
		\item آیا تحریم علمی جامع توصیه می‌شود؟
		\item تفکیک بین نهادهای علمی نظامی و غیرنظامی؟
		\item تأثیر بر دانشجویان و پژوهشگران بی‌تقصیر؟
		\item نقش دانشگاه‌های ایرانی در برنامه‌های تسلیحاتی؟
		\item آیا تبادل دانشجویی باید ادامه یابد
		(به‌عنوان ابزار نفوذ فرهنگی مثبت)؟
	\end{enumerate} \\
	\midrule
	\textbf{نمونه‌ی تاریخی} &
	\textbf{روسیه (از ۲۰۲۲):}
	بسیاری از دانشگاه‌های غربی همکاری با نهادهای
	روسی را قطع کردند اما بورسیه‌ی دانشجویان
	فرّاری روسی را ادامه دادند — مدل «تحریم نهاد،
	حمایت از فرد».
	
	\medskip
	\textbf{ایران (از ۲۰۰۶):}
	تحریم‌های \en{OFAC} همکاری علمی برخی
	دانشگاه‌های آمریکایی با ایران را محدود کرده —
	اما اثرات جانبی بر دانشجویان ایرانی در آمریکا
	نیز داشته. \\
	\bottomrule
\end{longtable}


\subsection{موضوع ۱.۶: تحریم‌های دیپلماتیک}

\begin{longtable}{r p{11cm}}
	\toprule
	\textbf{عنصر} & \textbf{شرح} \\
	\midrule
	\endhead
	\textbf{تعریف} &
	درخواست کاهش سطح روابط دیپلماتیک،
	اخراج دیپلمات‌ها، بستن سفارتخانه‌ها،
	یا فراخوان سفرا. \\
	\midrule
	\textbf{نوع پرونده} &
	نوع ب \\
	\midrule
	\textbf{پرسش‌های کلیدی} &
	\begin{enumerate}[nosep, label=\textcolor{cPrimary}{\roman*.}]
		\item آیا قطع کامل روابط دیپلماتیک مفید است
		یا بسته شدن کانال‌های ارتباطی زیان‌بار خواهد بود؟
		\item آیا سفارت‌ها و نمایندگی‌های ایران
		مراکز جاسوسی و عملیات ترور هستند؟ (شواهد؟)
		\item تأثیر بر خدمات کنسولی شهروندان ایرانی؟
		\item مدل بینابینی: حفظ روابط حداقلی برای
		کانال‌های بشردوستانه؟
	\end{enumerate} \\
	\midrule
	\textbf{نمونه‌ی تاریخی} &
	\textbf{آلبانی (۲۰۱۸):}
	اخراج سفیر و اولین دبیر سفارت ایران
	به دلیل طراحی عملیات تروریستی.
	
	\medskip
	\textbf{کانادا (۲۰۱۲):}
	قطع روابط دیپلماتیک با ایران و بستن سفارت. \\
	\bottomrule
\end{longtable}


% ============================================================
%  بخش دوم: مستندسازی جرایم
% ============================================================
\section{حوزه‌ی دوم: مستندسازی جرایم حکومتی}

\subsection{موضوع ۲.۱: کشتارهای جمعی}

\begin{longtable}{r p{11cm}}
	\toprule
	\textbf{عنصر} & \textbf{شرح} \\
	\midrule
	\endhead
	\textbf{تعریف} &
	مستندسازی نظام‌مند کشتارهای جمعی توسط
	نیروهای امنیتی جمهوری اسلامی شامل:
	کشتار ۶۷، آبان ۹۸، کشتار زاهدان ۱۴۰۱،
	و سرکوب جنبش ژینا. \\
	\midrule
	\textbf{نوع پرونده} &
	نوع الف (مستندسازی + کیفرخواست سایه) \\
	\midrule
	\textbf{پرسش‌های کلیدی} &
	\begin{enumerate}[nosep, label=\textcolor{cPrimary}{\roman*.}]
		\item تعداد دقیق قربانیان هر واقعه چقدر است؟
		\item زنجیره‌ی فرماندهی در هر واقعه چه بود؟
		\item آیا شواهد کافی برای احراز «جرم علیه بشریت» وجود دارد؟
		\item آیا شواهد کافی برای احراز «نسل‌کشی» وجود دارد؟
		(به‌ویژه در مورد اقلیت‌های قومی و مذهبی)
		\item کدام مقام‌ها باید در فهرست کیفرخواست قرار گیرند؟
	\end{enumerate} \\
	\midrule
	\textbf{خروجی مورد انتظار} &
	\begin{itemize}[nosep]
		\item کیفرخواست سایه برای هر واقعه
		\item فهرست مظنونین با مستندات
		\item بسته‌ی ارجاع به \en{ICC} و دادگاه‌های ملی
	\end{itemize} \\
	\bottomrule
\end{longtable}


\subsection{موضوع ۲.۲: شکنجه‌ی سیستماتیک}

\begin{longtable}{r p{11cm}}
	\toprule
	\textbf{عنصر} & \textbf{شرح} \\
	\midrule
	\endhead
	\textbf{تعریف} &
	مستندسازی شکنجه‌ی سیستماتیک در زندان‌ها و
	بازداشتگاه‌های جمهوری اسلامی شامل:
	شکنجه‌ی جسمی، تجاوز، انفرادی طولانی‌مدت،
	شکنجه‌ی روانی، و محرومیت از درمان. \\
	\midrule
	\textbf{نوع پرونده} &
	نوع الف \\
	\midrule
	\textbf{پرسش‌های کلیدی} &
	\begin{enumerate}[nosep, label=\textcolor{cPrimary}{\roman*.}]
		\item کدام بازداشتگاه‌ها مراکز اصلی شکنجه هستند؟
		\item کدام بازجوها و مقام‌ها عاملان مستقیم هستند؟
		\item آیا شکنجه «سیستماتیک» است یا «موردی»؟
		(معیار: آیا الگوی مشترک و سیاست مرکزی وجود دارد؟)
		\item آیا شواهد برای صلاحیت جهانی کافی است؟
	\end{enumerate} \\
	\bottomrule
\end{longtable}


\subsection{موضوع ۲.۳: اعدام‌های خودسرانه}

\begin{longtable}{r p{11cm}}
	\toprule
	\textbf{عنصر} & \textbf{شرح} \\
	\midrule
	\endhead
	\textbf{تعریف} &
	مستندسازی اعدام‌هایی که بدون رعایت استانداردهای
	حداقلی دادرسی عادلانه صورت گرفته‌اند شامل:
	اعدام‌های سیاسی، اعدام نوجوانان،
	اعدام‌های مرتبط با اعتراضات، و
	اعدام‌های مبتنی بر اعترافات اجباری. \\
	\midrule
	\textbf{نوع پرونده} &
	نوع الف \\
	\midrule
	\textbf{همکاری‌های موجود} &
	\en{Iran Human Rights} (نروژ) و
	\en{Abdorrahman Boroumand Center}
	بانک اطلاعاتی گسترده‌ای از اعدام‌ها دارند.
	دادگاه در سایه می‌تواند با این نهادها همکاری کند. \\
	\bottomrule
\end{longtable}


% ============================================================
%  بخش سوم: موضوعات سیاسی-اخلاقی
% ============================================================
\section{حوزه‌ی سوم: موضوعات سیاسی‌ـ‌اخلاقی}

\subsection{موضوع ۳.۱: مداخله‌ی بشردوستانه}

\begin{longtable}{r p{11cm}}
	\toprule
	\textbf{عنصر} & \textbf{شرح} \\
	\midrule
	\endhead
	\textbf{تعریف} &
	بررسی شرایط و حدود مداخله‌ی نظامی خارجی
	برای حمایت از مردم ایران بر اساس اصل
	«مسئولیت حمایت» (\en{R2P}). \\
	\midrule
	\textbf{نوع پرونده} &
	نوع ب + نوع ج \\
	\midrule
	\textbf{پرسش‌های کلیدی} &
	\begin{enumerate}[nosep, label=\textcolor{cPrimary}{\roman*.}]
		\item آیا شرایط فعلی ایران معیارهای \en{R2P} را
		برآورده می‌کند؟
		\item اگر بله، چه نوع مداخله‌ای مجاز و مناسب است؟
		(منطقه‌ی پرواز ممنوع؟ حمایت تسلیحاتی از مقاومت؟
		مداخله‌ی زمینی؟)
		\item هزینه‌های احتمالی مداخله بر مردم عادی چیست؟
		\item آیا جایگزین‌های غیرنظامی همه آزمایش شده‌اند؟
		\item موضع رسمی جنبش درباره‌ی مداخله چه باید باشد؟
	\end{enumerate} \\
	\midrule
	\textbf{ملاحظات ویژه} &
	\begin{itemize}[nosep]
		\item این حساس‌ترین موضوع سیاسی است.
		\item هر رأی باید با بیشترین احتیاط صادر شود.
		\item نظرات اقلیت حتماً ثبت شوند.
		\item جمهوری اسلامی از هر اظهارنظر درباره‌ی مداخله
		بهره‌برداری تبلیغاتی خواهد کرد.
	\end{itemize} \\
	\bottomrule
\end{longtable}


\subsection{موضوع ۳.۲: برخورد با وابستگان حکومتی در خارج}

\begin{longtable}{r p{11cm}}
	\toprule
	\textbf{عنصر} & \textbf{شرح} \\
	\midrule
	\endhead
	\textbf{تعریف} &
	بررسی سیاست مناسب در قبال فرزندان، همسران،
	و بستگان مقام‌های حکومتی که در کشورهای غربی
	زندگی می‌کنند. \\
	\midrule
	\textbf{نوع پرونده} &
	نوع ب \\
	\midrule
	\textbf{پرسش‌های کلیدی} &
	\begin{enumerate}[nosep, label=\textcolor{cPrimary}{\roman*.}]
		\item آیا درخواست اخراج فرزندان مقام‌ها مشروع است؟
		\item تفکیک بین «فرزند بالغ همراه» و «فرزند بی‌تقصیر»؟
		\item آیا مصادره‌ی اموال خانواده‌ها مجاز است؟
		\item معیار تفکیک: سن؟ میزان همراهی فعال؟ منبع ثروت؟
		\item آیا «شرمسارسازی عمومی» (\en{public shaming})
		ابزار مشروعی است؟
	\end{enumerate} \\
	\midrule
	\textbf{ملاحظات ویژه} &
	\begin{itemize}[nosep]
		\item اصل \keyword{مسئولیت فردی}: هیچ‌کس به‌خاطر
		جرم پدر یا مادرش مجازات نمی‌شود.
		\item اما: اگر فرزند از ثروت نامشروع پدر/مادر
		بهره‌مند است، ردیابی و مصادره‌ی آن ثروت مجاز است.
		\item معیار پیشنهادی: \keyword{آیا شخص فعالانه
			از نظام حمایت کرده یا صرفاً فرزند مقام است؟}
	\end{itemize} \\
	\bottomrule
\end{longtable}


\subsection{موضوع ۳.۳: خشونت در جنبش}

\begin{longtable}{r p{11cm}}
	\toprule
	\textbf{عنصر} & \textbf{شرح} \\
	\midrule
	\endhead
	\textbf{تعریف} &
	بررسی حدود مشروع «دفاع مشروع» و «مقاومت فعال»
	در برابر سرکوب حکومتی. \\
	\midrule
	\textbf{نوع پرونده} &
	نوع ب + نوع ج \\
	\midrule
	\textbf{پرسش‌های کلیدی} &
	\begin{enumerate}[nosep, label=\textcolor{cPrimary}{\roman*.}]
		\item آیا حمله به نیروهای سرکوب‌گر در لحظه‌ی سرکوب
		«دفاع مشروع» است؟
		\item آیا خرابکاری در تأسیسات حکومتی مجاز است؟
		\item خط قرمزها کجاست؟ (حمله به غیرنظامیان؟
		حمله به خانواده‌ی نیروهای امنیتی؟)
		\item آیا جنبش باید رسماً از «مقاومت مسلحانه» حمایت کند؟
		\item تفکیک بین «خشونت دفاعی» و «تروریسم»؟
	\end{enumerate} \\
	\midrule
	\textbf{ملاحظات ویژه} &
	\begin{itemize}[nosep]
		\item دادگاه در سایه نهاد «مهار خشونت» است، نه ترویج آن.
		\item هر رأی باید بر اساس حقوق بین‌الملل بشردوستانه
		و اصول جنگ عادلانه (\en{jus ad bellum} و
		\en{jus in bello}) باشد.
		\item نظرات اقلیت حتماً ثبت شوند.
	\end{itemize} \\
	\bottomrule
\end{longtable}


% ============================================================
%  جمع‌بندی پیوست ک
% ============================================================
\section{جمع‌بندی}

\begin{principleBox}[خلاصه‌ی پیوست ک]
	\begin{enumerate}[nosep, label=\textcolor{cPrimary}{\arabic*.}]
		\item این فهرست \textbf{۱۲ موضوع اصلی} در سه حوزه را
		پوشش می‌دهد: تحریم‌ها، مستندسازی جرایم،
		و موضوعات سیاسی‌ـ‌اخلاقی.
		\item فهرست \textbf{باز و قابل بسط} است.
		\item هر موضوع شامل: تعریف، نوع پرونده، پرسش‌های کلیدی،
		نمونه‌های تاریخی، و ملاحظات ویژه است.
		\item هدف: ایجاد \keyword{چارچوب ساختارمند}
		برای داوری مسائل پیچیده‌ی جنبش.
	\end{enumerate}
\end{principleBox}

\clearpage

% ============================================================
%  بخش سوم: مداخله‌ی بشردوستانه و R2P
% ============================================================
\section{حوزه‌ی سوم: مداخله‌ی بشردوستانه و \en{R2P}}

\begin{warningBox}[حساس‌ترین موضوع]
	درخواست مداخله‌ی نظامی خارجی \keyword{بحث‌برانگیزترین}
	موضوع در اپوزیسیون ایران است و بالاترین پتانسیل
	شکاف‌افکنی و بهره‌برداری تبلیغاتی حکومت را دارد.
	دادگاه در سایه باید این موضوع را با بیشترین دقت
	و کمترین شتاب بررسی کند.
\end{warningBox}

\subsection{موضوع ۳.۱: ارزیابی شرایط \en{R2P}}

\begin{longtable}{r p{11cm}}
	\toprule
	\textbf{عنصر} & \textbf{شرح} \\
	\midrule
	\endhead
	\textbf{تعریف} &
	ارزیابی اینکه آیا شرایط ایران معیارهای «مسئولیت حمایت»
	(\en{Responsibility to Protect}) مصوب اجلاس جهانی ۲۰۰۵
	سازمان ملل را برآورده می‌کند یا خیر.
	\en{R2P} شامل سه ستون است:
	(۱) مسئولیت دولت برای حمایت از شهروندان؛
	(۲) مسئولیت جامعه‌ی بین‌المللی برای کمک؛
	(۳) مسئولیت جامعه‌ی بین‌المللی برای مداخله
	در صورت ناکامی دولت. \\
	\midrule
	\textbf{نوع پرونده} &
	نوع ج (تحلیل گفتمانی) — سنگین‌ترین سطح \\
	\midrule
	\textbf{چارچوب ارزیابی} &
	\begin{principleBox}[شش معیار مداخله‌ی نظامی مشروع — بر اساس گزارش \en{ICISS} ۲۰۰۱]
		\begin{enumerate}[nosep, label=\textcolor{cAccent}{\arabic*.}]
			\item \textbf{علت موجه (\en{Just Cause}):}
			آیا نسل‌کشی، پاکسازی قومی، جرایم جنگی،
			یا جرایم علیه بشریت در مقیاس بزرگ
			در حال وقوع یا قریب‌الوقوع است؟
			
			\item \textbf{نیت درست (\en{Right Intention}):}
			آیا هدف اصلی مداخله، توقف رنج انسانی است
			یا منافع ژئوپلیتیک مداخله‌کنندگان؟
			
			\item \textbf{آخرین راه‌حل (\en{Last Resort}):}
			آیا تمام ابزارهای غیرنظامی (تحریم، دیپلماسی،
			ارجاع به \en{ICC}، \ldots) امتحان شده یا
			به‌وضوح ناکارآمد شناخته شده‌اند؟
			
			\item \textbf{ابزار متناسب (\en{Proportional Means}):}
			آیا مقیاس و مدت مداخله حداقل لازم است
			و آسیب جانبی قابل قبول؟
			
			\item \textbf{چشم‌انداز معقول (\en{Reasonable Prospects}):}
			آیا مداخله احتمال معقولی برای موفقیت دارد
			و وضعیت را بدتر نخواهد کرد؟
			
			\item \textbf{مرجع صالح (\en{Right Authority}):}
			آیا مجوز شورای امنیت یا یک ائتلاف مشروع
			بین‌المللی وجود دارد؟
		\end{enumerate}
	\end{principleBox} \\
	\midrule
	\textbf{پرسش‌های کلیدی} &
	\begin{enumerate}[nosep, label=\textcolor{cPrimary}{\roman*.}]
		\item آیا سطح سرکوب در ایران به آستانه‌ی \en{R2P} رسیده؟
		(مقایسه با لیبی ۲۰۱۱، سوریه ۲۰۱۱، رواندا ۱۹۹۴)
		\item آیا جنبش باید فعالانه از مداخله حمایت کند
		یا صرفاً آن را «ممانعت نکند»؟
		\item خطر تبدیل شدن ایران به سوریه‌ی دوم چقدر است؟
		\item اگر مداخله‌ی تمام‌عیار نظامی مطلوب نیست،
		چه طیفی از اقدامات «زیرِ آستانه‌ی جنگ»
		(\en{below-the-threshold measures}) پیشنهاد می‌شود؟
		\item نقش منطقه‌ای (\en{regional arrangements})
		طبق فصل هشتم منشور ملل متحد چیست؟
	\end{enumerate} \\
	\midrule
	\textbf{خروجی پیشنهادی} &
	\begin{itemize}[nosep]
		\item «ارزیابی وضعیت» (\en{Situation Assessment})
		مستند و مفصل (۵۰+ صفحه)
		\item نقشه‌ی حرارتی (\en{heat-map}) سطح سرکوب
		در استان‌ها و حوزه‌ها
		\item مقایسه‌ی تطبیقی با ۵ مورد تاریخی
		\item توصیه‌نامه‌ی محرمانه با سطح‌بندی سناریوها
	\end{itemize} \\
	\midrule
	\textbf{نمونه‌ی تاریخی الهام‌بخش} &
	\begin{itemize}[nosep]
		\item \textbf{لیبی ۲۰۱۱:} قطعنامه‌ی ۱۹۷۳ شورای امنیت —
		تنها مورد اجرای کامل \en{R2P} نظامی
		\item \textbf{کوزوو ۱۹۹۹:} مداخله‌ی ناتو بدون مجوز شورای امنیت —
		«غیرقانونی اما مشروع»
		(\en{illegal but legitimate})
		\item \textbf{رواندا ۱۹۹۴:} عدم مداخله و نسل‌کشی ۸۰۰,۰۰۰ نفر —
		نمونه‌ی «شکست اخلاقی» جامعه‌ی بین‌المللی
	\end{itemize} \\
	\bottomrule
\end{longtable}

% ─────────────────────────────────────────
\subsection{موضوع ۳.۲: درخواست ایجاد مناطق پرواز ممنوع و کریدور بشردوستانه}

\begin{longtable}{r p{11cm}}
	\toprule
	\textbf{عنصر} & \textbf{شرح} \\
	\midrule
	\endhead
	\textbf{تعریف} &
	بررسی امکان‌سنجی و تدوین درخواست رسمی برای ایجاد:
	\begin{itemize}[nosep]
		\item مناطق پرواز ممنوع (\en{No-Fly Zone})
		بر فراز مناطق با بالاترین شدت سرکوب
		(الگوی عراق شمالی ۱۹۹۱ و لیبی ۲۰۱۱)
		\item کریدورهای بشردوستانه
		(\en{Humanitarian Corridors})
		برای دسترسی سازمان‌های امدادی
		\item مناطق امن
		(\en{Safe Zones / Safe Havens})
		برای حفاظت از غیرنظامیان
	\end{itemize} \\
	\midrule
	\textbf{نوع پرونده} &
	نوع ج (تحلیل گفتمانی) — سنگین \\
	\midrule
	\textbf{سابقه‌ی تاریخی} &
	\begin{itemize}[nosep]
		\item \textbf{عراق شمالی ۱۹۹۱ (\en{Operation Provide Comfort}):}
		حفاظت از کردها پس از قیام ناموفق
		\item \textbf{بوسنی ۱۹۹۳–۱۹۹۵:} مناطق امن سازمان ملل
		(شکست سربرنیتسا — درس عبرت)
		\item \textbf{سوریه ۲۰۱۲–۲۰۲۰:} بحث‌های ناتمام
		درباره‌ی منطقه‌ی پرواز ممنوع
	\end{itemize} \\
	\midrule
	\textbf{ملاحظات کلیدی} &
	\begin{warningBox}[هشدار عملیاتی]
		تجربه‌ی سربرنیتسا (ژوئیه ۱۹۹۵) نشان داد که
		اعلام «منطقه‌ی امن» بدون ضمانت اجرایی
		و نیروی بازدارنده‌ی کافی، می‌تواند فاجعه‌ای
		بدتر بیافریند: تمرکز غیرنظامیان در مناطقی
		که عملاً محافظت نمی‌شوند.
	\end{warningBox} \\
	\bottomrule
\end{longtable}

% ─────────────────────────────────────────
\subsection{موضوع ۳.۳: درخواست ارجاع وضعیت ایران به دیوان کیفری بین‌المللی}

\begin{longtable}{r p{11cm}}
	\toprule
	\textbf{عنصر} & \textbf{شرح} \\
	\midrule
	\endhead
	\textbf{تعریف} &
	تدوین درخواست مستدل برای ارجاع وضعیت ایران
	(\en{Referral of the Situation in Iran})
	به دیوان کیفری بین‌المللی (\en{ICC})
	از طریق شورای امنیت (ماده‌ی ۱۳(ب) اساسنامه‌ی رم)
	— با توجه به اینکه ایران عضو اساسنامه‌ی رم نیست. \\
	\midrule
	\textbf{نوع پرونده} &
	نوع ب (اعلام‌دادخواهی بین‌المللی) — سنگین \\
	\midrule
	\textbf{ساختار درخواست} &
	\begin{enumerate}[nosep, label=\textcolor{cAccent}{\arabic*.}]
		\item شناسایی جرایم در صلاحیت \en{ICC}:
		\begin{itemize}[nosep]
			\item جرایم علیه بشریت (ماده‌ی ۷)
			\item احتمال نسل‌کشی فرهنگی (ماده‌ی ۶)
			\item شکنجه‌ی سیستماتیک
			\item تعقیب و آزار به دلایل سیاسی، جنسیتی،
			قومی و مذهبی
		\end{itemize}
		\item ناکامی یا عدم تمایل نظام قضایی داخلی
		(\en{unwillingness or inability}) طبق ماده‌ی ۱۷
		\item مستندسازی الگوی سیستماتیک (\en{chapeau elements})
		\item شناسایی مسئولیت فرماندهی (\en{command responsibility})
	\end{enumerate} \\
	\midrule
	\textbf{نمونه‌ی تاریخی} &
	\begin{itemize}[nosep]
		\item \textbf{سودان / دارفور ۲۰۰۵:}
		ارجاع شورای امنیت (قطعنامه‌ی ۱۵۹۳)
		— صدور حکم بازداشت عمر البشیر
		\item \textbf{لیبی ۲۰۱۱:}
		ارجاع شورای امنیت (قطعنامه‌ی ۱۹۷۰)
		— صدور حکم بازداشت معمر قذافی
		\item \textbf{میانمار ۲۰۱۸:}
		ارجاع توسط بنگلادش (عضو اساسنامه‌ی رم)
		به دلیل عبور قربانیان از مرز —
		الگوی جایگزین برای دولت‌های غیرعضو
	\end{itemize} \\
	\midrule
	\textbf{خروجی} &
	\begin{itemize}[nosep]
		\item متن رسمی درخواست ارجاع (۲۰+ صفحه)
		\item ضمیمه‌ی مستندات حمایتی
		\item فهرست امضاکنندگان (حقوقدانان، سازمان‌ها)
		\item نقشه‌ی لابی‌گری در شورای امنیت
	\end{itemize} \\
	\bottomrule
\end{longtable}

% ─────────────────────────────────────────
\subsection{موضوع ۳.۴: درخواست تأسیس دادگاه ویژه یا تشکیل ساختار هیبریدی}

\begin{longtable}{r p{11cm}}
	\toprule
	\textbf{عنصر} & \textbf{شرح} \\
	\midrule
	\endhead
	\textbf{تعریف} &
	بررسی و تدوین طرح پیشنهادی برای تأسیس
	یکی از اَشکال زیر:
	\begin{enumerate}[nosep, label=\textcolor{cPrimary}{\alph*.}]
		\item \textbf{دادگاه ویژه‌ی بین‌المللی}
		(\en{Special International Tribunal}) —
		به‌صورت \en{ad hoc} مانند ICTY یا ICTR
		\item \textbf{دادگاه هیبریدی / مختلط}
		(\en{Hybrid Tribunal}) —
		ترکیب عناصر بین‌المللی و داخلی
		(الگوی سیرالئون، کامبوج، لبنان)
		\item \textbf{دادگاه ایرانی در تبعید}
		(\en{Court in Exile}) —
		با قضات ایرانی و نظارت بین‌المللی
	\end{enumerate} \\
	\midrule
	\textbf{نوع پرونده} &
	نوع ج (تحلیل گفتمانی) — سنگین‌ترین سطح \\
	\midrule
	\textbf{مقایسه‌ی الگوها} &
	% ══════════════════════════════════════════════════════════════
	%  اصلاح: TikZ با متن فارسی — تمام متن‌ها در \RL{}
	%  جهت نودها برای RTL اصلاح شد
	% ══════════════════════════════════════════════════════════════
	\begin{tikzpicture}[
		node distance=0.3cm,
		every node/.style={
			font=\small,
			align=center,
			text width=3.5cm,
			minimum height=1.2cm,
			rounded corners=3pt
		}
		]
		\node[fill=cPrimary!15, draw=cPrimary] (icty)
		{\textbf{ICTY}\\\footnotesize \RL{یوگوسلاوی}\\
			\footnotesize \RL{کاملاً بین‌المللی}};
		\node[fill=cAccent!15, draw=cAccent, left=of icty] (scsl)
		{\textbf{SCSL}\\\footnotesize \RL{سیرالئون}\\
			\footnotesize \RL{هیبریدی}};
		\node[fill=cWarn!15, draw=cWarn, left=of scsl] (eccc)
		{\textbf{ECCC}\\\footnotesize \RL{کامبوج}\\
			\footnotesize \RL{هیبریدی / ملی‌تر}};
		\node[fill=green!15, draw=green!60!black, left=of eccc] (iran)
		{\textbf{\RL{ایران؟}}\\\footnotesize \RL{مدل پیشنهادی}\\
			\footnotesize \RL{← کدام الگو؟}};
		\draw[-{Stealth[length=4pt]}, thick, cPrimary] (icty) -- (scsl);
		\draw[-{Stealth[length=4pt]}, thick, cAccent] (scsl) -- (eccc);
		\draw[-{Stealth[length=4pt]}, thick, dashed, cWarn] (eccc) -- (iran);
	\end{tikzpicture} \\
	\midrule
	\textbf{نمونه‌های تاریخی مفصل} &
	\begin{enumerate}[nosep, label=\textcolor{cAccent}{\arabic*.}]
		\item \textbf{ICTY (یوگوسلاوی، ۱۹۹۳–۲۰۱۷):}
		\begin{itemize}[nosep]
			\item تأسیس با قطعنامه‌ی ۸۲۷ شورای امنیت
			\item ۱۶۱ متهم، ۹۰ محکومیت
			\item محاکمه‌ی میلوشویچ (فوت حین محاکمه)
			\item محاکمه‌ی کارادژیچ و ملادیچ
			\item هزینه‌ی کل: ≈ ۲.۴ میلیارد دلار
			\item \keyword{درس:} طولانی‌بودن دادرسی‌ها
		\end{itemize}
		
		\item \textbf{ICTR (رواندا، ۱۹۹۴–۲۰۱۵):}
		\begin{itemize}[nosep]
			\item تأسیس با قطعنامه‌ی ۹۵۵ شورای امنیت
			\item ۹۳ متهم، ۶۲ محکومیت
			\item اولین محکومیت برای نسل‌کشی در تاریخ
			حقوق بین‌الملل (پرونده‌ی آکایسو)
			\item اولین شناسایی تجاوز جنسی
			به‌عنوان ابزار نسل‌کشی
			\item \keyword{درس:} اهمیت مشارکت جامعه‌ی محلی
		\end{itemize}
		
		\item \textbf{SCSL (سیرالئون، ۲۰۰۲–۲۰۱۳):}
		\begin{itemize}[nosep]
			\item تأسیس با توافق سازمان ملل و دولت سیرالئون
			\item محاکمه‌ی چارلز تیلور (رئیس‌جمهور لیبریا)
			\item هزینه‌ی بسیار کمتر از ICTY/ICTR
			\item \keyword{درس:} الگوی هیبریدی مقرون‌به‌صرفه‌تر
		\end{itemize}
		
		\item \textbf{STL (لبنان، ۲۰۰۷–۲۰۲۳):}
		\begin{itemize}[nosep]
			\item تأسیس برای رسیدگی به ترور رفیق حریری
			\item محاکمه‌ی غیابی اعضای حزب‌الله
			\item \keyword{درس:} امکان محاکمه‌ی غیابی
			و محدودیت‌های آن
		\end{itemize}
	\end{enumerate} \\
	\bottomrule
\end{longtable}


% ============================================================
%  بخش چهارم: تحریم‌ها — طبقه‌بندی و ارزیابی
% ============================================================
\section{حوزه‌ی چهارم: تحریم‌ها — طبقه‌بندی، ارزیابی و پیشنهاد}

\begin{infoBox}[اهمیت موضوع]
	تحریم‌ها مهم‌ترین ابزار فشار بین‌المللی غیرنظامی
	هستند. دادگاه در سایه می‌تواند نقش مهمی در
	\keyword{ارزیابی اثربخشی}، \keyword{پیشنهاد تحریم‌های هدفمند}
	و \keyword{مستندسازی آسیب‌های جانبی} ایفا کند.
\end{infoBox}

\subsection{موضوع ۴.۱: تحریم‌های هدفمند فردی (\en{Targeted / Smart Sanctions})}

\begin{longtable}{r p{11cm}}
	\toprule
	\textbf{عنصر} & \textbf{شرح} \\
	\midrule
	\endhead
	\textbf{تعریف} &
	شناسایی، مستندسازی و درخواست تحریم علیه
	\keyword{افراد مشخص} مسئول نقض حقوق بشر.
	شامل:
	\begin{itemize}[nosep]
		\item مسدود کردن دارایی‌ها (\en{Asset Freeze})
		\item ممنوعیت سفر (\en{Travel Ban})
		\item ممنوعیت صدور ویزا
		\item محرومیت از خدمات بانکی بین‌المللی
	\end{itemize} \\
	\midrule
	\textbf{نوع پرونده} &
	نوع ب (اعلام‌دادخواهی بین‌المللی) \\
	\midrule
	\textbf{چارچوب‌های قانونی موجود} &
	\begin{enumerate}[nosep, label=\textcolor{cPrimary}{\arabic*.}]
		\item \textbf{قانون مگنیتسکی جهانی}
		(\en{Global Magnitsky Act}) — ایالات متحده
		\item \textbf{نظام تحریمی حقوق بشر اتحادیه‌ی اروپا}
		(\en{EU Global Human Rights Sanctions Regime}) — ۲۰۲۰
		\item \textbf{قانون مگنیتسکی کانادا}
		(\en{Justice for Victims of Corrupt Foreign
			Officials Act}) — ۲۰۱۷
		\item \textbf{قانون تحریم‌های مگنیتسکی بریتانیا}
		(\en{UK Sanctions and Anti-Money Laundering Act}) — ۲۰۱۸
		\item \textbf{قانون تحریمی استرالیا}
		(\en{Autonomous Sanctions Act}) — ۲۰۱۱
	\end{enumerate} \\
	\midrule
	\textbf{فعالیت دادگاه} &
	\begin{enumerate}[nosep, label=\textcolor{cAccent}{\roman*.}]
		\item تهیه‌ی «پرونده‌ی تحریمی»
		(\en{Sanctions Dossier}) برای هر فرد:
		\begin{itemize}[nosep]
			\item هویت کامل و مشخصات
			\item سمت و جایگاه در ساختار قدرت
			\item مستندات نقض حقوق بشر
			\item دارایی‌های شناسایی‌شده در خارج
			\item ارتباطات مالی بین‌المللی
		\end{itemize}
		\item ارسال رسمی به واحدهای \en{OFAC}
		(آمریکا)، \en{CFSP} (اروپا) و معادل‌ها
		\item پیگیری و ارزیابی دوره‌ای اثربخشی
	\end{enumerate} \\
	\midrule
	\textbf{دسته‌بندی اهداف تحریمی} &
	\begin{principleBox}[طبقه‌بندی سه‌لایه‌ی اهداف تحریمی]
		\begin{itemize}[nosep]
			\item \textbf{لایه‌ی ۱ — رهبری:}
			رهبر، اعضای مجمع تشخیص، شورای نگهبان،
			فرماندهان ارشد سپاه و بسیج
			\item \textbf{لایه‌ی ۲ — مجریان مستقیم:}
			فرماندهان عملیاتی سرکوب، رؤسای زندان‌ها،
			بازجویان، قضات شعب امنیتی
			\item \textbf{لایه‌ی ۳ — تسهیل‌کنندگان:}
			مدیران رسانه‌های حکومتی، مسئولان سانسور،
			دانشمندان هسته‌ای درگیر در برنامه‌ی نظامی
		\end{itemize}
	\end{principleBox} \\
	\bottomrule
\end{longtable}

% ─────────────────────────────────────────
\subsection{موضوع ۴.۲: تحریم‌های بخشی (\en{Sectoral Sanctions})}

\begin{longtable}{r p{11cm}}
	\toprule
	\textbf{عنصر} & \textbf{شرح} \\
	\midrule
	\endhead
	\textbf{تعریف} &
	بررسی و ارزیابی تحریم‌هایی که بخش‌های اقتصادی
	را هدف می‌گیرند. هدف: شناسایی \keyword{تحریم‌های مؤثر}
	بر حکومت و \keyword{کم‌آسیب} به مردم. \\
	\midrule
	\textbf{بخش‌های هدف} &
	\begin{enumerate}[nosep, label=\textcolor{cAccent}{\arabic*.}]
		\item \textbf{نفت و گاز:} — ستون فقرات مالی حکومت
		\item \textbf{پتروشیمی و فلزات:} — منبع دوم درآمد ارزی
		\item \textbf{بانکداری و مالی:} — قطع از \en{SWIFT}،
		محدودیت تراکنش‌های بین‌المللی
		\item \textbf{فناوری و ارتباطات:}
		\begin{itemize}[nosep]
			\item تحریم تجهیزات نظارت و سانسور
			(\en{surveillance technology})
			\item ممنوعیت فروش \en{DPI} و ابزار فیلترینگ
		\end{itemize}
		\item \textbf{تسلیحاتی و نظامی:}
		— تحریم تسلیحاتی جامع
		\item \textbf{هسته‌ای:}
		— محدودیت‌های فناوری دومنظوره
		\item \textbf{حمل‌ونقل:}
		— تحریم کشتیرانی و هواپیمایی
	\end{enumerate} \\
	\midrule
	\textbf{تحلیل آسیب جانبی} &
	\begin{warningBox}[وظیفه‌ی اخلاقی دادگاه]
		دادگاه در سایه \keyword{موظف} است برای هر
		پیشنهاد تحریم بخشی، «ارزیابی تأثیر
		بر جمعیت غیرنظامی»
		(\en{Civilian Impact Assessment})
		تهیه کند. تجربه‌ی تحریم‌های عراق (۱۹۹۰–۲۰۰۳)
		نشان داد که تحریم‌های فراگیر می‌توانند
		به فاجعه‌ی انسانی منجر شوند.
	\end{warningBox}
	\begin{itemize}[nosep]
		\item تحلیل آثار بر قیمت دارو و غذا
		\item تحلیل آثار بر اشتغال و معیشت
		\item پیشنهاد معافیت‌های بشردوستانه
		(\en{Humanitarian Exemptions})
		\item مکانیسم پایش مستمر (\en{Monitoring Mechanism})
	\end{itemize} \\
	\bottomrule
\end{longtable}

% ─────────────────────────────────────────
\subsection{موضوع ۴.۳: تحریم‌های ثانویه و فراسرزمینی (\en{Secondary / Extraterritorial Sanctions})}

\begin{longtable}{r p{11cm}}
	\toprule
	\textbf{عنصر} & \textbf{شرح} \\
	\midrule
	\endhead
	\textbf{تعریف} &
	ارزیابی تحریم‌هایی که اشخاص ثالث (شرکت‌ها و
	کشورهای غیرآمریکایی) را نیز هدف قرار می‌دهند
	اگر با ایران تجارت کنند. \\
	\midrule
	\textbf{پرسش‌های کلیدی} &
	\begin{enumerate}[nosep, label=\textcolor{cPrimary}{\roman*.}]
		\item آیا تحریم‌های ثانویه مؤثرتر از اولیه‌اند؟
		\item مشروعیت حقوق بین‌المللی آنها چیست؟
		\item آیا از آنها حمایت شود یا فقط پذیرفته شوند؟
		\item اثر بر متحدان بالقوه (چین، هند، ترکیه) چیست؟
	\end{enumerate} \\
	\midrule
	\textbf{نمونه} &
	\begin{itemize}[nosep]
		\item تحریم‌های ثانویه آمریکا علیه
		خریداران نفت ایران (۲۰۱۸–اکنون)
		\item \en{CAATSA} (تحریم روسیه، ایران، کره‌ی شمالی)
		\item قانون «تحکیم تحریم‌های ایران» (\en{ISA})
	\end{itemize} \\
	\bottomrule
\end{longtable}

% ─────────────────────────────────────────
\subsection{موضوع ۴.۴: ارزیابی شدت و تناسب تحریم‌ها}

\begin{longtable}{r p{11cm}}
	\toprule
	\textbf{عنصر} & \textbf{شرح} \\
	\midrule
	\endhead
	\textbf{تعریف} &
	توسعه‌ی یک \keyword{مقیاس سنجش شدت تحریم}
	(\en{Sanctions Severity Index})
	برای ارزیابی و مقایسه‌ی سطح فشار وارده. \\
	\midrule
	\textbf{معیارهای سنجش} &
	\begin{principleBox}[ماتریس ارزیابی تحریم]
		هر تحریم از ۵ بُعد ارزیابی می‌شود (هر کدام ۱ تا ۵):
		\begin{enumerate}[nosep, label=\textcolor{cAccent}{\arabic*.}]
			\item \textbf{اثرگذاری بر حکومت} (۱=ناچیز، ۵=فلج‌کننده)
			\item \textbf{آسیب به مردم عادی} (۱=ناچیز، ۵=فاجعه‌بار)
			\item \textbf{امکان دور زدن} (۱=غیرممکن، ۵=بسیار آسان)
			\item \textbf{مشروعیت بین‌المللی} (۱=جنجالی، ۵=اجماع کامل)
			\item \textbf{بازگشت‌پذیری} (۱=برگشت‌ناپذیر، ۵=سریعاً قابل رفع)
		\end{enumerate}
		\textbf{نسبت مطلوب:} حداکثر امتیاز بُعد ۱ + حداقل امتیاز بُعد ۲
	\end{principleBox} \\
	\midrule
	\textbf{نتیجه‌ی ارزیابی} &
	\begin{itemize}[nosep]
		\item تحریم‌های \textbf{با اولویت بالا}:
		شاخص بالای ۳.۵ (مؤثر بر حکومت، کم‌آسیب بر مردم)
		\item تحریم‌های \textbf{مشکل‌دار}:
		مؤثر ولی آسیب‌زای جانبی بالا — نیاز به اصلاح
		\item تحریم‌های \textbf{ناکارآمد}:
		بدون اثر بر حکومت، فقط مردم را آزار می‌دهند — درخواست رفع
	\end{itemize} \\
	\bottomrule
\end{longtable}


% ============================================================
%  بخش پنجم: مواجهه با عوامل، نفوذی‌ها و نیروهای نیابتی
% ============================================================
\section{حوزه‌ی پنجم: عوامل، نفوذی‌ها و دست‌نشاندگان رژیم}

\begin{warningBox}[حساسیت بسیار بالا]
	این حوزه \keyword{بیشترین خطر سوءاستفاده} را دارد.
	اتهام «نفوذی» و «عامل رژیم» به‌سادگی
	می‌تواند ابزار حذف مخالفان درون‌گروهی شود.
	بنابراین \keyword{بالاترین استانداردهای ادله}
	و \keyword{محکم‌ترین ضمانت‌های دادرسی منصفانه}
	در این حوزه ضروری است.
\end{warningBox}

\subsection{موضوع ۵.۱: تعریف و طبقه‌بندی}

\begin{longtable}{r p{11cm}}
	\toprule
	\textbf{عنصر} & \textbf{شرح} \\
	\midrule
	\endhead
	\textbf{دسته‌بندی} &
	\begin{principleBox}[طبقه‌بندی چهارگانه]
		\begin{enumerate}[nosep, label=\textcolor{cAccent}{\arabic*.}]
			\item \textbf{عامل اطلاعاتی مستقیم (\en{Agent}):}
			فردی که مستقیماً از نهادهای اطلاعاتی‌/امنیتی
			جمهوری اسلامی (وزارت اطلاعات، سازمان اطلاعات سپاه)
			مأموریت دریافت کرده و گزارش می‌دهد.
			\textit{بالاترین سطح خطر.}
			
			\item \textbf{نفوذی سازمان‌یافته (\en{Infiltrator}):}
			فردی که به‌صورت سیستماتیک در سازمان‌های
			اپوزیسیون یا رسانه‌ها نفوذ کرده تا
			اطلاعات جمع‌آوری یا اختلال ایجاد کند.
			\textit{ممکن است خود داوطلب یا اجیرشده باشد.}
			
			\item \textbf{دست‌نشانده‌ی سیاسی (\en{Proxy}):}
			فردی که بدون ارتباط اطلاعاتی مستقیم،
			مواضع و سیاست‌هایی منطبق بر منافع جمهوری اسلامی
			ترویج می‌کند — ممکن است از روی اعتقاد یا منفعت مالی باشد.
			\textit{خاکستری‌ترین حوزه.}
			
			\item \textbf{همسو (\en{Fellow Traveler}):}
			فردی که بدون ارتباط مستقیم یا مأموریت،
			گاه مواضعی موافق حکومت اتخاذ می‌کند —
			ممکن است صادقانه باشد.
			\textit{باید از اتهام‌زنی مصون بماند.}
		\end{enumerate}
	\end{principleBox} \\
	\midrule
	\textbf{خط قرمز} &
	\begin{warningBox}[اصل بنیادین]
		دادگاه در سایه تنها حق بررسی دسته‌های ۱ و ۲ را دارد.
		دسته‌ی ۳ فقط در صورت وجود مستندات مالی یا ارتباطی
		بررسی می‌شود. دسته‌ی ۴ \keyword{هرگز} موضوع بررسی قرار نمی‌گیرد.
		\textbf{اختلاف نظر سیاسی جرم نیست.}
	\end{warningBox} \\
	\bottomrule
\end{longtable}

\subsection{موضوع ۵.۲: معیارهای ادله و دادرسی}

\begin{longtable}{r p{11cm}}
	\toprule
	\textbf{عنصر} & \textbf{شرح} \\
	\midrule
	\endhead
	\textbf{سطح ادله} &
	\begin{itemize}[nosep]
		\item \textbf{حداقل:} «فراتر از شک معقول»
		(\en{beyond reasonable doubt}) — بالاترین استاندارد
		\item هیچ پرونده‌ای بر مبنای «شایعه»، «حدس»،
		«اعتراف اجباری» یا «گزارش ناشناس واحد»
		قابل پذیرش نیست
		\item \keyword{اصل برائت}: هر فرد بی‌گناه فرض می‌شود
		تا خلاف آن اثبات شود
	\end{itemize} \\
	\midrule
	\textbf{انواع ادله‌ی قابل قبول} &
	\begin{enumerate}[nosep, label=\textcolor{cAccent}{\arabic*.}]
		\item اسناد مالی (حواله‌ها، حساب‌های بانکی، قراردادها)
		\item مکاتبات مکتوب یا الکترونیکی (ایمیل، پیام‌رسان)
		\item شهادت حداقل دو شاهد مستقل و معتبر
		\item اسناد رسمی حکومتی (در صورت دسترسی)
		\item گزارش‌های تحقیقاتی معتبر
		(رسانه‌ها، سازمان‌های حقوق بشری)
		\item ادله‌ی دیجیتال با قابلیت تأیید فنی
		(\en{digital forensics})
	\end{enumerate} \\
	\midrule
	\textbf{ضمانت‌های دادرسی} &
	\begin{enumerate}[nosep, label=\textcolor{cPrimary}{\roman*.}]
		\item حق اطلاع از اتهام به‌صورت کامل و مکتوب
		\item حق داشتن وکیل مدافع (حتی تعیین‌شده از سوی دادگاه)
		\item حق دسترسی به کلیه‌ی ادله‌ی مورد استفاده
		\item حق پاسخ‌گویی و دفاع کامل
		\item حق تجدیدنظر
		\item حق جبران خسارت در صورت تبرئه
		\item \keyword{محرمانگی}: تا زمان صدور رأی، هویت متهم علنی نمی‌شود
	\end{enumerate} \\
	\bottomrule
\end{longtable}

\subsection{موضوع ۵.۳: سازوکارهای مواجهه و نتایج}

\begin{longtable}{r p{11cm}}
	\toprule
	\textbf{عنصر} & \textbf{شرح} \\
	\midrule
	\endhead
	\textbf{در صورت اثبات} &
	\begin{enumerate}[nosep, label=\textcolor{cAccent}{\arabic*.}]
		\item \textbf{اطلاع‌رسانی محرمانه} به سازمان‌ها و نهادهای مرتبط
		\item \textbf{اعلام عمومی} (فقط دسته‌ی ۱ با ادله‌ی محکم)
		\item \textbf{ارجاع به مقامات کشور محل اقامت}
		(در صورت نقض قوانین آن کشور: جاسوسی، تحریم‌شکنی و …)
		\item \textbf{درخواست اخراج دیپلماتیک} (\en{persona non grata})
		\item \textbf{ارائه‌ی مستندات به سازمان‌های بین‌المللی}
		\item \textbf{ثبت در بانک اطلاعاتی} برای دوره‌ی گذار
	\end{enumerate} \\
	\midrule
	\textbf{نمونه‌های تاریخی} &
	\begin{itemize}[nosep]
		\item \textbf{ترور میکونوس (۱۹۹۲):}
		دادگاه آلمان حکم به دست‌داشتن رهبران ایران داد
		— اخراج دیپلمات‌ها
		\item \textbf{بمب‌گذاری آمیا (آرژانتین ۱۹۹۴):}
		صدور حکم بازداشت بین‌المللی (\en{Red Notice})
		برای مقامات ایرانی
		\item \textbf{ترور فخری‌زاده (بحث‌های حقوقی):}
		مرز بین عملیات اطلاعاتی و حاکمیت دولت‌ها
		\item \textbf{شبکه‌ی اطلاعاتی شتازی (\en{Stasi})
			آلمان شرقی:}
		پس از اتحاد آلمان، میلیون‌ها پرونده
		بازگشایی و هزاران نفوذی شناسایی شدند —
		الگوی «گاوک‌بِهورده» (\en{Gauck-Behörde})
	\end{itemize} \\
	\bottomrule
\end{longtable}


% ============================================================
%  بخش ششم: اخراج و محدودسازی وابستگان سران
% ============================================================
\section{حوزه‌ی ششم: وابستگان و فرزندان سران نظام}

\begin{warningBox}[حساسیت حقوقی و اخلاقی]
	هرگونه درخواست مرتبط با اعضای خانواده‌ی مقامات
	باید با \keyword{اصل عدم مجازات جمعی}
	(\en{principle of non-collective punishment})
	و \keyword{حقوق فردی} سازگار باشد.
	دادگاه نمی‌تواند صرفاً به دلیل قرابت خانوادگی
	کسی را مجرم بشناسد — اما می‌تواند
	\keyword{شبکه‌ی منافع} و \keyword{جریان ثروت}
	را مستند و افشا کند.
\end{warningBox}

\subsection{موضوع ۶.۱: تعریف «سران» و «وابستگان»}

\begin{longtable}{r p{11cm}}
	\toprule
	\textbf{عنصر} & \textbf{شرح} \\
	\midrule
	\endhead
	\textbf{تعریف سران} &
	\begin{principleBox}[تعریف پیشنهادی «سران نظام»]
		فرد یا افرادی که در دوره‌ی مورد بررسی
		(۱۳۵۷ تا زمان گذار) یکی از سِمَت‌های زیر
		را داشته‌اند یا دارند:
		\begin{enumerate}[nosep, label=\textcolor{cAccent}{\arabic*.}]
			\item رهبری (ولایت فقیه)
			\item ریاست جمهوری
			\item ریاست قوه‌ی قضائیه
			\item ریاست مجلس شورای اسلامی
			\item عضویت شورای نگهبان
			\item عضویت مجمع تشخیص مصلحت نظام
			\item فرماندهی کل سپاه پاسداران و معاونان ارشد
			\item وزارت اطلاعات (وزیر و معاونان)
			\item فرماندهی نیروی انتظامی (فراجا)
			\item ریاست صداوسیما
			\item دبیری شورای عالی امنیت ملی
			\item دادستانی کل کشور و روسای دیوان عالی
			\item فرماندهان سرکوب مستقیم
			(فرمانده بسیج، فرمانده یگان ویژه و …)
		\end{enumerate}
	\end{principleBox} \\
	\midrule
	\textbf{تعریف وابستگان} &
	\begin{enumerate}[nosep, label=\textcolor{cPrimary}{\arabic*.}]
		\item \textbf{خویشاوندان درجه یک:}
		همسر، فرزندان، والدین
		\item \textbf{خویشاوندان درجه دو:}
		خواهر و برادر، داماد و عروس
		\item \textbf{وابستگان مالی (\en{financial associates}):}
		افرادی که بدون قرابت خونی، از طریق
		شرکت‌ها، بنیادها یا بنگاه‌های اقتصادی
		با سران مرتبط‌اند
		\item \textbf{وابستگان ایدئولوژیک:}
		\begin{warningBox}[هشدار]
			این دسته \keyword{باید با نهایت احتیاط} تعریف شود.
			صرف داشتن اعتقادات مذهبی یا سیاسی
			مشابه حکومت \keyword{دلیل بررسی نیست}.
			تنها در صورتی که فرد فعالانه
			در \keyword{ترویج سرکوب} یا \keyword{تطهیر نقض حقوق بشر}
			مشارکت داشته باشد.
		\end{warningBox}
	\end{enumerate} \\
	\bottomrule
\end{longtable}

\subsection{موضوع ۶.۲: اقدامات پیشنهادی}

\begin{longtable}{r p{11cm}}
	\toprule
	\textbf{اقدام} & \textbf{شرح و مبنای حقوقی} \\
	\midrule
	\endhead
	\textbf{۱. مستندسازی ثروت} &
	شناسایی و مستندسازی دارایی‌های خارجی سران
	و وابستگان:
	\begin{itemize}[nosep]
		\item املاک (کانادا، امارات، ترکیه، اروپا و …)
		\item حساب‌های بانکی
		\item شرکت‌های پوششی (\en{shell companies})
		\item سرمایه‌گذاری‌ها
	\end{itemize}
	\textit{مبنا: قوانین مبارزه با پول‌شویی
		(\en{AML}) و تأمین مالی تروریسم (\en{CFT})} \\
	\midrule
	\textbf{۲. درخواست لغو اقامت و ویزا} &
	\begin{itemize}[nosep]
		\item تهیه‌ی پرونده‌ی حقوقی برای درخواست
		لغو اقامت افرادی که از ثروت نامشروع
		برای اخذ اقامت استفاده کرده‌اند
		\item استناد به قوانین «ویزای طلایی»
		(\en{Golden Visa}) و مبارزه با فساد
		\item \textbf{نکته:} این اقدام باید فردی و مستند باشد،
		نه جمعی و خودسرانه
	\end{itemize} \\
	\midrule
	\textbf{۳. درخواست مسدودسازی دارایی} &
	\begin{itemize}[nosep]
		\item ارائه‌ی مستندات به واحدهای \en{FIU}
		(\en{Financial Intelligence Unit}) کشورهای مقصد
		\item درخواست مسدودسازی موقت (\en{freezing order})
		تا روشن شدن منبع ثروت
	\end{itemize} \\
	\midrule
	\textbf{۴. درخواست اخراج دیپلماتیک} &
	\begin{itemize}[nosep]
		\item فقط برای کسانی که دارای سمت دیپلماتیک
		یا شبه‌دیپلماتیک هستند
		\item مبنا: سوءاستفاده از مصونیت دیپلماتیک
		\item نمونه: اخراج دیپلمات‌های ایرانی
		پس از حکم میکونوس (۱۹۹۷)
	\end{itemize} \\
	\midrule
	\textbf{۵. درخواست اخراج فرزندان و بستگان} &
	\begin{warningBox}[محدودیت حقوقی مهم]
		در نظام‌های حقوقی غربی، «اخراج» صرفاً
		بر مبنای قرابت خانوادگی \keyword{امکان‌پذیر نیست}
		مگر:
		\begin{itemize}[nosep]
			\item فرد خود مرتکب تخلف شده باشد
			(اقامت با مدارک جعلی، فرار مالیاتی، پول‌شویی)
			\item فرد تهدید امنیتی باشد
			\item مجوز اقامت بر مبنای اطلاعات نادرست
			صادر شده باشد
		\end{itemize}
		\textbf{بنابراین:} مسیر حقوقی «اثبات تخلف فردی»
		مؤثرتر و مشروع‌تر از «اخراج جمعی» است.
	\end{warningBox} \\
	\midrule
  \textbf{۶. افشای عمومی} &
\begin{itemize}[nosep]
	\item انتشار «گزارش ثروت سران»
	(\en{Wealth of Rulers Report})
	به‌صورت سالانه
	\item مستندسازی تناقض بین شعارهای ساده‌زیستی
	و سبک زندگی مجلل فرزندان و بستگان
	\item اطلاع‌رسانی به رسانه‌ها و سازمان‌های
	بین‌المللی ضد فساد (\en{Transparency International})
\end{itemize} \\
\bottomrule
\end{longtable}


% ============================================================
%  بخش هفتم: همکاری با ساختارهای بین‌المللی
% ============================================================
\section{حوزه‌ی هفتم: همکاری با ساختارهای حقوق بشری بین‌المللی}

\subsection{موضوع ۷.۱: همکاری با نماینده‌ی ویژه‌ی سازمان ملل}

\begin{longtable}{r p{11cm}}
\toprule
\textbf{عنصر} & \textbf{شرح} \\
\midrule
\endhead
\textbf{تعریف} &
ایجاد کانال ارتباطی رسمی و نظام‌مند با
گزارشگر ویژه‌ی سازمان ملل در امور حقوق بشر ایران
(\en{Special Rapporteur on the situation of
	human rights in Iran}) و سایر رویه‌های ویژه
(\en{Special Procedures}). \\
\midrule
\textbf{اَشکال همکاری} &
\begin{enumerate}[nosep, label=\textcolor{cAccent}{\arabic*.}]
	\item \textbf{ارائه‌ی اطلاعات و مستندات:}
	ارسال منظم پرونده‌ها، گزارش‌ها و شهادت‌نامه‌ها
	به دفتر گزارشگر ویژه
	\item \textbf{ارائه‌ی شکایت فردی (\en{Individual Complaint}):}
	کمک به قربانیان برای ثبت شکایت
	در نظام‌های معاهده‌ای سازمان ملل
	\item \textbf{مشارکت در بررسی دوره‌ای جهانی
		(\en{Universal Periodic Review — UPR}):}
	تهیه و ارائه‌ی «گزارش سایه»
	(\en{Shadow Report}) برای جلسات \en{UPR} ایران
	\item \textbf{همکاری با هیئت حقیقت‌یاب:}
	در صورت تشکیل هیئت حقیقت‌یاب بین‌المللی
	(\en{Fact-Finding Mission})
	مشابه هیئت‌های میانمار و سوریه
	\item \textbf{ارائه‌ی بریفینگ} (\en{briefing})
	به شورای حقوق بشر و کمیته‌ی سوم مجمع عمومی
\end{enumerate} \\
\midrule
\textbf{گزارشگران ویژه‌ی مرتبط} &
\begin{itemize}[nosep]
	\item گزارشگر ویژه‌ی وضعیت حقوق بشر ایران
	(فعلی: جاوید رحمان / جانشین)
	\item گزارشگر ویژه‌ی شکنجه
	\item گزارشگر ویژه‌ی اعدام‌های فراقانونی
	\item گزارشگر ویژه‌ی آزادی بیان
	\item گزارشگر ویژه‌ی حقوق زنان
	\item گزارشگر ویژه‌ی مدافعان حقوق بشر
	\item گروه کاری بازداشت خودسرانه
	(\en{Working Group on Arbitrary Detention})
\end{itemize} \\
\bottomrule
\end{longtable}

\subsection{موضوع ۷.۲: همکاری با ساز‌و‌کار مستندسازی مستقل ایران (\en{IIMM})}

\begin{longtable}{r p{11cm}}
\toprule
\textbf{عنصر} & \textbf{شرح} \\
\midrule
\endhead
\textbf{تعریف} &
در نوامبر ۲۰۲۲، شورای حقوق بشر سازمان ملل
هیئت حقیقت‌یاب مستقل بین‌المللی
(\en{Independent International Fact-Finding Mission
	on Iran}) را تأسیس کرد.
در مارس ۲۰۲۳، مجمع عمومی خواستار تأسیس
سازوکار مستقل مستندسازی ایران شد
(الگوی \en{IIIM} سوریه). \\
\midrule
\textbf{اَشکال همکاری} &
\begin{enumerate}[nosep, label=\textcolor{cPrimary}{\arabic*.}]
	\item ارائه‌ی مستندات تأیید‌شده
	\item معرفی شهود و قربانیان (با رضایت آنها)
	\item همکاری فنی در مستندسازی دیجیتال
	\item مشارکت در حفظ زنجیره‌ی ادله
	(\en{chain of custody})
	\item تبادل اطلاعات با رعایت محرمانگی
\end{enumerate} \\
\midrule
\textbf{نمونه‌ی الهام‌بخش} &
\begin{itemize}[nosep]
	\item \textbf{IIIM سوریه (۲۰۱۶–اکنون):}
	سازوکار بین‌المللی بی‌طرف و مستقل
	برای جمع‌آوری ادله‌ی جرایم جنگی در سوریه.
	بیش از ۲ میلیون سند جمع‌آوری شده.
	\keyword{الگوی موفق} برای ایران.
	\item \textbf{UNITAD عراق (۲۰۱۷–اکنون):}
	تیم تحقیقاتی سازمان ملل برای
	مستندسازی جرایم داعش
\end{itemize} \\
\bottomrule
\end{longtable}

\subsection{موضوع ۷.۳: مشارکت در تدوین قطعنامه‌های بین‌المللی}

\begin{longtable}{r p{11cm}}
\toprule
\textbf{عنصر} & \textbf{شرح} \\
\midrule
\endhead
\textbf{تعریف} &
مشارکت فعال در فرآیند تدوین و اصلاح
قطعنامه‌های مربوط به ایران در:
\begin{itemize}[nosep]
	\item مجمع عمومی (قطعنامه‌ی سالانه‌ی حقوق بشر ایران)
	\item شورای حقوق بشر
	\item شورای امنیت (در صورت امکان)
	\item نهادهای معاهده‌ای (\en{Treaty Bodies})
\end{itemize} \\
\midrule
\textbf{فعالیت‌های مشخص} &
\begin{enumerate}[nosep, label=\textcolor{cAccent}{\arabic*.}]
	\item تهیه‌ی پیش‌نویس‌های پیشنهادی
	(\en{draft language}) برای درج در قطعنامه‌ها
	\item لابی‌گری در هیئت‌های نمایندگی کشورهای حامی
	(کانادا، سوئد، آلمان، فرانسه و …)
	\item تهیه‌ی «یادداشت مفهومی» (\en{concept note})
	برای رویدادهای جانبی (\en{side events})
	در حاشیه‌ی اجلاس‌ها
	\item ارائه‌ی شهادت شفاهی
	(\en{oral testimony}) در جلسات رسمی
	\item پیگیری اجرای مفاد قطعنامه‌های مصوب
	و گزارش‌دهی درباره‌ی عدم اجرا
\end{enumerate} \\
\midrule
\textbf{نمونه‌ی تاریخی} &
\begin{itemize}[nosep]
	\item \textbf{قطعنامه‌های سالانه‌ی مجمع عمومی
		درباره‌ی حقوق بشر ایران} (از ۱۹۸۵ تاکنون)
	\item \textbf{قطعنامه‌ی ۱۶/۱ شورای حقوق بشر (۲۰۱۱):}
	تعیین گزارشگر ویژه برای سوریه
	\item \textbf{قطعنامه‌ی ۳۴/۱ (میانمار ۲۰۱۷):}
	تأسیس هیئت حقیقت‌یاب
	\item نقش سازمان‌های غیردولتی
	(\en{Human Rights Watch}، \en{Amnesty International})
	در شکل‌دهی به زبان قطعنامه‌ها —
	الگویی برای دادگاه در سایه
\end{itemize} \\
\bottomrule
\end{longtable}

% ============================================================
%  بخش هشتم: افشا و مستندسازی جنایات
% ============================================================
\section{حوزه‌ی هشتم: افشا، مستندسازی و بایگانی جنایات}

\begin{infoBox}[فلسفه‌ی این حوزه]
	حتی اگر هیچ‌یک از اقدامات حقوقی بین‌المللی
	در کوتاه‌مدت به نتیجه نرسد،
	\keyword{مستندسازی دقیق و حرفه‌ای جنایات}
	ارزش ذاتی و بلندمدت دارد.
	سند تاریخی قابل انکار نیست.
	همان‌طور که آرشیو \en{Stasi}،
	اسناد هولوکاست و مستندات رواندا
	دهه‌ها بعد مبنای عدالت شدند.
\end{infoBox}

% FIXED: تضمین اینکه عنوان بدون حداقل ۶ خط فضا، به صفحه‌ی بعد برود
\needspace{6\baselineskip}%
\subsection{موضوع ۸.۱: مستندسازی نظام‌مند نقض حقوق بشر}
\nopagebreak[4]%            % FIXED: ممنوعیت شکست بین عنوان و جدول
\vspace{-1ex}%              % FIXED: کاهش فاصله‌ی عمودی اضافی
\noindent%                  % FIXED: حذف تورفتگی پاراگراف

{% FIXED: گروه‌بندی برای محدود کردن دامنه‌ی head/tail
	\tablefirsthead{%
		\toprule
		\textbf{عنصر} & \textbf{شرح} \\
		\midrule
	}
	
	\tablehead{%
	    \toprule
		\multicolumn{2}{r}{\small\itshape ادامه از صفحه‌ی قبل} \\
		\midrule
		\textbf{عنصر} & \textbf{شرح} \\
		\midrule
	}
	
	\tabletail{%
		\midrule
		\multicolumn{2}{l}{\small\itshape ادامه در صفحه‌ی بعد} \\
		\bottomrule
	}
	
	\tablelasttail{%
		\bottomrule
	}
	
	% FIXED: کمک به supertabular برای محاسبه‌ی دقیق‌تر عرض
	%\twocolumn[\setlength{\STableWidth}{\textwidth}]%  ← فقط اگر twocolumn دارید
	% اگر تک‌ستونی است، خط بالا را حذف کنید
	
	\begin{supertabular}{r p{11cm}}
		
		% ─────────────────────────────────────
		%  بخش ۱: تعریف
		% ─────────────────────────────────────
		\textbf{تعریف} &
		ایجاد یک \keyword{بانک اطلاعاتی جامع}%
		(\en{Comprehensive Violations Database})
		از کلیه‌ی موارد نقض حقوق بشر توسط جمهوری اسلامی
		از ۱۳۵۷ تا زمان حال. \\
		\midrule
		
		% ─────────────────────────────────────
		%  بخش ۲: دسته‌بندی جرایم
		%  FIXED: هر جرم در سطر مجزا → شکست صفحه ممکن
		% ─────────────────────────────────────
		\multicolumn{2}{c}{%
			\cellcolor{BgLightBlue}%
			\textbf{طبقه‌بندی ده‌گانه‌ی جرایم}%
		} \\
		\midrule
		
		% ── ۱. اعدام‌ها ──
		\textbf{دسته‌بندی جرایم} &
		\textcolor{cAccent}{\textbf{۱.}} \textbf{اعدام‌ها:}
		\begin{itemize}[nosep, topsep=2pt]
			\item اعدام‌های سیاسی (۱۳۶۰–۱۳۶۷ و پس از آن)
			\item قتل‌عام ۶۷ (کشتار زندانیان سیاسی)
			\item اعدام‌های علنی
			\item اعدام نوجوانان (نقض \en{CRC})
		\end{itemize} \\
		
		% ── ۲. شکنجه و بدرفتاری ──
		&
		\textcolor{cAccent}{\textbf{۲.}} \textbf{شکنجه و بدرفتاری:}
		\begin{itemize}[nosep, topsep=2pt]
			\item شکنجه‌ی جسمی (شلاق، آویزان‌کردن، ضرب و شتم)
			\item شکنجه‌ی روانی (سلول انفرادی طولانی،
			تهدید خانواده، اعدام‌های نمایشی)
			\item شکنجه‌ی جنسی و تجاوز در بازداشتگاه‌ها
		\end{itemize} \\
		
		% ── ۳. بازداشت خودسرانه ──
		&
		\textcolor{cAccent}{\textbf{۳.}} \textbf{بازداشت خودسرانه:}
		\begin{itemize}[nosep, topsep=2pt]
			\item بازداشت بدون حکم قضایی
			\item بازداشت‌های طولانی‌مدت بدون محاکمه
			\item بازداشت گروگان‌گیرانه (\en{hostage diplomacy})
		\end{itemize} \\
		
		% ── ۴. ناپدید‌سازی اجباری ──
		&
		\textcolor{cAccent}{\textbf{۴.}} \textbf{ناپدید‌سازی اجباری:}% FIXED: نیم‌فاصله
		\begin{itemize}[nosep, topsep=2pt]
			\item ناپدید‌شدگان دهه‌ی ۶۰% FIXED: نیم‌فاصله
			\item قتل‌های زنجیره‌ای (۱۳۷۷)
			\item گورهای جمعی (خاوران و \ldots)% FIXED: \ldots
		\end{itemize} \\
		
		% ── ۵. سرکوب آزادی بیان و رسانه ──
		&
		\textcolor{cAccent}{\textbf{۵.}} \textbf{سرکوب آزادی بیان و رسانه:}
		\begin{itemize}[nosep, topsep=2pt]
			\item زندانی‌کردن روزنامه‌نگاران
			\item توقیف نشریات
			\item سانسور اینترنت و فیلترینگ
			\item قطع اینترنت در اعتراضات
		\end{itemize} \\
		
		% ── ۶. سرکوب اقلیت‌های قومی و مذهبی ──
		&
		\textcolor{cAccent}{\textbf{۶.}} \textbf{سرکوب اقلیت‌های قومی و مذهبی:}
		\begin{itemize}[nosep, topsep=2pt]
			\item بهائیان (محرومیت سیستماتیک)
			\item سنّی‌ها (محدودیت عبادت در تهران)% FIXED: تشدید
			\item مسیحیان و یهودیان تازه‌مسلمان
			\item کردها، بلوچ‌ها، عرب‌ها، ترکمن‌ها
			\item دراویش گنابادی
		\end{itemize} \\
		
		% ── ۷. نقض حقوق زنان ──
		&
		\textcolor{cAccent}{\textbf{۷.}} \textbf{نقض حقوق زنان:}
		\begin{itemize}[nosep, topsep=2pt]
			\item حجاب اجباری و مجازات بدنی
			\item نابرابری حقوقی (ارث، شهادت، طلاق)
			\item گشت ارشاد و خشونت سازمان‌یافته
			\item مورد مهسا (ژینا) امینی
		\end{itemize} \\
		
		% ── ۸. نقض حقوق کارگران ──
		&
		\textcolor{cAccent}{\textbf{۸.}} \textbf{نقض حقوق کارگران:}
		\begin{itemize}[nosep, topsep=2pt]
			\item ممنوعیت سندیکاهای مستقل
			\item بازداشت فعالان کارگری
			\item عدم پرداخت دستمزد و بی‌تأمینی
		\end{itemize} \\
		
		% ── ۹. سرکوب اعتراضات مردمی ──
		&
		\textcolor{cAccent}{\textbf{۹.}} \textbf{سرکوب اعتراضات مردمی:}
		\begin{itemize}[nosep, topsep=2pt]
			\item ۱۸ تیر ۱۳۷۸ (کوی دانشگاه)
			\item جنبش سبز ۱۳۸۸
			\item آبان ۱۳۹۸ (کشتار بیش از ۱۵۰۰ نفر)% FIXED: حذف +
			\item جنبش ژن، ژیان، ئازادی ۱۴۰۱% FIXED: ویرگول فارسی
		\end{itemize} \\
		
		% ── ۱۰. تروریسم دولتی و ترور فراسرزمینی ──
		&
		\textcolor{cAccent}{\textbf{۱۰.}} \textbf{تروریسم دولتی و ترور فراسرزمینی:}
		\begin{itemize}[nosep, topsep=2pt]
			\item ترور مخالفان در خارج (قاسملو، شرفکندی، فرخ‌زاد و \ldots)% FIXED: \ldots
			\item بمب‌گذاری آمیا (آرژانتین)
			\item عملیات تروریستی در اروپا
			(رستوران میکونوس، ویلپنت و \ldots)% FIXED: \ldots
			\item گروگان‌گیری سفارت آمریکا (۱۳۵۸)
		\end{itemize} \\
		\midrule
		
		\\
		% ─────────────────────────────────────
		%  بخش ۳: استاندارد مستندسازی
		% ─────────────────────────────────────
		\textbf{استاندارد مستندسازی} &
		بر اساس \en{Berkeley Protocol on Digital
			Open Source Investigations} (۲۰۲۰)
		و \en{ICC Evidence Collection Guidelines}:
		\begin{enumerate}[nosep, topsep=2pt, label=\textcolor{cPrimary}{\roman*.}]
			\item هر واقعه: چه کسی، چه چیزی، کجا، کی، چگونه
			\item شناسایی قربانی (با رضایت یا ناشناس)
			\item شناسایی عاملان و زنجیره‌ی فرماندهی
			\item ادله‌ی فیزیکی و دیجیتال
			\item شهادت شاهدان و بازماندگان
			\item تأیید متقاطع (\en{cross-verification})
			\item حفظ زنجیره‌ی ادله (\en{chain of custody})
			\item فراداده (\en{metadata}) و هش‌گذاری دیجیتال% FIXED: فرادیتا → فراداده
		\end{enumerate} \\
		
	\end{supertabular}
}% پایان گروه

\noindent%                  % FIXED: حذف تورفتگی پاراگراف
\subsection{موضوع ۸.۲: شناسایی مسببین، عاملین و مباشران}

\begin{longtable}{r p{11cm}}
\toprule
\textbf{عنصر} & \textbf{شرح} \\
\midrule
\endhead
\textbf{تعریف} &
ایجاد \keyword{بانک اطلاعاتی مسئولان نقض حقوق بشر}
با سطح‌بندی مسئولیت بر اساس حقوق بین‌الملل کیفری. \\
\midrule
\textbf{سطوح مسئولیت} &
\begin{principleBox}[هرم مسئولیت — بر اساس اساسنامه‌ی رم]
	\begin{enumerate}[nosep, label=\textcolor{cAccent}{\arabic*.}]
		\item \textbf{مسئولیت فرماندهی (\en{Command Responsibility}):}
		مقامات ارشدی که دستور سرکوب را صادر کرده‌اند
		یا از آن آگاه بوده و جلوگیری نکرده‌اند
		(ماده‌ی ۲۸ اساسنامه‌ی رم)
		
		\item \textbf{مباشرت مستقیم (\en{Direct Perpetration}):}
		افرادی که شخصاً اقدام مجرمانه انجام داده‌اند
		(بازجو، تیرانداز، شکنجه‌گر)
		
		\item \textbf{مشارکت (\en{Co-perpetration}):}
		مشارکت مشترک در طرح مجرمانه
		(ماده‌ی ۲۵(۳)(الف))
		
		\item \textbf{معاونت و مساعدت (\en{Aiding and Abetting}):}
		کمک مادی یا معنوی به ارتکاب جرم
		(ماده‌ی ۲۵(۳)(ج))
		
		\item \textbf{تحریک (\en{Inducement / Solicitation}):}
		ترغیب و تشویق به ارتکاب جرم
		(مانند سخنرانی‌های تحریک‌آمیز مقامات)
		
		\item \textbf{سهم در اقدام گروهی
			(\en{Contribution to Group Action}):}
		مشارکت در فعالیت گروهی مجرمانه
		با آگاهی از نیت گروه
		(ماده‌ی ۲۵(۳)(د))
	\end{enumerate}
\end{principleBox} \\
\midrule
\textbf{ساختار پرونده‌ی فردی} &
برای هر فرد شناسایی‌شده:
\begin{enumerate}[nosep, label=\textcolor{cPrimary}{\roman*.}]
	\item مشخصات هویتی کامل (نام، سمت، عکس)
	\item سوابق خدمتی و جایگاه در ساختار قدرت
	\item \keyword{نمودار ارتباطی}
	(\en{link chart / network diagram})
	— نشان‌دهنده‌ی زنجیره‌ی فرماندهی
	\item فهرست جرایم منتسب با ادله
	\item سطح مسئولیت (از ۶ سطح فوق)
	\item وضعیت فعلی (محل اقامت، وضعیت تحریمی)
	\item ارزیابی امکان تعقیب قضایی
	(در کدام حوزه‌ی صلاحیتی)
\end{enumerate} \\
\midrule
\textbf{نمونه‌ی تاریخی} &
\begin{itemize}[nosep]
	\item \textbf{پروژه‌ی سزار (سوریه):}
	یک عکاس نظامی سوری (\en{Caesar}) ۵۵,۰۰۰
	عکس از اجساد شکنجه‌شدگان را قاچاق کرد —
	مبنای «قانون سزار» آمریکا و تعقیب
	مقامات سوری در اروپا
	\item \textbf{پرونده‌ی انور رسلان (آلمان ۲۰۲۲):}
	محاکمه‌ی رئیس بازجویی سوریه در کوبلنتس
	بر اساس صلاحیت جهانی —
	اولین محکومیت جرایم علیه بشریت
	مرتبط با سوریه
	\item \textbf{پرونده‌ی حمید نوری (سوئد ۲۰۲۲):}
	محاکمه‌ی دست‌یار دادیار اوین
	به اتهام مشارکت در کشتار ۶۷ —
	\keyword{مستقیماً مرتبط با ایران}
	\item \textbf{گاچاچا (رواندا):}
	دادگاه‌های محلی/سنتی رواندا که
	بیش از ۱.۲ میلیون پرونده‌ی
	مربوط به نسل‌کشی را رسیدگی کردند
\end{itemize} \\
\bottomrule
\end{longtable}

\subsection{موضوع ۸.۳: حفاظت از شهود و قربانیان}

\begin{longtable}{r p{11cm}}
\toprule
\textbf{عنصر} & \textbf{شرح} \\
\midrule
\endhead
\textbf{تعریف} &
ایجاد برنامه‌ی حمایت از شهود و قربانیانی که
اطلاعات یا شهادت ارائه می‌دهند. \\
\midrule
\textbf{سازوکارها} &
\begin{enumerate}[nosep, label=\textcolor{cAccent}{\arabic*.}]
	\item \textbf{ناشناس‌سازی هویت:}
	استفاده از نام مستعار و کدگذاری در تمام اسناد
	\item \textbf{شهادت از راه دور:}
	استفاده از ویدئوکنفرانس با تغییر صدا و تصویر
	\item \textbf{ارزیابی خطر فردی:}
	برای هر شاهد، تحلیل خطر
	(\en{risk assessment}) انجام شود
	\item \textbf{مشاوره‌ی حقوقی:}
	ارائه‌ی راهنمایی درباره‌ی حقوق و مسئولیت‌ها
	\item \textbf{حمایت روان‌شناختی:}
	ارجاع به متخصصان ترومای اختصاصی
	\item \textbf{حمایت مادی اضطراری:}
	در صورت تهدید جانی، کمک به جابه‌جایی
\end{enumerate} \\
\midrule
\textbf{نمونه} &
\begin{itemize}[nosep]
	\item برنامه‌ی حمایت شهود \en{ICC}
	(\en{Victims and Witnesses Unit})
	\item برنامه‌ی \en{ICTY} برای شهادت محرمانه
	\item تجربه‌ی کمیسیون حقیقت و آشتی آفریقای جنوبی
	(\en{TRC}) در حمایت از شهود
\end{itemize} \\
\bottomrule
\end{longtable}


% ============================================================
%  بخش نهم: صلاحیت جهانی و تعقیب در دادگاه‌های ملی
% ============================================================
\section{حوزه‌ی نهم: صلاحیت جهانی و تعقیب در دادگاه‌های ملی}

\subsection{موضوع ۹.۱: شناسایی فرصت‌های صلاحیت جهانی}

\begin{longtable}{r p{11cm}}
\toprule
\textbf{عنصر} & \textbf{شرح} \\
\midrule
\endhead
\textbf{تعریف} &
صلاحیت جهانی (\en{Universal Jurisdiction})
به معنای توانایی دادگاه‌های ملی یک کشور
برای رسیدگی به جرایم بین‌المللی (نسل‌کشی،
جرایم علیه بشریت، جرایم جنگی، شکنجه)
بدون توجه به محل وقوع جرم یا تابعیت
مرتکب و قربانی. \\
\midrule
\textbf{کشورهای دارای قوانین مناسب} &
\begin{principleBox}[کشورهای با بهترین چارچوب صلاحیت جهانی]
	\begin{enumerate}[nosep, label=\textcolor{cAccent}{\arabic*.}]
		\item \textbf{آلمان:}
		قانون جرایم علیه حقوق بین‌الملل
		(\en{Völkerstrafgesetzbuch — VStGB}) ۲۰۰۲
		— صلاحیت جهانی محض (بدون نیاز به ارتباط)
		— واحد ویژه‌ی \en{ZBKV} در پلیس فدرال
		— \keyword{موفق‌ترین الگو} (پرونده‌های سوریه و عراق)
		
		\item \textbf{سوئد:}
		پرونده‌ی حمید نوری — نقطه‌ی عطف برای ایران
		— واحد جرایم بین‌المللی در پلیس سوئد
		
		\item \textbf{فرانسه:}
		واحد تحقیق جرایم علیه بشریت
		(\en{Pôle crimes contre l'humanité})
		— پرونده‌های سوریه و روآندا
		
		\item \textbf{هلند:}
		واحد جرایم جنگی (\en{Team Internationale Misdrijven})
		
		\item \textbf{بلژیک:}
		پیشتاز صلاحیت جهانی (با محدودیت‌های ۲۰۰۳)
		
		\item \textbf{آرژانتین:}
		پرونده‌ی شکایت علیه فرانکو (اسپانیا)
		— پرونده‌ی باز علیه ایران (آمیا)
		
		\item \textbf{کانادا:}
		قانون جرایم علیه بشریت و جرایم جنگی (۲۰۰۰)
		— برنامه‌ی «بدون پناهگاه امن»
		(\en{No Safe Haven})
		
		\item \textbf{بریتانیا:}
		امکان تعقیب محدود (نیاز به رضایت دادستان کل)
	\end{enumerate}
\end{principleBox} \\
\midrule
\textbf{فعالیت دادگاه در سایه} &
\begin{enumerate}[nosep, label=\textcolor{cPrimary}{\roman*.}]
	\item شناسایی مظنونان مقیم هر کشور
	\item تهیه‌ی «پرونده‌ی آماده‌ی تعقیب»
	(\en{prosecution-ready dossier}) مطابق
	با قوانین آن کشور
	\item ارتباط با وکلای متخصص در هر حوزه‌ی صلاحیتی
	\item حمایت از شاکیان خصوصی در طرح شکایت
	\item پیگیری و گزارش‌دهی از روند پرونده‌ها
\end{enumerate} \\
\bottomrule
\end{longtable}

\subsection{موضوع ۹.۲: شکایت به دیوان بین‌المللی دادگستری (\en{ICJ})}

\begin{longtable}{r p{11cm}}
\toprule
\textbf{عنصر} & \textbf{شرح} \\
\midrule
\endhead
\textbf{تعریف} &
بررسی امکان طرح شکایت دولتی علیه ایران
در \en{ICJ} بر اساس معاهدات دوجانبه یا چندجانبه
که ایران عضو آنهاست. \\
\midrule
\textbf{مبانی صلاحیتی ممکن} &
\begin{enumerate}[nosep, label=\textcolor{cAccent}{\arabic*.}]
	\item \textbf{کنوانسیون منع نسل‌کشی (۱۹۴۸):}
	ایران عضو است — ماده‌ی ۹ حق ارجاع به \en{ICJ}
	(الگوی گامبیا علیه میانمار ۲۰۱۹)
	\item \textbf{کنوانسیون مبارزه با شکنجه (\en{CAT}):}
	ایران امضا کرده اما تصویب نکرده —
	محدودیت صلاحیتی
	\item \textbf{کنوانسیون رفع تبعیض نژادی (\en{CERD}):}
	ایران عضو است — سابقه‌ی ارجاع
	(اوکراین علیه روسیه ۲۰۲۲)
	\item \textbf{معاهده‌ی مودت ایران و آمریکا (۱۹۵۵):}
	مبنای پرونده‌های قبلی ایران در \en{ICJ}
\end{enumerate} \\
\midrule
\textbf{محدودیت} &
\begin{warningBox}[محدودیت اساسی]
	دادگاه در سایه \keyword{حق طرح دعوا} در \en{ICJ} ندارد
	(فقط دولت‌ها دارند). اما می‌تواند:
	\begin{itemize}[nosep]
		\item مستندات و تحلیل حقوقی تهیه کند
		\item دولت‌های ثالث را \keyword{ترغیب و حمایت} کند
		\item به‌عنوان \en{amicus curiae} (دوست دادگاه)
		یادداشت ارائه دهد
	\end{itemize}
	\textbf{الگوی الهام‌بخش:} گامبیا در شکایت علیه میانمار
	با حمایت سازمان همکاری اسلامی عمل کرد.
	آیا کشوری حاضر است با حمایت دادگاه در سایه
	علیه ایران در \en{ICJ} اقامه‌ی دعوا کند؟
\end{warningBox} \\
\bottomrule
\end{longtable}


% ============================================================
%  بخش دهم: موضوعات ویژه و نوآورانه
% ============================================================
\section{حوزه‌ی دهم: موضوعات ویژه و نوآورانه}

\subsection{موضوع ۱۰.۱: اکوساید و جرایم زیست‌محیطی}

\begin{longtable}{r p{11cm}}
\toprule
\textbf{عنصر} & \textbf{شرح} \\
\midrule
\endhead
\textbf{تعریف} &
بررسی و مستندسازی \keyword{تخریب عمدی
	یا ناشی از بی‌توجهی محیط زیست}
به‌عنوان جرم بین‌المللی نوظهور (\en{ecocide}).
نمونه‌ها:
\begin{itemize}[nosep]
	\item خشک‌شدن دریاچه‌ی ارومیه
	\item بحران آب خوزستان و اصفهان
	\item فرونشست زمین در دشت‌ها
	\item آلودگی هوای شدید کلان‌شهرها
	\item انتقال آب بین‌حوضه‌ای بدون ارزیابی محیطی
	\item تخریب جنگل‌ها و تالاب‌ها
\end{itemize} \\
\midrule
\textbf{مبنای حقوقی} &
\begin{itemize}[nosep]
	\item تلاش بین‌المللی برای افزودن «اکوساید»
	به‌عنوان جرم پنجم اساسنامه‌ی رم
	(پیشنهاد ۲۰۲۱ کارگروه \en{Stop Ecocide})
	\item ماده‌ی ۸(۲)(ب)(iv) اساسنامه‌ی رم:
	تخریب گسترده‌ی محیط زیست در مخاصمات
	\item رأی مشورتی \en{ICJ} درباره‌ی تسلیحات هسته‌ای (۱۹۹۶):
	تعهد عدم آسیب به محیط زیست
	\item اصل ۵۰ قانون اساسی جمهوری اسلامی
	(که خود حکومت آن را نقض می‌کند)
\end{itemize} \\
\midrule
\textbf{نمونه‌ی تاریخی} &
\begin{itemize}[nosep]
	\item \textbf{عامل نارنجی (ویتنام):}
	تخریب جنگل‌ها توسط آمریکا —
	بحث‌های حقوقی پس از جنگ
	\item \textbf{آتش‌زدن چاه‌های نفت کویت (۱۹۹۱):}
	محکومیت عراق و پرداخت خسارت
	\item \textbf{جنگل‌زدایی آمازون:}
	تلاش‌ها برای محاکمه‌ی مقامات برزیلی
\end{itemize} \\
\bottomrule
\end{longtable}

\subsection{موضوع ۱۰.۲: جرایم سایبری دولتی و قطع اینترنت}

\begin{longtable}{r p{11cm}}
\toprule
\textbf{عنصر} & \textbf{شرح} \\
\midrule
\endhead
\textbf{تعریف} &
بررسی و مستندسازی:
\begin{itemize}[nosep]
	\item قطع عمدی اینترنت در اعتراضات
	(\en{internet shutdowns})
	— به‌عنوان ابزار سرکوب
	\item حملات سایبری به زیرساخت‌های مدنی
	\item جاسوسی سایبری از شهروندان و اپوزیسیون
	\item نظارت دیجیتال فراگیر (\en{mass surveillance})
	\item استفاده از بدافزارهای دولتی
	(\en{state-sponsored malware})
\end{itemize} \\
\midrule
\textbf{مبنای حقوقی} &
\begin{itemize}[nosep]
	\item قطعنامه‌ی شورای حقوق بشر \en{A/HRC/RES/32/13} (۲۰۱۶):
	دسترسی به اینترنت حق بشری است
	\item گزارش ویژه‌ی دبیرکل سازمان ملل
	درباره‌ی قطع اینترنت (۲۰۲۲)
	\item راهنمای تالین (\en{Tallinn Manual}) درباره‌ی
	حقوق بین‌الملل قابل‌اعمال بر عملیات سایبری
\end{itemize} \\
\midrule
\textbf{مورد ایران} &
\begin{itemize}[nosep]
	\item \textbf{آبان ۱۳۹۸:} قطع کامل اینترنت به مدت ۷ روز
	— همزمان با کشتار ۱۵۰۰+ نفر
	— عمدی و برنامه‌ریزی‌شده
	\item \textbf{مهسا ۱۴۰۱:} قطع مکرر و محدودسازی شدید
	\item \textbf{شبکه‌ی ملی اطلاعات (\en{NIN}):}
	طرح جمهوری اسلامی برای جایگزینی
	اینترنت جهانی با شبکه‌ی داخلی قابل‌کنترل
\end{itemize} \\
\bottomrule
\end{longtable}

\subsection{موضوع ۱۰.۳: فساد سیستماتیک و غارت ثروت ملی}

\begin{longtable}{r p{11cm}}
\toprule
\textbf{عنصر} & \textbf{شرح} \\
\midrule
\endhead
\textbf{تعریف} &
مستندسازی فساد ساختاری نظام به‌عنوان
\keyword{جرم اقتصادی علیه ملت ایران}:
\begin{itemize}[nosep]
	\item بنیادهای اقتصادی نظامی (ستاد اجرایی فرمان امام،
	بنیاد مستضعفان، آستان قدس و …)
	\item اقتصاد سپاه و شرکت‌های تابعه
	(خاتم‌الانبیا و …)
	\item اختلاس‌های بزرگ (بابک زنجانی، پرونده‌ی بانکی و …)
	\item قاچاق سازمان‌یافته‌ی دولتی
	\item پول‌شویی بین‌المللی
\end{itemize} \\
\midrule
\textbf{مبنای حقوقی} &
\begin{itemize}[nosep]
	\item کنوانسیون سازمان ملل علیه فساد
	(\en{UNCAC}) — ایران عضو
	\item اصول بازگرداندن دارایی‌ها
	(\en{asset recovery}) در \en{UNCAC}
	\item اصول اقتصادی عدالت انتقالی
\end{itemize} \\
\midrule
\textbf{خروجی} &
\begin{itemize}[nosep]
	\item «گزارش فساد سالانه» (\en{Annual Corruption Report})
	\item نقشه‌ی ارتباطات مالی سران و نهادها
	\item تخمین حجم ثروت غارت‌شده
	\item پیشنهاد سازوکارهای بازگرداندن دارایی‌ها
	پس از گذار
\end{itemize} \\
\bottomrule
\end{longtable}

\subsection{موضوع ۱۰.۴: نقض حقوق اقلیت‌های جنسی و جنسیتی (\en{LGBTQ+})}

\begin{longtable}{r p{11cm}}
\toprule
\textbf{عنصر} & \textbf{شرح} \\
\midrule
\endhead
\textbf{تعریف} &
مستندسازی سرکوب سیستماتیک اقلیت‌های جنسی
و جنسیتی:
\begin{itemize}[nosep]
	\item مجازات اعدام برای هم‌جنس‌گرایی در قانون مجازات
	\item اعدام‌های ثبت‌شده و ثبت‌نشده
	\item فشار برای تغییر جنسیت اجباری
	\item آزار و شکنجه‌ی سیستماتیک
	\item محرومیت اجتماعی و اقتصادی
\end{itemize} \\
\midrule
\textbf{مبنای حقوقی} &
\begin{itemize}[nosep]
	\item اصول یوگیاکارتا (\en{Yogyakarta Principles}) ۲۰۰۶
	\item قطعنامه‌ی ۱۷/۱۹ شورای حقوق بشر (۲۰۱۱) —
	اولین قطعنامه‌ی سازمان ملل درباره‌ی حقوق \en{LGBTQ+}
	\item تفسیر ماده‌ی ۷ اساسنامه‌ی رم:
	«تعقیب و آزار» بر مبنای جنسیت
\end{itemize} \\
\bottomrule
\end{longtable}

\subsection{موضوع ۱۰.۵: گروگان‌گیری دولتی (\en{State Hostage-Taking})}

\begin{longtable}{r p{11cm}}
\toprule
\textbf{عنصر} & \textbf{شرح} \\
\midrule
\endhead
\textbf{تعریف} &
مستندسازی الگوی بازداشت شهروندان دوتابعیتی
و خارجی به‌عنوان ابزار سیاسی و اهرم فشار:
\begin{itemize}[nosep]
	\item نازنین زاغری (بریتانیا–ایران)
	\item سیامک نمازی (آمریکا–ایران)
	\item احمدرضا جلالی (سوئد–ایران)
	\item و ده‌ها مورد دیگر
\end{itemize} \\
\midrule
\textbf{مبنای حقوقی} &
\begin{itemize}[nosep]
	\item کنوانسیون بین‌المللی علیه گروگان‌گیری (۱۹۷۹)
	\item قطعنامه‌ی \en{A/RES/77/243} مجمع عمومی (۲۰۲۲):
	«مقابله با بازداشت خودسرانه‌ی
	دولتی در روابط بین دولت‌ها»
	\item ابتکار کانادا: «اعلامیه‌ی علیه
	بازداشت خودسرانه‌ی دولتی»
	(\en{Declaration Against Arbitrary
		Detention in State-to-State Relations})
	— ۲۰۲۱، با ۷۰+ کشور امضاکننده
\end{itemize} \\
\midrule
\textbf{خروجی} &
\begin{itemize}[nosep]
	\item بانک اطلاعاتی گروگان‌ها و بازداشت‌شدگان
	\item تحلیل الگوی گروگان‌گیری و اهرم‌سازی
	\item پیشنهاد سازوکار چندجانبه‌ی پاسخ‌دهی
	\item حمایت حقوقی از خانواده‌ها
\end{itemize} \\
\bottomrule
\end{longtable}

\subsection{موضوع ۱۰.۶: نقض حقوق آموزشی و آکادمیک}

\begin{longtable}{r p{11cm}}
\toprule
\textbf{عنصر} & \textbf{شرح} \\
\midrule
\endhead
\textbf{تعریف} &
\begin{itemize}[nosep]
	\item محرومیت بهائیان از تحصیلات دانشگاهی
	(\en{ستاره‌دار کردن})
	\item اخراج اساتید و دانشجویان به دلایل سیاسی
	\item «انقلاب فرهنگی» ۱۳۵۹ و تعطیلی دانشگاه‌ها
	\item سرکوب تشکل‌های دانشجویی
	\item کنترل ایدئولوژیک محتوای آموزشی
	\item ممانعت از تحصیل دختران در برخی رشته‌ها
\end{itemize} \\
\midrule
\textbf{مبنای حقوقی} &
\begin{itemize}[nosep]
	\item ماده‌ی ۲۶ اعلامیه‌ی جهانی حقوق بشر (حق آموزش)
	\item ماده‌ی ۱۳ میثاق حقوق اقتصادی، اجتماعی و فرهنگی
	\item توصیه‌نامه‌ی یونسکو درباره‌ی آزادی آکادمیک
\end{itemize} \\
\bottomrule
\end{longtable}


% ============================================================
%  بخش یازدهم: درس‌هایی از تجربه‌های تاریخی
% ============================================================
\section{حوزه‌ی یازدهم: جدول تطبیقی — درس‌هایی از تجربه‌های تاریخی}

\begin{infoBox}[هدف این بخش]
استخراج \keyword{موضوعات الهام‌بخش}
از تجربه‌های عدالت انتقالی و مقابله با فجایع
در دیگر کشورها، برای غنی‌سازی فهرست موضوعات
قابل‌بررسی دادگاه در سایه.
\end{infoBox}

% ══════════════════════════════════════════════════════════════
%  اصلاح اساسی: جدول landscape با متن فارسی
%  ─────────────────────────────────────────────
%  تمام متن‌های فارسی در سلول‌ها باید در \RL{} قرار گیرند
%  عرض ستون‌ها تنظیم شد برای جلوگیری از سرریز
% ══════════════════════════════════════════════════════════════

\begin{landscape}
\begin{longtable}{r r p{3.8cm} p{3.8cm} p{4cm}}
	\toprule
	\textbf{\RL{ردیف}} &
	\textbf{\RL{کشور / بحران}} &
	\textbf{\RL{موضوعات کلیدی}} &
	\textbf{\RL{سازوکار عدالت}} &
	\textbf{\RL{درس برای ایران}} \\
	\midrule
	\endhead
	
	۱ &
	{\RL{رواندا (۱۹۹۴)}} &
	{\RL{نسل‌کشی ۸۰۰,۰۰۰ توتسی؛
			نقش رسانه‌ی نفرت‌پراکن؛
			تجاوز سیستماتیک به‌عنوان ابزار جنگ؛
			نقش همسایگان}} &
	{\RL{\en{ICTR} (دادگاه بین‌المللی)؛
			دادگاه‌های گاچاچا (محلی)؛
			کمیسیون ملی وحدت و آشتی}} &
	{\RL{محاکمه‌ی رسانه‌ی نفرت‌پراکن
			→ قابل‌تعمیم به صداوسیمای ایران؛
			مدل چندلایه‌ی عدالت؛
			شناسایی تجاوز جنسی به‌عنوان نسل‌کشی}} \\
	\midrule
	
	۲ &
	{\RL{یوگوسلاوی سابق (۱۹۹۱–۱۹۹۹)}} &
	{\RL{پاکسازی قومی در بوسنی؛
			نسل‌کشی سربرنیتسا؛
			محاصره‌ی سارایوو؛
			جرایم در کوزوو؛
			تخریب میراث فرهنگی}} &
	{\RL{\en{ICTY} (۱۹۹۳–۲۰۱۷)؛
			\en{IRMCT} (سازوکار باقیمانده)؛
			دادگاه‌های ملی بوسنی}} &
	{\RL{محاکمه‌ی رئیس‌جمهور (میلوشویچ)
			→ سابقه‌ی تعقیب بالاترین مقامات؛
			مسئولیت فرماندهی؛
			تخریب میراث فرهنگی به‌عنوان جرم جنگی}} \\
	\midrule
	
	۳ &
	{\RL{میانمار / روهینگیا (۲۰۱۷)}} &
	{\RL{پاکسازی قومی ۷۰۰,۰۰۰+ روهینگیا؛
			تجاوز سیستماتیک؛
			آتش‌زدن روستاها؛
			انکار شهروندی}} &
	{\RL{شکایت گامبیا در \en{ICJ}؛
			تحقیق \en{ICC}؛
			هیئت حقیقت‌یاب سازمان ملل؛
			\en{IIMM}}} &
	{\RL{الگوی «دولت ثالث شاکی»
			→ آیا کشوری حاضر است علیه ایران شکایت کند؟
			قرار موقت \en{ICJ}؛
			سازوکار مستندسازی بلندمدت}} \\
	\midrule
	
  ۴ &
{\RL{سودان / دارفور (۲۰۰۳–اکنون)}} &
{\RL{نسل‌کشی قبایل فور، مساعد، زاغاوه؛
		شبه‌نظامیان جنجاوید؛
		جابه‌جایی اجباری ۲.۵ میلیون نفر}} &
{\RL{ارجاع شورای امنیت به \en{ICC}؛
		حکم بازداشت البشیر (۲۰۰۹)؛
		عدم همکاری کشورها در دستگیری}} &
{\RL{محدودیت‌های اجرای حکم \en{ICC}
		بدون همکاری دولت‌ها؛
		اهمیت فشار مستمر حتی بدون دستگیری؛
		تحول پس از سقوط البشیر}} \\
\midrule

۵ &
{\RL{اویغورها / چین (۲۰۱۷–اکنون)}} &
{\RL{بازداشتگاه‌های اجباری ۱+ میلیون نفر؛
		نسل‌کشی فرهنگی؛
		عقیم‌سازی اجباری؛
		کار اجباری؛
		نظارت دیجیتال فراگیر}} &
{\RL{دادگاه اویغور لندن ۲۰۲۱
		— الگوی دادگاه مردمی؛
		تحریم‌های هدفمند (مگنیتسکی)؛
		قطعنامه‌های پارلمانی}} &
{\RL{نزدیک‌ترین الگو به دادگاه در سایه:
		ساختار غیردولتی با اثرگذاری بین‌المللی؛
		نسل‌کشی فرهنگی → قابل‌تعمیم به اقوام ایرانی؛
		نظارت دیجیتال → شبکه‌ی ملی اطلاعات}} \\
\midrule

۶ &
{\RL{سوریه (۲۰۱۱–اکنون)}} &
{\RL{کشتار غیرنظامیان (۵۰۰,۰۰۰+)؛
		استفاده از سلاح شیمیایی؛
		شکنجه‌ی سیستماتیک؛
		آوارگی ۱۲+ میلیون نفر}} &
{\RL{\en{IIIM} (سازوکار مستندسازی)؛
		محاکمات صلاحیت جهانی (آلمان، سوئد)؛
		پرونده‌ی سزار؛
		کمیسیون تحقیق سازمان ملل}} &
{\RL{\en{IIIM} → الگوی سازوکار مستندسازی ایران؛
		محاکمات آلمان → نقشه‌ی راه تعقیب مقامات ایرانی؛
		نقش ایران در سوریه → مسئولیت مشترک؛
		پرونده‌ی سزار → الهام‌بخش}} \\
\midrule

۷ &
{\RL{کامبوج / خمرهای سرخ (۱۹۷۵–۱۹۷۹)}} &
{\RL{نسل‌کشی ۱.۷ میلیون نفر؛
		قتل نخبگان و تحصیل‌کرده‌ها؛
		اردوگاه‌های کار اجباری؛
		شکنجه‌گاه S-21}} &
{\RL{\en{ECCC} (دادگاه فوق‌العاده‌ی کامبوج)
		— تأسیس ۲۰۰۳ (۲۵ سال بعد!)؛
		ترکیب قضات ملی و بین‌المللی}} &
{\RL{درس مهم: دیر رسیدن عدالت
		بهتر از نرسیدن آن است؛
		قتل نخبگان → مشابه سرکوب نخبگان ایرانی؛
		مدل هیبریدی → ممکن برای دوره‌ی گذار ایران}} \\
\midrule

۸ &
{\RL{آفریقای جنوبی (آپارتاید)}} &
{\RL{جداسازی نژادی نظام‌مند؛
		سرکوب خشونت‌بار اعتراضات (سووتو ۱۹۷۶)؛
		بازداشت و شکنجه‌ی فعالان؛
		ترورهای دولتی}} &
{\RL{کمیسیون حقیقت و آشتی (\en{TRC})؛
		عفو مشروط در قبال اعتراف کامل؛
		تحریم‌های بین‌المللی}} &
{\RL{مدل \en{TRC} → یکی از گزینه‌های
		عدالت انتقالی ایران؛
		تحریم‌های ورزشی و فرهنگی → تجربه‌ی قابل‌تکرار؛
		تحول مسالمت‌آمیز → الهام‌بخش}} \\
\midrule

۹ &
{\RL{آرژانتین (حکومت نظامی ۱۹۷۶–۱۹۸۳)}} &
{\RL{ناپدیدسازی ۳۰,۰۰۰ نفر؛
		پروازهای مرگ؛
		سرقت نوزادان؛
		شکنجه‌ی سیستماتیک در \en{ESMA}}} &
{\RL{کمیسیون ملی ناپدیدشدگان (\en{CONADEP})؛
		محاکمه‌ی خونتا (۱۹۸۵)؛
		ابطال قوانین عفو (۲۰۰۵)؛
		ادامه‌ی محاکمات تا ۲۰۲۰+}} &
{\RL{مهم‌ترین درس: حتی قوانین عفو قابل ابطال‌اند؛
		مادران میدان مه‌یو → نقش مادران خاوران؛
		سرقت نوزادان → آزمایش \en{DNA} دهه‌ها بعد؛
		محاکمه‌ی خونتا → الگوی محاکمه‌ی رهبران}} \\
\midrule

۱۰ &
{\RL{شیلی (پینوشه ۱۹۷۳–۱۹۹۰)}} &
{\RL{کودتا علیه آلنده؛
		ناپدیدسازی و شکنجه (۳,۰۰۰+ کشته)؛
		عملیات کندور؛
		استادیوم ملی به‌عنوان بازداشتگاه}} &
{\RL{کمیسیون رتیگ (\en{Rettig}) ۱۹۹۱؛
		کمیسیون والچ (\en{Valech}) ۲۰۰۳؛
		بازداشت پینوشه در لندن (۱۹۹۸)؛
		محاکمات داخلی پس از ۲۰۰۰}} &
{\RL{بازداشت پینوشه: حتی دیکتاتور بازنشسته مصون نیست؛
		صلاحیت جهانی اسپانیا → الگوی فعال‌سازی؛
		عملیات کندور → شبکه‌ی ترور فراسرزمینی ایران؛
		مصونیت رئیس دولت سابق → بحث حقوقی مهم}} \\
\midrule

۱۱ &
{\RL{سیرالئون (جنگ داخلی ۱۹۹۱–۲۰۰۲)}} &
{\RL{قطع دست و پای غیرنظامیان؛
		سربازگیری کودکان؛
		خشونت جنسی گسترده؛
		الماس‌های خونین}} &
{\RL{دادگاه ویژه‌ی سیرالئون (\en{SCSL})؛
		کمیسیون حقیقت و آشتی؛
		محاکمه‌ی چارلز تیلور}} &
{\RL{محاکمه‌ی رئیس‌جمهور کشور دیگر
		→ مسئولیت مقامات ایرانی در سوریه، عراق و یمن؟
		سربازگیری کودکان → بسیج نوجوانان در جنگ ایران–عراق؛
		مدل هیبریدی مقرون‌به‌صرفه}} \\
\midrule

۱۲ &
{\RL{تیمور شرقی (۱۹۹۹)}} &
{\RL{خشونت پس از همه‌پرسی استقلال؛
		۱,۴۰۰ کشته و ۳۰۰,۰۰۰ آواره؛
		تخریب ۷۰٪ زیرساخت‌ها؛
		نقش شبه‌نظامیان اندونزیایی}} &
{\RL{هیئت‌های ویژه‌ی دادگاه ناحیه‌ی دیلی؛
		کمیسیون پذیرش، حقیقت و آشتی (\en{CAVR})؛
		کمیسیون حقیقت و دوستی مشترک با اندونزی}} &
{\RL{همه‌پرسی تحت نظارت سازمان ملل → الگوی گذار؛
		کمیسیون مشترک با کشور ناقض
		→ آیا ساختار مشترکی ممکن است؟
		ضعف اجرا → درس عبرت}} \\
\midrule

۱۳ &
{\RL{لیبی (۲۰۱۱)}} &
{\RL{سرکوب خشونت‌بار اعتراضات؛
		بمباران شهرها؛
		استفاده از مزدوران؛
		زندان ابوسلیم (قتل‌عام ۱۹۹۶)}} &
{\RL{ارجاع شورای امنیت به \en{ICC}؛
		مداخله‌ی نظامی ناتو (\en{R2P})؛
		حکم بازداشت قذافی}} &
{\RL{تنها اجرای کامل \en{R2P}
		→ هم الهام‌بخش و هم هشداردهنده؛
		هرج‌ومرج پس از سقوط → لزوم برنامه‌ریزی برای «روز بعد»؛
		حکم \en{ICC} بدون دستگیری → محدودیت‌های عملی}} \\
\midrule

۱۴ &
{\RL{عراق (دوره‌ی صدام و پس از آن)}} &
{\RL{انفال و نسل‌کشی کردها (۱۹۸۸)؛
		حلبچه (سلاح شیمیایی)؛
		سرکوب شیعیان (۱۹۹۱)؛
		تالاب‌خشکانی؛
		تحریم‌های فراگیر}} &
{\RL{دادگاه ویژه‌ی عراق؛
		محاکمه و اعدام صدام (۲۰۰۶)؛
		\en{UNITAD} (مستندسازی جرایم داعش)؛
		دادگاه‌های ملی}} &
{\RL{تحریم‌های فراگیر عراق →
		نمونه‌ی آنچه نباید تکرار شود؛
		محاکمه‌ی صدام → انتقاد از عدم دادرسی منصفانه؛
		انفال → نمونه‌ی نسل‌کشی قومی؛
		حلبچه → سلاح شیمیایی و مسئولیت فروشندگان}} \\
\midrule

۱۵ &
{\RL{گواتمالا (جنگ داخلی ۱۹۶۰–۱۹۹۶)}} &
{\RL{نسل‌کشی بومیان مایا؛
		۲۰۰,۰۰۰ کشته و ناپدیدشده؛
		ارتش و شبه‌نظامیان؛
		حمایت آمریکا از نظامیان}} &
{\RL{کمیسیون روشنگری تاریخی (\en{CEH})؛
		محاکمه‌ی ریوس مونت (۲۰۱۳) —
		اولین محاکمه‌ی نسل‌کشی در دادگاه ملی؛
		\en{CICIG} (کمیسیون ضد فساد سازمان ملل)}} &
{\RL{محاکمه‌ی رئیس‌جمهور سابق در دادگاه ملی
		→ الگوی عدالت انتقالی داخلی؛
		\en{CICIG} → الگوی همکاری بین‌المللی–ملی؛
		نقش جامعه‌ی مدنی در پیگیری عدالت ۳۰+ ساله}} \\

\bottomrule
\end{longtable}
\end{landscape}


% ============================================================
%  بخش دوازدهم: موضوعات الهام‌گرفته از تجربه‌های تاریخی
% ============================================================
\section{حوزه‌ی دوازدهم: موضوعات تکمیلی الهام‌گرفته از تجربه‌های جهانی}

\begin{infoBox}[هدف]
موضوعات زیر بر اساس بررسی تطبیقی بخش قبل
به فهرست اضافه می‌شوند.
\end{infoBox}

\subsection{موضوع ۱۲.۱: رسانه‌ی نفرت‌پراکن و مسئولیت کیفری}

\begin{longtable}{r p{11cm}}
\toprule
\textbf{عنصر} & \textbf{شرح} \\
\midrule
\endhead
\textbf{تعریف} &
بررسی مسئولیت کیفری رسانه‌هایی که
به‌طور سیستماتیک نفرت، تبعیض و خشونت
علیه گروه‌های خاص را ترویج کرده‌اند. \\
\midrule
\textbf{مصادیق ایران} &
\begin{itemize}[nosep]
\item صداوسیمای جمهوری اسلامی:
\begin{itemize}[nosep]
	\item پخش اعترافات اجباری
	\item تبلیغات ضد بهایی
	\item تحریک خشونت علیه معترضان
	\item تحقیر اقلیت‌های قومی
	\item تبلیغات جنگ‌طلبانه
\end{itemize}
\item روزنامه‌ی کیهان و رسانه‌های وابسته به سپاه
\item شبکه‌های مجازی حکومتی (سایبری سپاه و بسیج)
\item منابر و خطبه‌های جمعه‌ی رسمی
\end{itemize} \\
\midrule
\textbf{نمونه‌ی تاریخی} &
\begin{principleBox}[پرونده‌ی رسانه‌ی رواندا — \en{ICTR}]
در پرونده‌ی «رسانه‌ها» (\en{Media Case}) ۲۰۰۳،
\en{ICTR} سه مدیر رسانه‌ی رواندایی را
به اتهام تحریک مستقیم و علنی به نسل‌کشی
محکوم کرد:
\begin{itemize}[nosep]
	\item فردیناند ناهیمانا (مدیر رادیو \en{RTLM})
	\item ژان بوسکو باراگویتزا (مدیر رادیو \en{RTLM})
	\item حسن نگزه (مدیر روزنامه‌ی کانگورا)
\end{itemize}
\textbf{اصل تثبیت‌شده:}
«رسانه‌ای که آگاهانه و مستقیماً به
خشونت و نسل‌کشی تحریک کند، مسئولیت کیفری دارد.»
\end{principleBox} \\
\midrule
\textbf{فعالیت دادگاه} &
\begin{enumerate}[nosep, label=\textcolor{cAccent}{\arabic*.}]
\item آرشیو کامل برنامه‌های تحریک‌آمیز صداوسیما
\item تحلیل محتوای نفرت‌پراکنی
(\en{content analysis}) به روش علمی
\item شناسایی مسئولان تولید و پخش
\item تدوین کیفرخواست سایه برای مدیران ارشد
\item ارائه‌ی مستندات به \en{UNESCO} و سازمان‌های
آزادی رسانه (\en{RSF}، \en{CPJ})
\end{enumerate} \\
\bottomrule
\end{longtable}

\subsection{موضوع ۱۲.۲: سربازگیری و بسیج کودکان}

\begin{longtable}{r p{11cm}}
\toprule
\textbf{عنصر} & \textbf{شرح} \\
\midrule
\endhead
\textbf{تعریف} &
بررسی و مستندسازی استفاده از کودکان
و نوجوانان زیر ۱۸ سال در:
\begin{itemize}[nosep]
\item جنگ ایران–عراق (بسیج دانش‌آموزان،
عملیات‌های انتحاری، مین‌روبی انسانی)
\item نیروهای نیابتی (فاطمیون، زینبیون)
— اعزام نوجوانان افغان و پاکستانی به سوریه
\item بسیج دانش‌آموزی و آموزش‌های نظامی
\item تشکیل «دسته‌های خودجوش» در سرکوب اعتراضات
\end{itemize} \\
\midrule
\textbf{مبنای حقوقی} &
\begin{itemize}[nosep]
\item ماده‌ی ۸(۲)(ب)(xxvi) اساسنامه‌ی رم:
سربازگیری کودکان زیر ۱۵ سال جرم جنگی
\item پروتکل اختیاری کنوانسیون حقوق کودک
درباره‌ی مشارکت کودکان در مخاصمات مسلحانه
\item اصول پاریس (۲۰۰۷) — آزادسازی و
بازپذیری کودکان سرباز
\end{itemize} \\
\midrule
\textbf{نمونه‌ی تاریخی} &
\begin{itemize}[nosep]
\item \textbf{پرونده‌ی توماس لوبانگا (\en{ICC} ۲۰۱۲):}
اولین محکومیت \en{ICC} — سربازگیری کودکان
در کنگو → سابقه‌ساز
\item \textbf{پرونده‌ی دومینیک اونگون (\en{ICC} ۲۰۲۱):}
ارتش مقاومت خداوند (اوگاندا) — سربازگیری و
بردگی جنسی کودکان
\item \textbf{سیرالئون:} کودکان سرباز انقلابیون — \en{SCSL}
\end{itemize} \\
\bottomrule
\end{longtable}

\subsection{موضوع ۱۲.۳: تخریب و غارت میراث فرهنگی}

\begin{longtable}{r p{11cm}}
\toprule
\textbf{عنصر} & \textbf{شرح} \\
\midrule
\endhead
\textbf{تعریف} &
مستندسازی تخریب عمدی یا ناشی از بی‌توجهی
میراث فرهنگی ایران:
\begin{itemize}[nosep]
\item تخریب بافت‌های تاریخی شهری
\item آسیب به محوطه‌های باستان‌شناسی
به‌واسطه‌ی سدسازی و طرح‌های عمرانی
\item تخریب گورستان‌ها
(گورستان بهائیان، خاوران و …)
\item سرقت و قاچاق آثار باستانی
\item تغییر نام‌ها و محو هویت تاریخی مکان‌ها
\item تبدیل اماکن تاریخی به ساختمان‌های حکومتی/مذهبی
\end{itemize} \\
\midrule
\textbf{مبنای حقوقی} &
\begin{itemize}[nosep]
\item کنوانسیون لاهه (۱۹۵۴) برای حمایت از
اموال فرهنگی در مخاصمات مسلحانه
\item کنوانسیون میراث جهانی \en{UNESCO} (۱۹۷۲)
\item ماده‌ی ۸(۲)(ب)(ix) اساسنامه‌ی رم:
حمله‌ی عمدی به اماکن تاریخی و مذهبی
\item \textbf{پرونده‌ی المهدی (\en{ICC} ۲۰۱۶):}
تخریب مقابر تیمبوکتو (مالی)
— اولین محکومیت صرفاً برای تخریب میراث فرهنگی
\end{itemize} \\
\bottomrule
\end{longtable}

\subsection{موضوع ۱۲.۴: نقض حقوق بین‌المللی بشردوستانه در مخاصمات خارجی}

\begin{longtable}{r p{11cm}}
\toprule
\textbf{عنصر} & \textbf{شرح} \\
\midrule
\endhead
\textbf{تعریف} &
بررسی مسئولیت جمهوری اسلامی در نقض حقوق بشر
و حقوق بشردوستانه در \keyword{خارج از مرزهای ایران}
از طریق نیروهای نظامی یا نیابتی:
\begin{itemize}[nosep]
\item \textbf{سوریه:} نقش سپاه قدس و نیروهای
فاطمیون/زینبیون در حمایت از اسد
— مشارکت در محاصره و بمباران غیرنظامیان
\item \textbf{عراق:} نقش شبه‌نظامیان حشد الشعبی
وابسته به ایران — اتهامات پاکسازی قومی
\item \textbf{یمن:} حمایت از حوثی‌ها
— حملات موشکی به غیرنظامیان عربستان
\item \textbf{لبنان:} حزب‌الله به‌عنوان بازوی نظامی ایران
— جنگ ۲۰۰۶ و پس از آن
\item \textbf{فلسطین / غزه:} حمایت از حماس و جهاد اسلامی
— پیچیدگی‌های حقوقی
\end{itemize} \\
\midrule
\textbf{مبنای حقوقی} &
\begin{itemize}[nosep]
\item کنوانسیون‌های ژنو (۱۹۴۹) و پروتکل‌های الحاقی
\item مسئولیت دولت برای اقدامات خارج از سرزمین
(\en{extraterritorial obligations})
\item مسئولیت تسهیل (\en{complicity})
در جرایم بین‌المللی
\item ماده‌ی ۲۵(۳)(ج) اساسنامه‌ی رم:
معاونت و مساعدت
\end{itemize} \\
\midrule
\textbf{نمونه} &
\begin{itemize}[nosep]
\item \textbf{پرونده‌ی تیلور (\en{SCSL}):}
رئیس‌جمهور لیبریا به دلیل حمایت
از شورشیان سیرالئون محکوم شد
→ الگوی مسئولیت برای حمایت از گروه‌های مسلح
\item \textbf{پرونده‌ی نیکاراگوئه (\en{ICJ} ۱۹۸۶):}
مسئولیت آمریکا برای حمایت از کنتراها
→ حقوق بین‌الملل عرفی درباره‌ی مداخله
\end{itemize} \\
\bottomrule
\end{longtable}

\subsection{موضوع ۱۲.۵: عدالت جنسیتی و خشونت جنسی سیستماتیک}

\begin{longtable}{r p{11cm}}
\toprule
\textbf{عنصر} & \textbf{شرح} \\
\midrule
\endhead
\textbf{تعریف} &
مستندسازی خشونت جنسی و جنسیتی به‌عنوان
ابزار سرکوب سیاسی:
\begin{itemize}[nosep]
\item تجاوز در بازداشتگاه‌ها (کهریزک ۱۳۸۸ و …)
\item «ازدواج‌های اجباری» پیش از اعدام
(دهه‌ی ۶۰)
\item آزمایش بکارت اجباری
\item خشونت جنسی در سرکوب اعتراضات
(جنبش ژن ژیان ئازادی)
\item حجاب اجباری به‌عنوان خشونت جنسیتی نهادینه
\end{itemize} \\
\midrule
\textbf{مبنای حقوقی} &
\begin{itemize}[nosep]
\item \textbf{پرونده‌ی آکایسو (\en{ICTR} ۱۹۹۸):}
تجاوز به‌عنوان عنصر نسل‌کشی
\item \textbf{پرونده‌ی کونارچ (\en{ICTY}):}
تجاوز به‌عنوان جرم علیه بشریت
\item ماده‌ی ۷(۱)(ز) اساسنامه‌ی رم:
تجاوز، بردگی جنسی، فحشای اجباری
به‌عنوان جرایم علیه بشریت
\item قطعنامه‌ی ۱۳۲۵ شورای امنیت (۲۰۰۰):
زنان، صلح و امنیت
\item کنوانسیون رفع تبعیض علیه زنان (\en{CEDAW})
\end{itemize} \\
\midrule
\textbf{خروجی} &
\begin{enumerate}[nosep, label=\textcolor{cAccent}{\arabic*.}]
\item بانک اطلاعاتی اختصاصی خشونت جنسی
با بالاترین سطح محرمانگی
\item پروتکل مصاحبه‌ی حساس به جنسیت
(\en{gender-sensitive interview protocol})
\item حمایت روان‌شناختی تخصصی از شاهدان
\item تحلیل الگوی سیستماتیک
(\en{pattern analysis})
\item کیفرخواست سایه برای فرماندهان مسئول
\end{enumerate} \\
\bottomrule
\end{longtable}

\subsection{موضوع ۱۲.۶: حق تعیین سرنوشت و حقوق اقوام}

\begin{longtable}{r p{11cm}}
\toprule
\textbf{عنصر} & \textbf{شرح} \\
\midrule
\endhead
\textbf{تعریف} &
بررسی وضعیت حقوقی اقوام و ملیت‌های ساکن ایران
در چارچوب حقوق بین‌الملل:
\begin{itemize}[nosep]
\item کردها، بلوچ‌ها، عرب‌ها (اهوازی‌ها)،
آذری‌ها، ترکمن‌ها، لرها، گیلکی‌ها و …
\item حق تعیین سرنوشت داخلی
(\en{internal self-determination})
— نه لزوماً جدایی بلکه خودمختاری،
فدرالیسم یا عدم‌تمرکز
\item حقوق زبانی و فرهنگی
\item عدم تبعیض در توسعه‌ی اقتصادی مناطق
\end{itemize} \\
\midrule
\textbf{ملاحظات حساس} &
\begin{warningBox}[هشدار بسیار مهم]
این موضوع \keyword{بالاترین پتانسیل اختلاف‌زایی}
در اپوزیسیون ایرانی را دارد. دادگاه در سایه باید:
\begin{itemize}[nosep]
	\item از تبدیل‌شدن به ابزار هر جریان سیاسی خاص
	(تجزیه‌طلب یا مرکزگرا) بپرهیزد
	\item بر \keyword{حقوق فردی و جمعیِ مستند در حقوق بین‌الملل}
	تمرکز کند
	\item موضع صریح درباره‌ی تمامیت ارضی نگیرد
	(خارج از صلاحیت دادگاه)
	\item صرفاً \keyword{نقض حقوق} را مستند و ارزیابی کند
\end{itemize}
\end{warningBox} \\
\midrule
\textbf{مبنای حقوقی} &
\begin{itemize}[nosep]
\item ماده‌ی ۱ مشترک میثاقین (حق تعیین سرنوشت)
\item اعلامیه‌ی حقوق اشخاص متعلق به اقلیت‌ها (۱۹۹۲)
\item اعلامیه‌ی حقوق اقوام بومی (\en{UNDRIP}) ۲۰۰۷
\item ماده‌ی ۲۷ میثاق حقوق مدنی و سیاسی
(حقوق اقلیت‌ها)
\end{itemize} \\
\bottomrule
\end{longtable}

\subsection{موضوع ۱۲.۷: نقض آزادی مذهبی و عقیده}

\begin{longtable}{r p{11cm}}
\toprule
\textbf{عنصر} & \textbf{شرح} \\
\midrule
\endhead
\textbf{تعریف} &
مستندسازی نقض سیستماتیک آزادی مذهب و عقیده:
\begin{itemize}[nosep]
\item \textbf{بهائیان:} سرکوب‌شده‌ترین اقلیت مذهبی —
محرومیت از تحصیل، شغل، مالکیت؛
زندان؛ اعدام؛ تخریب اماکن مقدس
\item \textbf{مسیحیان تازه‌ایمان (\en{converts}):}
زندان و شکنجه به اتهام «ارتداد»
\item \textbf{سنی‌ها:} ممنوعیت ساخت مسجد در تهران؛
تبعیض در استخدام دولتی
\item \textbf{دراویش گنابادی:} سرکوب خشونت‌بار ۲۰۱۸
\item \textbf{یهودیان:} فشار مستمر و تبلیغات ضدسامی
\item \textbf{زرتشتیان:} محدودیت‌های اجتماعی و حقوقی
\item \textbf{بی‌دینان و ملحدان:} مجازات اعدام
برای «ارتداد» در قانون مجازات
\end{itemize} \\
\midrule
\textbf{مبنای حقوقی} &
\begin{itemize}[nosep]
\item ماده‌ی ۱۸ میثاق حقوق مدنی و سیاسی
\item تفسیر عمومی شماره‌ی ۲۲ کمیته‌ی حقوق بشر
\item گزارش‌های گزارشگر ویژه‌ی آزادی مذهب و عقیده
\item قطعنامه‌های مجمع عمومی درباره‌ی
مبارزه با عدم تحمل مذهبی
\end{itemize} \\
\bottomrule
\end{longtable}


% ============================================================
%  جمع‌بندی نهایی پیوست ک
% ============================================================
\section{جمع‌بندی نهایی پیوست ک}

\begin{principleBox}[خلاصه‌ی فهرست موضوعات]
این پیوست \textbf{۱۲ حوزه‌ی موضوعی} با بیش از
\textbf{۴۰ موضوع مشخص} برای بررسی دادگاه در سایه
شناسایی کرده است:

\begin{enumerate}[nosep, label=\textcolor{cPrimary}{\arabic*.}]
\item \textbf{تحریم‌ها} (۶ موضوع)
\item \textbf{مستندسازی جرایم} (۳ موضوع)
\item \textbf{مداخله‌ی بشردوستانه و \en{R2P}} (۴ موضوع)
\item \textbf{تحریم‌ها — طبقه‌بندی و ارزیابی} (۴ موضوع)
\item \textbf{عوامل و نفوذی‌ها} (۳ موضوع)
\item \textbf{وابستگان سران} (۲ موضوع)
\item \textbf{همکاری بین‌المللی} (۳ موضوع)
\item \textbf{افشا و مستندسازی} (۳ موضوع)
\item \textbf{صلاحیت جهانی} (۲ موضوع)
\item \textbf{موضوعات ویژه} (۶ موضوع)
\item \textbf{درس‌های تاریخی} (۱۵ مورد تطبیقی)
\item \textbf{موضوعات تکمیلی} (۷ موضوع)
\end{enumerate}

\medskip
\textbf{اصول راهنما:}
\begin{itemize}[nosep]
\item فهرست \keyword{باز و قابل بسط} است
\item هر موضوع دارای تعریف، مبنای حقوقی، پرسش‌های کلیدی،
و نمونه‌های تاریخی است
\item اولویت‌بندی موضوعات توسط هیئت عالی انجام می‌شود
\item برخی موضوعات (مداخله، اقوام) \keyword{حساسیت بالا} دارند
و نیازمند احتیاط ویژه هستند
\end{itemize}
\end{principleBox}

\clearpage

	
	% ============================================================
	%  پیوست: مواجهه با وابستگان، عوامل و افراد خاکستری
	%  نسخه‌ی پاک‌سازی‌شده — بدون تکرار

	% ============================================================
	\section{حوزه‌ی اول: مواجهه با وابستگان مقام‌ها}
	
	\begin{warningBox}[حساسیت موضوع]
		این حوزه از حساس‌ترین موضوعات داوری است.
		هرگونه اقدام علیه «وابستگان» باید از آزمون‌های
		سخت‌گیرانه‌ی \keyword{مسئولیت فردی}
		(\en{individual responsibility}) و
		\keyword{تناسب} عبور کند.
		مجازات خانوادگی (\en{collective punishment})
		ممنوع است.
	\end{warningBox}
	
	
	% ─────────────────────────────────────────
	\subsection{موضوع ۱.۱: اخراج فرزندان و بستگان سران}
	
	\begin{longtable}{r p{11cm}}
		\toprule
		\textbf{عنصر} & \textbf{شرح} \\
		\midrule
		\endhead
		\textbf{تعریف} &
		درخواست از دولت‌های ثالث (کانادا، بریتانیا، آلمان، …)
		برای لغو اقامت یا تابعیت و اخراج
		فرزندان و بستگان سران جمهوری اسلامی. \\
		\midrule
		\textbf{نوع پرونده} &
		نوع ب (داوری تصمیم سیاسی) +
		نوع الف (مستندسازی ثروت‌های نامشروع) \\
		\midrule
		\textbf{تعاریف لازم} &
		\begin{itemize}[nosep]
			\item \keyword{سران:} رهبر، رئیس‌جمهور، اعضای شورای نگهبان،
			فرماندهان سپاه و بسیج (از سطح سرتیپ به بالا)،
			رؤسای قوه‌ی قضائیه، دادستان‌ها و قضات حکم‌دهنده‌ی
			اعدام سیاسی، وزرای اطلاعات و کشور،
			فرمانداران استان‌ها در دوره‌های سرکوب.
			\item \keyword{وابستگان نسبی:} فرزندان، همسران، والدین،
			خواهران و برادران.
			\item \keyword{وابستگان سببی:} داماد، عروس، خویشاوندان
			سببی درجه‌ی اول.
			\item \keyword{وابستگان عملکردی:}
			اشخاصی که رابطه‌ی خویشاوندی ندارند اما
			از ثروت یا نفوذ مقام بهره‌مند شده‌اند
			(شرکای تجاری، واسطه‌های مالی).
		\end{itemize} \\
		\midrule
		\textbf{آزمون سه‌لایه‌ای پیشنهادی} &
		% ← FIX: added unbreakable to override breakable inside table cell
		\begin{principleBox}[آزمون ویژه‌ی وابستگان]
			\textbf{لایه‌ی ۱ — ذی‌نفعی فعال:}
			آیا شخص از ثروت‌های حاصل از فساد یا سرکوب
			بهره‌مند شده؟ (مثلاً تحصیل در دانشگاه گران‌قیمت
			خارجی با پول مشکوک)
			
			\medskip
			\textbf{لایه‌ی ۲ — مشارکت فعال:}
			آیا شخص خودش در فعالیت‌های حکومتی،
			تبلیغاتی، اطلاعاتی، یا اقتصادی حکومت نقش دارد؟
			
			\medskip
			\textbf{لایه‌ی ۳ — غرابت ایدئولوژیک:}
			آیا شخص فعالانه از ایدئولوژی حکومت دفاع
			یا تبلیغ می‌کند؟
			
			\medskip
			\textbf{نتیجه‌گیری:}
			\begin{itemize}[nosep]
				\item هر سه لایه مثبت → \textcolor{cAccent}{اخراج توصیه‌شده}
				\item لایه‌ی ۱ و ۲ مثبت → \textcolor{cPrimary}{تحریم مالی توصیه‌شده}
				\item فقط لایه‌ی ۱ مثبت → \textcolor{cGold}{بررسی بیشتر}
				\item هیچ لایه مثبت نیست → \textcolor{cGray}{اقدامی توصیه نمی‌شود}
			\end{itemize}
		\end{principleBox} \\
		\midrule
		\textbf{نمونه‌ی تاریخی} &
		\textbf{اوکراین (از ۲۰۲۲):}
		اتحادیه‌ی اروپا فرزندان و خویشاوندان بسیاری از
		الیگارش‌ها و مقام‌های روسی را تحریم کرد —
		اما فقط آن‌هایی که خود ذی‌نفع فعال بودند.
		\keyword{مدل ذی‌نفعی فعال.}
		
		\medskip
		\textbf{کره‌ی شمالی:}
		تحریم‌های شورای امنیت شامل اعضای خانواده‌ی
		رهبری نیز هست — اما این مدل بسیار استثنایی
		و بحث‌برانگیز است.
		
		\medskip
		\textbf{درس:} تحریم خانوادگی فقط بر اساس رابطه‌ی
		خویشاوندی — بدون اثبات ذی‌نفعی یا مشارکت فعال —
		\keyword{مجازات جمعی محسوب می‌شود}
		و با حقوق بین‌الملل مغایرت دارد. \\
		\midrule
		\textbf{ملاحظات ویژه} &
		\begin{itemize}[nosep]
			\item فرزندان خردسال \textbf{هرگز} مشمول نمی‌شوند.
			\item فرزندان بالغ فقط در صورت اثبات ذی‌نفعی فعال.
			\item حق دفاع: هر شخص نام‌بُرده حق دارد
			از طریق وکیل تسخیری دفاع شود.
			\item جمهوری اسلامی تلاش خواهد کرد از این موضوع
			بهره‌برداری تبلیغاتی کند —
			آزمون سه‌لایه‌ای مستدل، بهترین سپر دفاعی است.
		\end{itemize} \\
		\bottomrule
	\end{longtable}
	
	
	% ─────────────────────────────────────────
	\subsection{موضوع ۱.۲: عوامل، نفوذی‌ها و دست‌نشاندگان}
	
	\begin{longtable}{r p{11cm}}
		\toprule
		\textbf{عنصر} & \textbf{شرح} \\
		\midrule
		\endhead
		\textbf{تعریف} &
		شناسایی و مواجهه با اشخاصی که در خارج از کشور
		به نفع جمهوری اسلامی فعالیت می‌کنند:
		عوامل اطلاعاتی، لابیست‌ها، فعالان رسانه‌ای حکومتی،
		اعضای شبکه‌های نفوذ (\en{influence networks})،
		و مدیران مالی ثروت‌های حکومتی. \\
		\midrule
		\textbf{نوع پرونده} &
		نوع الف (مستندسازی) + نوع ب (توصیه‌ی اقدام) \\
		\midrule
		\textbf{طبقه‌بندی پیشنهادی} &
		\begin{principleBox}[سه سطح عاملیت]
			\textbf{سطح ۱ — عامل مستقیم (\en{Agent}):}
			شخصی که مستقیماً از نهاد اطلاعاتی‌ـ‌امنیتی
			جمهوری اسلامی دستور می‌گیرد
			(مأمور وزارت اطلاعات، سپاه قدس، …).
			\textcolor{cAccent}{بالاترین سطح خطر. توصیه: گزارش به
				مراجع قضایی و امنیتی کشور میزبان.}
			
			\medskip
			\textbf{سطح ۲ — نفوذی (\en{Influence Agent}):}
			شخصی که بدون رابطه‌ی رسمی اطلاعاتی،
			اما به‌صورت آگاهانه و مستمر منافع حکومت را
			در فضای عمومی، رسانه‌ای، یا سیاسی پیش می‌برد.
			\textcolor{cAccent}{توصیه: افشای عمومی مستند + درخواست
				ثبت به‌عنوان \en{foreign agent}.}
			
			\medskip
			\textbf{سطح ۳ — همراه ناآگاه (\en{Unwitting Asset}):}
			شخصی که بدون آگاهی کامل، مورد بهره‌برداری
			اطلاعاتی حکومت قرار گرفته.
			\textcolor{cPrimary}{توصیه: تماس محرمانه و هشدار.
				افشای عمومی فقط در صورت عدم تغییر رفتار.}
		\end{principleBox} \\
		\midrule
		\textbf{پرسش‌های کلیدی} &
		\begin{enumerate}[nosep, label=\textcolor{cPrimary}{\roman*.}]
			\item استاندارد اثبات عاملیت چقدر باید سخت‌گیرانه باشد؟
			(خطر اتهام ناروا بسیار بالاست)
			\item تفاوت بین «لابی‌گری مشروع» و «نفوذ اطلاعاتی» کجاست؟
			\item آیا اشخاص ایرانی‌تبار شاغل در رسانه‌های دولتی
			(\en{Press TV}، \en{HispanTV}) عامل محسوب می‌شوند؟
			\item سازوکار حفاظت از افراد ناحق متهم‌شده چیست؟
			\item آیا افشای عمومی در همه‌ی موارد مناسب است
			یا در برخی موارد اقدام محرمانه بهتر است؟
		\end{enumerate} \\
		\midrule
		\textbf{نمونه‌های تاریخی} &
		\textbf{قانون \en{FARA} آمریکا (۱۹۳۸):}
		قانون ثبت عوامل خارجی — الزام به افشای
		فعالیت‌های لابی‌گری به نفع دولت‌های خارجی.
		ایران یکی از کشورهایی است که عوامل آن
		بارها \en{FARA} را نقض کرده‌اند.
		
		\medskip
		\textbf{استرالیا (\en{FITS Act}, ۲۰۱۸):}
		قانون مشابه \en{FARA} اما سخت‌گیرانه‌تر —
		پس از کشف نفوذ گسترده‌ی چین
		در سیاست استرالیا تصویب شد.
		
		\medskip
		\textbf{بریتانیا (\en{National Security Act}, ۲۰۲۳):}
		جدیدترین قانون ضد جاسوسی و نفوذ خارجی —
		شامل جرم‌انگاری فعالیت اطلاعاتی پنهان
		و الزام ثبت فعالیت‌های سیاسی به نفع دولت‌های خارجی.
		
		\medskip
		\textbf{کانادا (\en{Foreign Influence Transparency Registry}, ۲۰۲۴):}
		قانون جدید شفافیت نفوذ خارجی —
		واکنش به نفوذ چین در سیاست کانادا.
		قابل استفاده به‌عنوان مبنای درخواست
		ثبت عوامل ایرانی نیز هست. \\
		\bottomrule
	\end{longtable}
	
	
	% ============================================================
	%  بخش دوم: افراد بینابینی — «خاکستری‌ها»
	% ============================================================
	\section{حوزه‌ی دوم: «خاکستری‌ها» — جایگاه افراد بینابینی}
	
	\begin{infoBox}[چرا این بخش مهم است؟]
		در هر فرایند گذار سیاسی، بزرگ‌ترین بخش جمعیت
		نه در اردوگاه سرکوب‌گران قرار دارد و نه
		در صف مقاومت فعال. این «اکثریت خاکستری»
		سرنوشت گذار را تعیین می‌کند.
		نحوه‌ی مواجهه با آن‌ها از مهم‌ترین
		تصمیمات راهبردی هر جنبش آزادی‌خواهانه است.
	\end{infoBox}
	
	
	% ─────────────────────────────────────────
	\subsection{موضوع ۲.۱: تعریف و شناخت «خاکستری‌ها»}
	
	\begin{longtable}{r p{11cm}}
		\toprule
		\textbf{عنصر} & \textbf{شرح} \\
		\midrule
		\endhead
		\textbf{تعریف} &
		افرادی که نه کاملاً حکومتی‌اند و نه کاملاً اپوزیسیون:
		\begin{itemize}[nosep]
			\item ورزشکارانی که پرچم حکومت را حمل می‌کنند اما
			اعتراض نمی‌کنند و شعار نمی‌دهند؛
			\item هنرمندانی که در نهادهای حکومتی فعالیت دارند اما
			در خفا از جنبش حمایت می‌کنند؛
			\item تاجرانی که با نهادهای حکومتی کار می‌کنند اما
			به مردم هم خدمت می‌رسانند؛
			\item اساتید دانشگاه‌های دولتی؛
			\item کارمندان دولت بدون نقش سرکوب‌گرانه.
		\end{itemize} \\
		\midrule
		\textbf{چالش اساسی} &
		\begin{warningBox}[معضل اخلاقی]
			آیا در شرایط انقلابی، سکوت مساوی با همراهی است؟
			یا باید شرایط فشار بر افراد داخل ایران را
			در نظر گرفت و پذیرفت که بسیاری
			\keyword{انتخاب‌های محدودی} دارند؟
			
			\medskip
			دادگاه در سایه باید بین دو افراط بپرهیزد:
			\begin{itemize}[nosep]
				\item \textbf{افراط ۱:} هر کس ساکت است، مقصر است
				→ \textcolor{cAccent}{مجازات جمعی و بی‌عدالتی}
				\item \textbf{افراط ۲:} هر کس سرکوب نکرده، بی‌گناه است
				→ \textcolor{cAccent}{نادیده‌گرفتن تسهیل‌کنندگان سرکوب}
			\end{itemize}
		\end{warningBox} \\
		\bottomrule
	\end{longtable}
	
	
	% ─────────────────────────────────────────
	\subsection{موضوع ۲.۲: «نردبان خاکستری» — طبقه‌بندی پنج‌پله‌ای}
	
	\begin{longtable}{r p{11cm}}
		\toprule
		\textbf{عنصر} & \textbf{شرح} \\
		\midrule
		\endhead
		\textbf{تعریف} &
		«نردبان خاکستری» یک طیف پیوسته است که
		از سکوت همدلانه تا مشارکت غیرمستقیم در سرکوب
		را در بر می‌گیرد. هدف: جایگزین‌کردن
		قضاوت سیاه‌وسفید با ارزیابی ساختارمند و عادلانه. \\
		\midrule
		\textbf{پله‌های نردبان} &
		\begin{principleBox}[نردبان خاکستری — پنج پله]
			\begin{enumerate}[nosep, label=\textcolor{cGray}{\arabic*.}]
				\item \textbf{خاکستری روشن — مقاوم خاموش:}
				کارمندان و متخصصانی که بدون مشارکت در سرکوب،
				به دلایل معیشتی یا وطن‌دوستانه در ایران مانده‌اند
				و در خفا با جنبش همدل‌اند.
				\textcolor{cGold}{→ هیچ اقدامی علیه آن‌ها توصیه نمی‌شود.
					در صورت امکان، حمایت شوند.}
				
				\item \textbf{خاکستری روشن — منفعل:}
				افرادی که نه مشارکت فعال در سرکوب دارند و نه
				حمایت پنهانی از جنبش. صرفاً سکوت کرده‌اند.
				\textcolor{cGold}{→ اقدامی توصیه نمی‌شود.
					اما انتظار همکاری آتی وجود دارد.}
				
				\item \textbf{خاکستری میانه — ذی‌نفع بدون مشارکت:}
				کسانی که از امکانات حکومت بهره‌مند شده‌اند
				(بورسیه‌ی حکومتی، قرارداد دولتی، …)
				بدون مشارکت مستقیم در سرکوب.
				\textcolor{cPrimary}{→ بررسی موردی.
					اولویت: شفاف‌سازی منابع مالی.}
				
				\item \textbf{خاکستری تیره — همراه فعال:}
				کسانی که فعالانه (اما نه در نقش سرکوب‌گر)
				با حکومت همکاری می‌کنند: تبلیغات، لابی‌گری،
				حضور در مراسم رسمی، دفاع رسانه‌ای.
				\textcolor{cAccent}{→ افشای مستند. درخواست ثبت
					به‌عنوان عامل خارجی در صورت فعالیت خارج از کشور.}
				
				\item \textbf{خاکستری بسیار تیره — مباشر غیرمستقیم:}
				کسانی که بدون حضور مستقیم در صحنه‌ی سرکوب،
				اما با تأمین مالی، فنی، یا اطلاعاتی
				به دستگاه سرکوب کمک کرده‌اند.
				\textcolor{cAccent}{→ مستندسازی + کیفرخواست سایه
					+ درخواست تحریم هدفمند.}
			\end{enumerate}
		\end{principleBox} \\
		\midrule
		\textbf{نمودار بصری} &
		% ← FIX: ALL Persian text inside TikZ nodes wrapped with \RL{}
		% ← FIX: each line of a multi-line node uses a separate \RL{}
		\begin{tikzpicture}[
			node distance=0.15cm,
			every node/.style={
				font=\small, align=center,
				text width=2.5cm, minimum height=1cm,
				rounded corners=3pt
				}
				]
				% ← FIX: Every Persian text in TikZ nodes wrapped with \RL{}
				\node[fill=green!15, draw=green!60!black] (p1)
				{\RL{۱}\\\RL{مقاوم خاموش}};
				\node[fill=yellow!15, draw=yellow!60!black, right=of p1] (p2)
				{\RL{۲}\\\RL{منفعل}};
				\node[fill=orange!15, draw=orange!60!black, right=of p2] (p3)
				{\RL{۳}\\\RL{ذی‌نفع}};
				\node[fill=red!15, draw=red!60!black, right=of p3] (p4)
				{\RL{۴}\\\RL{همراه فعال}};
				\node[fill=red!30, draw=red!80!black, right=of p4] (p5)
				{\RL{۵}\\\RL{مباشر غیرمستقیم}};
				
				\draw[-{Stealth}, thick, green!60!black] (p1) -- (p2);
				\draw[-{Stealth}, thick, yellow!60!black] (p2) -- (p3);
				\draw[-{Stealth}, thick, orange!60!black] (p3) -- (p4);
				\draw[-{Stealth}, thick, red!60!black] (p4) -- (p5);
				
				% ← FIX: labels below nodes also wrapped with \RL{}
				\node[below=0.3cm of p1, font=\tiny\color{green!60!black},
				text width=2cm, minimum height=0pt] {\RL{حمایت}};
				\node[below=0.3cm of p2, font=\tiny\color{yellow!60!black},
				text width=2cm, minimum height=0pt] {\RL{بی‌اقدام}};
				\node[below=0.3cm of p3, font=\tiny\color{orange!60!black},
				text width=2cm, minimum height=0pt] {\RL{بررسی}};
				\node[below=0.3cm of p4, font=\tiny\color{red!60!black},
				text width=2cm, minimum height=0pt] {\RL{افشا}};
				\node[below=0.3cm of p5, font=\tiny\color{red!80!black},
				text width=2cm, minimum height=0pt] {\RL{تعقیب}};
			\end{tikzpicture} \\
		\midrule
	\textbf{معیارهای تشخیص} &
\begin{enumerate}[nosep, label=\textcolor{cPrimary}{\roman*.}]
\item \textbf{منبع درآمد:} آیا از بودجه‌ی عمومی،
نهادهای حکومتی، یا شرکت‌های وابسته حقوق می‌گیرد؟
\item \textbf{فعالیت عمومی:} آیا در رسانه‌ها، مراسم رسمی،
یا فضای مجازی از حکومت دفاع می‌کند؟
\item \textbf{نقش در سرکوب:} آیا مستقیم یا غیرمستقیم
در عملیات سرکوب، سانسور، بازجویی،
یا قضاوت مشارکت دارد؟
\item \textbf{دسترسی به اطلاعات:} آیا به اطلاعات حساس
حکومتی دسترسی دارد و آن‌ها را فاش نمی‌کند؟
\item \textbf{ثروت نامتناسب:} آیا ثروتی بسیار فراتر از
درآمد رسمی خود دارد؟
\item \textbf{رفتار در لحظات بحرانی:} در جریان اعتراضات
۱۳۸۸، ۱۳۹۸ و ۱۴۰۱ چه موضعی گرفت؟
\end{enumerate} \\
\midrule
\textbf{ویژگی پویایی نردبان} &
\begin{infoBox}[اصل مهم: نردبان پویاست]
افراد می‌توانند بین پله‌ها \keyword{حرکت کنند}:
\begin{itemize}[nosep]
\item فردی در پله‌ی ۴ (همراه فعال) که استعفا دهد
و اطلاعات افشا کند → می‌تواند به پله‌ی ۱ برود
\item فردی در پله‌ی ۲ (منفعل) که قرارداد مالی
با نهاد سرکوب‌گر ببندد → به پله‌ی ۳ یا ۴ می‌رود
\item \textbf{ارزیابی دوره‌ای} (هر ۶ ماه) ضروری است
\end{itemize}
\end{infoBox} \\
\bottomrule
\end{longtable}


% ─────────────────────────────────────────
\subsection{موضوع ۲.۳: «خروج از خاکستری» — معیارها و مشوق‌ها}

\begin{longtable}{r p{11cm}}
\toprule
\textbf{عنصر} & \textbf{شرح} \\
\midrule
\endhead
\textbf{تعریف} &
«خروج از خاکستری» یعنی تغییر وضعیت از
سکوت یا همراهی منفعلانه به موضع‌گیری
فعالانه در برابر حکومت.
این فرایند می‌تواند تدریجی یا ناگهانی باشد. \\
\midrule
\textbf{معیارهای خروج} &
\begin{enumerate}[nosep, label=\textcolor{cAccent}{\arabic*.}]
\item \textbf{اعلام موضع عمومی:}
بیانیه‌ی رسمی، مصاحبه، حضور در تجمعات
(حتی اگر کوتاه باشد).
\item \textbf{استعفا از سمت حکومتی:}
استعفای علنی از هرگونه موقعیت در
نهادهای دولتی، قضایی، یا نظامی.
\item \textbf{همکاری با نهادهای مستقل:}
ارائه‌ی اطلاعات، کمک مالی، یا
مشارکت در پروژه‌های مستقل.
\item \textbf{حمایت از زندانیان سیاسی:}
کمک به خانواده‌ی زندانیان، امضای
نامه‌های آزادی، یا حمایت رسانه‌ای.
\item \textbf{تحریم شخصی:}
تحریم خودخواسته از هرگونه همکاری
با حکومت (حتی در سطح خرد).
\item \textbf{سوت‌زنی (\en{whistleblowing}):}
افشای اطلاعات داخلی درباره‌ی فساد
یا سرکوب — بالاترین ارزش خروج.
\end{enumerate} \\
\midrule
\textbf{مشوق‌های خروج} &
\begin{principleBox}[بسته‌ی تشویقی خروج از خاکستری]
دادگاه در سایه باید نه فقط «تنبیه خاکستری‌ها»
بلکه \keyword{تشویق خروج} را هم در دستور کار داشته باشد:
\begin{enumerate}[nosep, label=\textcolor{cPrimary}{\arabic*.}]
\item \textbf{عدم تعقیب گذشته:}
افرادی که پیش از نقطه‌ی مشخصی
از حکومت جدا شوند، از بررسی فعالیت‌های
گذشته‌شان (در پله‌های ۱–۳) معاف شوند
\item \textbf{حمایت امنیتی:}
حفاظت هویت سوت‌زن‌ها و افشاگران
\item \textbf{حمایت حقوقی:}
مشاوره‌ی رایگان برای کسانی که
به‌دنبال استعفا یا فرار از ایران‌اند
\item \textbf{به‌رسمیت‌شناسی عمومی:}
تقدیر از افرادی که با شجاعت
از حکومت جدا شده‌اند (با رضایت خودشان)
\item \textbf{نقش در دوره‌ی گذار:}
وعده‌ی مشارکت در ساختارهای آینده
برای متخصصانی که به‌موقع جدا شوند
\end{enumerate}
\end{principleBox} \\
\midrule
\textbf{پرسش‌های کلیدی} &
\begin{enumerate}[nosep, label=\textcolor{cPrimary}{\roman*.}]
\item آیا «خروج از خاکستری» باید توسط خود فرد
اعلام شود یا توسط نهادهای مستقل تشخیص داده شود؟
\item چه مدت زمانی باید از «خروج» بگذرد
تا فرد از وضعیت خاکستری خارج شود؟
\item آیا «خروج» یک‌باره است یا تدریجی؟
\item آیا «بازگشت به خاکستری» پس از خروج ممکن است؟
(مثلاً فردی که پس از اعلام مخالفت،
دوباره با حکومت همکاری کند)
\item آیا «خروج» باید مستند شود
یا صرف ادعای شفاهی کافی است؟
\item آیا برای افراد داخل ایران
(که خطر دستگیری دارند)
معیارهای متفاوتی باید تعریف شود؟
\end{enumerate} \\
\bottomrule
\end{longtable}


% ─────────────────────────────────────────
\subsection{موضوع ۲.۴: نمونه‌های تاریخی مرتبط}

\begin{longtable}{r p{4.5cm} p{6cm}}
\toprule
\textbf{نمونه} &
\textbf{شرح} &
\textbf{درس برای ایران} \\
\midrule
\endhead

\textbf{نازی‌زدایی آلمان}
\newline (\en{Denazification})
\newline ۱۹۴۵–۱۹۵۰ &
جمعیت به پنج دسته تقسیم شد:
(۱) مجرمان اصلی (\en{Major Offenders})،
(۲) مجرمان (\en{Offenders})،
(۳) مجرمان جزئی (\en{Lesser Offenders})،
(۴) پیروان (\en{Followers})،
(۵) تبرئه‌شدگان (\en{Exonerated}).
بیش از ۸ میلیون نفر بررسی شدند. &
\textbf{درس مثبت:} طبقه‌بندی ساختارمند
از عدالت کور جلوگیری کرد.
\newline
\textbf{درس منفی:} فرایند بسیار کُند بود
و بسیاری از مجرمان واقعی
از لابه‌لای طبقه‌بندی فرار کردند.
«نردبان خاکستری» دادگاه در سایه
باید ساده‌تر و عملیاتی‌تر باشد. \\
\midrule

\textbf{گاچاچا — رواندا}
\newline ۱۹۹۴–۲۰۱۲ &
دادگاه‌های محلی جمعیت را
به سه دسته تقسیم کردند بر اساس
سطح مشارکت در نسل‌کشی.
افراد دسته‌ی سوم (پیروان)
با اعتراف و خدمات اجتماعی آزاد شدند.
بیش از ۱.۲ میلیون پرونده رسیدگی شد. &
\keyword{مدل عدالت ترمیمی محلی.}
اعتراف در ازای تخفیف مجازات —
آیا مدلی مشابه برای «خاکستری‌ها»ی
ایرانی طراحی‌شدنی است؟ \\
\midrule

\textbf{جنبش ضد آپارتاید}
\newline آفریقای جنوبی &
افراد بر اساس میزان همکاری با
رژیم آپارتاید طبقه‌بندی می‌شدند:
از همکاران فعال تا مخالفان خاموش.
بازیکنان فوتبال، هنرمندان، و دانشگاهیانی
که از رژیم جدا شدند، نقش نمادین مهمی
در تضعیف مشروعیت حکومت داشتند. &
تأثیر نمادین «خروج» افراد سرشناس
بسیار بیشتر از تأثیر عملی آن است.
\newline
کمیسیون حقیقت و آشتی (\en{TRC})
بستری برای «خروج رسمی از خاکستری»
فراهم کرد: اعتراف کامل در ازای عفو. \\
\midrule

\textbf{همبستگی (\en{Solidarność})}
\newline لهستان ۱۹۸۰–۱۹۸۹ &
کارگران در کارخانه‌ها به‌صورت تدریجی
از همراهی با رژیم خارج شدند:
(۱) اعتراضات محلی،
(۲) تشکیل سندیکای مستقل،
(۳) اعتصاب سراسری،
(۴) مذاکرات با حکومت. &
\keyword{الگوی خروج تدریجی:}
لازم نیست خروج ناگهانی و قهرمانانه باشد.
مراحل میانی (مانند عدم شرکت در
انتخابات فرمایشی، عدم حضور در مراسم رسمی)
هم ارزشمند است. \\
\midrule

\textbf{مفهوم «همدلی پنهان»}
\newline (\en{Hidden Sympathizers}) &
در بسیاری از جنبش‌های مقاومت
(فرانسه‌ی اشغالی، اروپای شرقی،
آفریقای جنوبی)، افرادی وجود داشتند
که در ظاهر همراه رژیم بودند
اما در خفا از مقاومت حمایت می‌کردند. &
\textbf{درس کلیدی:}
نباید فرض شود که هر کس در دستگاه
دولت حضور دارد، دشمن است.
بسیاری از تحولات بزرگ تاریخی
با «خیانت نخبگان حکومتی»
(\en{elite defection}) آغاز شده‌اند
— و دادگاه در سایه باید
بستری برای تسهیل این فرایند باشد. \\

\bottomrule
\end{longtable}


% ─────────────────────────────────────────
\subsection{موضوع ۲.۵: ملاحظات اخلاقی و خطوط قرمز}

\begin{warningBox}[اصول بنیادین مواجهه با خاکستری‌ها]
\begin{enumerate}[nosep, label=\textcolor{cAccent}{\arabic*.}]
\item \textbf{اصل تفاوت بین داخل و خارج:}
افراد داخل ایران تحت فشار شدید زندگی می‌کنند
و انتخاب‌های محدودی دارند.
معیارهای ارزیابی آن‌ها باید
\keyword{ملایم‌تر} از معیارهای ارزیابی
افراد مقیم خارج باشد.

\item \textbf{اصل بار اثبات:}
برای هر ادعایی درباره‌ی «خاکستری تیره» بودن فرد،
بار اثبات بر عهده‌ی \keyword{مدعی} است
— نه بر عهده‌ی فرد متهم.

\item \textbf{اصل تناسب پیامدها:}
پیامدهای ارزیابی هر پله باید
متناسب با سطح مشارکت باشد.
هیچ‌گاه از ابزارهای پله‌ی ۵ (کیفرخواست)
برای فرد پله‌ی ۲ (منفعل) استفاده نشود.

\item \textbf{اصل تشویق، نه تنبیه:}
هدف اصلی نردبان خاکستری
\keyword{تشویق خروج} از حکومت است
— نه «تنبیه ساکتان».

\item \textbf{اصل حق دفاع:}
هر فرد ارزیابی‌شده حق دارد
از طریق مکانیسم اعتراض رسمی،
رتبه‌بندی خود را به چالش بکشد.

\item \textbf{اصل عدم انتشار عمومی:}
رتبه‌بندی افراد در پله‌های ۱ تا ۳
\keyword{هرگز} به‌صورت عمومی منتشر نمی‌شود.
فقط پله‌های ۴ و ۵ (با ادله‌ی محکم)
ممکن است موضوع افشای عمومی شوند.

\item \textbf{اصل محرمانگی داخل:}
هرگونه اطلاعات درباره‌ی «مقاومان خاموش»
(پله‌ی ۱) باید با بالاترین سطح امنیتی
محافظت شود — افشای هویت آن‌ها
مستقیماً جان آن‌ها را به خطر می‌اندازد.
\end{enumerate}
\end{warningBox}


% ============================================================
%  جمع‌بندی بصری — نمودار کلی
% ============================================================
\subsection*{جمع‌بندی بصری: نقشه‌ی مفهومی حوزه‌ی خاکستری‌ها}

% ← FIX: ALL Persian text inside mindmap TikZ nodes wrapped with \RL{}
% ← FIX: multi-line nodes split into separate \RL{} per line
% ← FIX: level 2 concept text width increased slightly for \RL{} content
\begin{tikzpicture}[
mindmap,
every node/.style={concept, minimum size=1.5cm, text width=2.8cm,
font=\small\bfseries, align=center},
root concept/.append style={concept color=cPrimary!30},
level 1 concept/.append style={sibling angle=72},
level 2 concept/.append style={sibling angle=45, font=\tiny,
text width=2.2cm},  % ← FIX: slightly wider for RL text
]
\node[root concept] {\RL{خاکستری‌ها}}
child[concept color=green!30] {
node {\RL{مقاوم خاموش}\\\RL{(پله ۱)}}
child { node {\RL{حمایت}} }
child { node {\RL{حفاظت هویت}} }
}
child[concept color=yellow!30] {
node {\RL{منفعل}\\\RL{(پله ۲)}}
child { node {\RL{بی‌اقدام}} }
child { node {\RL{انتظار آتی}} }
}
child[concept color=orange!30] {
node {\RL{ذی‌نفع}\\\RL{(پله ۳)}}
child { node {\RL{بررسی مالی}} }
child { node {\RL{شفاف‌سازی}} }
}
child[concept color=red!20] {
node {\RL{همراه فعال}\\\RL{(پله ۴)}}
child { node {\RL{افشای مستند}} }
child { node {\RL{ثبت عامل خارجی}} }
}
child[concept color=red!40] {
node {\RL{مباشر غیرمستقیم}\\\RL{(پله ۵)}}
child { node {\RL{کیفرخواست سایه}} }
child { node {\RL{تحریم هدفمند}} }
};
\end{tikzpicture}


\begin{center}
\rule{0.6\textwidth}{0.5pt}

\vspace{1em}
{\Large\color{cPrimary}
\textbf{پایان بخش — مواجهه با وابستگان، عوامل و خاکستری‌ها}}

\vspace{0.5em}
{\small\color{gray}
نسخه‌ی پاک‌سازی‌شده — بدون تکرار.}
\end{center}

% ============================================================
%  جمع‌بندی فهرست موضوعات — ماتریس نهایی
% ============================================================
\section{جمع‌بندی: ماتریس کلان موضوعات قابل‌بررسی}

\begin{infoBox}[توضیح]
جدول زیر \keyword{فهرست یکپارچه و نهایی}
تمام موضوعاتی است که دادگاه در سایه
می‌تواند در دستور کار قرار دهد.
هر موضوع با اولویت، نوع پرونده و سطح
حساسیت مشخص شده است.
\end{infoBox}

\begin{landscape}
% ← FIX: ALL \rotatebox{90}{...} Persian text wrapped with \RL{}
% ← FIX: column spec adjusted for landscape readability
\begin{longtable}{r r p{5cm} c c c}
\toprule
\textbf{ردیف} &
\textbf{حوزه} &
\textbf{موضوع} &
\textbf{نوع پرونده} &
\textbf{اولویت} &
\textbf{حساسیت} \\
\midrule
\endhead

% --- حوزه‌ی ۱: حقوق بشر مدنی و سیاسی ---
۱ & \multirow{6}{*}{\rotatebox{90}{\RL{\textbf{حقوق بشر}}}}
& اعدام‌های سیاسی و کشتار ۶۷ & الف & ★★★★★ & متوسط \\
۲ & & شکنجه و بدرفتاری سیستماتیک & الف & ★★★★★ & متوسط \\
۳ & & بازداشت خودسرانه و ناپدیدسازی & الف/ب & ★★★★★ & متوسط \\
۴ & & سرکوب اعتراضات (آبان ۹۸، ژینا ۱۴۰۱) & الف & ★★★★★ & کم \\
۵ & & سرکوب آزادی بیان و رسانه & الف/ب & ★★★★☆ & کم \\
۶ & & نقض حقوق اقلیت‌های قومی و مذهبی & الف & ★★★★★ & بالا \\
\midrule

% --- حوزه‌ی ۲: حقوق زنان و اقلیت‌های جنسی ---
۷ & \multirow{4}{*}{\rotatebox{90}{\RL{\textbf{جنسیتی}}}}
& حجاب اجباری و نابرابری حقوقی & الف/ج & ★★★★★ & کم \\
۸ & & خشونت جنسی سیستماتیک در بازداشتگاه‌ها & الف & ★★★★★ & بالا \\
۹ & & سرکوب اقلیت‌های جنسی (\en{LGBTQ+}) & الف & ★★★★☆ & بالا \\
۱۰ & & نقض حقوق آموزشی زنان & الف & ★★★☆☆ & کم \\
\midrule

% --- حوزه‌ی ۳: مداخله و R2P ---
۱۱ & \multirow{4}{*}{\rotatebox{90}{\textbf{\en{R2P}}}}
& ارزیابی شرایط \en{R2P} & ج & ★★★☆☆ & بسیار بالا \\
۱۲ & & مناطق پرواز ممنوع و کریدور بشردوستانه & ج & ★★☆☆☆ & بسیار بالا \\
۱۳ & & ارجاع به \en{ICC} & ب & ★★★★★ & بالا \\
۱۴ & & تأسیس دادگاه ویژه / هیبریدی & ج & ★★★★☆ & بالا \\
\midrule

% --- حوزه‌ی ۴: تحریم‌ها ---
۱۵ & \multirow{4}{*}{\rotatebox{90}{\RL{\textbf{تحریم‌ها}}}}
& تحریم‌های هدفمند فردی (مگنیتسکی) & ب & ★★★★★ & کم \\
۱۶ & & تحریم‌های بخشی — ارزیابی و پیشنهاد & ج & ★★★★☆ & متوسط \\
۱۷ & & تحریم‌های ثانویه و فراسرزمینی & ج & ★★★☆☆ & بالا \\
۱۸ & & ارزیابی شدت و تناسب — شاخص‌سازی & ج & ★★★★☆ & متوسط \\
\midrule

% --- حوزه‌ی ۵: عوامل و نفوذی‌ها ---
۱۹ & \multirow{3}{*}{\rotatebox{90}{\RL{\textbf{نفوذ}}}}
& شناسایی عوامل اطلاعاتی و نفوذی‌ها & الف & ★★★★☆ & بسیار بالا \\
۲۰ & & معیارهای ادله و دادرسی منصفانه & ج & ★★★★★ & بالا \\
۲۱ & & سازوکارهای مواجهه و نتایج & الف & ★★★☆☆ & بسیار بالا \\
\midrule

% --- حوزه‌ی ۶: وابستگان سران ---
۲۲ & \multirow{3}{*}{\rotatebox{90}{\RL{\textbf{سران}}}}
& تعریف سران و وابستگان & ج & ★★★★☆ & بالا \\
۲۳ & & مستندسازی ثروت و درخواست مسدودسازی & ب & ★★★★★ & متوسط \\
۲۴ & & درخواست لغو اقامت و اخراج (فردی/مستند) & ب & ★★★★☆ & بالا \\
\midrule

% --- حوزه‌ی ۷: همکاری بین‌المللی ---
۲۵ & \multirow{3}{*}{\rotatebox{90}{\RL{\textbf{بین‌الملل}}}}
& همکاری با گزارشگران ویژه‌ی سازمان ملل & ب & ★★★★★ & کم \\
۲۶ & & همکاری با \en{IIMM} و هیئت‌های حقیقت‌یاب & ب & ★★★★★ & کم \\
۲۷ & & مشارکت در تدوین قطعنامه‌ها & ب & ★★★★☆ & کم \\
\midrule

% --- حوزه‌ی ۸: مستندسازی ---
۲۸ & \multirow{3}{*}{\rotatebox{90}{\RL{\textbf{مستند}}}}
& بانک اطلاعاتی جامع نقض حقوق بشر & الف & ★★★★★ & متوسط \\
۲۹ & & شناسایی مسببین و زنجیره‌ی فرماندهی & الف & ★★★★★ & بالا \\
۳۰ & & حفاظت از شهود و قربانیان & — & ★★★★★ & بسیار بالا \\
\midrule

% --- حوزه‌ی ۹: صلاحیت جهانی ---
۳۱ & \multirow{2}{*}{\rotatebox{90}{\RL{\textbf{صلاحیت}}}}
& تعقیب از طریق صلاحیت جهانی & ب & ★★★★★ & متوسط \\
۳۲ & & شکایت در \en{ICJ} از طریق دولت ثالث & ب & ★★★★☆ & بالا \\
\midrule

% --- حوزه‌ی ۱۰: موضوعات ویژه ---
۳۳ & \multirow{6}{*}{\rotatebox{90}{\RL{\textbf{ویژه}}}}
& اکوساید و جرایم زیست‌محیطی & الف/ج & ★★★☆☆ & متوسط \\
۳۴ & & جرایم سایبری دولتی و قطع اینترنت & الف & ★★★★☆ & کم \\
۳۵ & & فساد سیستماتیک و غارت ثروت ملی & الف/ب & ★★★★★ & متوسط \\
۳۶ & & گروگان‌گیری دولتی & الف/ب & ★★★★★ & کم \\
۳۷ & & نقض حقوق آموزشی و آکادمیک & الف & ★★★☆☆ & کم \\
۳۸ & & نقض آزادی مذهب و عقیده & الف & ★★★★★ & متوسط \\
\midrule

% --- حوزه‌ی ۱۲: الهام از تجربه‌ها ---
۳۹ & \multirow{5}{*}{\rotatebox{90}{\RL{\textbf{تکمیلی}}}}
& رسانه‌ی نفرت‌پراکن و مسئولیت کیفری & الف & ★★★★★ & متوسط \\
۴۰ & & سربازگیری و بسیج کودکان & الف & ★★★★☆ & متوسط \\
۴۱ & & تخریب و غارت میراث فرهنگی & الف & ★★★☆☆ & کم \\
۴۲ & & نقض حقوق بشردوستانه در مخاصمات خارجی & الف & ★★★★☆ & بالا \\
۴۳ & & عدالت جنسیتی و خشونت جنسی سیستماتیک & الف & ★★★★★ & بسیار بالا \\
۴۴ & & حق تعیین سرنوشت و حقوق اقوام & ج & ★★★☆☆ & بسیار بالا \\

\bottomrule
\caption{ماتریس کلان ۴۴ موضوع قابل‌بررسی در دادگاه در سایه}
\end{longtable}
\end{landscape}

\begin{infoBox}[راهنمای جدول]
\begin{itemize}[nosep]
\item \textbf{نوع پرونده:}
الف = اعلام‌دادخواهی داخلی؛
ب = اعلام‌دادخواهی بین‌المللی؛
ج = تحلیل گفتمانی
\item \textbf{اولویت:}
★★★★★ = فوری و حیاتی؛
★☆☆☆☆ = بلندمدت
\item \textbf{حساسیت:}
بسیار بالا = بالاترین خطر سوءاستفاده و اختلاف‌زایی؛
کم = اجماع گسترده
\end{itemize}
\end{infoBox}


% ============================================================
%  نقشه‌ی راه زمانی — اولویت‌بندی مراحل
% ============================================================
\section{نقشه‌ی راه زمانی: اولویت‌بندی موضوعات}

\begin{figure}[H]
	\noindent\makebox[\textwidth][c]{%
\begin{tikzpicture}[
node distance=0.4cm,
phase/.style={
draw=#1, fill=#1!8, rounded corners=5pt,
text width=13cm, minimum height=1.2cm,
font=\small, align=right}]

\node[phase=cPrimary] (p1) {
\RL{\textbf{فاز ۱ — ماه‌های ۱–۶: بنیاد‌گذاری}}\\[3pt]
\RL{موضوعات: ۲۸ (بانک اطلاعاتی)، ۳۰ (حفاظت شهود)،
۲۰ (معیارهای ادله)، ۲۵–۲۶ (همکاری بین‌المللی)}
};

\node[phase=cAccent, below=of p1] (p2) {
\RL{\textbf{فاز ۲ — ماه‌های ۷–۱۲: پرونده‌سازی اولیه}}\\[3pt]
\RL{موضوعات: ۱–۴ (اعدام، شکنجه، بازداشت، سرکوب اعتراضات)،
۱۵ (تحریم‌های مگنیتسکی)، ۲۹ (شناسایی مسببین)،
۳۱ (صلاحیت جهانی)، ۳۹ (رسانه‌ی نفرت‌پراکن)}
};

\node[phase=cWarn, below=of p2] (p3) {
\RL{\textbf{فاز ۳ — ماه‌های ۱۳–۱۸: گسترش دامنه}}\\[3pt]
\RL{موضوعات: ۶–۹ (اقلیت‌ها، زنان، \en{LGBTQ+})،
۱۳ (ارجاع به \en{ICC})، ۲۲–۲۴ (سران و وابستگان)،
۳۵–۳۶ (فساد، گروگان‌گیری)، ۴۲–۴۳ (خشونت جنسی، مخاصمات خارجی)}
};

\node[phase=green!60!black, below=of p3] (p4) {
\RL{\textbf{فاز ۴ — ماه‌های ۱۹–۲۴: موضوعات پیچیده و حساس}}\\[3pt]
\RL{موضوعات: ۱۱–۱۲ (\en{R2P}، مناطق امن)، ۱۴ (دادگاه ویژه)،
۱۹–۲۱ (عوامل و نفوذی‌ها)، ۳۲ (شکایت \en{ICJ})، ۴۴ (حقوق اقوام)}
};

% ← FIX: رنگ فاز ۵ از orange!60 به cConsolid تغییر یافت
\node[phase=cConsolid, below=of p4] (p5) {
	\RL{\textbf{فاز ۵ — سال سوم به بعد: تثبیت و تداوم}}\\[3pt]
	\RL{موضوعات: ۱۶–۱۸ (ارزیابی تحریم‌ها)، ۲۷ (قطعنامه‌ها)،
		۳۳–۳۴ (اکوساید، سایبری)، ۴۰–۴۱ (کودکان سرباز، میراث فرهنگی)}\\[1pt]
	\RL{+ به‌روزرسانی مستمر تمام حوزه‌ها}
};

% فلش‌ها
\draw[-{Stealth}, thick, cPrimary!70] (p1.south) -- (p2.north);
\draw[-{Stealth}, thick, cAccent!70] (p2.south) -- (p3.north);
\draw[-{Stealth}, thick, cWarn!70] (p3.south) -- (p4.north);
\draw[-{Stealth}, thick, cConsolid!70] (p4.south) -- (p5.north);

\end{tikzpicture}}
\caption{نقشه‌ی راه زمانی: اولویت‌بندی موضوعات}
\end{figure}

% ============================================================
%  یادداشت پایانی
% ============================================================
\section*{یادداشت پایانی}

\begin{principleBox}[جمع‌بندی]
این فهرست شامل \textbf{۴۴ موضوع اصلی}
در \textbf{۱۲ حوزه‌ی موضوعی} است که
دادگاه در سایه می‌تواند آنها را
در دستور کار خود قرار دهد.

\medskip
فهرست بر اساس:
\begin{enumerate}[nosep, label=\textcolor{cAccent}{\arabic*.}]
\item موضوعات مطرح‌شده توسط پیشنهاددهنده
\item تجربه‌ی ۱۵ کشور و بحران تاریخی
\item اصول و موازین حقوق بین‌الملل
\item واقعیت‌های سیاسی و عملیاتی
\end{enumerate}
تدوین شده و قابل بازنگری و توسعه است.

\medskip
\textbf{اصول حاکم بر انتخاب و اولویت‌بندی موضوعات:}
\begin{itemize}[nosep]
\item فوریت و شدت نقض حقوق بشر
\item دسترسی به ادله و امکان مستندسازی
\item اثرگذاری بر فشار بین‌المللی
\item اجماع‌پذیری در اپوزیسیون
\item پتانسیل بهره‌برداری منفی حکومت
\item ظرفیت عملیاتی دادگاه در هر مرحله
\end{itemize}
\end{principleBox}

\begin{center}
\rule{0.6\textwidth}{0.5pt}

\vspace{1em}
{\Large\color{cPrimary}
\textbf{پایان پیوست — فهرست موضوعات قابل‌بررسی}}

\vspace{0.5em}
{\small\color{gray}
این پیوست بخشی از سند جامع
«دادگاه در سایه — بسته‌ی پیشنهادی نخبگانی» است.}
\end{center}

% ////////////////////////////////////////////////////////////
%  آخرین صفحه — پس از تمام پیوست‌ها
% ////////////////////////////////////////////////////////////
\clearpage
\thispagestyle{empty}
\vspace*{\fill}
\begin{center}

% ← FIX: Persian text in final TikZ logo wrapped with \RL{}
\begin{tikzpicture}
\draw[cPrimary, very thick] (0,0) circle (2cm);
\node[align=center, font=\large\bfseries, color=cPrimary] at (0,0.3)
{\RL{دادگاه در سایه}};
\node[align=center, font=\small, color=cGold] at (0,-0.4)
{\RL{وجدان منفصل مسئول}};
\end{tikzpicture}

\vspace{24pt}

{\color{cGray}\large
«عدالت تأخیری، عدالت انکاری است.\\[4pt]
اما عدالت شتاب‌زده، عدالت نیست.\\[4pt]
ما راهِ سوم را پیشنهاد می‌کنیم:\\[4pt]
\textbf{عدالت اندیشیده.}»}

\vspace{36pt}

{\small\color{cGray}
تمام حقوق این سند متعلق به هیچ‌کس و همه‌کس است.\\
مهدی سالم؛ ریچموند هیل\\[14pt]
\en{Creative Commons BY-SA 4.0}\\[8pt]
نسخه‌ی ۰.۶ — \today\\[4pt]
برای بحث، نقد، و بهبود.}

\end{center}
\vspace*{\fill}
\end{document}